% Auto-generated by extract_catalog.py — do not edit by hand.

% =====================================================================
\chapter{Comprehensive lemma catalog}
\label{app:catalog}
% =====================================================================

This appendix lists \emph{every} named result in the active build
chain --- 1000 entries across 40 files. Each entry shows the kind
(\texttt{Definition}, \texttt{Lemma}, \dots), the file:line
location (clickable to GitHub), and the \texttt{coqdoc}
docstring extracted from the source. Entries narrated in the
main chapters of this document carry a chapter/section badge.

This catalog is generated mechanically from the \texttt{.v}
files; see \texttt{docs/formal/extract\_catalog.py}. Adding a
docstring to a \texttt{.v} file improves both the catalog and
the deployed coqdoc HTML at the same time.

% --------------------------------------------
\section*{\texttt{altsub.v}}
\addcontentsline{toc}{section}{\texttt{altsub.v}}

\begin{description}
\item[\href{https://github.com/LLM4Rocq/mathcomp-eulerian/blob/main/\detokenize{altsub.v}\#L43}{\texttt{\detokenize{is_turn}}}] \emph{Definition, line 43}. \textsf{[Ch 7 §7.1]} \\
  \texttt{\detokenize{is_turn s i}} is the boolean indicator that interior position     \texttt{\detokenize{i : 'I_n}} is a turning point of \texttt{s : \{perm 'I\_n.+2\}}: the descent     indicator at slot \texttt{\detokenize{widen i}} differs (XOR) from the indicator at slot     \texttt{\detokenize{lift ord0 i}}.  These two slots straddle position \texttt{\detokenize{i}}.

\item[\href{https://github.com/LLM4Rocq/mathcomp-eulerian/blob/main/\detokenize{altsub.v}\#L49}{\texttt{\detokenize{turn_count}}}] \emph{Definition, line 49}. \textsf{[Ch 7 §7.1]} \\
  \texttt{\detokenize{turn_count s}} is the number of turning points of \texttt{\detokenize{s}}, i.e., the     cardinality of \texttt{\detokenize{[set i : 'I_n | is_turn s i]}}.  Stanley \textbackslash{}S\{\}1.6.2: the     "turn count" governing the longest alternating subsequence length.

\item[\href{https://github.com/LLM4Rocq/mathcomp-eulerian/blob/main/\detokenize{altsub.v}\#L69}{\texttt{\detokenize{alt_aux}}}] \emph{Fixpoint, line 69}. \textsf{[aux]} \\
  \texttt{\detokenize{alt_aux b x xs}} checks alternation of the seq \texttt{\detokenize{x :: xs}} given the     expected direction \texttt{\detokenize{b}} of the first comparison (\texttt{\detokenize{true}} = up, \texttt{\detokenize{false}}     = down).  Successive directions flip at each step.

\item[\href{https://github.com/LLM4Rocq/mathcomp-eulerian/blob/main/\detokenize{altsub.v}\#L79}{\texttt{\detokenize{is_alt}}}] \emph{Definition, line 79}. \textsf{[Ch 7 §7.2]} \\
  \texttt{\detokenize{is_alt xs}} holds when \texttt{\detokenize{xs}} is alternating: comparisons of consecutive     elements strictly alternate between \texttt{\detokenize{<}} and \texttt{\detokenize{>}}.  Two starts are     possible (up-down or down-up); both are accepted by disjunction.

\item[\href{https://github.com/LLM4Rocq/mathcomp-eulerian/blob/main/\detokenize{altsub.v}\#L93}{\texttt{\detokenize{perm_seq}}}] \emph{Definition, line 93}. \textsf{[Ch 7 §7.2]} \\
  \texttt{\detokenize{perm_seq s}} is the one-line word \texttt{\detokenize{[s 0; s 1; ...; s (n+1)]}} of     \texttt{s : \{perm 'I\_n.+2\}}, as a \texttt{\detokenize{seq nat}}.

\item[\href{https://github.com/LLM4Rocq/mathcomp-eulerian/blob/main/\detokenize{altsub.v}\#L97}{\texttt{\detokenize{size_perm_seq}}}] \emph{Lemma, line 97}. \textsf{[aux]} \\
  \texttt{\detokenize{perm_seq s}} always has length \texttt{\detokenize{n.+2}}.

\item[\href{https://github.com/LLM4Rocq/mathcomp-eulerian/blob/main/\detokenize{altsub.v}\#L108}{\texttt{\detokenize{pick_seq}}}] \emph{Definition, line 108}. \textsf{[Ch 7 §7.2]} \\
  \texttt{\detokenize{pick_seq s I}} is the subseq of \texttt{\detokenize{perm_seq s}} indexed by the position     set \texttt{I : \{set 'I\_n.+2\}} in ascending order: sort \texttt{\detokenize{I}} by underlying     \texttt{\detokenize{val}}, then read off the \texttt{\detokenize{s}}-values.

\item[\href{https://github.com/LLM4Rocq/mathcomp-eulerian/blob/main/\detokenize{altsub.v}\#L114}{\texttt{\detokenize{as_perm_max}}}] \emph{Definition, line 114}. \textsf{[Ch 7 §7.3]} \\
  \texttt{\detokenize{as_perm_max s}} is the bijective definition of \texttt{\detokenize{as(s)}} from Stanley     \textbackslash{}S\{\}1.6.2: the maximum cardinality of a position set \texttt{\detokenize{I}} whose ordered     image \texttt{\detokenize{pick_seq s I}} is alternating.

\item[\href{https://github.com/LLM4Rocq/mathcomp-eulerian/blob/main/\detokenize{altsub.v}\#L120}{\texttt{\detokenize{as_perm}}}] \emph{Definition, line 120}. \textsf{[Ch 7 §7.3]} \\
  \texttt{\detokenize{as_perm s}} is the DIRECT definition (Path X) of \texttt{\detokenize{as(s)}} as     \texttt{\detokenize{(turn_count s).+2}}.  The headline \texttt{\detokenize{as_perm_max_eq}} proves these     two definitions agree.

\item[\href{https://github.com/LLM4Rocq/mathcomp-eulerian/blob/main/\detokenize{altsub.v}\#L164}{\texttt{\detokenize{is_alt_cons2}}}] \emph{Lemma, line 164}. \textsf{[aux]} \\
  Defining equation of \texttt{\detokenize{is_alt}} on a 2-element prefix \texttt{\detokenize{x :: y :: xs}}:     either start with an ascent or a descent.  Used as a rewrite.

\item[\href{https://github.com/LLM4Rocq/mathcomp-eulerian/blob/main/\detokenize{altsub.v}\#L186}{\texttt{\detokenize{turn_count_le}}}] \emph{Lemma, line 186}. \textsf{[aux]} \\
  Cardinality bound: the turn-set of \texttt{s : \{perm 'I\_n.+2\}} has at     most \texttt{\detokenize{n}} turning points.

\item[\href{https://github.com/LLM4Rocq/mathcomp-eulerian/blob/main/\detokenize{altsub.v}\#L192}{\texttt{\detokenize{as_permE}}}] \emph{Lemma, line 192}. \textsf{[aux]} \\
  Definitional unfolding \texttt{\detokenize{as_perm s = (turn_count s).+2}}; rewrite rule.

\item[\href{https://github.com/LLM4Rocq/mathcomp-eulerian/blob/main/\detokenize{altsub.v}\#L200}{\texttt{\detokenize{turn_count_id}}}] \emph{Lemma, line 200}. \textsf{[aux]} \\
  The identity permutation has no descents, hence no turning points:     \texttt{\detokenize{turn_count 1 = 0}}.

\item[\href{https://github.com/LLM4Rocq/mathcomp-eulerian/blob/main/\detokenize{altsub.v}\#L231}{\texttt{\detokenize{turn_count_alt_desc}}}] \emph{Lemma, line 231}. \textsf{[aux]} \\
  When \texttt{\detokenize{s}} has the alternating descent pattern, every interior     position is a turning point: \texttt{\detokenize{turn_count s = n}}.

\item[\href{https://github.com/LLM4Rocq/mathcomp-eulerian/blob/main/\detokenize{altsub.v}\#L245}{\texttt{\detokenize{as_perm_alt_desc}}}] \emph{Lemma, line 245}. \textsf{[aux]} \\
  When \texttt{\detokenize{s}} has the alternating descent pattern, \texttt{\detokenize{as_perm s}} attains     its maximum value \texttt{\detokenize{n.+2}}.

\item[\href{https://github.com/LLM4Rocq/mathcomp-eulerian/blob/main/\detokenize{altsub.v}\#L346}{\texttt{\detokenize{sign_seq}}}] \emph{Definition, line 346}. \textsf{[aux]} \\
  \texttt{\detokenize{sign_seq xs}} is the boolean direction sequence at each adjacent     pair: \texttt{\detokenize{xs[i] < xs[i+1]}} for \texttt{\detokenize{i = 0, ..., size xs - 2}}.  Empty for     \texttt{\detokenize{xs = nil}}.

\item[\href{https://github.com/LLM4Rocq/mathcomp-eulerian/blob/main/\detokenize{altsub.v}\#L352}{\texttt{\detokenize{size_sign_seq}}}] \emph{Lemma, line 352}. \textsf{[aux]} \\
  Size of \texttt{\detokenize{sign_seq xs}}: one less than \texttt{\detokenize{size xs}}.

\item[\href{https://github.com/LLM4Rocq/mathcomp-eulerian/blob/main/\detokenize{altsub.v}\#L361}{\texttt{\detokenize{flip_count}}}] \emph{Definition, line 361}. \textsf{[aux]} \\
  \texttt{\detokenize{flip_count ss}} is the number of adjacent disagreements (XOR     flips) in a boolean seq \texttt{\detokenize{ss}}.  Counts indices \texttt{\detokenize{i}} with     \texttt{\detokenize{ss[i] != ss[i+1]}}.

\item[\href{https://github.com/LLM4Rocq/mathcomp-eulerian/blob/main/\detokenize{altsub.v}\#L370}{\texttt{\detokenize{turn_count_seq}}}] \emph{Definition, line 370}. \textsf{[aux]} \\
  \texttt{\detokenize{turn_count_seq xs}} is the seq-level turn count: number of direction     flips in \texttt{\detokenize{xs}}, i.e., \texttt{\detokenize{flip_count (sign_seq xs)}}.  The seq analogue     of \texttt{\detokenize{turn_count}} for permutations.

\item[\href{https://github.com/LLM4Rocq/mathcomp-eulerian/blob/main/\detokenize{altsub.v}\#L428}{\texttt{\detokenize{size_sign_seq_perm}}}] \emph{Lemma, line 428}. \textsf{[aux]} \\
  The sign sequence of \texttt{\detokenize{perm_seq s}} has length \texttt{\detokenize{n.+1}}, one slot per     pair of consecutive positions.

\item[\href{https://github.com/LLM4Rocq/mathcomp-eulerian/blob/main/\detokenize{altsub.v}\#L435}{\texttt{\detokenize{enum_ord_split}}}] \emph{Lemma, line 435}. \textsf{[aux]} \\
  Decomposition of \texttt{\detokenize{enum 'I_n.+2}} as \texttt{\detokenize{ord0}} followed by the \texttt{\detokenize{lift     ord0}}-image of \texttt{\detokenize{enum 'I_n.+1}}.  Used to align indexing of \texttt{\detokenize{perm_seq}}     with \texttt{\detokenize{sign_seq}}.

\item[\href{https://github.com/LLM4Rocq/mathcomp-eulerian/blob/main/\detokenize{altsub.v}\#L454}{\texttt{\detokenize{not_is_descentE}}}] \emph{Lemma, line 454}. \textsf{[aux]} \\
  Negated descent indicator as a strict less-than: \texttt{\detokenize{~~ is_descent s i}}     iff \texttt{\detokenize{val (s (widen i)) < val (s (lift ord0 i))}}.

\item[\href{https://github.com/LLM4Rocq/mathcomp-eulerian/blob/main/\detokenize{altsub.v}\#L472}{\texttt{\detokenize{sign_seq_perm_seq}}}] \emph{Lemma, line 472}. \textsf{[aux]} \\
  Bridge identifying \texttt{\detokenize{sign_seq (perm_seq s)}} with the (negated)     descent indicator over \texttt{\detokenize{enum 'I_n.+1}}.  Connects the seq-level     flip-count machinery to the perm-level descent/turn structure.

\item[\href{https://github.com/LLM4Rocq/mathcomp-eulerian/blob/main/\detokenize{altsub.v}\#L513}{\texttt{\detokenize{as_perm_max_le_size}}}] \emph{Lemma, line 513}. \textsf{[aux]} \\
  Trivial size bound: \texttt{\detokenize{as_perm_max s <= n.+2}}, the size of the     underlying \texttt{\detokenize{perm_seq s}}.

\item[\href{https://github.com/LLM4Rocq/mathcomp-eulerian/blob/main/\detokenize{altsub.v}\#L530}{\texttt{\detokenize{sorted_val_enum_ord}}}] \emph{Lemma, line 530}. \textsf{[aux]} \\
  \texttt{\detokenize{enum 'I_m}} is sorted by \texttt{\detokenize{val}} ascending.  Used to identify     \texttt{\detokenize{pick_seq s [set: ...]}} with \texttt{\detokenize{perm_seq s}}.

\item[\href{https://github.com/LLM4Rocq/mathcomp-eulerian/blob/main/\detokenize{altsub.v}\#L542}{\texttt{\detokenize{pick_seq_setT}}}] \emph{Lemma, line 542}. \textsf{[aux]} \\
  Picking the full position set yields the full one-line word:     \texttt{\detokenize{pick_seq s [set: 'I_n.+2] = perm_seq s}}.

\item[\href{https://github.com/LLM4Rocq/mathcomp-eulerian/blob/main/\detokenize{altsub.v}\#L563}{\texttt{\detokenize{is_alt_three}}}] \emph{Lemma, line 563}. \textsf{[aux]} \\
  An alternating seq of length at least \texttt{\detokenize{3}} cannot have its first     three entries strictly monotone in either direction.

\item[\href{https://github.com/LLM4Rocq/mathcomp-eulerian/blob/main/\detokenize{altsub.v}\#L599}{\texttt{\detokenize{slot_iv}}}] \emph{Definition, line 599}. \textsf{[aux]} \\
  \texttt{\detokenize{slot_iv i j}} is the set of descent slots \texttt{\detokenize{k : 'I_n.+1}} whose     underlying \texttt{\detokenize{nat}} satisfies \texttt{\detokenize{val i <= val k < val j}}.  Indexes the     "slot interval" between positions \texttt{\detokenize{i}} and \texttt{\detokenize{j}}.

\item[\href{https://github.com/LLM4Rocq/mathcomp-eulerian/blob/main/\detokenize{altsub.v}\#L603}{\texttt{\detokenize{mem_slot_iv}}}] \emph{Lemma, line 603}. \textsf{[aux]} \\
  Membership in \texttt{\detokenize{slot_iv i j}} unfolded to the \texttt{\detokenize{val}} inequalities.

\item[\href{https://github.com/LLM4Rocq/mathcomp-eulerian/blob/main/\detokenize{altsub.v}\#L619}{\texttt{\detokenize{slot_descent_const}}}] \emph{Lemma, line 619}. \textsf{[aux]} \\
  Constancy of the descent indicator on a turn-free slot interval:     if no turning point of \texttt{\detokenize{s}} lies in slot range [\texttt{\detokenize{val i, val j)}}, then     \texttt{\detokenize{is_descent s k1 = is_descent s k2}} for all slots \texttt{\detokenize{k1, k2}} in the     range.  Used to derive monotone behavior of \texttt{\detokenize{s}} on turn-free zones.

\item[\href{https://github.com/LLM4Rocq/mathcomp-eulerian/blob/main/\detokenize{altsub.v}\#L655}{\texttt{\detokenize{constant_descent_monotone}}}] \emph{Lemma, line 655}. \textsf{[aux]} \\
  Constant-descent implies monotonicity: if \texttt{\detokenize{is_descent s k = b}} for     all slots \texttt{\detokenize{k}} in [\texttt{\detokenize{val p, val q)}}, then \texttt{\detokenize{s}} is monotone on positions     \texttt{\detokenize{(p, q)}}: ascending if \texttt{\detokenize{b = false}}, descending if \texttt{\detokenize{b = true}}.

\item[\href{https://github.com/LLM4Rocq/mathcomp-eulerian/blob/main/\detokenize{altsub.v}\#L727}{\texttt{\detokenize{inter_turn_monotone}}}] \emph{Lemma, line 727}. \textsf{[aux]} \\
  Direction transport across a turn-free slot range: if no turning     point lies in [\texttt{\detokenize{val i, val j)}}, then the \texttt{\detokenize{s}}-comparison direction     on any sub-pair \texttt{\detokenize{(p, q)}} inside \texttt{\detokenize{[i, j]}} matches the boundary     direction at \texttt{\detokenize{(i, j)}}.

\item[\href{https://github.com/LLM4Rocq/mathcomp-eulerian/blob/main/\detokenize{altsub.v}\#L834}{\texttt{\detokenize{turn_iv}}}] \emph{Definition, line 834}. \textsf{[aux]} \\
  \texttt{\detokenize{turn_iv s a b}} is the set of turning points \texttt{\detokenize{t : 'I_n}} of \texttt{\detokenize{s}}     whose witness position \texttt{\detokenize{(val t).+1}} lies in the half-open interval     [\texttt{\detokenize{a, b)}}.

\item[\href{https://github.com/LLM4Rocq/mathcomp-eulerian/blob/main/\detokenize{altsub.v}\#L839}{\texttt{\detokenize{turn_iv_subset}}}] \emph{Lemma, line 839}. \textsf{[aux]} \\
  \texttt{\detokenize{turn_iv s a b}} is a subset of the global turn set of \texttt{\detokenize{s}}.

\item[\href{https://github.com/LLM4Rocq/mathcomp-eulerian/blob/main/\detokenize{altsub.v}\#L923}{\texttt{\detokenize{bool_triangle}}}] \emph{Lemma, line 923}. \textsf{[aux]} \\
  Hamming-style triangle inequality on three booleans:     \texttt{\detokenize{a XOR c <= (a XOR b) + (b XOR c)}}.  Building block for the     flip-count monotonicity arguments.

\item[\href{https://github.com/LLM4Rocq/mathcomp-eulerian/blob/main/\detokenize{altsub.v}\#L931}{\texttt{\detokenize{triangle_xor_nat}}}] \emph{Lemma, line 931}. \textsf{[aux]} \\
  Mirror variant of \texttt{\detokenize{triangle_xor_nat_g}} with the middle position on     the left.  Used for "insert at front" reductions in flip count.

\item[\href{https://github.com/LLM4Rocq/mathcomp-eulerian/blob/main/\detokenize{altsub.v}\#L950}{\texttt{\detokenize{flip_count_pairmap_insert}}}] \emph{Lemma, line 950}. \textsf{[aux]} \\
  Front-insertion monotonicity for \texttt{\detokenize{flip_count}} on a \texttt{\detokenize{pairmap}}:     inserting one element at the front never decreases the flip count.     Step 1 building block for \texttt{\detokenize{flip_count_pairmap_insert_anywhere}}.

\item[\href{https://github.com/LLM4Rocq/mathcomp-eulerian/blob/main/\detokenize{altsub.v}\#L967}{\texttt{\detokenize{flip_step_helper}}}] \emph{Lemma, line 967}. \textsf{[aux]} \\
  Inductive step helper for \texttt{\detokenize{flip_count_pairmap_insert_anywhere}}:     combines the IH bound \texttt{\detokenize{s1 <= ... + s2}} with a structural matching     condition to handle the four \texttt{\detokenize{(x < z) = (y < z)}} cases uniformly.

\item[\href{https://github.com/LLM4Rocq/mathcomp-eulerian/blob/main/\detokenize{altsub.v}\#L1002}{\texttt{\detokenize{flip_count_pairmap_insert_anywhere}}}] \emph{Lemma, line 1002}. \textsf{[aux]} \\
  Insertion-anywhere monotonicity for \texttt{\detokenize{flip_count}} on a \texttt{\detokenize{pairmap}}:     inserting one element at any position never decreases the flip     count.  Generalizes \texttt{\detokenize{flip_count_pairmap_insert}} to interior     positions.

\item[\href{https://github.com/LLM4Rocq/mathcomp-eulerian/blob/main/\detokenize{altsub.v}\#L1027}{\texttt{\detokenize{flip_count_sign_seq_insert_front}}}] \emph{Lemma, line 1027}. \textsf{[aux]} \\
  Sign-seq front-insertion monotonicity: inserting \texttt{\detokenize{y}} at the front     of \texttt{\detokenize{xs}} never decreases \texttt{\detokenize{flip_count (sign_seq _)}}.

\item[\href{https://github.com/LLM4Rocq/mathcomp-eulerian/blob/main/\detokenize{altsub.v}\#L1080}{\texttt{\detokenize{subseq_drop_extra}}}] \emph{Lemma, line 1080}. \textsf{[aux]} \\
  Strict-subseq dropping: if \texttt{\detokenize{xs}} is a proper subseq of \texttt{\detokenize{ys}}, some     index \texttt{\detokenize{i}} of \texttt{\detokenize{ys}} can be removed while still containing \texttt{\detokenize{xs}}.

\item[\href{https://github.com/LLM4Rocq/mathcomp-eulerian/blob/main/\detokenize{altsub.v}\#L1100}{\texttt{\detokenize{size_take_drop_skip}}}] \emph{Lemma, line 1100}. \textsf{[aux]} \\
  Size of \texttt{\detokenize{ys}} with element at position \texttt{\detokenize{i}} removed: \texttt{\detokenize{(size ys).-1}}.

\item[\href{https://github.com/LLM4Rocq/mathcomp-eulerian/blob/main/\detokenize{altsub.v}\#L1114}{\texttt{\detokenize{flip_count_pairmap_le_subseq}}}] \emph{Lemma, line 1114}. \textsf{[aux]} \\
  Subseq monotonicity for \texttt{\detokenize{flip_count}} on a \texttt{\detokenize{pairmap}}: if \texttt{\detokenize{xs}} is a     subseq of \texttt{\detokenize{ys}}, its flip count is at most that of \texttt{\detokenize{ys}}.  Iterates     \texttt{\detokenize{flip_count_pairmap_insert_anywhere}} across the size difference.

\item[\href{https://github.com/LLM4Rocq/mathcomp-eulerian/blob/main/\detokenize{altsub.v}\#L1151}{\texttt{\detokenize{flip_count_sign_seq_le_subseq}}}] \emph{Lemma, line 1151}. \textsf{[aux]} \\
  Subseq monotonicity specialized to \texttt{\detokenize{sign_seq}}: if \texttt{\detokenize{xs}} is a subseq     of \texttt{\detokenize{ys}}, then \texttt{\detokenize{flip_count (sign_seq xs) <= flip_count (sign_seq ys)}}.     Used in \texttt{\detokenize{as_perm_max_upper}} to bound subseq alternation by the global     flip count.

\item[\href{https://github.com/LLM4Rocq/mathcomp-eulerian/blob/main/\detokenize{altsub.v}\#L1183}{\texttt{\detokenize{is_alt_bool_aux}}}] \emph{Fixpoint, line 1183}. \textsf{[aux]} \\
  \texttt{\detokenize{is_alt_bool_aux a xs}} checks alternation of a boolean seq with a     leading "previous" element \texttt{\detokenize{a}}: every adjacent pair must XOR to true.

\item[\href{https://github.com/LLM4Rocq/mathcomp-eulerian/blob/main/\detokenize{altsub.v}\#L1191}{\texttt{\detokenize{is_alt_bool}}}] \emph{Definition, line 1191}. \textsf{[aux]} \\
  \texttt{\detokenize{is_alt_bool xs}} holds when the boolean seq \texttt{\detokenize{xs}} is fully     alternating: every adjacent pair differs.

\item[\href{https://github.com/LLM4Rocq/mathcomp-eulerian/blob/main/\detokenize{altsub.v}\#L1199}{\texttt{\detokenize{flip_count_is_alt_bool}}}] \emph{Lemma, line 1199}. \textsf{[aux]} \\
  Fully alternating boolean seqs maximize flip count:     \texttt{\detokenize{flip_count xs = (size xs).-1}} when \texttt{\detokenize{is_alt_bool xs}} holds.

\item[\href{https://github.com/LLM4Rocq/mathcomp-eulerian/blob/main/\detokenize{altsub.v}\#L1214}{\texttt{\detokenize{sign_seq_is_alt}}}] \emph{Lemma, line 1214}. \textsf{[aux]} \\
  The sign seq of an \texttt{\detokenize{is_alt}} nat seq is itself fully alternating     as a boolean seq: every adjacent direction pair differs.

\item[\href{https://github.com/LLM4Rocq/mathcomp-eulerian/blob/main/\detokenize{altsub.v}\#L1243}{\texttt{\detokenize{is_alt_flip_count}}}] \emph{Lemma, line 1243}. \textsf{[aux]} \\
  Flip-count formula for alternating seqs: an \texttt{\detokenize{is_alt}} seq of size     at least \texttt{\detokenize{2}} has \texttt{\detokenize{flip_count (sign_seq xs) = (size xs).-2}}.  Used     in \texttt{\detokenize{as_perm_max_upper}} to reduce alternation to flip-count bounds.

\item[\href{https://github.com/LLM4Rocq/mathcomp-eulerian/blob/main/\detokenize{altsub.v}\#L1263}{\texttt{\detokenize{sort_enum_strict_sorted}}}] \emph{Lemma, line 1263}. \textsf{[aux]} \\
  Sorting the enumeration of a \texttt{\{set 'I\_n.+2\}} by \texttt{\detokenize{val}} gives a     strictly sorted seq, since the underlying enumeration is     duplicate-free.

\item[\href{https://github.com/LLM4Rocq/mathcomp-eulerian/blob/main/\detokenize{altsub.v}\#L1307}{\texttt{\detokenize{turn_inj}}}] \emph{Definition, line 1307}. \textsf{[aux]} \\
  \texttt{\detokenize{turn_inj t}} injects an interior position \texttt{\detokenize{t : 'I_n}} into     \texttt{\{ 'I\_n.+2 \}} at the offset \texttt{\detokenize{(val t).+1}}.  Used to construct the     witness set \texttt{\detokenize{turn_witness}} for the lower-bound construction.

\item[\href{https://github.com/LLM4Rocq/mathcomp-eulerian/blob/main/\detokenize{altsub.v}\#L1311}{\texttt{\detokenize{val_turn_inj}}}] \emph{Lemma, line 1311}. \textsf{[aux]} \\
  Underlying value of \texttt{\detokenize{turn_inj t}}: \texttt{\detokenize{(val t).+1}}.

\item[\href{https://github.com/LLM4Rocq/mathcomp-eulerian/blob/main/\detokenize{altsub.v}\#L1315}{\texttt{\detokenize{turn_inj_inj}}}] \emph{Lemma, line 1315}. \textsf{[aux]} \\
  \texttt{\detokenize{turn_inj}} is injective.

\item[\href{https://github.com/LLM4Rocq/mathcomp-eulerian/blob/main/\detokenize{altsub.v}\#L1325}{\texttt{\detokenize{turn_witness}}}] \emph{Definition, line 1325}. \textsf{[aux]} \\
  \texttt{\detokenize{turn_witness s}} is the explicit alternating-subseq witness:     \texttt{\{ord0; ord\_max\} \$\textbackslash\{\}cup\$ \{turn\_inj t | t turning point of s\}}.  Has     cardinality \texttt{\detokenize{(turn_count s).+2}}; its \texttt{\detokenize{pick_seq}} is the alternating     subsequence of length \texttt{\detokenize{(turn_count s).+2}}.

\item[\href{https://github.com/LLM4Rocq/mathcomp-eulerian/blob/main/\detokenize{altsub.v}\#L1329}{\texttt{\detokenize{card_turn_image}}}] \emph{Lemma, line 1329}. \textsf{[aux]} \\
  Image of \texttt{\detokenize{turn_inj}} on the turn set has cardinality \texttt{\detokenize{turn_count s}}.

\item[\href{https://github.com/LLM4Rocq/mathcomp-eulerian/blob/main/\detokenize{altsub.v}\#L1335}{\texttt{\detokenize{turn_image_lt_max}}}] \emph{Lemma, line 1335}. \textsf{[aux]} \\
  Range bound: any element of the \texttt{\detokenize{turn_inj}}-image of the turn set     has \texttt{\detokenize{0 < val x <= n}}; in particular, neither \texttt{\detokenize{ord0}} nor \texttt{\detokenize{ord_max}}.

\item[\href{https://github.com/LLM4Rocq/mathcomp-eulerian/blob/main/\detokenize{altsub.v}\#L1345}{\texttt{\detokenize{ord0_notin_turn_image}}}] \emph{Lemma, line 1345}. \textsf{[aux]} \\
  \texttt{\detokenize{ord0}} is not in the \texttt{\detokenize{turn_inj}}-image of the turn set.

\item[\href{https://github.com/LLM4Rocq/mathcomp-eulerian/blob/main/\detokenize{altsub.v}\#L1353}{\texttt{\detokenize{ord_max_notin_turn_image}}}] \emph{Lemma, line 1353}. \textsf{[aux]} \\
  \texttt{\detokenize{ord_max}} is not in the \texttt{\detokenize{turn_inj}}-image of the turn set.

\item[\href{https://github.com/LLM4Rocq/mathcomp-eulerian/blob/main/\detokenize{altsub.v}\#L1363}{\texttt{\detokenize{ord0_neq_ord_max}}}] \emph{Lemma, line 1363}. \textsf{[aux]} \\
  Trivial distinction: \texttt{\detokenize{ord0}} and \texttt{\detokenize{ord_max}} differ in \texttt{\{'I\_n.+2\}}.

\item[\href{https://github.com/LLM4Rocq/mathcomp-eulerian/blob/main/\detokenize{altsub.v}\#L1368}{\texttt{\detokenize{card_turn_witness}}}] \emph{Lemma, line 1368}. \textsf{[aux]} \\
  Cardinality of the lower-bound witness set:     \texttt{\detokenize{#|turn_witness s| = (turn_count s).+2}}.

\item[\href{https://github.com/LLM4Rocq/mathcomp-eulerian/blob/main/\detokenize{altsub.v}\#L1381}{\texttt{\detokenize{slot_descent_const_strict}}}] \emph{Lemma, line 1381}. \textsf{[aux]} \\
  Strict-left variant of \texttt{\detokenize{slot_descent_const}}: hypothesis only excludes     turns with witness STRICTLY greater than \texttt{\detokenize{val i}}.  Used in the     lower-bound construction where \texttt{\detokenize{i}} itself may be a turn position.

\item[\href{https://github.com/LLM4Rocq/mathcomp-eulerian/blob/main/\detokenize{altsub.v}\#L1416}{\texttt{\detokenize{inter_turn_monotone_strict}}}] \emph{Lemma, line 1416}. \textsf{[aux]} \\
  Strict-left variant of \texttt{\detokenize{inter_turn_monotone}}: monotone direction     transport with a strict lower-bound hypothesis on turn witnesses.

\item[\href{https://github.com/LLM4Rocq/mathcomp-eulerian/blob/main/\detokenize{altsub.v}\#L1454}{\texttt{\detokenize{dir_left_of_turn}}}] \emph{Lemma, line 1454}. \textsf{[aux]} \\
  Direction LEFT of a turn witness: when \texttt{\detokenize{b = turn_inj t}} and no     interior turn lies in \texttt{\detokenize{(val a, val b)}}, the direction \texttt{\detokenize{s a}} vs     \texttt{\detokenize{s b}} equals \texttt{\detokenize{~~ is_descent s (widen_ord t)}} (the "left half"     descent indicator at slot \texttt{\detokenize{t}}).

\item[\href{https://github.com/LLM4Rocq/mathcomp-eulerian/blob/main/\detokenize{altsub.v}\#L1484}{\texttt{\detokenize{dir_right_of_turn}}}] \emph{Lemma, line 1484}. \textsf{[aux]} \\
  Direction RIGHT of a turn witness: dual of \texttt{\detokenize{dir_left_of_turn}}; the     direction \texttt{\detokenize{s b}} vs \texttt{\detokenize{s c}} equals \texttt{\detokenize{~~ is_descent s (lift ord0 t)}}     (the "right half" descent indicator at slot \texttt{\detokenize{t}}).

\item[\href{https://github.com/LLM4Rocq/mathcomp-eulerian/blob/main/\detokenize{altsub.v}\#L1517}{\texttt{\detokenize{sign_flip_at_turn}}}] \emph{Lemma, line 1517}. \textsf{[aux]} \\
  Key adjacency lemma for the lower bound: at an interior witness     \texttt{\detokenize{b = turn_inj t}} flanked by neighbors \texttt{\detokenize{a, c}} with no interior turns     in the half-open intervals, the directions across \texttt{\detokenize{(a, b)}} and \texttt{\detokenize{(b, c)}}     DIFFER (since \texttt{\detokenize{t}} is a turning point).  Used in \texttt{\detokenize{triple_flip_pos_seq}}     to verify the alternating property of \texttt{\detokenize{pick_seq s (turn_witness s)}}.

\item[\href{https://github.com/LLM4Rocq/mathcomp-eulerian/blob/main/\detokenize{altsub.v}\#L1543}{\texttt{\detokenize{uniq_adj}}}] \emph{Definition, line 1543}. \textsf{[aux]} \\
  \texttt{\detokenize{uniq_adj xs}} holds when every adjacent pair in \texttt{\detokenize{xs}} consists of     distinct elements.  Weaker than \texttt{\detokenize{uniq}}, used to characterize when     \texttt{\detokenize{is_alt}} follows from sign-seq alternation.

\item[\href{https://github.com/LLM4Rocq/mathcomp-eulerian/blob/main/\detokenize{altsub.v}\#L1552}{\texttt{\detokenize{uniq_adj_head_neq}}}] \emph{Lemma, line 1552}. \textsf{[aux]} \\
  Head distinctness extracted from \texttt{\detokenize{uniq_adj}}.

\item[\href{https://github.com/LLM4Rocq/mathcomp-eulerian/blob/main/\detokenize{altsub.v}\#L1561}{\texttt{\detokenize{alt_aux_from_sign}}}] \emph{Lemma, line 1561}. \textsf{[aux]} \\
  Key reverse direction: an alternating boolean sign seq plus adjacent     distinctness implies the \texttt{\detokenize{alt_aux}} predicate.  Building block for     \texttt{\detokenize{is_alt_from_sign}}.

\item[\href{https://github.com/LLM4Rocq/mathcomp-eulerian/blob/main/\detokenize{altsub.v}\#L1590}{\texttt{\detokenize{sign_seq_alt_of_triple_flip}}}] \emph{Lemma, line 1590}. \textsf{[aux]} \\
  Triple-flip sufficient condition for sign-seq alternation: if every     adjacent triple \texttt{\detokenize{(js[i], js[i+1], js[i+2])}} has flipping     \texttt{\detokenize{s}}-direction, then the sign seq of the \texttt{\detokenize{s}}-image is alternating     as a boolean seq.

\item[\href{https://github.com/LLM4Rocq/mathcomp-eulerian/blob/main/\detokenize{altsub.v}\#L1613}{\texttt{\detokenize{uniq_adj_of_uniq}}}] \emph{Lemma, line 1613}. \textsf{[aux]} \\
  \texttt{\detokenize{uniq_adj}} follows from full \texttt{\detokenize{uniq}} for nat seqs.

\item[\href{https://github.com/LLM4Rocq/mathcomp-eulerian/blob/main/\detokenize{altsub.v}\#L1625}{\texttt{\detokenize{is_alt_from_sign}}}] \emph{Lemma, line 1625}. \textsf{[aux]} \\
  Main reverse characterization: \texttt{\detokenize{is_alt xs}} follows from     \texttt{\detokenize{is_alt_bool (sign_seq xs)}} together with adjacent distinctness.     Used to construct alternating subseqs from sign-flip witnesses.

\item[\href{https://github.com/LLM4Rocq/mathcomp-eulerian/blob/main/\detokenize{altsub.v}\#L1648}{\texttt{\detokenize{pos_seq}}}] \emph{Definition, line 1648}. \textsf{[aux]} \\
  \texttt{\detokenize{pos_seq s}} is the \texttt{\detokenize{turn_witness s}} enumerated in ascending \texttt{\detokenize{val}}     order: the sorted seq of position indices in the lower-bound     witness construction.

\item[\href{https://github.com/LLM4Rocq/mathcomp-eulerian/blob/main/\detokenize{altsub.v}\#L1652}{\texttt{\detokenize{size_pos_seq}}}] \emph{Lemma, line 1652}. \textsf{[aux]} \\
  \texttt{\detokenize{pos_seq s}} has length \texttt{\detokenize{(turn_count s).+2}}.

\item[\href{https://github.com/LLM4Rocq/mathcomp-eulerian/blob/main/\detokenize{altsub.v}\#L1656}{\texttt{\detokenize{pos_seq_uniq}}}] \emph{Lemma, line 1656}. \textsf{[aux]} \\
  \texttt{\detokenize{pos_seq s}} is duplicate-free.

\item[\href{https://github.com/LLM4Rocq/mathcomp-eulerian/blob/main/\detokenize{altsub.v}\#L1661}{\texttt{\detokenize{pick_seq_pos_seq}}}] \emph{Lemma, line 1661}. \textsf{[aux]} \\
  Definitional unfolding: \texttt{\detokenize{pick_seq s (turn_witness s)}} equals the     map of \texttt{\detokenize{pos_seq s}} under \texttt{\detokenize{val (s _)}}.

\item[\href{https://github.com/LLM4Rocq/mathcomp-eulerian/blob/main/\detokenize{altsub.v}\#L1666}{\texttt{\detokenize{pos_seq_strict_sorted}}}] \emph{Lemma, line 1666}. \textsf{[aux]} \\
  \texttt{\detokenize{pos_seq s}} is strictly sorted by \texttt{\detokenize{val}}.

\item[\href{https://github.com/LLM4Rocq/mathcomp-eulerian/blob/main/\detokenize{altsub.v}\#L1671}{\texttt{\detokenize{mem_pos_seq}}}] \emph{Lemma, line 1671}. \textsf{[aux]} \\
  Membership transfer: \texttt{x \textbackslash\{\}in pos\_seq s} iff \texttt{x \textbackslash\{\}in turn\_witness s}.

\item[\href{https://github.com/LLM4Rocq/mathcomp-eulerian/blob/main/\detokenize{altsub.v}\#L1677}{\texttt{\detokenize{pos_seq_val_lt}}}] \emph{Lemma, line 1677}. \textsf{[aux]} \\
  Strict-sorted property: index-order in \texttt{\detokenize{pos_seq s}} matches     \texttt{\detokenize{val}}-order.

\item[\href{https://github.com/LLM4Rocq/mathcomp-eulerian/blob/main/\detokenize{altsub.v}\#L1689}{\texttt{\detokenize{pos_seq_val_lt_inv}}}] \emph{Lemma, line 1689}. \textsf{[aux]} \\
  Converse strict-sorted property: \texttt{\detokenize{val}}-order forces index-order     in \texttt{\detokenize{pos_seq s}}.

\item[\href{https://github.com/LLM4Rocq/mathcomp-eulerian/blob/main/\detokenize{altsub.v}\#L1702}{\texttt{\detokenize{no_inner_in_pos_seq}}}] \emph{Lemma, line 1702}. \textsf{[aux]} \\
  Gap property: no element of \texttt{\detokenize{pos_seq s}} has \texttt{\detokenize{val}} strictly between     consecutive entries \texttt{\detokenize{pos_seq s i}} and \texttt{\detokenize{pos_seq s i.+1}}.

\item[\href{https://github.com/LLM4Rocq/mathcomp-eulerian/blob/main/\detokenize{altsub.v}\#L1719}{\texttt{\detokenize{pos_seq_nth0}}}] \emph{Lemma, line 1719}. \textsf{[aux]} \\
  First element of \texttt{\detokenize{pos_seq s}} is \texttt{\detokenize{ord0}}.

\item[\href{https://github.com/LLM4Rocq/mathcomp-eulerian/blob/main/\detokenize{altsub.v}\#L1734}{\texttt{\detokenize{pos_seq_last_val}}}] \emph{Lemma, line 1734}. \textsf{[aux]} \\
  Last element of \texttt{\detokenize{pos_seq s}} has value \texttt{\detokenize{n.+1}} (i.e., is \texttt{\detokenize{ord_max}}).

\item[\href{https://github.com/LLM4Rocq/mathcomp-eulerian/blob/main/\detokenize{altsub.v}\#L1759}{\texttt{\detokenize{interior_is_turn_inj}}}] \emph{Lemma, line 1759}. \textsf{[aux]} \\
  Interior elements of \texttt{\detokenize{pos_seq s}} (indices \texttt{\detokenize{0 < i < size - 1}}) are     \texttt{\detokenize{turn_inj}}-images of turning points: they correspond to actual     turns of \texttt{\detokenize{s}}.

\item[\href{https://github.com/LLM4Rocq/mathcomp-eulerian/blob/main/\detokenize{altsub.v}\#L1792}{\texttt{\detokenize{triple_flip_pos_seq}}}] \emph{Lemma, line 1792}. \textsf{[aux]} \\
  Triple-flip property for adjacent triples in \texttt{\detokenize{pos_seq s}}: directions     across consecutive \texttt{\detokenize{s}}-comparisons differ.  The interior witness in     each triple is a turning point, so \texttt{\detokenize{sign_flip_at_turn}} applies.

\item[\href{https://github.com/LLM4Rocq/mathcomp-eulerian/blob/main/\detokenize{altsub.v}\#L1846}{\texttt{\detokenize{is_alt_pick_turn_witness}}}] \emph{Theorem, line 1846}. \textsf{[aux]} \\
  Witness alternation: \texttt{\detokenize{pick_seq s (turn_witness s)}} is alternating.     Combines \texttt{\detokenize{triple_flip_pos_seq}} with the sign-seq characterization     of \texttt{\detokenize{is_alt}}.

\item[\href{https://github.com/LLM4Rocq/mathcomp-eulerian/blob/main/\detokenize{altsub.v}\#L1863}{\texttt{\detokenize{as_perm_max_lower}}}] \emph{Theorem, line 1863}. \textsf{[aux]} \\
  Headline lower bound for \texttt{\detokenize{as_perm_max}}:     \texttt{\detokenize{(turn_count s).+2 <= as_perm_max s}}.  Witnessed by \texttt{\detokenize{turn_witness s}},     which has \texttt{\detokenize{(turn_count s).+2}} elements and induces an alternating     \texttt{\detokenize{pick_seq}} (Stanley \textbackslash{}S\{\}1.6.2).

\item[\href{https://github.com/LLM4Rocq/mathcomp-eulerian/blob/main/\detokenize{altsub.v}\#L1883}{\texttt{\detokenize{sort_enum_subseq_enum}}}] \emph{Lemma, line 1883}. \textsf{[aux]} \\
  Sorted enumeration of a subset is a subseq of the sorted full     enumeration: \texttt{\detokenize{sort by val (enum I)}} is a subseq of \texttt{\detokenize{enum 'I_n.+2}}.

\item[\href{https://github.com/LLM4Rocq/mathcomp-eulerian/blob/main/\detokenize{altsub.v}\#L1904}{\texttt{\detokenize{pick_seq_subseq_perm_seq}}}] \emph{Lemma, line 1904}. \textsf{[aux]} \\
  \texttt{\detokenize{pick_seq s I}} is a subseq of \texttt{\detokenize{perm_seq s}} for any position set \texttt{\detokenize{I}}.

\item[\href{https://github.com/LLM4Rocq/mathcomp-eulerian/blob/main/\detokenize{altsub.v}\#L1913}{\texttt{\detokenize{flip_count_pick_le_perm}}}] \emph{Lemma, line 1913}. \textsf{[aux]} \\
  Flip-count comparison for \texttt{\detokenize{pick_seq}} vs \texttt{\detokenize{perm_seq}}: any pick has     \texttt{\detokenize{flip_count}} at most that of the full perm seq.

\item[\href{https://github.com/LLM4Rocq/mathcomp-eulerian/blob/main/\detokenize{altsub.v}\#L1923}{\texttt{\detokenize{flip_count_as_sum}}}] \emph{Lemma, line 1923}. \textsf{[aux]} \\
  \texttt{\detokenize{flip_count}} expressed as an indexed sum of consecutive XOR-pairs.     Used to align with the \texttt{\detokenize{turn_count}} sum over \texttt{\detokenize{is_turn}}.

\item[\href{https://github.com/LLM4Rocq/mathcomp-eulerian/blob/main/\detokenize{altsub.v}\#L1941}{\texttt{\detokenize{neg_add_neg}}}] \emph{Lemma, line 1941}. \textsf{[aux]} \\
  XOR is invariant under double negation: \texttt{\detokenize{~~ a (+) ~~ b = a (+) b}}.

\item[\href{https://github.com/LLM4Rocq/mathcomp-eulerian/blob/main/\detokenize{altsub.v}\#L1948}{\texttt{\detokenize{flip_count_perm_seq_eq_turn_count}}}] \emph{Lemma, line 1948}. \textsf{[aux]} \\
  Bridge identity: \texttt{\detokenize{flip_count (sign_seq (perm_seq s)) = turn_count s}}.     The seq-level flip count of the full word equals the perm-level turn     count, via \texttt{\detokenize{sign_seq_perm_seq}} and the descent-XOR characterization     of \texttt{\detokenize{is_turn}}.

\item[\href{https://github.com/LLM4Rocq/mathcomp-eulerian/blob/main/\detokenize{altsub.v}\#L1986}{\texttt{\detokenize{as_perm_max_upper}}}] \emph{Lemma, line 1986}. \textsf{[aux]} \\
  Headline upper bound for \texttt{\detokenize{as_perm_max}}:     \texttt{\detokenize{as_perm_max s <= (turn_count s).+2}}.  Combines the alternating     \texttt{\detokenize{flip_count}} formula \texttt{\detokenize{is_alt_flip_count}} with the subseq monotonicity     \texttt{\detokenize{flip_count_pick_le_perm}} and the bridge     \texttt{\detokenize{flip_count_perm_seq_eq_turn_count}}.

\item[\href{https://github.com/LLM4Rocq/mathcomp-eulerian/blob/main/\detokenize{altsub.v}\#L2009}{\texttt{\detokenize{as_perm_max_eq}}}] \emph{Theorem, line 2009}. \textsf{[Ch 7 §7.4]} \\
  HEADLINE THEOREM (Stanley \textbackslash{}S\{\}1.6.2): the longest alternating     subsequence count of \texttt{s : \{perm 'I\_n.+2\}} equals \texttt{\detokenize{(turn_count s).+2}}.     That is, \texttt{\detokenize{as(w) = (turn_count w).+2}}, proving the two definitions     \texttt{\detokenize{as_perm_max}} (bijective max) and \texttt{\detokenize{as_perm}} (turn count) coincide.     Combines \texttt{\detokenize{as_perm_max_upper}} and \texttt{\detokenize{as_perm_max_lower}}.

\end{description}

% --------------------------------------------
\section*{\texttt{beta.v}}
\addcontentsline{toc}{section}{\texttt{beta.v}}

\begin{description}
\item[\href{https://github.com/LLM4Rocq/mathcomp-eulerian/blob/main/\detokenize{beta.v}\#L24}{\texttt{\detokenize{beta}}}] \emph{Definition, line 24}. \textsf{[Ch 5 §5.6]} \\
  \texttt{\detokenize{beta n D}} is the set-refined descent count: the number of permutations of     \texttt{\detokenize{'I_n.+1}} whose descent set equals exactly \texttt{\detokenize{D}}. Stanley's \texttt{\detokenize{beta_n(D)}}     (Stanley EC1, S1.4); summing \texttt{\detokenize{beta D}} over \texttt{\detokenize{|D| = k}} recovers     \texttt{\detokenize{eulerian n k}}.

\item[\href{https://github.com/LLM4Rocq/mathcomp-eulerian/blob/main/\detokenize{beta.v}\#L33}{\texttt{\detokenize{descent_set_rev_perm}}}] \emph{Lemma, line 33}. \textsf{[aux]} \\
  Descent set of the reversal: \texttt{\detokenize{D(rev_perm s)}} is the \texttt{\detokenize{rev_ord}}-image of the     complement of \texttt{\detokenize{D(s)}}. Set-level refinement of \texttt{\detokenize{is_descent_rev}}.

\item[\href{https://github.com/LLM4Rocq/mathcomp-eulerian/blob/main/\detokenize{beta.v}\#L47}{\texttt{\detokenize{descent_set_rev_perm_ord}}}] \emph{Lemma, line 47}. \textsf{[aux]} \\
  The reversal permutation \texttt{\detokenize{n,n-1,...,0}} has descent set equal to all     positions.

\item[\href{https://github.com/LLM4Rocq/mathcomp-eulerian/blob/main/\detokenize{beta.v}\#L61}{\texttt{\detokenize{imset_rev_ord_inv}}}] \emph{Lemma, line 61}. \textsf{[aux]} \\
  \texttt{\detokenize{rev_ord}} is involutive on subsets: applying \texttt{\detokenize{imset rev_ord}} twice is     the identity.

\item[\href{https://github.com/LLM4Rocq/mathcomp-eulerian/blob/main/\detokenize{beta.v}\#L71}{\texttt{\detokenize{imset_rev_ord_compl}}}] \emph{Lemma, line 71}. \textsf{[aux]} \\
  \texttt{\detokenize{rev_ord}}-image commutes with set complement.

\item[\href{https://github.com/LLM4Rocq/mathcomp-eulerian/blob/main/\detokenize{beta.v}\#L91}{\texttt{\detokenize{beta0}}}] \emph{Lemma, line 91}. \textsf{[aux]} \\
  \texttt{\detokenize{beta_n(emptyset) = 1}}: only the identity permutation has empty descent     set.

\item[\href{https://github.com/LLM4Rocq/mathcomp-eulerian/blob/main/\detokenize{beta.v}\#L102}{\texttt{\detokenize{descent_set_full_rev}}}] \emph{Lemma, line 102}. \textsf{[aux]} \\
  Only the strictly-decreasing permutation \texttt{\detokenize{rev_perm_ord}} has descent set     equal to all positions.

\item[\href{https://github.com/LLM4Rocq/mathcomp-eulerian/blob/main/\detokenize{beta.v}\#L118}{\texttt{\detokenize{beta_full}}}] \emph{Lemma, line 118}. \textsf{[aux]} \\
  \texttt{\detokenize{beta_n(top) = 1}}: only the reversal permutation has full descent set.

\item[\href{https://github.com/LLM4Rocq/mathcomp-eulerian/blob/main/\detokenize{beta.v}\#L132}{\texttt{\detokenize{sum_beta_eq_fact}}}] \emph{Lemma, line 132}. \textsf{[aux]} \\
  \texttt{\detokenize{sum_D beta_n(D) = (n+1)!}}: the \texttt{\detokenize{beta}} partition of \texttt{\{perm 'I\_n.+1\}} by     descent set covers all permutations.

\item[\href{https://github.com/LLM4Rocq/mathcomp-eulerian/blob/main/\detokenize{beta.v}\#L147}{\texttt{\detokenize{beta_eulerian}}}] \emph{Lemma, line 147}. \textsf{[Ch 5 §5.6]} \\
  Connection to Eulerian numbers: summing \texttt{\detokenize{beta D}} over all \texttt{\detokenize{D}} of     cardinality \texttt{\detokenize{k}} recovers the Eulerian number \texttt{\detokenize{A(n+1, k)}}     (Stanley EC1, S1.4).

\item[\href{https://github.com/LLM4Rocq/mathcomp-eulerian/blob/main/\detokenize{beta.v}\#L166}{\texttt{\detokenize{beta_rev}}}] \emph{Lemma, line 166}. \textsf{[Ch 5 §5.6]} \\
  Reversal-complement symmetry: \texttt{\detokenize{beta_n(D) = beta_n(rev_ord(complement D))}}.     Set-level refinement of \texttt{\detokenize{eulerian_symm}}, proved via the \texttt{\detokenize{rev_perm}}     involution.

\end{description}

% --------------------------------------------
\section*{\texttt{beta\_bridge.v}}
\addcontentsline{toc}{section}{\texttt{beta\_bridge.v}}

\begin{description}
\item[\href{https://github.com/LLM4Rocq/mathcomp-eulerian/blob/main/\detokenize{beta_bridge.v}\#L29}{\texttt{\detokenize{set_to_seq}}}] \emph{Definition, line 29}. \textsf{[aux]} \\
  \texttt{\detokenize{set_to_seq D}} -- convert the finset \texttt{D : \{set 'I\_n\}} to the     sorted \texttt{\detokenize{seq nat}} of underlying positions; bridges the finset     descent world to the seq-based \texttt{\detokenize{psi_cdindex}} machinery.

\item[\href{https://github.com/LLM4Rocq/mathcomp-eulerian/blob/main/\detokenize{beta_bridge.v}\#L35}{\texttt{\detokenize{mem_set_to_seq}}}] \emph{Lemma, line 35}. \textsf{[aux]} \\
  \texttt{\detokenize{mem_set_to_seq}} -- membership in \texttt{\detokenize{set_to_seq D}} equals membership     in the unsorted underlying value list (sorting is order-preserving     on memberships).

\item[\href{https://github.com/LLM4Rocq/mathcomp-eulerian/blob/main/\detokenize{beta_bridge.v}\#L41}{\texttt{\detokenize{mem_set_to_seq_iff}}}] \emph{Lemma, line 41}. \textsf{[aux]} \\
  \texttt{\detokenize{mem_set_to_seq_iff}} -- alternative form of membership in     \texttt{\detokenize{set_to_seq D}} as a \texttt{\detokenize{has}} over \texttt{\detokenize{enum D}}.

\item[\href{https://github.com/LLM4Rocq/mathcomp-eulerian/blob/main/\detokenize{beta_bridge.v}\#L53}{\texttt{\detokenize{mem_set_to_seq_ord}}}] \emph{Lemma, line 53}. \textsf{[aux]} \\
  \texttt{\detokenize{mem_set_to_seq_ord}} -- ordinal-level membership bridge: \texttt{\detokenize{val i}}     is in \texttt{\detokenize{set_to_seq D}} iff \texttt{\detokenize{i}} is in \texttt{\detokenize{D}}.

\item[\href{https://github.com/LLM4Rocq/mathcomp-eulerian/blob/main/\detokenize{beta_bridge.v}\#L67}{\texttt{\detokenize{uniq_set_to_seq}}}] \emph{Lemma, line 67}. \textsf{[aux]} \\
  \texttt{\detokenize{uniq_set_to_seq}} -- the seq-of-positions \texttt{\detokenize{set_to_seq D}} has no     duplicates.

\item[\href{https://github.com/LLM4Rocq/mathcomp-eulerian/blob/main/\detokenize{beta_bridge.v}\#L83}{\texttt{\detokenize{omega_seq_local}}}] \emph{Definition, line 83}. \textsf{[aux]} \\
  \texttt{\detokenize{omega_seq_local s}} -- the omega map at the seq level (mirrors     \texttt{\detokenize{psi_cdindex.omega_seq}}): the list of \texttt{\detokenize{k}} in \texttt{\detokenize{iota 0 (max s).+1}}     such that exactly one of \texttt{\detokenize{k}}, \texttt{\detokenize{k+1}} belongs to \texttt{\detokenize{s}}.

\item[\href{https://github.com/LLM4Rocq/mathcomp-eulerian/blob/main/\detokenize{beta_bridge.v}\#L94}{\texttt{\detokenize{foldr_maxn_set_to_seq_lb}}}] \emph{Lemma, line 94}. \textsf{[aux]} \\
  \texttt{\detokenize{foldr_maxn_set_to_seq_lb}} -- any element of \texttt{\detokenize{set_to_seq D}} is     bounded above by \texttt{\detokenize{foldr maxn 0}} over the same list (used to ensure     the \texttt{\detokenize{iota}} bound in \texttt{\detokenize{omega_seq_local}} covers the relevant range).

\item[\href{https://github.com/LLM4Rocq/mathcomp-eulerian/blob/main/\detokenize{beta_bridge.v}\#L106}{\texttt{\detokenize{omega_set_seq_local_bridge}}}] \emph{Lemma, line 106}. \textsf{[aux]} \\
  \texttt{\detokenize{omega_set_seq_local_bridge}} -- key bridge: \texttt{k \textbackslash\{\}in omega\_set D}     (finset side) iff \texttt{val k \textbackslash\{\}in omega\_seq\_local (set\_to\_seq D)} (seq     side). Connects the two formulations of Stanley's omega map.

\end{description}

% --------------------------------------------
\section*{\texttt{beta\_omega.v}}
\addcontentsline{toc}{section}{\texttt{beta\_omega.v}}

\begin{description}
\item[\href{https://github.com/LLM4Rocq/mathcomp-eulerian/blob/main/\detokenize{beta_omega.v}\#L28}{\texttt{\detokenize{sym_diff}}}] \emph{Definition, line 28}. \textsf{[Ch 8 §8.1]} \\
  \texttt{\detokenize{sym_diff D E}} -- symmetric difference \texttt{(D \textbackslash\{\} E) U (E \textbackslash\{\} D)} of two     finsets over \texttt{\{set 'I\_n\}}.

\item[\href{https://github.com/LLM4Rocq/mathcomp-eulerian/blob/main/\detokenize{beta_omega.v}\#L33}{\texttt{\detokenize{toggle_at}}}] \emph{Definition, line 33}. \textsf{[Ch 8 §8.1]} \\
  \texttt{\detokenize{toggle_at D i}} -- single-position toggle: flips membership of \texttt{\detokenize{i}}     in the descent set \texttt{\detokenize{D}} by symmetric difference with \texttt{\detokenize{[set i]}}.

\item[\href{https://github.com/LLM4Rocq/mathcomp-eulerian/blob/main/\detokenize{beta_omega.v}\#L38}{\texttt{\detokenize{toggle_at_in}}}] \emph{Lemma, line 38}. \textsf{[aux]} \\
  \texttt{\detokenize{toggle_at_in}} -- membership in \texttt{\detokenize{toggle_at D i}} is xor of     \texttt{\detokenize{(i == j)}} and \texttt{(j \textbackslash\{\}in D)}.

\item[\href{https://github.com/LLM4Rocq/mathcomp-eulerian/blob/main/\detokenize{beta_omega.v}\#L47}{\texttt{\detokenize{toggle_atK}}}] \emph{Lemma, line 47}. \textsf{[aux]} \\
  \texttt{\detokenize{toggle_atK}} -- toggling at the same position twice is the     identity.

\item[\href{https://github.com/LLM4Rocq/mathcomp-eulerian/blob/main/\detokenize{beta_omega.v}\#L56}{\texttt{\detokenize{toggle_at_self}}}] \emph{Lemma, line 56}. \textsf{[aux]} \\
  \texttt{\detokenize{toggle_at_self}} -- at the toggled position, membership is     flipped: \texttt{i \textbackslash\{\}in toggle\_at D i} iff \texttt{i \textbackslash\{\}notin D}.

\item[\href{https://github.com/LLM4Rocq/mathcomp-eulerian/blob/main/\detokenize{beta_omega.v}\#L62}{\texttt{\detokenize{toggle_at_other}}}] \emph{Lemma, line 62}. \textsf{[aux]} \\
  \texttt{\detokenize{toggle_at_other}} -- away from the toggled position \texttt{\detokenize{i}},     membership in \texttt{\detokenize{D}} is preserved.

\item[\href{https://github.com/LLM4Rocq/mathcomp-eulerian/blob/main/\detokenize{beta_omega.v}\#L73}{\texttt{\detokenize{omega_set}}}] \emph{Definition, line 73}. \textsf{[Ch 8 §8.2]} \\
  \texttt{\detokenize{omega_set D}} -- Stanley's omega-map (Stanley EC1 \textbackslash{}S\{\}1.6) at the     finset level: \texttt{k \textbackslash\{\}in omega\_set D} iff exactly one of \texttt{\detokenize{k}} and \texttt{\detokenize{k+1}}     belongs to \texttt{\detokenize{D}} (xor on consecutive positions).

\item[\href{https://github.com/LLM4Rocq/mathcomp-eulerian/blob/main/\detokenize{beta_omega.v}\#L78}{\texttt{\detokenize{mem_omega_set}}}] \emph{Lemma, line 78}. \textsf{[aux]} \\
  \texttt{\detokenize{mem_omega_set}} -- unfolds membership of \texttt{\detokenize{k}} in \texttt{\detokenize{omega_set D}} to     the underlying xor of consecutive \texttt{\detokenize{D}}-memberships.

\item[\href{https://github.com/LLM4Rocq/mathcomp-eulerian/blob/main/\detokenize{beta_omega.v}\#L86}{\texttt{\detokenize{toggle_at_omega_bit_i_new}}}] \emph{Lemma, line 86}. \textsf{[aux]} \\
  \texttt{\detokenize{toggle_at_omega_bit_i_new}} -- when consecutive \texttt{\detokenize{i,j}} are both in     \texttt{\detokenize{D}}, toggling \texttt{\detokenize{i}} turns the omega-bit at the connecting position     \texttt{\detokenize{k}} from off to on; produces the witness ordinal \texttt{\detokenize{k}}.

\item[\href{https://github.com/LLM4Rocq/mathcomp-eulerian/blob/main/\detokenize{beta_omega.v}\#L111}{\texttt{\detokenize{toggle_at_omega_strict_superset}}}] \emph{Lemma, line 111}. \textsf{[aux]} \\
  \texttt{\detokenize{toggle_at_omega_strict_superset}} -- key omega-monotonicity step     (left toggle): if \texttt{\detokenize{i,j}} both in \texttt{\detokenize{D}} with \texttt{\detokenize{j = i+1}} and the     predecessor of \texttt{\detokenize{i}} also lies in \texttt{\detokenize{D}} when defined, then toggling     out \texttt{\detokenize{i}} strictly enlarges \texttt{\detokenize{omega_set}}. Feeds into \texttt{\detokenize{beta_swap}} case     analysis.

\item[\href{https://github.com/LLM4Rocq/mathcomp-eulerian/blob/main/\detokenize{beta_omega.v}\#L155}{\texttt{\detokenize{toggle_at_j_omega_bit_i_new}}}] \emph{Lemma, line 155}. \textsf{[aux]} \\
  \texttt{\detokenize{toggle_at_j_omega_bit_i_new}} -- right-toggle analogue of     \texttt{\detokenize{toggle_at_omega_bit_i_new}}: toggling \texttt{\detokenize{j}} in the consecutive pair     \texttt{\detokenize{i,j}} turns the omega-bit at \texttt{\detokenize{k}} from off to on.

\item[\href{https://github.com/LLM4Rocq/mathcomp-eulerian/blob/main/\detokenize{beta_omega.v}\#L181}{\texttt{\detokenize{toggle_at_j_omega_strict_superset}}}] \emph{Lemma, line 181}. \textsf{[aux]} \\
  \texttt{\detokenize{toggle_at_j_omega_strict_superset}} -- right-toggle omega     monotonicity (case A of beta-swap): if \texttt{\detokenize{i,j}} both in \texttt{\detokenize{D}} with     \texttt{\detokenize{j = i+1}} and the successor of \texttt{\detokenize{j}} is also in \texttt{\detokenize{D}} (when defined),     then toggling \texttt{\detokenize{j}} strictly enlarges \texttt{\detokenize{omega_set}}.

\item[\href{https://github.com/LLM4Rocq/mathcomp-eulerian/blob/main/\detokenize{beta_omega.v}\#L225}{\texttt{\detokenize{left_ok}}}] \emph{Definition, line 225}. \textsf{[aux]} \\
  \texttt{\detokenize{left_ok s i l}} -- candidate left endpoint test: \texttt{\detokenize{l <= i}} and every     position between \texttt{\detokenize{l}} and \texttt{\detokenize{i}} is a descent of \texttt{\detokenize{s}}.

\item[\href{https://github.com/LLM4Rocq/mathcomp-eulerian/blob/main/\detokenize{beta_omega.v}\#L231}{\texttt{\detokenize{right_ok}}}] \emph{Definition, line 231}. \textsf{[aux]} \\
  \texttt{\detokenize{right_ok s i r}} -- candidate right endpoint test: \texttt{\detokenize{i <= r}} and     every position between \texttt{\detokenize{i}} and \texttt{\detokenize{r}} is a descent of \texttt{\detokenize{s}}.

\item[\href{https://github.com/LLM4Rocq/mathcomp-eulerian/blob/main/\detokenize{beta_omega.v}\#L237}{\texttt{\detokenize{left_ok_self}}}] \emph{Lemma, line 237}. \textsf{[aux]} \\
  \texttt{\detokenize{left_ok_self}} -- a descent position \texttt{\detokenize{i}} is its own valid left     endpoint candidate.

\item[\href{https://github.com/LLM4Rocq/mathcomp-eulerian/blob/main/\detokenize{beta_omega.v}\#L246}{\texttt{\detokenize{right_ok_self}}}] \emph{Lemma, line 246}. \textsf{[aux]} \\
  \texttt{\detokenize{right_ok_self}} -- a descent position \texttt{\detokenize{i}} is its own valid right     endpoint candidate.

\item[\href{https://github.com/LLM4Rocq/mathcomp-eulerian/blob/main/\detokenize{beta_omega.v}\#L256}{\texttt{\detokenize{block_left}}}] \emph{Definition, line 256}. \textsf{[aux]} \\
  \texttt{\detokenize{block_left s i}} -- left endpoint of the maximal Foata descent     block containing \texttt{\detokenize{i}}: if \texttt{\detokenize{i}} is a descent, the smallest \texttt{\detokenize{l <= i}}     such that all positions in \texttt{\detokenize{[l..i]}} are descents; otherwise \texttt{\detokenize{i}}.

\item[\href{https://github.com/LLM4Rocq/mathcomp-eulerian/blob/main/\detokenize{beta_omega.v}\#L263}{\texttt{\detokenize{block_right}}}] \emph{Definition, line 263}. \textsf{[aux]} \\
  \texttt{\detokenize{block_right s i}} -- right endpoint of the maximal Foata descent     block containing \texttt{\detokenize{i}}: if \texttt{\detokenize{i}} is a descent, the largest \texttt{\detokenize{r >= i}}     such that all positions in \texttt{\detokenize{[i..r]}} are descents; otherwise \texttt{\detokenize{i}}.

\item[\href{https://github.com/LLM4Rocq/mathcomp-eulerian/blob/main/\detokenize{beta_omega.v}\#L269}{\texttt{\detokenize{block_left_left_ok}}}] \emph{Lemma, line 269}. \textsf{[aux]} \\
  \texttt{\detokenize{block_left_left_ok}} -- the chosen \texttt{\detokenize{block_left s i}} satisfies the     \texttt{\detokenize{left_ok}} predicate (when \texttt{\detokenize{i}} is a descent).

\item[\href{https://github.com/LLM4Rocq/mathcomp-eulerian/blob/main/\detokenize{beta_omega.v}\#L277}{\texttt{\detokenize{block_right_right_ok}}}] \emph{Lemma, line 277}. \textsf{[aux]} \\
  \texttt{\detokenize{block_right_right_ok}} -- the chosen \texttt{\detokenize{block_right s i}} satisfies     the \texttt{\detokenize{right_ok}} predicate (when \texttt{\detokenize{i}} is a descent).

\item[\href{https://github.com/LLM4Rocq/mathcomp-eulerian/blob/main/\detokenize{beta_omega.v}\#L286}{\texttt{\detokenize{block_left_min}}}] \emph{Lemma, line 286}. \textsf{[aux]} \\
  \texttt{\detokenize{block_left_min}} -- minimality of \texttt{\detokenize{block_left s i}} among all     valid left-endpoint candidates.

\item[\href{https://github.com/LLM4Rocq/mathcomp-eulerian/blob/main/\detokenize{beta_omega.v}\#L295}{\texttt{\detokenize{block_right_max}}}] \emph{Lemma, line 295}. \textsf{[aux]} \\
  \texttt{\detokenize{block_right_max}} -- maximality of \texttt{\detokenize{block_right s i}} among all     valid right-endpoint candidates.

\item[\href{https://github.com/LLM4Rocq/mathcomp-eulerian/blob/main/\detokenize{beta_omega.v}\#L304}{\texttt{\detokenize{block_left_le}}}] \emph{Lemma, line 304}. \textsf{[aux]} \\
  \texttt{\detokenize{block_left_le}} -- the left endpoint of the block does not exceed     \texttt{\detokenize{i}}.

\item[\href{https://github.com/LLM4Rocq/mathcomp-eulerian/blob/main/\detokenize{beta_omega.v}\#L310}{\texttt{\detokenize{block_right_ge}}}] \emph{Lemma, line 310}. \textsf{[aux]} \\
  \texttt{\detokenize{block_right_ge}} -- the right endpoint of the block is at least     \texttt{\detokenize{i}}.

\item[\href{https://github.com/LLM4Rocq/mathcomp-eulerian/blob/main/\detokenize{beta_omega.v}\#L316}{\texttt{\detokenize{block_left_descent}}}] \emph{Lemma, line 316}. \textsf{[aux]} \\
  \texttt{\detokenize{block_left_descent}} -- every position from \texttt{\detokenize{block_left s i}} up     to \texttt{\detokenize{i}} is a descent of \texttt{\detokenize{s}}.

\item[\href{https://github.com/LLM4Rocq/mathcomp-eulerian/blob/main/\detokenize{beta_omega.v}\#L326}{\texttt{\detokenize{block_right_descent}}}] \emph{Lemma, line 326}. \textsf{[aux]} \\
  \texttt{\detokenize{block_right_descent}} -- every position from \texttt{\detokenize{i}} up to     \texttt{\detokenize{block_right s i}} is a descent of \texttt{\detokenize{s}}.

\item[\href{https://github.com/LLM4Rocq/mathcomp-eulerian/blob/main/\detokenize{beta_omega.v}\#L336}{\texttt{\detokenize{block_descent_chain}}}] \emph{Lemma, line 336}. \textsf{[aux]} \\
  \texttt{\detokenize{block_descent_chain}} -- every position in the full block     \texttt{\detokenize{[block_left s i .. block_right s i]}} is a descent of \texttt{\detokenize{s}}.

\item[\href{https://github.com/LLM4Rocq/mathcomp-eulerian/blob/main/\detokenize{beta_omega.v}\#L358}{\texttt{\detokenize{block_chain_step}}}] \emph{Lemma, line 358}. \textsf{[aux]} \\
  \texttt{\detokenize{block_chain_step}} -- single-step value comparison: for any \texttt{\detokenize{k}} in     the block, \texttt{\detokenize{s}} strictly decreases from position \texttt{\detokenize{k}} to \texttt{\detokenize{k+1}}.

\item[\href{https://github.com/LLM4Rocq/mathcomp-eulerian/blob/main/\detokenize{beta_omega.v}\#L367}{\texttt{\detokenize{block_chain_values}}}] \emph{Lemma, line 367}. \textsf{[aux]} \\
  \texttt{\detokenize{block_chain_values}} -- chain-of-values: the values \texttt{\detokenize{s p}} are     strictly decreasing across the entire Foata block (from     \texttt{\detokenize{block_left s i}} through to \texttt{\detokenize{block_right s i + 1}}).

\end{description}

% --------------------------------------------
\section*{\texttt{beta\_swap.v}}
\addcontentsline{toc}{section}{\texttt{beta\_swap.v}}

\begin{description}
\item[\href{https://github.com/LLM4Rocq/mathcomp-eulerian/blob/main/\detokenize{beta_swap.v}\#L34}{\texttt{\detokenize{alt_desc_set}}}] \emph{Definition, line 34}. \textsf{[Ch 8 §8.3]} \\
  \texttt{\detokenize{alt_desc_set n}} -- the canonical alternating descent set on     \texttt{\{set 'I\_n\}}: the set of even positions \texttt{\{i : \textasciitilde\{\}\textasciitilde\{\} odd i\}}. The     descent set that maximises \texttt{\detokenize{beta}} (Stanley Cor 1.6.5).

\item[\href{https://github.com/LLM4Rocq/mathcomp-eulerian/blob/main/\detokenize{beta_swap.v}\#L39}{\texttt{\detokenize{mem_alt_desc_set}}}] \emph{Lemma, line 39}. \textsf{[aux]} \\
  \texttt{\detokenize{mem_alt_desc_set}} -- membership in \texttt{\detokenize{alt_desc_set n}} reduces to     parity: \texttt{i \textbackslash\{\}in alt\_desc\_set n} iff \texttt{\detokenize{~~ odd i}}.

\item[\href{https://github.com/LLM4Rocq/mathcomp-eulerian/blob/main/\detokenize{beta_swap.v}\#L45}{\texttt{\detokenize{set_is_alt}}}] \emph{Definition, line 45}. \textsf{[Ch 8 §8.3]} \\
  \texttt{\detokenize{set_is_alt D}} -- decidable predicate: \texttt{\detokenize{D}} is "set-alternating"     iff every consecutive ordinal pair has differing \texttt{\detokenize{D}}-memberships.

\item[\href{https://github.com/LLM4Rocq/mathcomp-eulerian/blob/main/\detokenize{beta_swap.v}\#L53}{\texttt{\detokenize{set_not_altP}}}] \emph{Lemma, line 53}. \textsf{[aux]} \\
  \texttt{\detokenize{set_not_altP}} -- reflection: if \texttt{\detokenize{D}} is not set-alternating then     there exist consecutive \texttt{\detokenize{i,j}} with the same \texttt{\detokenize{D}}-membership.

\item[\href{https://github.com/LLM4Rocq/mathcomp-eulerian/blob/main/\detokenize{beta_swap.v}\#L69}{\texttt{\detokenize{toggle_at_compl}}}] \emph{Lemma, line 69}. \textsf{[aux]} \\
  \texttt{\detokenize{toggle_at_compl}} -- the \texttt{\detokenize{toggle_at}} operation commutes with     complementation: \texttt{\detokenize{toggle_at (~: D) i = ~: toggle_at D i}}.

\item[\href{https://github.com/LLM4Rocq/mathcomp-eulerian/blob/main/\detokenize{beta_swap.v}\#L99}{\texttt{\detokenize{compl_perm}}}] \emph{Definition, line 99}. \textsf{[Ch 8 §8.4]} \\
  \texttt{\detokenize{compl_perm s}} -- value-complement permutation: post-compose \texttt{\detokenize{s}}     with \texttt{\detokenize{rev_perm_ord n}} so that \texttt{\detokenize{compl_perm s i = rev_ord (s i)}}.     Bijection used to prove \texttt{\detokenize{beta_compl}} (set complement invariance).

\item[\href{https://github.com/LLM4Rocq/mathcomp-eulerian/blob/main/\detokenize{beta_swap.v}\#L104}{\texttt{\detokenize{compl_permE}}}] \emph{Lemma, line 104}. \textsf{[aux]} \\
  \texttt{\detokenize{compl_permE}} -- evaluation rule for \texttt{\detokenize{compl_perm s}}: it sends \texttt{\detokenize{i}}     to \texttt{\detokenize{rev_ord (s i)}}.

\item[\href{https://github.com/LLM4Rocq/mathcomp-eulerian/blob/main/\detokenize{beta_swap.v}\#L110}{\texttt{\detokenize{compl_perm_inj}}}] \emph{Lemma, line 110}. \textsf{[aux]} \\
  \texttt{\detokenize{compl_perm_inj}} -- \texttt{\detokenize{compl_perm}} is injective on permutations     (right-cancellation by \texttt{\detokenize{rev_perm_ord}}).

\item[\href{https://github.com/LLM4Rocq/mathcomp-eulerian/blob/main/\detokenize{beta_swap.v}\#L115}{\texttt{\detokenize{compl_perm_involutive}}}] \emph{Lemma, line 115}. \textsf{[aux]} \\
  \texttt{\detokenize{compl_perm_involutive}} -- \texttt{\detokenize{compl_perm}} is an involution:     applying it twice is the identity (since \texttt{\detokenize{rev_perm_ord}} is).

\item[\href{https://github.com/LLM4Rocq/mathcomp-eulerian/blob/main/\detokenize{beta_swap.v}\#L126}{\texttt{\detokenize{is_descent_compl}}}] \emph{Lemma, line 126}. \textsf{[aux]} \\
  \texttt{\detokenize{is_descent_compl}} -- value-complementing flips every descent     bit: \texttt{\detokenize{is_descent (compl_perm s) i = ~~ is_descent s i}}.

\item[\href{https://github.com/LLM4Rocq/mathcomp-eulerian/blob/main/\detokenize{beta_swap.v}\#L141}{\texttt{\detokenize{descent_set_compl}}}] \emph{Lemma, line 141}. \textsf{[Ch 8 §8.4]} \\
  \texttt{\detokenize{descent_set_compl}} -- value-complementing complements the     descent set: \texttt{\detokenize{descent_set (compl_perm s) = ~: descent_set s}}.

\item[\href{https://github.com/LLM4Rocq/mathcomp-eulerian/blob/main/\detokenize{beta_swap.v}\#L150}{\texttt{\detokenize{beta_compl}}}] \emph{Lemma, line 150}. \textsf{[Ch 8 §8.4]} \\
  \texttt{\detokenize{beta_compl}} -- the count \texttt{\detokenize{beta D}} is invariant under set     complementation: \texttt{\detokenize{beta D = beta (~: D)}}. Proven via the     value-complement involution \texttt{\detokenize{compl_perm}}.

\item[\href{https://github.com/LLM4Rocq/mathcomp-eulerian/blob/main/\detokenize{beta_swap.v}\#L170}{\texttt{\detokenize{set_is_alt_classify}}}] \emph{Lemma, line 170}. \textsf{[Ch 8 §8.4]} \\
  \texttt{\detokenize{set_is_alt_classify}} -- classification of set-alternating sets     (Stanley \textbackslash{}S\{\}1.6): exactly two such sets exist, \texttt{\detokenize{alt_desc_set n}} and     its complement. Together with \texttt{\detokenize{beta_compl}} this implies all     set-alternating sets share the same beta value.

\item[\href{https://github.com/LLM4Rocq/mathcomp-eulerian/blob/main/\detokenize{beta_swap.v}\#L206}{\texttt{\detokenize{beta_set_is_alt_eq}}}] \emph{Lemma, line 206}. \textsf{[aux]} \\
  \texttt{\detokenize{beta_set_is_alt_eq}} -- any set-alternating \texttt{\detokenize{D'}} has the same     beta as \texttt{\detokenize{alt_desc_set n}} (consequence of \texttt{\detokenize{set_is_alt_classify}}     and \texttt{\detokenize{beta_compl}}).

\item[\href{https://github.com/LLM4Rocq/mathcomp-eulerian/blob/main/\detokenize{beta_swap.v}\#L219}{\texttt{\detokenize{not_set_is_alt_n_ge2}}}] \emph{Lemma, line 219}. \textsf{[aux]} \\
  \texttt{\detokenize{not_set_is_alt_n_ge2}} -- if \texttt{\detokenize{D}} fails \texttt{\detokenize{set_is_alt}} then \texttt{\detokenize{n >= 2}}     (a violating consecutive pair requires at least two positions).

\item[\href{https://github.com/LLM4Rocq/mathcomp-eulerian/blob/main/\detokenize{beta_swap.v}\#L231}{\texttt{\detokenize{omega_set_alt_full}}}] \emph{Lemma, line 231}. \textsf{[aux]} \\
  \texttt{\detokenize{omega_set_alt_full}} -- the omega-set of \texttt{\detokenize{alt_desc_set m.+1}} is     the full set \texttt{\detokenize{[set: 'I_m]}}: every consecutive position pair has     differing memberships, so every omega-bit is on.

\item[\href{https://github.com/LLM4Rocq/mathcomp-eulerian/blob/main/\detokenize{beta_swap.v}\#L249}{\texttt{\detokenize{not_set_is_alt_omega_not_full}}}] \emph{Lemma, line 249}. \textsf{[aux]} \\
  \texttt{\detokenize{not_set_is_alt_omega_not_full}} -- a non-set-alternating \texttt{\detokenize{D}} has     \texttt{\detokenize{omega_set D}} strictly smaller than the full set; the offending     consecutive pair witnesses an "off" omega bit.

\end{description}

% --------------------------------------------
\section*{\texttt{cycles.v}}
\addcontentsline{toc}{section}{\texttt{cycles.v}}

\begin{description}
\item[\href{https://github.com/LLM4Rocq/mathcomp-eulerian/blob/main/\detokenize{cycles.v}\#L24}{\texttt{\detokenize{cycle_count}}}] \emph{Definition, line 24}. \textsf{[Ch 2 §2.1]} \\
  \texttt{\detokenize{cycle_count s}} is the number of cycles in the disjoint-cycle     decomposition of \texttt{\detokenize{s}}, equivalently the cardinality of \texttt{\detokenize{porbits s}},     the set of orbits of \texttt{\detokenize{s}} acting on \texttt{\detokenize{T}} by iteration.  Stanley's \texttt{\detokenize{c(w)}}.

\item[\href{https://github.com/LLM4Rocq/mathcomp-eulerian/blob/main/\detokenize{cycles.v}\#L34}{\texttt{\detokenize{stirling_c}}}] \emph{Definition, line 34}. \textsf{[Ch 2 §2.2]} \\
  \texttt{\detokenize{stirling_c n k}} is the number of permutations of \texttt{\detokenize{'I_n}} with exactly     \texttt{\detokenize{k}} cycles in their disjoint-cycle decomposition.  This is Stanley's     \texttt{\detokenize{c(n, k)}}, the signless Stirling number of the first kind (Stanley     EC1 \textbackslash{}S\{\}1.3.2).

\item[\href{https://github.com/LLM4Rocq/mathcomp-eulerian/blob/main/\detokenize{cycles.v}\#L45}{\texttt{\detokenize{cycle_count_le}}}] \emph{Lemma, line 45}. \textsf{[Ch 2 §2.1]} \\
  Each cycle is non-empty, so \texttt{\detokenize{cycle_count s <= #|T|}}.  Cycle analogue     of \texttt{\detokenize{des_le}} for descents.

\item[\href{https://github.com/LLM4Rocq/mathcomp-eulerian/blob/main/\detokenize{cycles.v}\#L54}{\texttt{\detokenize{porbit_id_singleton}}}] \emph{Lemma, line 54}. \textsf{[aux]} \\
  Under the identity permutation each \texttt{\detokenize{porbit}} is the singleton \texttt{\{x\}}.     Used in \texttt{\detokenize{cycle_count_id}}.

\item[\href{https://github.com/LLM4Rocq/mathcomp-eulerian/blob/main/\detokenize{cycles.v}\#L65}{\texttt{\detokenize{cycle_count_id}}}] \emph{Lemma, line 65}. \textsf{[Ch 2 §2.1]} \\
  The identity permutation on \texttt{\detokenize{T}} has exactly \texttt{\detokenize{#|T|}} cycles, one     singleton per element (Stanley \textbackslash{}S\{\}1.3.1: each fixed point is a 1-cycle).

\item[\href{https://github.com/LLM4Rocq/mathcomp-eulerian/blob/main/\detokenize{cycles.v}\#L86}{\texttt{\detokenize{stirling_c_row_sum}}}] \emph{Lemma, line 86}. \textsf{[aux]} \\
  Row sum: summing \texttt{\detokenize{stirling_c n k}} over all cycle counts \texttt{\detokenize{k <= n}}     recovers the total number of permutations of \texttt{\detokenize{'I_n}}.  Mirrors     \texttt{\detokenize{eulerian_row_sum}} in eulerian.v.

\item[\href{https://github.com/LLM4Rocq/mathcomp-eulerian/blob/main/\detokenize{cycles.v}\#L100}{\texttt{\detokenize{stirling_c_row_sum_fact}}}] \emph{Lemma, line 100}. \textsf{[Ch 2 §2.2]} \\
  Row sum, factorial form: \texttt{\textbackslash\{\}sum\_k c(n, k) = n!}.  Stanley's identity     derived from \texttt{\detokenize{stirling_c_row_sum}} via \texttt{\detokenize{card_Sn}}.

\end{description}

% --------------------------------------------
\section*{\texttt{cycles\_rec.v}}
\addcontentsline{toc}{section}{\texttt{cycles\_rec.v}}

\begin{description}
\item[\href{https://github.com/LLM4Rocq/mathcomp-eulerian/blob/main/\detokenize{cycles_rec.v}\#L55}{\texttt{\detokenize{lift_perm_id_iter}}}] \emph{Lemma, line 55}. \textsf{[aux]} \\
  Iterating \texttt{\detokenize{perm.lift_perm ord_max ord_max s}} at \texttt{\detokenize{ord_max}} stays at     \texttt{\detokenize{ord_max}}.  Used in \texttt{\detokenize{porbit_lift_perm_id_max}}.

\item[\href{https://github.com/LLM4Rocq/mathcomp-eulerian/blob/main/\detokenize{cycles_rec.v}\#L64}{\texttt{\detokenize{porbit_lift_perm_id_max}}}] \emph{Lemma, line 64}. \textsf{[aux]} \\
  The \texttt{\detokenize{porbit}} of \texttt{\detokenize{ord_max}} under \texttt{\detokenize{perm.lift_perm ord_max ord_max s}} is     the singleton \texttt{\{ord\_max\}}.

\item[\href{https://github.com/LLM4Rocq/mathcomp-eulerian/blob/main/\detokenize{cycles_rec.v}\#L76}{\texttt{\detokenize{lift_perm_id_iter_lift}}}] \emph{Lemma, line 76}. \textsf{[aux]} \\
  Iteration on lifted points commutes with \texttt{\detokenize{lift ord_max}}: applying     \texttt{\detokenize{(perm.lift_perm ord_max ord_max s)^+i}} at \texttt{\detokenize{lift ord_max k}} is the     lift of \texttt{\detokenize{(s^+i) k}}.

\item[\href{https://github.com/LLM4Rocq/mathcomp-eulerian/blob/main/\detokenize{cycles_rec.v}\#L86}{\texttt{\detokenize{porbit_lift_perm_id_lift}}}] \emph{Lemma, line 86}. \textsf{[aux]} \\
  The \texttt{\detokenize{porbit}} of \texttt{\detokenize{lift ord_max k}} under the lifted perm is the     \texttt{\detokenize{lift ord_max}}-image of the \texttt{\detokenize{porbit}} of \texttt{\detokenize{k}} under \texttt{\detokenize{s}}.

\item[\href{https://github.com/LLM4Rocq/mathcomp-eulerian/blob/main/\detokenize{cycles_rec.v}\#L100}{\texttt{\detokenize{porbits_lift_perm_id}}}] \emph{Lemma, line 100}. \textsf{[aux]} \\
  Decomposition of \texttt{\detokenize{porbits (perm.lift_perm ord_max ord_max s)}} into     the singleton orbit \texttt{\{ord\_max\}} and the \texttt{\detokenize{lift ord_max}}-images of     the orbits of \texttt{\detokenize{s}}.  Used in \texttt{\detokenize{cycle_count_lift_perm_id}}.

\item[\href{https://github.com/LLM4Rocq/mathcomp-eulerian/blob/main/\detokenize{cycles_rec.v}\#L124}{\texttt{\detokenize{cycle_count_lift_perm_id}}}] \emph{Lemma, line 124}. \textsf{[aux]} \\
  H1 proper: extending \texttt{\detokenize{s}} with \texttt{\detokenize{ord_max}} as a fixed point adds exactly     one cycle.  Building block for the \texttt{\detokenize{j = ord_max}} class of the     Stirling fiber bijection (see \texttt{\detokenize{stirling_fiber.v}}).

\end{description}

% --------------------------------------------
\section*{\texttt{descent.v}}
\addcontentsline{toc}{section}{\texttt{descent.v}}

\begin{description}
\item[\href{https://github.com/LLM4Rocq/mathcomp-eulerian/blob/main/\detokenize{descent.v}\#L20}{\texttt{\detokenize{is_descent}}}] \emph{Definition, line 20}. \textsf{[Ch 5 §5.1]} \\
  \texttt{\detokenize{is_descent s i}} holds when \texttt{\detokenize{s}} descends at position \texttt{\detokenize{i}}, i.e. when     \texttt{\detokenize{s i > s (i+1)}}. Stanley's predicate "\texttt{\detokenize{i}} is a descent of \texttt{\detokenize{w}}"     (Stanley EC1, S1.4).

\item[\href{https://github.com/LLM4Rocq/mathcomp-eulerian/blob/main/\detokenize{descent.v}\#L25}{\texttt{\detokenize{descent_set}}}] \emph{Definition, line 25}. \textsf{[Ch 5 §5.1]} \\
  \texttt{\detokenize{descent_set s}} is the descent set of \texttt{\detokenize{s}}, Stanley's \texttt{\detokenize{D(w)}}     (Stanley EC1, S1.4).

\item[\href{https://github.com/LLM4Rocq/mathcomp-eulerian/blob/main/\detokenize{descent.v}\#L29}{\texttt{\detokenize{des}}}] \emph{Definition, line 29}. \textsf{[Ch 5 §5.1]} \\
  \texttt{\detokenize{des s}} is the descent number of \texttt{\detokenize{s}}, Stanley's \texttt{\detokenize{d(w) = #D(w)}}     (Stanley EC1, S1.4).

\item[\href{https://github.com/LLM4Rocq/mathcomp-eulerian/blob/main/\detokenize{descent.v}\#L32}{\texttt{\detokenize{asc}}}] \emph{Definition, line 32}. \textsf{[Ch 5 §5.1]} \\
  \texttt{\detokenize{asc s = n - des s}} is the ascent number of \texttt{\detokenize{s}}.

\item[\href{https://github.com/LLM4Rocq/mathcomp-eulerian/blob/main/\detokenize{descent.v}\#L37}{\texttt{\detokenize{mem_descent_set}}}] \emph{Lemma, line 37}. \textsf{[aux]} \\
  Membership in \texttt{\detokenize{descent_set s}} is just \texttt{\detokenize{is_descent s i}}; rewrites \texttt{\detokenize{inE}}     form.

\item[\href{https://github.com/LLM4Rocq/mathcomp-eulerian/blob/main/\detokenize{descent.v}\#L41}{\texttt{\detokenize{des_le}}}] \emph{Lemma, line 41}. \textsf{[Ch 5 §5.1]} \\
  Descent number is bounded by \texttt{\detokenize{n}}, the number of consecutive positions.

\item[\href{https://github.com/LLM4Rocq/mathcomp-eulerian/blob/main/\detokenize{descent.v}\#L45}{\texttt{\detokenize{des_add_asc}}}] \emph{Lemma, line 45}. \textsf{[aux]} \\
  Descents and ascents partition the \texttt{\detokenize{n}} consecutive positions.

\item[\href{https://github.com/LLM4Rocq/mathcomp-eulerian/blob/main/\detokenize{descent.v}\#L49}{\texttt{\detokenize{des_id}}}] \emph{Lemma, line 49}. \textsf{[Ch 5 §5.1]} \\
  The identity permutation has no descents.

\item[\href{https://github.com/LLM4Rocq/mathcomp-eulerian/blob/main/\detokenize{descent.v}\#L78}{\texttt{\detokenize{rev_perm_ord}}}] \emph{Definition, line 78}. \textsf{[Ch 5 §5.2]} \\
  \texttt{\detokenize{rev_perm_ord}} is the reversal permutation \texttt{\detokenize{i |-> n - i}} of \texttt{\detokenize{'I_n.+1}}.

\item[\href{https://github.com/LLM4Rocq/mathcomp-eulerian/blob/main/\detokenize{descent.v}\#L81}{\texttt{\detokenize{rev_perm_ordE}}}] \emph{Lemma, line 81}. \textsf{[aux]} \\
  Equivalence with the underlying \texttt{\detokenize{rev_ord}} for rewriting.

\item[\href{https://github.com/LLM4Rocq/mathcomp-eulerian/blob/main/\detokenize{descent.v}\#L87}{\texttt{\detokenize{rev_perm}}}] \emph{Definition, line 87}. \textsf{[Ch 5 §5.2]} \\
  \texttt{\detokenize{rev_perm s}} is the reversal of \texttt{\detokenize{s}} (compose with \texttt{\detokenize{rev_perm_ord}} on     the left). It corresponds to reading the one-line notation of \texttt{\detokenize{s}}     backwards.

\item[\href{https://github.com/LLM4Rocq/mathcomp-eulerian/blob/main/\detokenize{descent.v}\#L90}{\texttt{\detokenize{rev_permE}}}] \emph{Lemma, line 90}. \textsf{[aux]} \\
  Defining equation: \texttt{\detokenize{rev_perm s j = s (rev_ord j)}}.

\item[\href{https://github.com/LLM4Rocq/mathcomp-eulerian/blob/main/\detokenize{descent.v}\#L95}{\texttt{\detokenize{is_descent_rev}}}] \emph{Lemma, line 95}. \textsf{[Ch 5 §5.2]} \\
  Reversal swaps descents and ascents pointwise: \texttt{\detokenize{i}} is a descent of     \texttt{\detokenize{rev_perm s}} iff \texttt{\detokenize{rev_ord i}} is not a descent of \texttt{\detokenize{s}}.

\item[\href{https://github.com/LLM4Rocq/mathcomp-eulerian/blob/main/\detokenize{descent.v}\#L113}{\texttt{\detokenize{des_rev_perm}}}] \emph{Lemma, line 113}. \textsf{[aux]} \\
  Descent count under reversal: \texttt{\detokenize{des (rev_perm s) = n - des s}}.     This is the involution underlying the symmetry     \texttt{\detokenize{eulerian n k = eulerian n (n-k)}}.

\end{description}

% --------------------------------------------
\section*{\texttt{eulerian.v}}
\addcontentsline{toc}{section}{\texttt{eulerian.v}}

\begin{description}
\item[\href{https://github.com/LLM4Rocq/mathcomp-eulerian/blob/main/\detokenize{eulerian.v}\#L16}{\texttt{\detokenize{eulerian}}}] \emph{Definition, line 16}. \textsf{[Ch 5 §5.3]} \\
  \texttt{\detokenize{eulerian n k}} is the Eulerian number \texttt{\detokenize{A(n+1, k)}}: the number of     permutations of \texttt{\detokenize{'I_n.+1}} with exactly \texttt{\detokenize{k}} descents. Stanley's     \texttt{\detokenize{A(n+1, k)}} (Stanley EC1, S1.4).

\item[\href{https://github.com/LLM4Rocq/mathcomp-eulerian/blob/main/\detokenize{eulerian.v}\#L27}{\texttt{\detokenize{eulerian_row_sum}}}] \emph{Lemma, line 27}. \textsf{[aux]} \\
  Row sum: summing \texttt{\detokenize{eulerian n k}} over \texttt{\detokenize{k <= n}} gives the cardinality of     \texttt{\{perm 'I\_n.+1\}}. Partition of permutations by descent count.

\item[\href{https://github.com/LLM4Rocq/mathcomp-eulerian/blob/main/\detokenize{eulerian.v}\#L38}{\texttt{\detokenize{eulerian_row_sum_fact}}}] \emph{Lemma, line 38}. \textsf{[Ch 5 §5.3]} \\
  Row sum equals \texttt{\detokenize{(n+1)!}}: combines \texttt{\detokenize{eulerian_row_sum}} with     \texttt{|S\_\{n+1\}| = (n+1)!}.

\item[\href{https://github.com/LLM4Rocq/mathcomp-eulerian/blob/main/\detokenize{eulerian.v}\#L43}{\texttt{\detokenize{eulerian_out_of_range}}}] \emph{Lemma, line 43}. \textsf{[Ch 5 §5.3]} \\
  Out-of-range vanishing: \texttt{\detokenize{eulerian n k = 0}} when \texttt{\detokenize{k > n}}, since no     permutation has more than \texttt{\detokenize{n}} descents.

\item[\href{https://github.com/LLM4Rocq/mathcomp-eulerian/blob/main/\detokenize{eulerian.v}\#L51}{\texttt{\detokenize{des0_id}}}] \emph{Lemma, line 51}. \textsf{[Ch 5 §5.3]} \\
  Only the identity has zero descents. Used to prove \texttt{\detokenize{eulerian n 0 = 1}}.

\item[\href{https://github.com/LLM4Rocq/mathcomp-eulerian/blob/main/\detokenize{eulerian.v}\#L81}{\texttt{\detokenize{eulerian_n_0}}}] \emph{Lemma, line 81}. \textsf{[aux]} \\
  Boundary value: \texttt{\detokenize{eulerian n 0 = 1}}, counting only the identity.

\item[\href{https://github.com/LLM4Rocq/mathcomp-eulerian/blob/main/\detokenize{eulerian.v}\#L95}{\texttt{\detokenize{rev_perm_inj}}}] \emph{Lemma, line 95}. \textsf{[aux]} \\
  Reversal is injective on permutations (left cancellation by     \texttt{\detokenize{rev_perm_ord}}).

\item[\href{https://github.com/LLM4Rocq/mathcomp-eulerian/blob/main/\detokenize{eulerian.v}\#L99}{\texttt{\detokenize{rev_perm_involutive}}}] \emph{Lemma, line 99}. \textsf{[aux]} \\
  Reversal is an involution: \texttt{\detokenize{rev_perm (rev_perm s) = s}}.

\item[\href{https://github.com/LLM4Rocq/mathcomp-eulerian/blob/main/\detokenize{eulerian.v}\#L104}{\texttt{\detokenize{eulerian_symm}}}] \emph{Lemma, line 104}. \textsf{[aux]} \\
  Symmetry of Eulerian numbers: \texttt{\detokenize{A(n+1, k) = A(n+1, n-k)}} (Stanley EC1, S1.4).     Proved via the reversal involution on \texttt{\{perm 'I\_n.+1\}}.

\item[\href{https://github.com/LLM4Rocq/mathcomp-eulerian/blob/main/\detokenize{eulerian.v}\#L117}{\texttt{\detokenize{eulerian_n_n}}}] \emph{Lemma, line 117}. \textsf{[aux]} \\
  Boundary value: \texttt{\detokenize{eulerian n n = 1}}, counting only the reversal     permutation.

\item[\href{https://github.com/LLM4Rocq/mathcomp-eulerian/blob/main/\detokenize{eulerian.v}\#L137}{\texttt{\detokenize{widenSn_inj}}}] \emph{Lemma, line 137}. \textsf{[aux]} \\
  \texttt{\detokenize{widen_ord (leqnSn n.+1) : 'I_n.+1 -> 'I_n.+2}} is injective.

\item[\href{https://github.com/LLM4Rocq/mathcomp-eulerian/blob/main/\detokenize{eulerian.v}\#L141}{\texttt{\detokenize{widenSn_neq_ord_max}}}] \emph{Lemma, line 141}. \textsf{[aux]} \\
  Widening from \texttt{\detokenize{'I_n.+1}} to \texttt{\detokenize{'I_n.+2}} never reaches \texttt{\detokenize{ord_max}}.

\item[\href{https://github.com/LLM4Rocq/mathcomp-eulerian/blob/main/\detokenize{eulerian.v}\#L156}{\texttt{\detokenize{insert_max_fun}}}] \emph{Definition, line 156}. \textsf{[Ch 5 §5.4]} \\
  Underlying function of \texttt{\detokenize{insert_max_perm t p}}: sends \texttt{\detokenize{p}} to \texttt{\detokenize{ord_max}} and     elsewhere lifts \texttt{\detokenize{t}} via \texttt{\detokenize{widen_ord}}.

\item[\href{https://github.com/LLM4Rocq/mathcomp-eulerian/blob/main/\detokenize{eulerian.v}\#L163}{\texttt{\detokenize{insert_max_fun_p}}}] \emph{Lemma, line 163}. \textsf{[aux]} \\
  \texttt{\detokenize{insert_max_fun}} sends position \texttt{\detokenize{p}} to \texttt{\detokenize{ord_max}}.

\item[\href{https://github.com/LLM4Rocq/mathcomp-eulerian/blob/main/\detokenize{eulerian.v}\#L167}{\texttt{\detokenize{insert_max_fun_lift}}}] \emph{Lemma, line 167}. \textsf{[aux]} \\
  \texttt{\detokenize{insert_max_fun}} on lifted positions widens \texttt{\detokenize{t}}.

\item[\href{https://github.com/LLM4Rocq/mathcomp-eulerian/blob/main/\detokenize{eulerian.v}\#L172}{\texttt{\detokenize{insert_max_fun_inj}}}] \emph{Lemma, line 172}. \textsf{[aux]} \\
  \texttt{\detokenize{insert_max_fun}} is injective: combine injectivity of \texttt{\detokenize{widen}} and \texttt{\detokenize{t}}.

\item[\href{https://github.com/LLM4Rocq/mathcomp-eulerian/blob/main/\detokenize{eulerian.v}\#L185}{\texttt{\detokenize{insert_max_perm}}}] \emph{Definition, line 185}. \textsf{[Ch 5 §5.4]} \\
  \texttt{\detokenize{insert_max_perm t p}} is the permutation of \texttt{\detokenize{'I_n.+2}} obtained from     \texttt{t : \{perm 'I\_n.+1\}} by inserting the value \texttt{\detokenize{ord_max}} at position \texttt{\detokenize{p}}.     Underlies the bijection \texttt{\{perm 'I\_n.+2\} \textasciitilde\{\}= \{perm 'I\_n.+1\} x 'I\_n.+2}.

\item[\href{https://github.com/LLM4Rocq/mathcomp-eulerian/blob/main/\detokenize{eulerian.v}\#L188}{\texttt{\detokenize{insert_max_permE}}}] \emph{Lemma, line 188}. \textsf{[aux]} \\
  Equivalence with the underlying \texttt{\detokenize{insert_max_fun}} for rewriting.

\item[\href{https://github.com/LLM4Rocq/mathcomp-eulerian/blob/main/\detokenize{eulerian.v}\#L192}{\texttt{\detokenize{insert_max_perm_at_p}}}] \emph{Lemma, line 192}. \textsf{[aux]} \\
  \texttt{\detokenize{insert_max_perm t p}} sends \texttt{\detokenize{p}} to \texttt{\detokenize{ord_max}}, the inserted maximum value.

\item[\href{https://github.com/LLM4Rocq/mathcomp-eulerian/blob/main/\detokenize{eulerian.v}\#L196}{\texttt{\detokenize{insert_max_perm_lift}}}] \emph{Lemma, line 196}. \textsf{[aux]} \\
  \texttt{\detokenize{insert_max_perm t p}} on lifted positions widens \texttt{\detokenize{t}}.

\item[\href{https://github.com/LLM4Rocq/mathcomp-eulerian/blob/main/\detokenize{eulerian.v}\#L212}{\texttt{\detokenize{extract_max_ne}}}] \emph{Fact, line 212}. \textsf{[aux]} \\
  Off-position values of \texttt{\detokenize{s}} avoid \texttt{\detokenize{ord_max}} when \texttt{\detokenize{s p = ord_max}}, by     injectivity.

\item[\href{https://github.com/LLM4Rocq/mathcomp-eulerian/blob/main/\detokenize{eulerian.v}\#L221}{\texttt{\detokenize{extract_max_fun}}}] \emph{Definition, line 221}. \textsf{[Ch 5 §5.4]} \\
  Underlying function of \texttt{\detokenize{extract_max_perm}}: removes the value \texttt{\detokenize{ord_max}}     (located at position \texttt{\detokenize{p}}) and reindexes the remaining values onto     \texttt{\detokenize{'I_n.+1}}.

\item[\href{https://github.com/LLM4Rocq/mathcomp-eulerian/blob/main/\detokenize{eulerian.v}\#L225}{\texttt{\detokenize{extract_max_funE}}}] \emph{Lemma, line 225}. \textsf{[aux]} \\
  Defining equation: widening \texttt{\detokenize{extract_max_fun j}} recovers \texttt{\detokenize{s (lift p j)}}.

\item[\href{https://github.com/LLM4Rocq/mathcomp-eulerian/blob/main/\detokenize{eulerian.v}\#L235}{\texttt{\detokenize{extract_max_fun_inj}}}] \emph{Lemma, line 235}. \textsf{[aux]} \\
  \texttt{\detokenize{extract_max_fun}} is injective.

\item[\href{https://github.com/LLM4Rocq/mathcomp-eulerian/blob/main/\detokenize{eulerian.v}\#L244}{\texttt{\detokenize{extract_max_perm}}}] \emph{Definition, line 244}. \textsf{[Ch 5 §5.4]} \\
  \texttt{\detokenize{extract_max_perm sp}} is the permutation of \texttt{\detokenize{'I_n.+1}} obtained from     \texttt{s : \{perm 'I\_n.+2\}} (with \texttt{\detokenize{s p = ord_max}}) by deleting the value     \texttt{\detokenize{ord_max}}.

\item[\href{https://github.com/LLM4Rocq/mathcomp-eulerian/blob/main/\detokenize{eulerian.v}\#L247}{\texttt{\detokenize{extract_max_permE}}}] \emph{Lemma, line 247}. \textsf{[aux]} \\
  Equivalence with the underlying \texttt{\detokenize{extract_max_fun}} for rewriting.

\item[\href{https://github.com/LLM4Rocq/mathcomp-eulerian/blob/main/\detokenize{eulerian.v}\#L251}{\texttt{\detokenize{extract_max_widen}}}] \emph{Lemma, line 251}. \textsf{[aux]} \\
  Defining equation: widening \texttt{\detokenize{extract_max_perm j}} recovers \texttt{\detokenize{s (lift p j)}}.

\item[\href{https://github.com/LLM4Rocq/mathcomp-eulerian/blob/main/\detokenize{eulerian.v}\#L268}{\texttt{\detokenize{insert_extract_max}}}] \emph{Lemma, line 268}. \textsf{[aux]} \\
  Right inverse: inserting the extracted max back at position \texttt{\detokenize{p}}     recovers \texttt{\detokenize{s}}.

\item[\href{https://github.com/LLM4Rocq/mathcomp-eulerian/blob/main/\detokenize{eulerian.v}\#L288}{\texttt{\detokenize{des_insert_max_ord0}}}] \emph{Lemma, line 288}. \textsf{[aux]} \\
  Inserting \texttt{\detokenize{ord_max}} at position \texttt{\detokenize{0}} adds exactly one descent (a fresh     descent is created at position \texttt{\detokenize{0}} from the new max value).

\item[\href{https://github.com/LLM4Rocq/mathcomp-eulerian/blob/main/\detokenize{eulerian.v}\#L317}{\texttt{\detokenize{des_insert_max_ord_max}}}] \emph{Lemma, line 317}. \textsf{[aux]} \\
  Inserting \texttt{\detokenize{ord_max}} at position \texttt{\detokenize{ord_max}} (the rightmost position) does not     change the descent count.

\item[\href{https://github.com/LLM4Rocq/mathcomp-eulerian/blob/main/\detokenize{eulerian.v}\#L362}{\texttt{\detokenize{des_insert_max_interior}}}] \emph{Lemma, line 362}. \textsf{[aux]} \\
  Inserting \texttt{\detokenize{ord_max}} at an interior position above \texttt{\detokenize{j}} adds a descent     iff \texttt{\detokenize{j}} was an ascent in \texttt{\detokenize{t}}. Key combinatorial lemma for the     Eulerian recurrence.

\item[\href{https://github.com/LLM4Rocq/mathcomp-eulerian/blob/main/\detokenize{eulerian.v}\#L434}{\texttt{\detokenize{insert_max_perm_fiber}}}] \emph{Lemma, line 434}. \textsf{[aux]} \\
  The inverse of \texttt{\detokenize{insert_max_perm t p}} sends \texttt{\detokenize{ord_max}} back to \texttt{\detokenize{p}}:     identifies the fiber of position-of-max under insertion.

\item[\href{https://github.com/LLM4Rocq/mathcomp-eulerian/blob/main/\detokenize{eulerian.v}\#L450}{\texttt{\detokenize{sum_descent}}}] \emph{Lemma, line 450}. \textsf{[aux]} \\
  The descent count \texttt{\detokenize{des t}} equals the sum of descent indicators     over positions.

\item[\href{https://github.com/LLM4Rocq/mathcomp-eulerian/blob/main/\detokenize{eulerian.v}\#L493}{\texttt{\detokenize{eulerian_rec}}}] \emph{Lemma, line 493}. \textsf{[Ch 5 §5.5]} \\
  Eulerian recurrence (Stanley EC1, S1.4):     \texttt{\detokenize{A(n+2, k+1) = (k+2) A(n+1, k+1) + (n+1-k) A(n+1, k)}}.     Proved via the insert-max bijection.

\item[\href{https://github.com/LLM4Rocq/mathcomp-eulerian/blob/main/\detokenize{eulerian.v}\#L538}{\texttt{\detokenize{worpitzky_binom_id}}}] \emph{Lemma, line 538}. \textsf{[aux]} \\
  Key algebraic identity for Worpitzky's induction step (valid for \texttt{\detokenize{k <= n}}):     \texttt{\detokenize{x * C(x+k, n+1) = (k+1) C(x+k, n+2) + (n+1-k) C(x+k+1, n+2)}}.

\item[\href{https://github.com/LLM4Rocq/mathcomp-eulerian/blob/main/\detokenize{eulerian.v}\#L555}{\texttt{\detokenize{worpitzky}}}] \emph{Lemma, line 555}. \textsf{[aux]} \\
  Worpitzky's identity (Stanley EC1, S1.4):     \texttt{\detokenize{x^(n+1) = sum_(k < n+1) A(n+1, k) * C(x+k, n+1)}}.     Proved by induction on \texttt{\detokenize{n}} using \texttt{\detokenize{eulerian_rec}} and \texttt{\detokenize{worpitzky_binom_id}}.

\item[\href{https://github.com/LLM4Rocq/mathcomp-eulerian/blob/main/\detokenize{eulerian.v}\#L593}{\texttt{\detokenize{binS}}}] \emph{Lemma, line 593}. \textsf{[aux]} \\
  Signed Pascal extension: a saturated-arithmetic-friendly form of Pascal's     rule, valid for all \texttt{\detokenize{t}} (including \texttt{\detokenize{t = 0}}, where \texttt{\detokenize{t.-1 = 0}}).

\item[\href{https://github.com/LLM4Rocq/mathcomp-eulerian/blob/main/\detokenize{eulerian.v}\#L598}{\texttt{\detokenize{aux_id_step}}}] \emph{Lemma, line 598}. \textsf{[aux]} \\
  Key recurrence for the alternating binomial sum: \texttt{\detokenize{g(N+1, u) = g(N, u.-1)}}.     Proved via Pascal applied to both binomial factors plus a reindex.

\item[\href{https://github.com/LLM4Rocq/mathcomp-eulerian/blob/main/\detokenize{eulerian.v}\#L626}{\texttt{\detokenize{aux_id}}}] \emph{Lemma, line 626}. \textsf{[aux]} \\
  Alternating binomial convolution identity:     \texttt{\detokenize{sum_j (-1)^j C(n+2, j) C(t-j, n+1) = [t == n+1]}}.     Used in the inversion step of \texttt{\detokenize{eulerian_explicit}}.

\item[\href{https://github.com/LLM4Rocq/mathcomp-eulerian/blob/main/\detokenize{eulerian.v}\#L669}{\texttt{\detokenize{eulerian_explicit}}}] \emph{Lemma, line 669}. \textsf{[aux]} \\
  Eulerian explicit formula (Stanley EC1, S1.4):     \texttt{\detokenize{A(n+1, k) = sum_(j <= k) (-1)^j C(n+2, j) (k+1-j)^(n+1)}}.     Proved by Worpitzky inversion against \texttt{\detokenize{aux_id}}.

\end{description}

% --------------------------------------------
\section*{\texttt{foata.v}}
\addcontentsline{toc}{section}{\texttt{foata.v}}

\begin{description}
\item[\href{https://github.com/LLM4Rocq/mathcomp-eulerian/blob/main/\detokenize{foata.v}\#L48}{\texttt{\detokenize{cyc_last_to_front}}}] \emph{Definition, line 48}. \textsf{[aux]} \\
  \texttt{\detokenize{cyc_last_to_front s}} moves the last letter of \texttt{\detokenize{s}} to the front.     Empty sequence stays empty.  Building block for \texttt{\detokenize{foata_step}}.

\item[\href{https://github.com/LLM4Rocq/mathcomp-eulerian/blob/main/\detokenize{foata.v}\#L55}{\texttt{\detokenize{split_blocks_aux}}}] \emph{Fixpoint, line 55}. \textsf{[aux]} \\
  \texttt{\detokenize{split_blocks_aux P cur s}} is the tail-recursive worker for     \texttt{\detokenize{split_blocks}}, threading the in-progress block \texttt{\detokenize{cur}}.  A boundary     occurs after each \texttt{\detokenize{P}}-letter.

\item[\href{https://github.com/LLM4Rocq/mathcomp-eulerian/blob/main/\detokenize{foata.v}\#L67}{\texttt{\detokenize{split_blocks}}}] \emph{Definition, line 67}. \textsf{[aux]} \\
  \texttt{\detokenize{split_blocks P s}} cuts \texttt{\detokenize{s}} into maximal blocks each ending with a     \texttt{\detokenize{P}}-letter (except possibly the last block).  Concatenating the     blocks recovers \texttt{\detokenize{s}}.

\item[\href{https://github.com/LLM4Rocq/mathcomp-eulerian/blob/main/\detokenize{foata.v}\#L74}{\texttt{\detokenize{foata_step}}}] \emph{Definition, line 74}. \textsf{[Ch 4 §4.1]} \\
  \texttt{\detokenize{foata_step a u}} is one step of the Foata bijection: compare the last     letter of \texttt{\detokenize{u}} to \texttt{\detokenize{a}}, split \texttt{\detokenize{u}} into blocks via \texttt{\detokenize{split_blocks}} (with     predicate \texttt{\detokenize{< a}} or \texttt{\detokenize{a <}} accordingly), cyclically rotate each block     via \texttt{\detokenize{cyc_last_to_front}}, and append \texttt{\detokenize{a}}.  See Stanley EC1 \textbackslash{}S\{\}1.3.4.

\item[\href{https://github.com/LLM4Rocq/mathcomp-eulerian/blob/main/\detokenize{foata.v}\#L86}{\texttt{\detokenize{foata}}}] \emph{Definition, line 86}. \textsf{[Ch 4 §4.1]} \\
  \texttt{\detokenize{foata w}} is Foata's first fundamental bijection at the word level:     iterate \texttt{\detokenize{foata_step}} left-to-right starting from the empty word.     Satisfies \texttt{\detokenize{inv_seq (foata w) = maj_seq w}} (Theorem \texttt{\detokenize{foata_inv_eq_maj}}).

\item[\href{https://github.com/LLM4Rocq/mathcomp-eulerian/blob/main/\detokenize{foata.v}\#L95}{\texttt{\detokenize{is_desc_seq}}}] \emph{Definition, line 95}. \textsf{[Ch 4 §4.1]} \\
  \texttt{\detokenize{is_desc_seq w k}} holds iff position \texttt{\detokenize{k}} is a descent of \texttt{\detokenize{w}}, i.e.     \texttt{w\_k > w\_\{k+1\}}.

\item[\href{https://github.com/LLM4Rocq/mathcomp-eulerian/blob/main/\detokenize{foata.v}\#L100}{\texttt{\detokenize{maj_seq}}}] \emph{Definition, line 100}. \textsf{[Ch 4 §4.1]} \\
  \texttt{\detokenize{maj_seq w}} is the major index of \texttt{\detokenize{w}} at the seq level: sum of     1-indexed descent positions \texttt{\detokenize{k+1}} over \texttt{\detokenize{is_desc_seq w k}}.

\item[\href{https://github.com/LLM4Rocq/mathcomp-eulerian/blob/main/\detokenize{foata.v}\#L105}{\texttt{\detokenize{inv_seq}}}] \emph{Definition, line 105}. \textsf{[Ch 4 §4.1]} \\
  \texttt{\detokenize{inv_seq w}} is the inversion count of \texttt{\detokenize{w}} at the seq level: number of     pairs \texttt{\detokenize{(i, j)}} with \texttt{\detokenize{i < j < size w}} and \texttt{\detokenize{w_i > w_j}}.

\item[\href{https://github.com/LLM4Rocq/mathcomp-eulerian/blob/main/\detokenize{foata.v}\#L110}{\texttt{\detokenize{count_gt}}}] \emph{Definition, line 110}. \textsf{[aux]} \\
  \texttt{\detokenize{count_gt a w}} is the number of letters of \texttt{\detokenize{w}} strictly greater than \texttt{\detokenize{a}}.

\item[\href{https://github.com/LLM4Rocq/mathcomp-eulerian/blob/main/\detokenize{foata.v}\#L143}{\texttt{\detokenize{cyc_last_to_front_perm_eq}}}] \emph{Lemma, line 143}. \textsf{[aux]} \\
  \texttt{\detokenize{cyc_last_to_front}} preserves the multiset of letters.

\item[\href{https://github.com/LLM4Rocq/mathcomp-eulerian/blob/main/\detokenize{foata.v}\#L155}{\texttt{\detokenize{split_blocks_aux_flatten}}}] \emph{Lemma, line 155}. \textsf{[aux]} \\
  Flattening the blocks of \texttt{\detokenize{split_blocks_aux}} yields \texttt{\detokenize{cur ++ s}}; the     accumulator-aware version of \texttt{\detokenize{split_blocks_flatten}}.

\item[\href{https://github.com/LLM4Rocq/mathcomp-eulerian/blob/main/\detokenize{foata.v}\#L166}{\texttt{\detokenize{split_blocks_flatten}}}] \emph{Lemma, line 166}. \textsf{[aux]} \\
  Concatenating the blocks of \texttt{\detokenize{split_blocks P s}} recovers \texttt{\detokenize{s}}.

\item[\href{https://github.com/LLM4Rocq/mathcomp-eulerian/blob/main/\detokenize{foata.v}\#L172}{\texttt{\detokenize{perm_eq_flatten_map_cyc}}}] \emph{Lemma, line 172}. \textsf{[aux]} \\
  Rotating each block via \texttt{\detokenize{cyc_last_to_front}} preserves the multiset     of the flattened sequence.

\item[\href{https://github.com/LLM4Rocq/mathcomp-eulerian/blob/main/\detokenize{foata.v}\#L181}{\texttt{\detokenize{foata_step_perm_eq}}}] \emph{Lemma, line 181}. \textsf{[aux]} \\
  \texttt{\detokenize{foata_step a u}} is a permutation of \texttt{\detokenize{rcons u a}}: the step rearranges     letters but preserves the multiset (with \texttt{\detokenize{a}} inserted at the end).

\item[\href{https://github.com/LLM4Rocq/mathcomp-eulerian/blob/main/\detokenize{foata.v}\#L202}{\texttt{\detokenize{foata_perm_eq}}}] \emph{Lemma, line 202}. \textsf{[aux]} \\
  \texttt{\detokenize{foata w}} is a permutation of \texttt{\detokenize{w}}: the bijection preserves the     multiset of letters.

\item[\href{https://github.com/LLM4Rocq/mathcomp-eulerian/blob/main/\detokenize{foata.v}\#L217}{\texttt{\detokenize{foata_size}}}] \emph{Lemma, line 217}. \textsf{[aux]} \\
  \texttt{\detokenize{foata}} preserves length.

\item[\href{https://github.com/LLM4Rocq/mathcomp-eulerian/blob/main/\detokenize{foata.v}\#L221}{\texttt{\detokenize{foata_uniq}}}] \emph{Lemma, line 221}. \textsf{[aux]} \\
  \texttt{\detokenize{foata}} preserves uniqueness.

\item[\href{https://github.com/LLM4Rocq/mathcomp-eulerian/blob/main/\detokenize{foata.v}\#L232}{\texttt{\detokenize{inv_seq_rcons}}}] \emph{Lemma, line 232}. \textsf{[aux]} \\
  Classical recursion: appending letter \texttt{\detokenize{a}} at the end of \texttt{\detokenize{w}} increases     \texttt{\detokenize{inv_seq}} by \texttt{\detokenize{count_gt a w}}, the number of letters of \texttt{\detokenize{w}} above \texttt{\detokenize{a}}.

\item[\href{https://github.com/LLM4Rocq/mathcomp-eulerian/blob/main/\detokenize{foata.v}\#L262}{\texttt{\detokenize{count_gt_perm_eq}}}] \emph{Lemma, line 262}. \textsf{[aux]} \\
  \texttt{\detokenize{count_gt}} only depends on the multiset (preserved under \texttt{\detokenize{perm_eq}}).

\item[\href{https://github.com/LLM4Rocq/mathcomp-eulerian/blob/main/\detokenize{foata.v}\#L278}{\texttt{\detokenize{maj_seq_rcons}}}] \emph{Lemma, line 278}. \textsf{[aux]} \\
  Appending letter \texttt{\detokenize{a}} to nonempty \texttt{\detokenize{w}} adds \texttt{\detokenize{size w}} to \texttt{\detokenize{maj_seq}} iff     the previous last letter exceeds \texttt{\detokenize{a}} (creating a new descent at position     \texttt{\detokenize{size w - 1}}, with 1-indexed contribution \texttt{\detokenize{size w}}); otherwise \texttt{\detokenize{maj_seq}}     is unchanged.

\item[\href{https://github.com/LLM4Rocq/mathcomp-eulerian/blob/main/\detokenize{foata.v}\#L311}{\texttt{\detokenize{foata_step_last}}}] \emph{Lemma, line 311}. \textsf{[aux]} \\
  \texttt{\detokenize{foata_step a u}} always ends with \texttt{\detokenize{a}}.

\item[\href{https://github.com/LLM4Rocq/mathcomp-eulerian/blob/main/\detokenize{foata.v}\#L321}{\texttt{\detokenize{foata_rcons}}}] \emph{Lemma, line 321}. \textsf{[aux]} \\
  \texttt{\detokenize{foata}} commutes with right-extension: \texttt{\detokenize{foata (rcons w a) =     foata_step a (foata w)}}.  Used for induction on \texttt{\detokenize{w}} from the right.

\item[\href{https://github.com/LLM4Rocq/mathcomp-eulerian/blob/main/\detokenize{foata.v}\#L331}{\texttt{\detokenize{foata_last}}}] \emph{Lemma, line 331}. \textsf{[aux]} \\
  Last letter of \texttt{\detokenize{foata (a :: w)}} equals last letter of \texttt{\detokenize{a :: w}}     (i.e. \texttt{\detokenize{last a w}}).

\item[\href{https://github.com/LLM4Rocq/mathcomp-eulerian/blob/main/\detokenize{foata.v}\#L350}{\texttt{\detokenize{foata_last_eq}}}] \emph{Lemma, line 350}. \textsf{[aux]} \\
  General form for the rcons induction: when \texttt{\detokenize{w}} is non-empty,     \texttt{\detokenize{last (foata w) = last w}} for any default.

\item[\href{https://github.com/LLM4Rocq/mathcomp-eulerian/blob/main/\detokenize{foata.v}\#L364}{\texttt{\detokenize{cross_inv}}}] \emph{Definition, line 364}. \textsf{[aux]} \\
  \texttt{\detokenize{cross_inv s1 s2}} counts cross-inversions between \texttt{\detokenize{s1}} and \texttt{\detokenize{s2}}:     pairs \texttt{\detokenize{(i, j)}} with \texttt{\detokenize{s1_i > s2_j}}.  Equals the inversions of     \texttt{\detokenize{s1 ++ s2}} that straddle the join point.

\item[\href{https://github.com/LLM4Rocq/mathcomp-eulerian/blob/main/\detokenize{foata.v}\#L368}{\texttt{\detokenize{cross_inv_nil_r}}}] \emph{Lemma, line 368}. \textsf{[aux]} \\
  \texttt{\detokenize{cross_inv}} is zero on the right when \texttt{\detokenize{s2}} is empty.

\item[\href{https://github.com/LLM4Rocq/mathcomp-eulerian/blob/main/\detokenize{foata.v}\#L373}{\texttt{\detokenize{cross_inv_rcons}}}] \emph{Lemma, line 373}. \textsf{[aux]} \\
  \texttt{\detokenize{cross_inv}} recursion when an element is appended on the right.

\item[\href{https://github.com/LLM4Rocq/mathcomp-eulerian/blob/main/\detokenize{foata.v}\#L380}{\texttt{\detokenize{inv_seq_cat}}}] \emph{Lemma, line 380}. \textsf{[aux]} \\
  Decomposition of \texttt{\detokenize{inv_seq}} over concatenation: left inversions, right     inversions, and cross-inversions.

\item[\href{https://github.com/LLM4Rocq/mathcomp-eulerian/blob/main/\detokenize{foata.v}\#L398}{\texttt{\detokenize{count_gt_cat}}}] \emph{Lemma, line 398}. \textsf{[aux]} \\
  \texttt{\detokenize{count_gt}} is additive under concatenation.

\item[\href{https://github.com/LLM4Rocq/mathcomp-eulerian/blob/main/\detokenize{foata.v}\#L405}{\texttt{\detokenize{cross_inv_perm_eq_r}}}] \emph{Lemma, line 405}. \textsf{[aux]} \\
  \texttt{\detokenize{cross_inv}} only depends on the multiset of its right argument.

\item[\href{https://github.com/LLM4Rocq/mathcomp-eulerian/blob/main/\detokenize{foata.v}\#L415}{\texttt{\detokenize{cross_inv_perm_eq_l}}}] \emph{Lemma, line 415}. \textsf{[aux]} \\
  \texttt{\detokenize{cross_inv}} only depends on the multiset of its left argument     (since \texttt{\detokenize{count_gt}} does).

\item[\href{https://github.com/LLM4Rocq/mathcomp-eulerian/blob/main/\detokenize{foata.v}\#L424}{\texttt{\detokenize{count_lt}}}] \emph{Definition, line 424}. \textsf{[aux]} \\
  \texttt{\detokenize{count_lt a w}} is the number of letters of \texttt{\detokenize{w}} strictly less than \texttt{\detokenize{a}}.

\item[\href{https://github.com/LLM4Rocq/mathcomp-eulerian/blob/main/\detokenize{foata.v}\#L431}{\texttt{\detokenize{perm_eq_flatten_map_pred}}}] \emph{Lemma, line 431}. \textsf{[aux]} \\
  If each block is \texttt{\detokenize{perm_eq}} to its image under \texttt{\detokenize{f}}, the flattened     images are \texttt{\detokenize{perm_eq}} to the flattened originals.

\item[\href{https://github.com/LLM4Rocq/mathcomp-eulerian/blob/main/\detokenize{foata.v}\#L444}{\texttt{\detokenize{inv_seq_flatten_swap_eq}}}] \emph{Lemma, line 444}. \textsf{[aux]} \\
  Block-rotation accounting: if every block is replaced by a \texttt{\detokenize{perm_eq}}     image, the change in \texttt{\detokenize{inv_seq}} of the flatten equals the sum of per-block     changes (cross-inversions are unaffected since they only see multisets).

\item[\href{https://github.com/LLM4Rocq/mathcomp-eulerian/blob/main/\detokenize{foata.v}\#L477}{\texttt{\detokenize{cyc_last_to_front_rcons}}}] \emph{Lemma, line 477}. \textsf{[aux]} \\
  Computational form: \texttt{\detokenize{cyc_last_to_front (rcons b' l) = l :: b'}}.

\item[\href{https://github.com/LLM4Rocq/mathcomp-eulerian/blob/main/\detokenize{foata.v}\#L488}{\texttt{\detokenize{inv_seq_cons_eq_rcons_shift}}}] \emph{Lemma, line 488}. \textsf{[aux]} \\
  Inversions shift when \texttt{\detokenize{l}} moves from the end to the front of \texttt{\detokenize{b'}}:     \texttt{\detokenize{inv_seq (l :: b') - inv_seq (rcons b' l)       = count_lt l b' - count_gt l b'}}.

\item[\href{https://github.com/LLM4Rocq/mathcomp-eulerian/blob/main/\detokenize{foata.v}\#L511}{\texttt{\detokenize{cyc_diff_block_lt}}}] \emph{Lemma, line 511}. \textsf{[aux]} \\
  Per-block inversion change in the \texttt{\detokenize{last u < a}} case: rotating drops     \texttt{\detokenize{size b'}} inversions when the small letter \texttt{\detokenize{l < a}} moves to the front     past the \texttt{\detokenize{b' > a}} block.  Used by \texttt{\detokenize{foata_step_inv_lt}}.

\item[\href{https://github.com/LLM4Rocq/mathcomp-eulerian/blob/main/\detokenize{foata.v}\#L535}{\texttt{\detokenize{cyc_diff_block_gt}}}] \emph{Lemma, line 535}. \textsf{[aux]} \\
  Per-block inversion change in the \texttt{\detokenize{a < last u}} case: rotating adds     \texttt{\detokenize{size b'}} inversions when the large letter \texttt{\detokenize{a < l}} moves to the front     past the \texttt{\detokenize{b' < a}} block.  Used by \texttt{\detokenize{foata_step_inv_gt}}.

\item[\href{https://github.com/LLM4Rocq/mathcomp-eulerian/blob/main/\detokenize{foata.v}\#L558}{\texttt{\detokenize{split_blocks_aux_all_nonempty}}}] \emph{Lemma, line 558}. \textsf{[aux]} \\
  Every block produced by \texttt{\detokenize{split_blocks_aux}} is non-empty.

\item[\href{https://github.com/LLM4Rocq/mathcomp-eulerian/blob/main/\detokenize{foata.v}\#L583}{\texttt{\detokenize{wf_block}}}] \emph{Definition, line 583}. \textsf{[aux]} \\
  \texttt{\detokenize{wf_block P b}} holds when \texttt{\detokenize{b}} is non-empty, its last letter satisfies     \texttt{\detokenize{P}}, and all non-last letters do not.  Captures the well-formedness     of blocks output by \texttt{\detokenize{split_blocks}} (except possibly the last).

\item[\href{https://github.com/LLM4Rocq/mathcomp-eulerian/blob/main/\detokenize{foata.v}\#L591}{\texttt{\detokenize{wf_block_rcons}}}] \emph{Lemma, line 591}. \textsf{[aux]} \\
  \texttt{\detokenize{wf_block P (rcons b' l)}} is equivalent to \texttt{\detokenize{P l && all (~~ P) b'}}.

\item[\href{https://github.com/LLM4Rocq/mathcomp-eulerian/blob/main/\detokenize{foata.v}\#L602}{\texttt{\detokenize{split_blocks_aux_wf}}}] \emph{Lemma, line 602}. \textsf{[aux]} \\
  Structural invariant of \texttt{\detokenize{split_blocks_aux}}: when \texttt{\detokenize{cur}} is all not-\texttt{\detokenize{P}}     and the last letter of \texttt{\detokenize{s}} satisfies \texttt{\detokenize{P}}, every produced block is a     \texttt{\detokenize{wf_block}}.

\item[\href{https://github.com/LLM4Rocq/mathcomp-eulerian/blob/main/\detokenize{foata.v}\#L631}{\texttt{\detokenize{split_blocks_wf}}}] \emph{Lemma, line 631}. \textsf{[aux]} \\
  Every block of \texttt{\detokenize{split_blocks P s}} is a \texttt{\detokenize{wf_block}} when \texttt{\detokenize{s}} is non-empty     and its last letter satisfies \texttt{\detokenize{P}}.

\item[\href{https://github.com/LLM4Rocq/mathcomp-eulerian/blob/main/\detokenize{foata.v}\#L638}{\texttt{\detokenize{split_blocks_aux_size_when_last_P}}}] \emph{Lemma, line 638}. \textsf{[aux]} \\
  Number of blocks of \texttt{\detokenize{split_blocks_aux P cur s}} equals \texttt{\detokenize{count P s}}, when     \texttt{\detokenize{s}} is non-empty and ends with a \texttt{\detokenize{P}}-letter.

\item[\href{https://github.com/LLM4Rocq/mathcomp-eulerian/blob/main/\detokenize{foata.v}\#L666}{\texttt{\detokenize{split_blocks_size_eq}}}] \emph{Lemma, line 666}. \textsf{[aux]} \\
  Number of blocks of \texttt{\detokenize{split_blocks P s}} equals \texttt{\detokenize{count P s}} (assuming     last letter of \texttt{\detokenize{s}} satisfies \texttt{\detokenize{P}}).

\item[\href{https://github.com/LLM4Rocq/mathcomp-eulerian/blob/main/\detokenize{foata.v}\#L673}{\texttt{\detokenize{wf_block_decomp}}}] \emph{Lemma, line 673}. \textsf{[aux]} \\
  Decomposition: any \texttt{\detokenize{wf_block P b}} is of the form \texttt{\detokenize{rcons b' l}} with     \texttt{\detokenize{P l}} and \texttt{\detokenize{b'}} all not-\texttt{\detokenize{P}}.

\item[\href{https://github.com/LLM4Rocq/mathcomp-eulerian/blob/main/\detokenize{foata.v}\#L688}{\texttt{\detokenize{size_count_lt_gt}}}] \emph{Lemma, line 688}. \textsf{[aux]} \\
  For \texttt{\detokenize{uniq u}} with \texttt{a \textbackslash\{\}notin u}, size of \texttt{\detokenize{u}} splits as     \texttt{\detokenize{count_lt a u + count_gt a u}} (no letter equals \texttt{\detokenize{a}}).

\item[\href{https://github.com/LLM4Rocq/mathcomp-eulerian/blob/main/\detokenize{foata.v}\#L705}{\texttt{\detokenize{sum_size_belast_wf}}}] \emph{Lemma, line 705}. \textsf{[aux]} \\
  Sum of \texttt{\detokenize{(size b).-1}} over \texttt{\detokenize{wf_block}}s equals the total length minus     the number of blocks.

\item[\href{https://github.com/LLM4Rocq/mathcomp-eulerian/blob/main/\detokenize{foata.v}\#L735}{\texttt{\detokenize{sum_inv_cyc_lt_blocks}}}] \emph{Lemma, line 735}. \textsf{[aux]} \\
  Aggregated per-block inversion change in the \texttt{\detokenize{last u < a}} case     (\texttt{\detokenize{P = (y < a)}}): summed cyc-rotation drop equals total internal sizes.     Used by \texttt{\detokenize{foata_step_inv_lt}}.

\item[\href{https://github.com/LLM4Rocq/mathcomp-eulerian/blob/main/\detokenize{foata.v}\#L772}{\texttt{\detokenize{sum_inv_cyc_gt_blocks}}}] \emph{Lemma, line 772}. \textsf{[aux]} \\
  Aggregated per-block inversion change in the \texttt{\detokenize{a < last u}} case     (\texttt{\detokenize{P = (a < y)}}): summed cyc-rotation gain equals total internal sizes.     Used by \texttt{\detokenize{foata_step_inv_gt}}.

\item[\href{https://github.com/LLM4Rocq/mathcomp-eulerian/blob/main/\detokenize{foata.v}\#L808}{\texttt{\detokenize{split_blocks_lt_strict}}}] \emph{Lemma, line 808}. \textsf{[aux]} \\
  Strong form of \texttt{\detokenize{split_blocks_wf}} in the \texttt{\detokenize{< a}} case: under \texttt{\detokenize{uniq u}}     and \texttt{a \textbackslash\{\}notin u}, every non-last letter in every block is strictly     \texttt{\detokenize{> a}} (not just not \texttt{\detokenize{< a}}).

\item[\href{https://github.com/LLM4Rocq/mathcomp-eulerian/blob/main/\detokenize{foata.v}\#L844}{\texttt{\detokenize{split_blocks_gt_strict}}}] \emph{Lemma, line 844}. \textsf{[aux]} \\
  Strong form of \texttt{\detokenize{split_blocks_wf}} in the \texttt{\detokenize{a <}} case: every non-last     letter is strictly \texttt{\detokenize{< a}}.

\item[\href{https://github.com/LLM4Rocq/mathcomp-eulerian/blob/main/\detokenize{foata.v}\#L880}{\texttt{\detokenize{foata_step_inv_lt}}}] \emph{Lemma, line 880}. \textsf{[aux]} \\
  Case \texttt{\detokenize{last u < a}} of \texttt{\detokenize{foata_step_inv}}: appending \texttt{\detokenize{a}} when last     letter is smaller leaves \texttt{\detokenize{inv_seq}} unchanged.

\item[\href{https://github.com/LLM4Rocq/mathcomp-eulerian/blob/main/\detokenize{foata.v}\#L950}{\texttt{\detokenize{foata_step_inv_gt}}}] \emph{Lemma, line 950}. \textsf{[aux]} \\
  Case \texttt{\detokenize{a < last u}} of \texttt{\detokenize{foata_step_inv}}: appending \texttt{\detokenize{a}} when last letter     is larger increases \texttt{\detokenize{inv_seq}} by \texttt{\detokenize{size u}}.

\item[\href{https://github.com/LLM4Rocq/mathcomp-eulerian/blob/main/\detokenize{foata.v}\#L1014}{\texttt{\detokenize{foata_step_inv}}}] \emph{Lemma, line 1014}. \textsf{[aux]} \\
  Combined per-step invariant: \texttt{\detokenize{foata_step}} adds \texttt{\detokenize{size u}} to \texttt{\detokenize{inv_seq}}     iff a new descent appears (i.e. \texttt{\detokenize{a < last u}}).  Mirrors \texttt{\detokenize{maj_seq_rcons}}.

\item[\href{https://github.com/LLM4Rocq/mathcomp-eulerian/blob/main/\detokenize{foata.v}\#L1038}{\texttt{\detokenize{foata_inv_eq_maj}}}] \emph{Theorem, line 1038}. \textsf{[aux]} \\
  Headline equidistribution at the seq level: \texttt{\detokenize{inv_seq (foata w) =     maj_seq w}} for \texttt{\detokenize{uniq w}}.  Proven by induction on \texttt{\detokenize{w}} from the right     using \texttt{\detokenize{foata_step_inv}} and \texttt{\detokenize{maj_seq_rcons}}.

\item[\href{https://github.com/LLM4Rocq/mathcomp-eulerian/blob/main/\detokenize{foata.v}\#L1070}{\texttt{\detokenize{cyc_first_to_back}}}] \emph{Definition, line 1070}. \textsf{[aux]} \\
  \texttt{\detokenize{cyc_first_to_back s}} moves the first letter of \texttt{\detokenize{s}} to the back.     Inverse of \texttt{\detokenize{cyc_last_to_front}}.

\item[\href{https://github.com/LLM4Rocq/mathcomp-eulerian/blob/main/\detokenize{foata.v}\#L1078}{\texttt{\detokenize{cyc_first_to_back_perm_eq}}}] \emph{Lemma, line 1078}. \textsf{[aux]} \\
  \texttt{\detokenize{cyc_first_to_back}} preserves the multiset.

\item[\href{https://github.com/LLM4Rocq/mathcomp-eulerian/blob/main/\detokenize{foata.v}\#L1087}{\texttt{\detokenize{cyc_first_to_backK}}}] \emph{Lemma, line 1087}. \textsf{[aux]} \\
  \texttt{\detokenize{cyc_first_to_back}} is the left inverse of \texttt{\detokenize{cyc_last_to_front}}.

\item[\href{https://github.com/LLM4Rocq/mathcomp-eulerian/blob/main/\detokenize{foata.v}\#L1097}{\texttt{\detokenize{split_blocks_inv_aux}}}] \emph{Fixpoint, line 1097}. \textsf{[aux]} \\
  \texttt{\detokenize{split_blocks_inv_aux P cur s}} is the tail-recursive worker for     \texttt{\detokenize{split_blocks_inv}}: cuts before each \texttt{\detokenize{P}}-letter rather than after.

\item[\href{https://github.com/LLM4Rocq/mathcomp-eulerian/blob/main/\detokenize{foata.v}\#L1112}{\texttt{\detokenize{split_blocks_inv}}}] \emph{Definition, line 1112}. \textsf{[aux]} \\
  \texttt{\detokenize{split_blocks_inv P s}} cuts \texttt{\detokenize{s}} into blocks each STARTING with a     \texttt{\detokenize{P}}-letter (cutting right before each \texttt{\detokenize{P}}-letter after the first).     Inverse-direction analogue of \texttt{\detokenize{split_blocks}}.

\item[\href{https://github.com/LLM4Rocq/mathcomp-eulerian/blob/main/\detokenize{foata.v}\#L1116}{\texttt{\detokenize{split_blocks_inv_aux_flatten}}}] \emph{Lemma, line 1116}. \textsf{[aux]} \\
  Accumulator-aware flatten property for \texttt{\detokenize{split_blocks_inv_aux}}.

\item[\href{https://github.com/LLM4Rocq/mathcomp-eulerian/blob/main/\detokenize{foata.v}\#L1162}{\texttt{\detokenize{split_blocks_inv_aux_app_notP}}}] \emph{Lemma, line 1162}. \textsf{[aux]} \\
  When a prefix \texttt{\detokenize{nots}} is all not-\texttt{\detokenize{P}}, it gets absorbed into the current     block accumulator.

\item[\href{https://github.com/LLM4Rocq/mathcomp-eulerian/blob/main/\detokenize{foata.v}\#L1189}{\texttt{\detokenize{split_blocks_inv_cyc_wf}}}] \emph{Lemma, line 1189}. \textsf{[aux]} \\
  Big cancellation: applying \texttt{\detokenize{split_blocks_inv P}} to the flatten of     rotated \texttt{\detokenize{wf_block}}s recovers the rotated blocks themselves.  Key step     in proving \texttt{\detokenize{foata_step_undoK}}.

\item[\href{https://github.com/LLM4Rocq/mathcomp-eulerian/blob/main/\detokenize{foata.v}\#L1224}{\texttt{\detokenize{foata_step_undo}}}] \emph{Definition, line 1224}. \textsf{[aux]} \\
  \texttt{\detokenize{foata_step_undo s}} inverts a single \texttt{\detokenize{foata_step}}: peels the trailing     letter \texttt{\detokenize{a}} and rotates each \texttt{\detokenize{split_blocks_inv}} block back via     \texttt{\detokenize{cyc_first_to_back}}.

\item[\href{https://github.com/LLM4Rocq/mathcomp-eulerian/blob/main/\detokenize{foata.v}\#L1235}{\texttt{\detokenize{split_blocks_inv_aux_eq}}}] \emph{Lemma, line 1235}. \textsf{[aux]} \\
  \texttt{\detokenize{split_blocks_inv_aux}} only depends on the predicate extensionally.

\item[\href{https://github.com/LLM4Rocq/mathcomp-eulerian/blob/main/\detokenize{foata.v}\#L1245}{\texttt{\detokenize{split_blocks_inv_eq}}}] \emph{Lemma, line 1245}. \textsf{[aux]} \\
  \texttt{\detokenize{split_blocks_inv}} only depends on the predicate extensionally.

\item[\href{https://github.com/LLM4Rocq/mathcomp-eulerian/blob/main/\detokenize{foata.v}\#L1254}{\texttt{\detokenize{foata_step_undoK}}}] \emph{Lemma, line 1254}. \textsf{[aux]} \\
  Cancellation: \texttt{\detokenize{foata_step_undo}} is the left inverse of \texttt{\detokenize{foata_step}}     on uniq inputs with a fresh letter.

\item[\href{https://github.com/LLM4Rocq/mathcomp-eulerian/blob/main/\detokenize{foata.v}\#L1360}{\texttt{\detokenize{foata_inv_aux}}}] \emph{Fixpoint, line 1360}. \textsf{[aux]} \\
  \texttt{\detokenize{foata_inv_aux n s}} inverts \texttt{\detokenize{n}} iterations of \texttt{\detokenize{foata_step}} by     repeatedly applying \texttt{\detokenize{foata_step_undo}}.

\item[\href{https://github.com/LLM4Rocq/mathcomp-eulerian/blob/main/\detokenize{foata.v}\#L1370}{\texttt{\detokenize{foata_inv}}}] \emph{Definition, line 1370}. \textsf{[aux]} \\
  \texttt{\detokenize{foata_inv s}} is the seq-level inverse of \texttt{\detokenize{foata}}: repeatedly peels     \texttt{\detokenize{foata_step}}s for \texttt{\detokenize{size s}} iterations.

\item[\href{https://github.com/LLM4Rocq/mathcomp-eulerian/blob/main/\detokenize{foata.v}\#L1375}{\texttt{\detokenize{foata_invK_aux}}}] \emph{Lemma, line 1375}. \textsf{[aux]} \\
  Length-indexed cancellation: \texttt{\detokenize{foata_inv_aux n (foata w) = w}} for     \texttt{\detokenize{size w = n}} and \texttt{\detokenize{uniq w}}.

\item[\href{https://github.com/LLM4Rocq/mathcomp-eulerian/blob/main/\detokenize{foata.v}\#L1404}{\texttt{\detokenize{foata_invK}}}] \emph{Lemma, line 1404}. \textsf{[aux]} \\
  \texttt{\detokenize{foata_inv}} is the left inverse of \texttt{\detokenize{foata}} on uniq sequences.

\item[\href{https://github.com/LLM4Rocq/mathcomp-eulerian/blob/main/\detokenize{foata.v}\#L1413}{\texttt{\detokenize{foata_inj_uniq}}}] \emph{Lemma, line 1413}. \textsf{[aux]} \\
  \texttt{\detokenize{foata}} is injective on uniq sequences (immediate from \texttt{\detokenize{foata_invK}}).

\item[\href{https://github.com/LLM4Rocq/mathcomp-eulerian/blob/main/\detokenize{foata.v}\#L1429}{\texttt{\detokenize{maj_eq_maj_seq}}}] \emph{Lemma, line 1429}. \textsf{[aux]} \\
  Bridge: perm-level \texttt{\detokenize{maj}} equals seq-level \texttt{\detokenize{maj_seq}} of \texttt{\detokenize{perm_to_seq}}.

\item[\href{https://github.com/LLM4Rocq/mathcomp-eulerian/blob/main/\detokenize{foata.v}\#L1450}{\texttt{\detokenize{foata_perm_to_seq_size}}}] \emph{Lemma, line 1450}. \textsf{[aux]} \\
  \texttt{\detokenize{foata (perm_to_seq s)}} has size \texttt{\detokenize{n.+1}}; needed to build \texttt{\detokenize{foata_perm}}.

\item[\href{https://github.com/LLM4Rocq/mathcomp-eulerian/blob/main/\detokenize{foata.v}\#L1455}{\texttt{\detokenize{foata_perm_to_seq_uniq}}}] \emph{Lemma, line 1455}. \textsf{[aux]} \\
  \texttt{\detokenize{foata (perm_to_seq s)}} is uniq; needed to build \texttt{\detokenize{foata_perm}}.

\item[\href{https://github.com/LLM4Rocq/mathcomp-eulerian/blob/main/\detokenize{foata.v}\#L1461}{\texttt{\detokenize{foata_perm_to_seq_bnd}}}] \emph{Lemma, line 1461}. \textsf{[aux]} \\
  \texttt{\detokenize{foata (perm_to_seq s)}} is bounded by \texttt{\detokenize{n.+1}}; needed to build     \texttt{\detokenize{foata_perm}}.

\item[\href{https://github.com/LLM4Rocq/mathcomp-eulerian/blob/main/\detokenize{foata.v}\#L1471}{\texttt{\detokenize{foata_perm}}}] \emph{Definition, line 1471}. \textsf{[Ch 4 §4.2]} \\
  \texttt{\detokenize{foata_perm s}} is the Foata bijection lifted to \texttt{\{perm 'I\_n.+1\}}:     apply \texttt{\detokenize{foata}} at the seq level, then re-package as a permutation.

\item[\href{https://github.com/LLM4Rocq/mathcomp-eulerian/blob/main/\detokenize{foata.v}\#L1477}{\texttt{\detokenize{perm_to_seq_foata_perm}}}] \emph{Lemma, line 1477}. \textsf{[aux]} \\
  Commutation: \texttt{\detokenize{perm_to_seq}} of \texttt{\detokenize{foata_perm s}} is \texttt{\detokenize{foata}} of     \texttt{\detokenize{perm_to_seq s}}.

\item[\href{https://github.com/LLM4Rocq/mathcomp-eulerian/blob/main/\detokenize{foata.v}\#L1484}{\texttt{\detokenize{inv_eq_inv_seq}}}] \emph{Lemma, line 1484}. \textsf{[aux]} \\
  Bridge: perm-level \texttt{\detokenize{inv}} equals seq-level \texttt{\detokenize{inv_seq}} of \texttt{\detokenize{perm_to_seq}}.

\item[\href{https://github.com/LLM4Rocq/mathcomp-eulerian/blob/main/\detokenize{foata.v}\#L1516}{\texttt{\detokenize{foata_perm_inv_maj}}}] \emph{Lemma, line 1516}. \textsf{[Ch 4 §4.3]} \\
  Headline perm-level identity: \texttt{\detokenize{inv (foata_perm s) = maj s}}, i.e.     \texttt{\detokenize{foata_perm}} sends maj-classes to inv-classes.  Direct lift of     \texttt{\detokenize{foata_inv_eq_maj}} via \texttt{\detokenize{maj_eq_maj_seq}} and \texttt{\detokenize{inv_eq_inv_seq}}.

\item[\href{https://github.com/LLM4Rocq/mathcomp-eulerian/blob/main/\detokenize{foata.v}\#L1527}{\texttt{\detokenize{foata_perm_inj}}}] \emph{Lemma, line 1527}. \textsf{[Ch 4 §4.3]} \\
  \texttt{\detokenize{foata_perm}} is injective; follows from \texttt{\detokenize{foata_inj_uniq}} (cancellation     via \texttt{\detokenize{foata_invK}}).  On the finite set \texttt{\{perm 'I\_n.+1\}} this implies     surjective, used to prove \texttt{\detokenize{inv_maj_equidistr}}.

\item[\href{https://github.com/LLM4Rocq/mathcomp-eulerian/blob/main/\detokenize{foata.v}\#L1542}{\texttt{\detokenize{inv_maj_equidistr}}}] \emph{Theorem, line 1542}. \textsf{[Ch 4 §4.3]} \\
  MacMahon's equidistribution theorem (Stanley EC1 \textbackslash{}S\{\}1.3.4): the     statistics \texttt{\detokenize{inv}} and \texttt{\detokenize{maj}} are equidistributed on \texttt{\{perm 'I\_n.+1\}}.     Proved via the bijection \texttt{\detokenize{foata_perm}}, which satisfies     \texttt{\detokenize{inv (foata_perm s) = maj s}} and is injective hence bijective.

\end{description}

% --------------------------------------------
\section*{\texttt{fps/fps.v}}
\addcontentsline{toc}{section}{\texttt{fps/fps.v}}

\begin{description}
\item[\href{https://github.com/LLM4Rocq/mathcomp-eulerian/blob/main/\detokenize{fps/fps.v}\#L76}{\texttt{\detokenize{fpsP}}}] \emph{Lemma, line 76}. \textsf{[aux]} \\
  Coefficientwise equality is equality (uses functional extensionality).

\item[\href{https://github.com/LLM4Rocq/mathcomp-eulerian/blob/main/\detokenize{fps/fps.v}\#L92}{\texttt{\detokenize{fps_zero}}}] \emph{Definition, line 92}. \textsf{[aux]} \\
  (no docstring)

\item[\href{https://github.com/LLM4Rocq/mathcomp-eulerian/blob/main/\detokenize{fps/fps.v}\#L93}{\texttt{\detokenize{fps_add}}}] \emph{Definition, line 93}. \textsf{[aux]} \\
  (no docstring)

\item[\href{https://github.com/LLM4Rocq/mathcomp-eulerian/blob/main/\detokenize{fps/fps.v}\#L94}{\texttt{\detokenize{fps_opp}}}] \emph{Definition, line 94}. \textsf{[aux]} \\
  (no docstring)

\item[\href{https://github.com/LLM4Rocq/mathcomp-eulerian/blob/main/\detokenize{fps/fps.v}\#L96}{\texttt{\detokenize{fps_addA}}}] \emph{Fact, line 96}. \textsf{[aux]} \\
  (no docstring)

\item[\href{https://github.com/LLM4Rocq/mathcomp-eulerian/blob/main/\detokenize{fps/fps.v}\#L99}{\texttt{\detokenize{fps_addC}}}] \emph{Fact, line 99}. \textsf{[aux]} \\
  (no docstring)

\item[\href{https://github.com/LLM4Rocq/mathcomp-eulerian/blob/main/\detokenize{fps/fps.v}\#L102}{\texttt{\detokenize{fps_add0f}}}] \emph{Fact, line 102}. \textsf{[aux]} \\
  (no docstring)

\item[\href{https://github.com/LLM4Rocq/mathcomp-eulerian/blob/main/\detokenize{fps/fps.v}\#L105}{\texttt{\detokenize{fps_addNf}}}] \emph{Fact, line 105}. \textsf{[aux]} \\
  (no docstring)

\item[\href{https://github.com/LLM4Rocq/mathcomp-eulerian/blob/main/\detokenize{fps/fps.v}\#L111}{\texttt{\detokenize{coef_fps0}}}] \emph{Lemma, line 111}. \textsf{[aux]} \\
  (no docstring)

\item[\href{https://github.com/LLM4Rocq/mathcomp-eulerian/blob/main/\detokenize{fps/fps.v}\#L114}{\texttt{\detokenize{coef_fpsD}}}] \emph{Lemma, line 114}. \textsf{[aux]} \\
  (no docstring)

\item[\href{https://github.com/LLM4Rocq/mathcomp-eulerian/blob/main/\detokenize{fps/fps.v}\#L117}{\texttt{\detokenize{coef_fpsN}}}] \emph{Lemma, line 117}. \textsf{[aux]} \\
  (no docstring)

\item[\href{https://github.com/LLM4Rocq/mathcomp-eulerian/blob/main/\detokenize{fps/fps.v}\#L120}{\texttt{\detokenize{coef_fpsB}}}] \emph{Lemma, line 120}. \textsf{[aux]} \\
  (no docstring)

\item[\href{https://github.com/LLM4Rocq/mathcomp-eulerian/blob/main/\detokenize{fps/fps.v}\#L123}{\texttt{\detokenize{coef_fps_sum}}}] \emph{Lemma, line 123}. \textsf{[aux]} \\
  (no docstring)

\item[\href{https://github.com/LLM4Rocq/mathcomp-eulerian/blob/main/\detokenize{fps/fps.v}\#L138}{\texttt{\detokenize{fps_one}}}] \emph{Definition, line 138}. \textsf{[aux]} \\
  (no docstring)

\item[\href{https://github.com/LLM4Rocq/mathcomp-eulerian/blob/main/\detokenize{fps/fps.v}\#L140}{\texttt{\detokenize{fps_mul}}}] \emph{Definition, line 140}. \textsf{[aux]} \\
  (no docstring)

\item[\href{https://github.com/LLM4Rocq/mathcomp-eulerian/blob/main/\detokenize{fps/fps.v}\#L144}{\texttt{\detokenize{fps_trunc}}}] \emph{Definition, line 144}. \textsf{[aux]} \\
  Polynomial truncation: the first m+1 coefficients.

\item[\href{https://github.com/LLM4Rocq/mathcomp-eulerian/blob/main/\detokenize{fps/fps.v}\#L146}{\texttt{\detokenize{coef_fps_trunc}}}] \emph{Lemma, line 146}. \textsf{[aux]} \\
  (no docstring)

\item[\href{https://github.com/LLM4Rocq/mathcomp-eulerian/blob/main/\detokenize{fps/fps.v}\#L152}{\texttt{\detokenize{coef_fpsM_trunc}}}] \emph{Lemma, line 152}. \textsf{[aux]} \\
  (no docstring)

\item[\href{https://github.com/LLM4Rocq/mathcomp-eulerian/blob/main/\detokenize{fps/fps.v}\#L163}{\texttt{\detokenize{coefM_eqr}}}] \emph{Lemma, line 163}. \textsf{[aux]} \\
  (no docstring)

\item[\href{https://github.com/LLM4Rocq/mathcomp-eulerian/blob/main/\detokenize{fps/fps.v}\#L169}{\texttt{\detokenize{coefM_eql}}}] \emph{Lemma, line 169}. \textsf{[aux]} \\
  (no docstring)

\item[\href{https://github.com/LLM4Rocq/mathcomp-eulerian/blob/main/\detokenize{fps/fps.v}\#L175}{\texttt{\detokenize{fps_mulA}}}] \emph{Fact, line 175}. \textsf{[aux]} \\
  (no docstring)

\item[\href{https://github.com/LLM4Rocq/mathcomp-eulerian/blob/main/\detokenize{fps/fps.v}\#L187}{\texttt{\detokenize{fps_mulC}}}] \emph{Fact, line 187}. \textsf{[aux]} \\
  (no docstring)

\item[\href{https://github.com/LLM4Rocq/mathcomp-eulerian/blob/main/\detokenize{fps/fps.v}\#L193}{\texttt{\detokenize{fps_trunc_one}}}] \emph{Lemma, line 193}. \textsf{[aux]} \\
  (no docstring)

\item[\href{https://github.com/LLM4Rocq/mathcomp-eulerian/blob/main/\detokenize{fps/fps.v}\#L200}{\texttt{\detokenize{fps_mul1f}}}] \emph{Fact, line 200}. \textsf{[aux]} \\
  (no docstring)

\item[\href{https://github.com/LLM4Rocq/mathcomp-eulerian/blob/main/\detokenize{fps/fps.v}\#L207}{\texttt{\detokenize{fps_mulDl}}}] \emph{Fact, line 207}. \textsf{[aux]} \\
  (no docstring)

\item[\href{https://github.com/LLM4Rocq/mathcomp-eulerian/blob/main/\detokenize{fps/fps.v}\#L213}{\texttt{\detokenize{fps_oner_neq0}}}] \emph{Fact, line 213}. \textsf{[aux]} \\
  (no docstring)

\item[\href{https://github.com/LLM4Rocq/mathcomp-eulerian/blob/main/\detokenize{fps/fps.v}\#L223}{\texttt{\detokenize{coef_fps1}}}] \emph{Lemma, line 223}. \textsf{[aux]} \\
  (no docstring)

\item[\href{https://github.com/LLM4Rocq/mathcomp-eulerian/blob/main/\detokenize{fps/fps.v}\#L226}{\texttt{\detokenize{coef_fpsM}}}] \emph{Lemma, line 226}. \textsf{[aux]} \\
  (no docstring)

\item[\href{https://github.com/LLM4Rocq/mathcomp-eulerian/blob/main/\detokenize{fps/fps.v}\#L230}{\texttt{\detokenize{coef0_fpsM}}}] \emph{Lemma, line 230}. \textsf{[aux]} \\
  Coefficient 0 is a ring morphism.

\item[\href{https://github.com/LLM4Rocq/mathcomp-eulerian/blob/main/\detokenize{fps/fps.v}\#L237}{\texttt{\detokenize{fps_scale}}}] \emph{Definition, line 237}. \textsf{[aux]} \\
  (no docstring)

\item[\href{https://github.com/LLM4Rocq/mathcomp-eulerian/blob/main/\detokenize{fps/fps.v}\#L239}{\texttt{\detokenize{fps_scaleA}}}] \emph{Fact, line 239}. \textsf{[aux]} \\
  (no docstring)

\item[\href{https://github.com/LLM4Rocq/mathcomp-eulerian/blob/main/\detokenize{fps/fps.v}\#L242}{\texttt{\detokenize{fps_scale1f}}}] \emph{Fact, line 242}. \textsf{[aux]} \\
  (no docstring)

\item[\href{https://github.com/LLM4Rocq/mathcomp-eulerian/blob/main/\detokenize{fps/fps.v}\#L245}{\texttt{\detokenize{fps_scaleDr}}}] \emph{Fact, line 245}. \textsf{[aux]} \\
  (no docstring)

\item[\href{https://github.com/LLM4Rocq/mathcomp-eulerian/blob/main/\detokenize{fps/fps.v}\#L248}{\texttt{\detokenize{fps_scaleDl}}}] \emph{Fact, line 248}. \textsf{[aux]} \\
  (no docstring)

\item[\href{https://github.com/LLM4Rocq/mathcomp-eulerian/blob/main/\detokenize{fps/fps.v}\#L255}{\texttt{\detokenize{coef_fpsZ}}}] \emph{Lemma, line 255}. \textsf{[aux]} \\
  (no docstring)

\item[\href{https://github.com/LLM4Rocq/mathcomp-eulerian/blob/main/\detokenize{fps/fps.v}\#L258}{\texttt{\detokenize{fps_scaleAl}}}] \emph{Fact, line 258}. \textsf{[aux]} \\
  (no docstring)

\item[\href{https://github.com/LLM4Rocq/mathcomp-eulerian/blob/main/\detokenize{fps/fps.v}\#L267}{\texttt{\detokenize{fps_scaleAr}}}] \emph{Fact, line 267}. \textsf{[aux]} \\
  (no docstring)

\item[\href{https://github.com/LLM4Rocq/mathcomp-eulerian/blob/main/\detokenize{fps/fps.v}\#L286}{\texttt{\detokenize{poly_fps}}}] \emph{Definition, line 286}. \textsf{[aux]} \\
  (no docstring)

\item[\href{https://github.com/LLM4Rocq/mathcomp-eulerian/blob/main/\detokenize{fps/fps.v}\#L288}{\texttt{\detokenize{coef_poly_fps}}}] \emph{Lemma, line 288}. \textsf{[aux]} \\
  (no docstring)

\item[\href{https://github.com/LLM4Rocq/mathcomp-eulerian/blob/main/\detokenize{fps/fps.v}\#L291}{\texttt{\detokenize{poly_fps_inj}}}] \emph{Lemma, line 291}. \textsf{[aux]} \\
  (no docstring)

\item[\href{https://github.com/LLM4Rocq/mathcomp-eulerian/blob/main/\detokenize{fps/fps.v}\#L294}{\texttt{\detokenize{poly_fps_is_zmod_morphism}}}] \emph{Fact, line 294}. \textsf{[aux]} \\
  (no docstring)

\item[\href{https://github.com/LLM4Rocq/mathcomp-eulerian/blob/main/\detokenize{fps/fps.v}\#L297}{\texttt{\detokenize{poly_fps_is_monoid_morphism}}}] \emph{Fact, line 297}. \textsf{[aux]} \\
  (no docstring)

\item[\href{https://github.com/LLM4Rocq/mathcomp-eulerian/blob/main/\detokenize{fps/fps.v}\#L310}{\texttt{\detokenize{poly_fpsC}}}] \emph{Lemma, line 310}. \textsf{[aux]} \\
  (no docstring)

\item[\href{https://github.com/LLM4Rocq/mathcomp-eulerian/blob/main/\detokenize{fps/fps.v}\#L331}{\texttt{\detokenize{inv_pref}}}] \emph{Fixpoint, line 331}. \textsf{[aux]} \\
  (no docstring)

\item[\href{https://github.com/LLM4Rocq/mathcomp-eulerian/blob/main/\detokenize{fps/fps.v}\#L337}{\texttt{\detokenize{size_inv_pref}}}] \emph{Lemma, line 337}. \textsf{[aux]} \\
  (no docstring)

\item[\href{https://github.com/LLM4Rocq/mathcomp-eulerian/blob/main/\detokenize{fps/fps.v}\#L340}{\texttt{\detokenize{fps_inv_coef}}}] \emph{Definition, line 340}. \textsf{[aux]} \\
  (no docstring)

\item[\href{https://github.com/LLM4Rocq/mathcomp-eulerian/blob/main/\detokenize{fps/fps.v}\#L342}{\texttt{\detokenize{inv_prefE}}}] \emph{Lemma, line 342}. \textsf{[aux]} \\
  (no docstring)

\item[\href{https://github.com/LLM4Rocq/mathcomp-eulerian/blob/main/\detokenize{fps/fps.v}\#L350}{\texttt{\detokenize{fps_inv_coef0}}}] \emph{Lemma, line 350}. \textsf{[aux]} \\
  (no docstring)

\item[\href{https://github.com/LLM4Rocq/mathcomp-eulerian/blob/main/\detokenize{fps/fps.v}\#L353}{\texttt{\detokenize{fps_inv_coefS}}}] \emph{Lemma, line 353}. \textsf{[aux]} \\
  (no docstring)

\item[\href{https://github.com/LLM4Rocq/mathcomp-eulerian/blob/main/\detokenize{fps/fps.v}\#L361}{\texttt{\detokenize{fps_unit}}}] \emph{Definition, line 361}. \textsf{[aux]} \\
  (no docstring)

\item[\href{https://github.com/LLM4Rocq/mathcomp-eulerian/blob/main/\detokenize{fps/fps.v}\#L363}{\texttt{\detokenize{fps_inv}}}] \emph{Definition, line 363}. \textsf{[aux]} \\
  (no docstring)

\item[\href{https://github.com/LLM4Rocq/mathcomp-eulerian/blob/main/\detokenize{fps/fps.v}\#L367}{\texttt{\detokenize{fps_mulVf}}}] \emph{Fact, line 367}. \textsf{[aux]} \\
  (no docstring)

\item[\href{https://github.com/LLM4Rocq/mathcomp-eulerian/blob/main/\detokenize{fps/fps.v}\#L383}{\texttt{\detokenize{fps_unitP}}}] \emph{Fact, line 383}. \textsf{[aux]} \\
  (no docstring)

\item[\href{https://github.com/LLM4Rocq/mathcomp-eulerian/blob/main/\detokenize{fps/fps.v}\#L389}{\texttt{\detokenize{fps_inv_out}}}] \emph{Fact, line 389}. \textsf{[aux]} \\
  (no docstring)

\item[\href{https://github.com/LLM4Rocq/mathcomp-eulerian/blob/main/\detokenize{fps/fps.v}\#L401}{\texttt{\detokenize{fps_unitE}}}] \emph{Lemma, line 401}. \textsf{[aux]} \\
  The headline characterization of units.

\item[\href{https://github.com/LLM4Rocq/mathcomp-eulerian/blob/main/\detokenize{fps/fps.v}\#L414}{\texttt{\detokenize{fps_geom}}}] \emph{Definition, line 414}. \textsf{[aux]} \\
  The all-ones series 1 + x + x\textasciicircum{}2 + ...

\item[\href{https://github.com/LLM4Rocq/mathcomp-eulerian/blob/main/\detokenize{fps/fps.v}\#L416}{\texttt{\detokenize{coef_fps_onesubX}}}] \emph{Lemma, line 416}. \textsf{[aux]} \\
  (no docstring)

\item[\href{https://github.com/LLM4Rocq/mathcomp-eulerian/blob/main/\detokenize{fps/fps.v}\#L428}{\texttt{\detokenize{onesubX_mul_geom}}}] \emph{Lemma, line 428}. \textsf{[aux]} \\
  (no docstring)

\item[\href{https://github.com/LLM4Rocq/mathcomp-eulerian/blob/main/\detokenize{fps/fps.v}\#L440}{\texttt{\detokenize{fps_onesubX_unit}}}] \emph{Lemma, line 440}. \textsf{[aux]} \\
  (no docstring)

\item[\href{https://github.com/LLM4Rocq/mathcomp-eulerian/blob/main/\detokenize{fps/fps.v}\#L446}{\texttt{\detokenize{fps_inv_onesubX}}}] \emph{Lemma, line 446}. \textsf{[aux]} \\
  The acceptance test: the inverse of 1 - x is the geometric series.

\end{description}

% --------------------------------------------
\section*{\texttt{fps/fps\_deriv.v}}
\addcontentsline{toc}{section}{\texttt{fps/fps\_deriv.v}}

\begin{description}
\item[\href{https://github.com/LLM4Rocq/mathcomp-eulerian/blob/main/\detokenize{fps/fps_deriv.v}\#L37}{\texttt{\detokenize{fps_deriv}}}] \emph{Definition, line 37}. \textsf{[aux]} \\
  (no docstring)

\item[\href{https://github.com/LLM4Rocq/mathcomp-eulerian/blob/main/\detokenize{fps/fps_deriv.v}\#L39}{\texttt{\detokenize{coef_fps_deriv}}}] \emph{Lemma, line 39}. \textsf{[aux]} \\
  (no docstring)

\item[\href{https://github.com/LLM4Rocq/mathcomp-eulerian/blob/main/\detokenize{fps/fps_deriv.v}\#L42}{\texttt{\detokenize{fps_derivD}}}] \emph{Lemma, line 42}. \textsf{[aux]} \\
  (no docstring)

\item[\href{https://github.com/LLM4Rocq/mathcomp-eulerian/blob/main/\detokenize{fps/fps_deriv.v}\#L45}{\texttt{\detokenize{fps_derivN}}}] \emph{Lemma, line 45}. \textsf{[aux]} \\
  (no docstring)

\item[\href{https://github.com/LLM4Rocq/mathcomp-eulerian/blob/main/\detokenize{fps/fps_deriv.v}\#L48}{\texttt{\detokenize{fps_derivB}}}] \emph{Lemma, line 48}. \textsf{[aux]} \\
  (no docstring)

\item[\href{https://github.com/LLM4Rocq/mathcomp-eulerian/blob/main/\detokenize{fps/fps_deriv.v}\#L51}{\texttt{\detokenize{fps_derivZ}}}] \emph{Lemma, line 51}. \textsf{[aux]} \\
  (no docstring)

\item[\href{https://github.com/LLM4Rocq/mathcomp-eulerian/blob/main/\detokenize{fps/fps_deriv.v}\#L54}{\texttt{\detokenize{fps_deriv0}}}] \emph{Lemma, line 54}. \textsf{[aux]} \\
  (no docstring)

\item[\href{https://github.com/LLM4Rocq/mathcomp-eulerian/blob/main/\detokenize{fps/fps_deriv.v}\#L57}{\texttt{\detokenize{fps_deriv1}}}] \emph{Lemma, line 57}. \textsf{[aux]} \\
  (no docstring)

\item[\href{https://github.com/LLM4Rocq/mathcomp-eulerian/blob/main/\detokenize{fps/fps_deriv.v}\#L60}{\texttt{\detokenize{fps_derivC}}}] \emph{Lemma, line 60}. \textsf{[aux]} \\
  (no docstring)

\item[\href{https://github.com/LLM4Rocq/mathcomp-eulerian/blob/main/\detokenize{fps/fps_deriv.v}\#L63}{\texttt{\detokenize{fps_deriv_sum}}}] \emph{Lemma, line 63}. \textsf{[aux]} \\
  (no docstring)

\item[\href{https://github.com/LLM4Rocq/mathcomp-eulerian/blob/main/\detokenize{fps/fps_deriv.v}\#L71}{\texttt{\detokenize{fps_derivM}}}] \emph{Lemma, line 71}. \textsf{[aux]} \\
  The Leibniz rule.

\item[\href{https://github.com/LLM4Rocq/mathcomp-eulerian/blob/main/\detokenize{fps/fps_deriv.v}\#L95}{\texttt{\detokenize{fps_deriv_poly}}}] \emph{Lemma, line 95}. \textsf{[aux]} \\
  Compatibility with the polynomial derivative.

\item[\href{https://github.com/LLM4Rocq/mathcomp-eulerian/blob/main/\detokenize{fps/fps_deriv.v}\#L101}{\texttt{\detokenize{fps_derivX}}}] \emph{Lemma, line 101}. \textsf{[aux]} \\
  (no docstring)

\item[\href{https://github.com/LLM4Rocq/mathcomp-eulerian/blob/main/\detokenize{fps/fps_deriv.v}\#L115}{\texttt{\detokenize{natf_neq0_char0}}}] \emph{Lemma, line 115}. \textsf{[aux]} \\
  (no docstring)

\item[\href{https://github.com/LLM4Rocq/mathcomp-eulerian/blob/main/\detokenize{fps/fps_deriv.v}\#L118}{\texttt{\detokenize{natSf_neq0}}}] \emph{Lemma, line 118}. \textsf{[aux]} \\
  (no docstring)

\item[\href{https://github.com/LLM4Rocq/mathcomp-eulerian/blob/main/\detokenize{fps/fps_deriv.v}\#L121}{\texttt{\detokenize{fps_prim}}}] \emph{Definition, line 121}. \textsf{[aux]} \\
  (no docstring)

\item[\href{https://github.com/LLM4Rocq/mathcomp-eulerian/blob/main/\detokenize{fps/fps_deriv.v}\#L124}{\texttt{\detokenize{coef_fps_prim0}}}] \emph{Lemma, line 124}. \textsf{[aux]} \\
  (no docstring)

\item[\href{https://github.com/LLM4Rocq/mathcomp-eulerian/blob/main/\detokenize{fps/fps_deriv.v}\#L127}{\texttt{\detokenize{coef_fps_primS}}}] \emph{Lemma, line 127}. \textsf{[aux]} \\
  (no docstring)

\item[\href{https://github.com/LLM4Rocq/mathcomp-eulerian/blob/main/\detokenize{fps/fps_deriv.v}\#L131}{\texttt{\detokenize{fps_derivK}}}] \emph{Lemma, line 131}. \textsf{[aux]} \\
  Differentiating a primitive recovers the series.

\item[\href{https://github.com/LLM4Rocq/mathcomp-eulerian/blob/main/\detokenize{fps/fps_deriv.v}\#L138}{\texttt{\detokenize{fps_deriv_eq0}}}] \emph{Lemma, line 138}. \textsf{[aux]} \\
  A series with zero derivative is its constant coefficient.

\item[\href{https://github.com/LLM4Rocq/mathcomp-eulerian/blob/main/\detokenize{fps/fps_deriv.v}\#L149}{\texttt{\detokenize{fps_deriv_inj0}}}] \emph{Lemma, line 149}. \textsf{[aux]} \\
  (no docstring)

\end{description}

% --------------------------------------------
\section*{\texttt{fps/fps\_egf.v}}
\addcontentsline{toc}{section}{\texttt{fps/fps\_egf.v}}

\begin{description}
\item[\href{https://github.com/LLM4Rocq/mathcomp-eulerian/blob/main/\detokenize{fps/fps_egf.v}\#L35}{\texttt{\detokenize{egf}}}] \emph{Definition, line 35}. \textsf{[aux]} \\
  (no docstring)

\item[\href{https://github.com/LLM4Rocq/mathcomp-eulerian/blob/main/\detokenize{fps/fps_egf.v}\#L37}{\texttt{\detokenize{coef_egf}}}] \emph{Lemma, line 37}. \textsf{[aux]} \\
  (no docstring)

\item[\href{https://github.com/LLM4Rocq/mathcomp-eulerian/blob/main/\detokenize{fps/fps_egf.v}\#L40}{\texttt{\detokenize{factf_neq0}}}] \emph{Lemma, line 40}. \textsf{[aux]} \\
  (no docstring)

\item[\href{https://github.com/LLM4Rocq/mathcomp-eulerian/blob/main/\detokenize{fps/fps_egf.v}\#L43}{\texttt{\detokenize{binf_neq0}}}] \emph{Lemma, line 43}. \textsf{[aux]} \\
  (no docstring)

\item[\href{https://github.com/LLM4Rocq/mathcomp-eulerian/blob/main/\detokenize{fps/fps_egf.v}\#L47}{\texttt{\detokenize{egf_eq}}}] \emph{Lemma, line 47}. \textsf{[aux]} \\
  egf is injective, coefficientwise.

\item[\href{https://github.com/LLM4Rocq/mathcomp-eulerian/blob/main/\detokenize{fps/fps_egf.v}\#L54}{\texttt{\detokenize{eq_egf}}}] \emph{Lemma, line 54}. \textsf{[aux]} \\
  egf respects pointwise equality (via functional extensionality).

\item[\href{https://github.com/LLM4Rocq/mathcomp-eulerian/blob/main/\detokenize{fps/fps_egf.v}\#L57}{\texttt{\detokenize{egfD}}}] \emph{Lemma, line 57}. \textsf{[aux]} \\
  (no docstring)

\item[\href{https://github.com/LLM4Rocq/mathcomp-eulerian/blob/main/\detokenize{fps/fps_egf.v}\#L60}{\texttt{\detokenize{egfZ}}}] \emph{Lemma, line 60}. \textsf{[aux]} \\
  (no docstring)

\item[\href{https://github.com/LLM4Rocq/mathcomp-eulerian/blob/main/\detokenize{fps/fps_egf.v}\#L64}{\texttt{\detokenize{egf_dirac}}}] \emph{Lemma, line 64}. \textsf{[aux]} \\
  The egf of the Kronecker delta is 1.

\item[\href{https://github.com/LLM4Rocq/mathcomp-eulerian/blob/main/\detokenize{fps/fps_egf.v}\#L71}{\texttt{\detokenize{egf_mul}}}] \emph{Lemma, line 71}. \textsf{[aux]} \\
  THE workhorse: products of egfs are egfs of binomial convolutions.

\item[\href{https://github.com/LLM4Rocq/mathcomp-eulerian/blob/main/\detokenize{fps/fps_egf.v}\#L87}{\texttt{\detokenize{egf_deriv}}}] \emph{Lemma, line 87}. \textsf{[aux]} \\
  Differentiation of an egf is the index shift.

\end{description}

% --------------------------------------------
\section*{\texttt{fps/fps\_ode.v}}
\addcontentsline{toc}{section}{\texttt{fps/fps\_ode.v}}

\begin{description}
\item[\href{https://github.com/LLM4Rocq/mathcomp-eulerian/blob/main/\detokenize{fps/fps_ode.v}\#L32}{\texttt{\detokenize{fps_quad_ode_uniq}}}] \emph{Lemma, line 32}. \textsf{[aux]} \\
  Solutions of 2 y' = c + y\textasciicircum{}2 are determined by their constant term.

\end{description}

% --------------------------------------------
\section*{\texttt{fps/fps\_ogf.v}}
\addcontentsline{toc}{section}{\texttt{fps/fps\_ogf.v}}

\begin{description}
\item[\href{https://github.com/LLM4Rocq/mathcomp-eulerian/blob/main/\detokenize{fps/fps_ogf.v}\#L30}{\texttt{\detokenize{sum_bin_stair}}}] \emph{Lemma, line 30}. \textsf{[aux]} \\
  \textbackslash{}sum\_(j <= m) 'C(m - j + k, k) = 'C(m + k + 1, k + 1).

\item[\href{https://github.com/LLM4Rocq/mathcomp-eulerian/blob/main/\detokenize{fps/fps_ogf.v}\#L48}{\texttt{\detokenize{coef_fps_geomXn}}}] \emph{Lemma, line 48}. \textsf{[aux]} \\
  The negative binomial expansion: 1/(1-x)\textasciicircum{}(k+1) = sum\_m C(m+k,k) x\textasciicircum{}m.

\item[\href{https://github.com/LLM4Rocq/mathcomp-eulerian/blob/main/\detokenize{fps/fps_ogf.v}\#L59}{\texttt{\detokenize{onesubXn_mul_geomn}}}] \emph{Lemma, line 59}. \textsf{[aux]} \\
  Powers of 1 - x and of the geometric series are mutually inverse.

\end{description}

% --------------------------------------------
\section*{\texttt{fps/fps\_trig.v}}
\addcontentsline{toc}{section}{\texttt{fps/fps\_trig.v}}

\begin{description}
\item[\href{https://github.com/LLM4Rocq/mathcomp-eulerian/blob/main/\detokenize{fps/fps_trig.v}\#L34}{\texttt{\detokenize{expf}}}] \emph{Definition, line 34}. \textsf{[aux]} \\
  (no docstring)

\item[\href{https://github.com/LLM4Rocq/mathcomp-eulerian/blob/main/\detokenize{fps/fps_trig.v}\#L35}{\texttt{\detokenize{sinf}}}] \emph{Definition, line 35}. \textsf{[aux]} \\
  (no docstring)

\item[\href{https://github.com/LLM4Rocq/mathcomp-eulerian/blob/main/\detokenize{fps/fps_trig.v}\#L37}{\texttt{\detokenize{cosf}}}] \emph{Definition, line 37}. \textsf{[aux]} \\
  (no docstring)

\item[\href{https://github.com/LLM4Rocq/mathcomp-eulerian/blob/main/\detokenize{fps/fps_trig.v}\#L39}{\texttt{\detokenize{secf}}}] \emph{Definition, line 39}. \textsf{[aux]} \\
  (no docstring)

\item[\href{https://github.com/LLM4Rocq/mathcomp-eulerian/blob/main/\detokenize{fps/fps_trig.v}\#L40}{\texttt{\detokenize{tanf}}}] \emph{Definition, line 40}. \textsf{[aux]} \\
  (no docstring)

\item[\href{https://github.com/LLM4Rocq/mathcomp-eulerian/blob/main/\detokenize{fps/fps_trig.v}\#L46}{\texttt{\detokenize{expf0}}}] \emph{Lemma, line 46}. \textsf{[aux]} \\
  (no docstring)

\item[\href{https://github.com/LLM4Rocq/mathcomp-eulerian/blob/main/\detokenize{fps/fps_trig.v}\#L49}{\texttt{\detokenize{sinf0}}}] \emph{Lemma, line 49}. \textsf{[aux]} \\
  (no docstring)

\item[\href{https://github.com/LLM4Rocq/mathcomp-eulerian/blob/main/\detokenize{fps/fps_trig.v}\#L52}{\texttt{\detokenize{cosf0}}}] \emph{Lemma, line 52}. \textsf{[aux]} \\
  (no docstring)

\item[\href{https://github.com/LLM4Rocq/mathcomp-eulerian/blob/main/\detokenize{fps/fps_trig.v}\#L55}{\texttt{\detokenize{cosf_unit}}}] \emph{Lemma, line 55}. \textsf{[aux]} \\
  (no docstring)

\item[\href{https://github.com/LLM4Rocq/mathcomp-eulerian/blob/main/\detokenize{fps/fps_trig.v}\#L58}{\texttt{\detokenize{cos_mul_sec}}}] \emph{Lemma, line 58}. \textsf{[aux]} \\
  (no docstring)

\item[\href{https://github.com/LLM4Rocq/mathcomp-eulerian/blob/main/\detokenize{fps/fps_trig.v}\#L61}{\texttt{\detokenize{secf0}}}] \emph{Lemma, line 61}. \textsf{[aux]} \\
  (no docstring)

\item[\href{https://github.com/LLM4Rocq/mathcomp-eulerian/blob/main/\detokenize{fps/fps_trig.v}\#L67}{\texttt{\detokenize{tanf0}}}] \emph{Lemma, line 67}. \textsf{[aux]} \\
  (no docstring)

\item[\href{https://github.com/LLM4Rocq/mathcomp-eulerian/blob/main/\detokenize{fps/fps_trig.v}\#L74}{\texttt{\detokenize{expf_deriv}}}] \emph{Lemma, line 74}. \textsf{[aux]} \\
  (no docstring)

\item[\href{https://github.com/LLM4Rocq/mathcomp-eulerian/blob/main/\detokenize{fps/fps_trig.v}\#L77}{\texttt{\detokenize{sinf_deriv}}}] \emph{Lemma, line 77}. \textsf{[aux]} \\
  (no docstring)

\item[\href{https://github.com/LLM4Rocq/mathcomp-eulerian/blob/main/\detokenize{fps/fps_trig.v}\#L85}{\texttt{\detokenize{cosf_deriv}}}] \emph{Lemma, line 85}. \textsf{[aux]} \\
  (no docstring)

\item[\href{https://github.com/LLM4Rocq/mathcomp-eulerian/blob/main/\detokenize{fps/fps_trig.v}\#L97}{\texttt{\detokenize{sin2cos2}}}] \emph{Lemma, line 97}. \textsf{[aux]} \\
  (no docstring)

\item[\href{https://github.com/LLM4Rocq/mathcomp-eulerian/blob/main/\detokenize{fps/fps_trig.v}\#L108}{\texttt{\detokenize{sec2f}}}] \emph{Lemma, line 108}. \textsf{[aux]} \\
  sec\textasciicircum{}2 = 1 + tan\textasciicircum{}2.

\item[\href{https://github.com/LLM4Rocq/mathcomp-eulerian/blob/main/\detokenize{fps/fps_trig.v}\#L121}{\texttt{\detokenize{secf_deriv}}}] \emph{Lemma, line 121}. \textsf{[aux]} \\
  (no docstring)

\item[\href{https://github.com/LLM4Rocq/mathcomp-eulerian/blob/main/\detokenize{fps/fps_trig.v}\#L132}{\texttt{\detokenize{tanf_deriv}}}] \emph{Lemma, line 132}. \textsf{[aux]} \\
  (no docstring)

\item[\href{https://github.com/LLM4Rocq/mathcomp-eulerian/blob/main/\detokenize{fps/fps_trig.v}\#L142}{\texttt{\detokenize{sectan0}}}] \emph{Lemma, line 142}. \textsf{[aux]} \\
  (no docstring)

\item[\href{https://github.com/LLM4Rocq/mathcomp-eulerian/blob/main/\detokenize{fps/fps_trig.v}\#L145}{\texttt{\detokenize{sectan_ode}}}] \emph{Theorem, line 145}. \textsf{[aux]} \\
  (no docstring)

\end{description}

% --------------------------------------------
\section*{\texttt{inversions.v}}
\addcontentsline{toc}{section}{\texttt{inversions.v}}

\begin{description}
\item[\href{https://github.com/LLM4Rocq/mathcomp-eulerian/blob/main/\detokenize{inversions.v}\#L27}{\texttt{\detokenize{is_inv}}}] \emph{Definition, line 27}. \textsf{[Ch 3 §3.1]} \\
  \texttt{\detokenize{is_inv s i j}} holds when the pair \texttt{\detokenize{(i, j)}} is an inversion of \texttt{\detokenize{s}},     i.e. \texttt{\detokenize{i < j}} but \texttt{\detokenize{s i > s j}}.  Stanley EC1 \textbackslash{}S\{\}1.3.3.

\item[\href{https://github.com/LLM4Rocq/mathcomp-eulerian/blob/main/\detokenize{inversions.v}\#L32}{\texttt{\detokenize{inv_set}}}] \emph{Definition, line 32}. \textsf{[Ch 3 §3.1]} \\
  \texttt{\detokenize{inv_set s}} is the set of inversions of \texttt{\detokenize{s}}: pairs \texttt{\detokenize{(i, j)}} with     \texttt{\detokenize{i < j}} but \texttt{\detokenize{s i > s j}}.

\item[\href{https://github.com/LLM4Rocq/mathcomp-eulerian/blob/main/\detokenize{inversions.v}\#L37}{\texttt{\detokenize{inv}}}] \emph{Definition, line 37}. \textsf{[Ch 3 §3.1]} \\
  \texttt{\detokenize{inv s}} is the inversion count of \texttt{\detokenize{s}}, i.e. \texttt{\detokenize{#|inv_set s|}}.     Stanley's \texttt{\detokenize{inv(w)}}.

\item[\href{https://github.com/LLM4Rocq/mathcomp-eulerian/blob/main/\detokenize{inversions.v}\#L41}{\texttt{\detokenize{inv_id}}}] \emph{Lemma, line 41}. \textsf{[Ch 3 §3.1]} \\
  The identity permutation has no inversions: \texttt{\detokenize{inv 1 = 0}}.

\item[\href{https://github.com/LLM4Rocq/mathcomp-eulerian/blob/main/\detokenize{inversions.v}\#L51}{\texttt{\detokenize{inv_double_sum}}}] \emph{Lemma, line 51}. \textsf{[aux]} \\
  Bridge lemma: \texttt{\detokenize{inv s}} expressed as the nested double-sum     \texttt{\textbackslash\{\}sum\_j \textbackslash\{\}sum\_(i < j) (s j < s i)}.  Bridges the \texttt{\detokenize{inv_set}}     cardinality form to the seq-level double-sum form of \texttt{\detokenize{inv_seq}},     used in \texttt{\detokenize{foata.v}} to prove \texttt{\detokenize{inv_eq_inv_seq}}.

\item[\href{https://github.com/LLM4Rocq/mathcomp-eulerian/blob/main/\detokenize{inversions.v}\#L82}{\texttt{\detokenize{maj}}}] \emph{Definition, line 82}. \textsf{[Ch 3 §3.2]} \\
  \texttt{\detokenize{maj s}} is the major index: the sum of the (1-indexed) descent     positions of \texttt{\detokenize{s}}.  Stanley EC1 \textbackslash{}S\{\}1.3.3.  Our \texttt{\detokenize{descent_set s}} is     0-indexed in \texttt{\detokenize{'I_n}}; we shift by 1 to recover Stanley's     1-indexed positions in \texttt{\{1, ..., n\}}.

\item[\href{https://github.com/LLM4Rocq/mathcomp-eulerian/blob/main/\detokenize{inversions.v}\#L85}{\texttt{\detokenize{maj_id}}}] \emph{Lemma, line 85}. \textsf{[aux]} \\
  The identity permutation has zero major index: \texttt{\detokenize{maj 1 = 0}}.

\item[\href{https://github.com/LLM4Rocq/mathcomp-eulerian/blob/main/\detokenize{inversions.v}\#L101}{\texttt{\detokenize{card_ord_pair}}}] \emph{Lemma, line 101}. \textsf{[aux]} \\
  The number of strictly-ordered pairs in \texttt{\detokenize{'I_m * 'I_m}} is the binomial     coefficient \texttt{\detokenize{'C(m, 2)}}.  Used to bound \texttt{\detokenize{inv}} and \texttt{\detokenize{maj}}.

\item[\href{https://github.com/LLM4Rocq/mathcomp-eulerian/blob/main/\detokenize{inversions.v}\#L142}{\texttt{\detokenize{inv_le}}}] \emph{Lemma, line 142}. \textsf{[Ch 3 §3.1]} \\
  Bound: \texttt{\detokenize{inv s <= 'C(n+1, 2)}} for \texttt{s : \{perm 'I\_n.+1\}}, since     inversions are a subset of all strictly-ordered pairs.

\item[\href{https://github.com/LLM4Rocq/mathcomp-eulerian/blob/main/\detokenize{inversions.v}\#L154}{\texttt{\detokenize{sum_succ_ord_eq_bin2}}}] \emph{Lemma, line 154}. \textsf{[aux]} \\
  Helper: \texttt{\textbackslash\{\}sum\_(i : 'I\_n) (val i).+1 = 'C(n+1, 2)}.  Used to bound     \texttt{\detokenize{maj}}.

\item[\href{https://github.com/LLM4Rocq/mathcomp-eulerian/blob/main/\detokenize{inversions.v}\#L168}{\texttt{\detokenize{maj_le}}}] \emph{Lemma, line 168}. \textsf{[Ch 3 §3.2]} \\
  Bound: \texttt{\detokenize{maj s <= 'C(n+1, 2)}} for \texttt{s : \{perm 'I\_n.+1\}}, by comparing     descent-position sum to the full sum of \texttt{\detokenize{(val i).+1}} over \texttt{\detokenize{'I_n}}.

\item[\href{https://github.com/LLM4Rocq/mathcomp-eulerian/blob/main/\detokenize{inversions.v}\#L180}{\texttt{\detokenize{coinv_set}}}] \emph{Definition, line 180}. \textsf{[Ch 3 §3.1]} \\
  \texttt{\detokenize{coinv_set s}} is the co-inversion set of \texttt{\detokenize{s}}: pairs \texttt{\detokenize{(i, j)}} with     \texttt{\detokenize{i < j}} and \texttt{\detokenize{s i < s j}}.  Since \texttt{\detokenize{s}} is a permutation, every     strictly-ordered pair is either an inversion or a co-inversion.

\item[\href{https://github.com/LLM4Rocq/mathcomp-eulerian/blob/main/\detokenize{inversions.v}\#L185}{\texttt{\detokenize{rev_ord_lt}}}] \emph{Lemma, line 185}. \textsf{[aux]} \\
  Order-reversal property of \texttt{\detokenize{rev_ord}}: \texttt{\detokenize{rev_ord a < rev_ord b}} iff     \texttt{\detokenize{b < a}}.  Used in the \texttt{\detokenize{inv}} reversal identity.

\item[\href{https://github.com/LLM4Rocq/mathcomp-eulerian/blob/main/\detokenize{inversions.v}\#L194}{\texttt{\detokenize{card_split_ord_pair}}}] \emph{Lemma, line 194}. \textsf{[aux]} \\
  Split: every strictly-ordered pair is either an inversion or a     co-inversion, so \texttt{\detokenize{inv s + #|coinv_set s| = 'C(n+1, 2)}}.

\item[\href{https://github.com/LLM4Rocq/mathcomp-eulerian/blob/main/\detokenize{inversions.v}\#L214}{\texttt{\detokenize{inv_rev_perm_eq_coinv}}}] \emph{Lemma, line 214}. \textsf{[aux]} \\
  The inversions of \texttt{\detokenize{rev_perm s}} are in bijection with the     co-inversions of \texttt{\detokenize{s}} via \texttt{\detokenize{(i, j) |-> (rev_ord j, rev_ord i)}}.

\item[\href{https://github.com/LLM4Rocq/mathcomp-eulerian/blob/main/\detokenize{inversions.v}\#L238}{\texttt{\detokenize{inv_rev_perm}}}] \emph{Lemma, line 238}. \textsf{[Ch 3 §3.1]} \\
  Reversal identity for \texttt{\detokenize{inv}}: \texttt{\detokenize{inv (rev_perm s) = 'C(n+1, 2) - inv s}}.     Stanley EC1 \textbackslash{}S\{\}1.3.3.

\item[\href{https://github.com/LLM4Rocq/mathcomp-eulerian/blob/main/\detokenize{inversions.v}\#L260}{\texttt{\detokenize{maj_rev_perm}}}] \emph{Lemma, line 260}. \textsf{[aux]} \\
  Reversal identity for \texttt{\detokenize{maj}}:     \texttt{\detokenize{maj (rev_perm s) + (n+1) * des s = 'C(n+1, 2) + maj s}}.     The naive \texttt{\detokenize{maj (rev_perm s) = 'C(n+1, 2) - maj s}} is FALSE; the     correct identity carries a \texttt{\detokenize{(n+1) * des s}} offset (Stanley EC1).

\end{description}

% --------------------------------------------
\section*{\texttt{mmtree.v}}
\addcontentsline{toc}{section}{\texttt{mmtree.v}}

\begin{description}
\item[\href{https://github.com/LLM4Rocq/mathcomp-eulerian/blob/main/\detokenize{mmtree.v}\#L39}{\texttt{\detokenize{mmtree}}}] \emph{Inductive, line 39}. \textsf{[aux]} \\
  \texttt{\detokenize{mmtree T}} is the labelled binary tree datatype underlying Stanley's     min-max tree construction \texttt{\detokenize{M(w)}} for sequences with labels in \texttt{\detokenize{T}}.

\item[\href{https://github.com/LLM4Rocq/mathcomp-eulerian/blob/main/\detokenize{mmtree.v}\#L48}{\texttt{\detokenize{mmtree_to_seq}}}] \emph{Fixpoint, line 48}. \textsf{[aux]} \\
  \texttt{\detokenize{mmtree_to_seq t}} is the in-order traversal of \texttt{\detokenize{t}}: left subtree, then     root, then right subtree.

\item[\href{https://github.com/LLM4Rocq/mathcomp-eulerian/blob/main/\detokenize{mmtree.v}\#L56}{\texttt{\detokenize{min_pos}}}] \emph{Definition, line 56}. \textsf{[aux]} \\
  \texttt{\detokenize{min_pos s}} is the least index \texttt{\detokenize{j}} in \texttt{\detokenize{s}} at which the minimum of \texttt{\detokenize{s}}     occurs; returns 0 on the empty sequence (vacuous case).

\item[\href{https://github.com/LLM4Rocq/mathcomp-eulerian/blob/main/\detokenize{mmtree.v}\#L61}{\texttt{\detokenize{min_in}}}] \emph{Lemma, line 61}. \textsf{[aux]} \\
  \texttt{\detokenize{min_in}} states that \texttt{\detokenize{foldr minn a s}} always lies in the sequence     \texttt{\detokenize{a :: s}}; used to bound \texttt{\detokenize{min_pos}} within range.

\item[\href{https://github.com/LLM4Rocq/mathcomp-eulerian/blob/main/\detokenize{mmtree.v}\#L74}{\texttt{\detokenize{min_pos_lt}}}] \emph{Lemma, line 74}. \textsf{[aux]} \\
  \texttt{\detokenize{min_pos_lt}} : on a nonempty sequence the split index \texttt{\detokenize{min_pos s}} is in     range, ensuring the recursion in \texttt{\detokenize{mmtree_of_seq_fuel}} is well-defined.

\item[\href{https://github.com/LLM4Rocq/mathcomp-eulerian/blob/main/\detokenize{mmtree.v}\#L85}{\texttt{\detokenize{mmtree_of_seq_fuel}}}] \emph{Fixpoint, line 85}. \textsf{[aux]} \\
  \texttt{\detokenize{mmtree_of_seq_fuel fuel s}} is the fuel-bounded M1 tree construction:     splits \texttt{\detokenize{s}} at \texttt{\detokenize{min_pos s}}, recursing on the take/drop halves.  Fuel     equal to \texttt{\detokenize{size s}} suffices since each recursion strictly shrinks.

\item[\href{https://github.com/LLM4Rocq/mathcomp-eulerian/blob/main/\detokenize{mmtree.v}\#L101}{\texttt{\detokenize{mmtree_of_seq}}}] \emph{Definition, line 101}. \textsf{[aux]} \\
  \texttt{\detokenize{mmtree_of_seq s}} runs \texttt{\detokenize{mmtree_of_seq_fuel}} with fuel \texttt{\detokenize{size s}}, the     M1 (minimum-only) variant of Stanley's min-max tree construction.

\item[\href{https://github.com/LLM4Rocq/mathcomp-eulerian/blob/main/\detokenize{mmtree.v}\#L107}{\texttt{\detokenize{mmtree_of_seq_fuel_correct}}}] \emph{Lemma, line 107}. \textsf{[aux]} \\
  \texttt{\detokenize{mmtree_of_seq_fuel_correct}} : the fuel-bounded construction round-trips     in-order, i.e. \texttt{\detokenize{mmtree_to_seq}} inverts \texttt{\detokenize{mmtree_of_seq_fuel}} when fuel     bounds the sequence size.

\item[\href{https://github.com/LLM4Rocq/mathcomp-eulerian/blob/main/\detokenize{mmtree.v}\#L136}{\texttt{\detokenize{mmtree_of_seqK}}}] \emph{Theorem, line 136}. \textsf{[aux]} \\
  \texttt{\detokenize{mmtree_of_seqK}} is the M1 round-trip theorem: in-order traversal     inverts \texttt{\detokenize{mmtree_of_seq}} for every input sequence.

\item[\href{https://github.com/LLM4Rocq/mathcomp-eulerian/blob/main/\detokenize{mmtree.v}\#L148}{\texttt{\detokenize{ex_seq}}}] \emph{Definition, line 148}. \textsf{[aux]} \\
  \texttt{\detokenize{ex_seq}} is the concrete test sequence used to demonstrate that the     construction yields a genuinely branching tree.

\end{description}

% --------------------------------------------
\section*{\texttt{ordinal\_reindex.v}}
\addcontentsline{toc}{section}{\texttt{ordinal\_reindex.v}}

\begin{description}
\item[\href{https://github.com/LLM4Rocq/mathcomp-eulerian/blob/main/\detokenize{ordinal_reindex.v}\#L18}{\texttt{\detokenize{leq_lift}}}] \emph{Lemma, line 18}. \textsf{[aux]} \\
  Weak monotonicity of \texttt{\detokenize{lift k : 'I_n -> 'I_n.+1}}: lifting preserves \texttt{\detokenize{<=}}.

\item[\href{https://github.com/LLM4Rocq/mathcomp-eulerian/blob/main/\detokenize{ordinal_reindex.v}\#L22}{\texttt{\detokenize{ltn_lift}}}] \emph{Lemma, line 22}. \textsf{[aux]} \\
  Strict monotonicity of \texttt{\detokenize{lift k : 'I_n -> 'I_n.+1}}: lifting preserves \texttt{\detokenize{<}}.

\end{description}

% --------------------------------------------
\section*{\texttt{perm\_compress.v}}
\addcontentsline{toc}{section}{\texttt{perm\_compress.v}}

\begin{description}
\item[\href{https://github.com/LLM4Rocq/mathcomp-eulerian/blob/main/\detokenize{perm_compress.v}\#L23}{\texttt{\detokenize{drop_fun_ne}}}] \emph{Fact, line 23}. \textsf{[aux]} \\
  \texttt{\detokenize{s}} sends \texttt{\detokenize{ord_max}} and \texttt{\detokenize{lift ord_max i}} to distinct values, by     injectivity.

\item[\href{https://github.com/LLM4Rocq/mathcomp-eulerian/blob/main/\detokenize{perm_compress.v}\#L29}{\texttt{\detokenize{drop_fun}}}] \emph{Definition, line 29}. \textsf{[aux]} \\
  \texttt{\detokenize{drop_fun i}} is the unique \texttt{\detokenize{j : 'I_n}} such that     \texttt{\detokenize{s (lift ord_max i) = lift (s ord_max) j}}. This is \texttt{\detokenize{s}} restricted to     \texttt{'I\_n.+1 \textbackslash\{\} \{ord\_max\}} and reindexed onto \texttt{\detokenize{'I_n}}.

\item[\href{https://github.com/LLM4Rocq/mathcomp-eulerian/blob/main/\detokenize{perm_compress.v}\#L33}{\texttt{\detokenize{drop_funE}}}] \emph{Lemma, line 33}. \textsf{[aux]} \\
  Defining equation for \texttt{\detokenize{drop_fun}}: \texttt{\detokenize{s (lift ord_max i)}} equals     \texttt{\detokenize{lift (s ord_max) (drop_fun i)}}.

\item[\href{https://github.com/LLM4Rocq/mathcomp-eulerian/blob/main/\detokenize{perm_compress.v}\#L37}{\texttt{\detokenize{drop_fun_inj}}}] \emph{Lemma, line 37}. \textsf{[aux]} \\
  \texttt{\detokenize{drop_fun}} is injective, inheriting injectivity from \texttt{\detokenize{s}}.

\item[\href{https://github.com/LLM4Rocq/mathcomp-eulerian/blob/main/\detokenize{perm_compress.v}\#L45}{\texttt{\detokenize{drop_perm}}}] \emph{Definition, line 45}. \textsf{[aux]} \\
  \texttt{\detokenize{drop_perm s}} is the permutation of \texttt{\detokenize{'I_n}} obtained by deleting position     \texttt{\detokenize{ord_max}} from \texttt{\detokenize{s}} and collapsing around the value \texttt{\detokenize{s ord_max}}.

\item[\href{https://github.com/LLM4Rocq/mathcomp-eulerian/blob/main/\detokenize{perm_compress.v}\#L48}{\texttt{\detokenize{drop_permE}}}] \emph{Lemma, line 48}. \textsf{[aux]} \\
  Equivalence with the underlying \texttt{\detokenize{drop_fun}} for rewriting.

\item[\href{https://github.com/LLM4Rocq/mathcomp-eulerian/blob/main/\detokenize{perm_compress.v}\#L53}{\texttt{\detokenize{lift_drop_permE}}}] \emph{Lemma, line 53}. \textsf{[aux]} \\
  Defining equation: lifting \texttt{\detokenize{drop_perm}} back through \texttt{\detokenize{s ord_max}} recovers \texttt{\detokenize{s}}     on lifted positions. Used to relate descents of \texttt{\detokenize{s}} and \texttt{\detokenize{drop_perm s}}.

\item[\href{https://github.com/LLM4Rocq/mathcomp-eulerian/blob/main/\detokenize{perm_compress.v}\#L68}{\texttt{\detokenize{lift_fun}}}] \emph{Definition, line 68}. \textsf{[aux]} \\
  Underlying function of \texttt{\detokenize{lift_perm k t}}: sends \texttt{\detokenize{ord_max}} to \texttt{\detokenize{k}} and lifts     \texttt{\detokenize{t}} elsewhere via \texttt{\detokenize{lift k}}.

\item[\href{https://github.com/LLM4Rocq/mathcomp-eulerian/blob/main/\detokenize{perm_compress.v}\#L75}{\texttt{\detokenize{lift_fun_ord_max}}}] \emph{Lemma, line 75}. \textsf{[aux]} \\
  \texttt{\detokenize{lift_fun}} sends \texttt{\detokenize{ord_max}} to \texttt{\detokenize{k}}.

\item[\href{https://github.com/LLM4Rocq/mathcomp-eulerian/blob/main/\detokenize{perm_compress.v}\#L79}{\texttt{\detokenize{lift_fun_lift}}}] \emph{Lemma, line 79}. \textsf{[aux]} \\
  \texttt{\detokenize{lift_fun}} on lifted positions equals \texttt{\detokenize{lift k (t .)}}.

\item[\href{https://github.com/LLM4Rocq/mathcomp-eulerian/blob/main/\detokenize{perm_compress.v}\#L83}{\texttt{\detokenize{lift_fun_inj}}}] \emph{Lemma, line 83}. \textsf{[aux]} \\
  \texttt{\detokenize{lift_fun}} is injective: combine injectivity of \texttt{\detokenize{lift}} and \texttt{\detokenize{t}}.

\item[\href{https://github.com/LLM4Rocq/mathcomp-eulerian/blob/main/\detokenize{perm_compress.v}\#L95}{\texttt{\detokenize{lift_perm}}}] \emph{Definition, line 95}. \textsf{[aux]} \\
  \texttt{\detokenize{lift_perm k t}} is the permutation of \texttt{\detokenize{'I_n.+1}} obtained from     \texttt{t : \{perm 'I\_n\}} by inserting a new maximum position \texttt{\detokenize{ord_max}} sent to     \texttt{\detokenize{k}}. Inverse of \texttt{\detokenize{drop_perm}}.

\item[\href{https://github.com/LLM4Rocq/mathcomp-eulerian/blob/main/\detokenize{perm_compress.v}\#L98}{\texttt{\detokenize{lift_permE}}}] \emph{Lemma, line 98}. \textsf{[aux]} \\
  Equivalence with the underlying \texttt{\detokenize{lift_fun}} for rewriting.

\item[\href{https://github.com/LLM4Rocq/mathcomp-eulerian/blob/main/\detokenize{perm_compress.v}\#L102}{\texttt{\detokenize{lift_perm_ord_max}}}] \emph{Lemma, line 102}. \textsf{[aux]} \\
  \texttt{\detokenize{lift_perm k t}} sends \texttt{\detokenize{ord_max}} to \texttt{\detokenize{k}}.

\item[\href{https://github.com/LLM4Rocq/mathcomp-eulerian/blob/main/\detokenize{perm_compress.v}\#L106}{\texttt{\detokenize{lift_perm_lift}}}] \emph{Lemma, line 106}. \textsf{[aux]} \\
  \texttt{\detokenize{lift_perm k t}} on lifted positions equals \texttt{\detokenize{lift k (t .)}}.

\item[\href{https://github.com/LLM4Rocq/mathcomp-eulerian/blob/main/\detokenize{perm_compress.v}\#L125}{\texttt{\detokenize{lift_perm_drop_perm}}}] \emph{Lemma, line 125}. \textsf{[aux]} \\
  Right inverse: lifting a dropped permutation at the dropped value     recovers the original. Establishes the bijection     \texttt{\{perm 'I\_n.+1\}} \textasciitilde{}= \texttt{\detokenize{'I_n.+1}} x \texttt{\{perm 'I\_n\}}.

\end{description}

% --------------------------------------------
\section*{\texttt{perm\_seq\_basics.v}}
\addcontentsline{toc}{section}{\texttt{perm\_seq\_basics.v}}

\begin{description}
\item[\href{https://github.com/LLM4Rocq/mathcomp-eulerian/blob/main/\detokenize{perm_seq_basics.v}\#L35}{\texttt{\detokenize{is_descent_seq}}}] \emph{Definition, line 35}. \textsf{[aux]} \\
  \texttt{\detokenize{is_descent_seq w k}} is the descent predicate: \texttt{\detokenize{k}} is a descent of     \texttt{\detokenize{w}} iff \texttt{w\_k > w\_\{k+1\}}.  Sequence-level mirror of Stanley's     \texttt{Des(w) = \{ i : w\_i > w\_\{i+1\} \}}.  Moved here from     \texttt{\detokenize{psi_descent_v2.v}} so \textbackslash{}S\{\}1.4 files can use it without importing the     cd-index machinery.

\item[\href{https://github.com/LLM4Rocq/mathcomp-eulerian/blob/main/\detokenize{perm_seq_basics.v}\#L44}{\texttt{\detokenize{perm_to_seq}}}] \emph{Definition, line 44}. \textsf{[aux]} \\
  \texttt{\detokenize{perm_to_seq s}} -- bridge from \texttt{\{perm 'I\_n\}} to \texttt{\detokenize{seq nat}}: the     list \texttt{\detokenize{[val (s 0); val (s 1); ...; val (s (n-1))]}}.

\item[\href{https://github.com/LLM4Rocq/mathcomp-eulerian/blob/main/\detokenize{perm_seq_basics.v}\#L49}{\texttt{\detokenize{perm_to_seq_size}}}] \emph{Lemma, line 49}. \textsf{[aux]} \\
  \texttt{\detokenize{perm_to_seq_size}} -- the seq view of a permutation has length     \texttt{\detokenize{n}}.

\item[\href{https://github.com/LLM4Rocq/mathcomp-eulerian/blob/main/\detokenize{perm_seq_basics.v}\#L54}{\texttt{\detokenize{perm_to_seq_uniq}}}] \emph{Lemma, line 54}. \textsf{[aux]} \\
  \texttt{\detokenize{perm_to_seq_uniq}} -- the seq view of a permutation has no     duplicates (since \texttt{\detokenize{s}} is injective).

\item[\href{https://github.com/LLM4Rocq/mathcomp-eulerian/blob/main/\detokenize{perm_seq_basics.v}\#L62}{\texttt{\detokenize{nth_perm_to_seq}}}] \emph{Lemma, line 62}. \textsf{[aux]} \\
  \texttt{\detokenize{nth_perm_to_seq}} -- accessing \texttt{\detokenize{perm_to_seq s}} at index \texttt{\detokenize{k}}     yields \texttt{\detokenize{val (s (Ordinal Hk))}}.

\item[\href{https://github.com/LLM4Rocq/mathcomp-eulerian/blob/main/\detokenize{perm_seq_basics.v}\#L72}{\texttt{\detokenize{perm_to_seq_inj}}}] \emph{Lemma, line 72}. \textsf{[aux]} \\
  \texttt{\detokenize{perm_to_seq_inj}} -- the \texttt{\detokenize{perm_to_seq}} map is injective: distinct     permutations yield distinct value lists.

\item[\href{https://github.com/LLM4Rocq/mathcomp-eulerian/blob/main/\detokenize{perm_seq_basics.v}\#L90}{\texttt{\detokenize{is_descent_perm_seq}}}] \emph{Lemma, line 90}. \textsf{[aux]} \\
  \texttt{\detokenize{is_descent_perm_seq}} -- descent equivalence: \texttt{\detokenize{is_descent s i}}     (perm side) equals \texttt{\detokenize{is_descent_seq (perm_to_seq s) (val i)}} (seq     side).

\item[\href{https://github.com/LLM4Rocq/mathcomp-eulerian/blob/main/\detokenize{perm_seq_basics.v}\#L107}{\texttt{\detokenize{descent_to_bvec}}}] \emph{Definition, line 107}. \textsf{[aux]} \\
  \texttt{\detokenize{descent_to_bvec D}} -- boolean-vector representation of a descent     set, indexed by \texttt{\detokenize{enum 'I_n}}: the \texttt{\detokenize{i}}-th entry is \texttt{(i \textbackslash\{\}in D)}.

\item[\href{https://github.com/LLM4Rocq/mathcomp-eulerian/blob/main/\detokenize{perm_seq_basics.v}\#L111}{\texttt{\detokenize{size_descent_to_bvec}}}] \emph{Lemma, line 111}. \textsf{[aux]} \\
  \texttt{\detokenize{size_descent_to_bvec}} -- the boolean vector has length \texttt{\detokenize{n}}.

\item[\href{https://github.com/LLM4Rocq/mathcomp-eulerian/blob/main/\detokenize{perm_seq_basics.v}\#L117}{\texttt{\detokenize{nth_descent_to_bvec}}}] \emph{Lemma, line 117}. \textsf{[aux]} \\
  \texttt{\detokenize{nth_descent_to_bvec}} -- accessing \texttt{\detokenize{descent_to_bvec D}} at \texttt{\detokenize{k}}     yields \texttt{(Ordinal Hk \textbackslash\{\}in D)}.

\item[\href{https://github.com/LLM4Rocq/mathcomp-eulerian/blob/main/\detokenize{perm_seq_basics.v}\#L127}{\texttt{\detokenize{descent_to_bvec_inj}}}] \emph{Lemma, line 127}. \textsf{[aux]} \\
  \texttt{\detokenize{descent_to_bvec_inj}} -- two descent sets giving the same     boolean vector must be equal.

\item[\href{https://github.com/LLM4Rocq/mathcomp-eulerian/blob/main/\detokenize{perm_seq_basics.v}\#L145}{\texttt{\detokenize{desc_positions_bvec}}}] \emph{Lemma, line 145}. \textsf{[aux]} \\
  \texttt{\detokenize{desc_positions_bvec}} -- the descent positions extracted from the     boolean vector \texttt{\detokenize{descent_to_bvec D}} coincide with \texttt{\detokenize{set_to_seq D}}     (both sorted-asc lists of positions of \texttt{\detokenize{D}}).

\item[\href{https://github.com/LLM4Rocq/mathcomp-eulerian/blob/main/\detokenize{perm_seq_basics.v}\#L183}{\texttt{\detokenize{seq_nth_bound}}}] \emph{Lemma, line 183}. \textsf{[aux]} \\
  \texttt{\detokenize{seq_nth_bound}} -- under the section hypotheses (uniqness, bound,     size), every entry of \texttt{\detokenize{w}} read at an ordinal index is itself     bounded by \texttt{\detokenize{n}}.

\item[\href{https://github.com/LLM4Rocq/mathcomp-eulerian/blob/main/\detokenize{perm_seq_basics.v}\#L191}{\texttt{\detokenize{seq_to_fun}}}] \emph{Definition, line 191}. \textsf{[aux]} \\
  \texttt{\detokenize{seq_to_fun i}} -- the underlying function \texttt{\detokenize{i |-> nth 0 w (val i)}}     on ordinals, used to lift a uniq bounded seq to a permutation.

\item[\href{https://github.com/LLM4Rocq/mathcomp-eulerian/blob/main/\detokenize{perm_seq_basics.v}\#L195}{\texttt{\detokenize{seq_to_fun_inj}}}] \emph{Lemma, line 195}. \textsf{[aux]} \\
  \texttt{\detokenize{seq_to_fun_inj}} -- \texttt{\detokenize{seq_to_fun}} is injective (uses uniqueness of     \texttt{\detokenize{w}}).

\item[\href{https://github.com/LLM4Rocq/mathcomp-eulerian/blob/main/\detokenize{perm_seq_basics.v}\#L208}{\texttt{\detokenize{seq_to_perm}}}] \emph{Definition, line 208}. \textsf{[aux]} \\
  \texttt{\detokenize{seq_to_perm}} -- the permutation built from \texttt{\detokenize{seq_to_fun}}; inverse     of \texttt{\detokenize{perm_to_seq}} for uniq, \texttt{\detokenize{n}}-bounded seqs of length \texttt{\detokenize{n}}.

\item[\href{https://github.com/LLM4Rocq/mathcomp-eulerian/blob/main/\detokenize{perm_seq_basics.v}\#L218}{\texttt{\detokenize{perm_to_seq_bnd}}}] \emph{Lemma, line 218}. \textsf{[aux]} \\
  \texttt{\detokenize{perm_to_seq_bnd}} -- every entry of \texttt{\detokenize{perm_to_seq s}} is strictly     less than \texttt{\detokenize{n}}.

\item[\href{https://github.com/LLM4Rocq/mathcomp-eulerian/blob/main/\detokenize{perm_seq_basics.v}\#L225}{\texttt{\detokenize{perm_to_seq_seq_to_perm}}}] \emph{Lemma, line 225}. \textsf{[aux]} \\
  \texttt{\detokenize{perm_to_seq_seq_to_perm}} -- round-trip identity:     \texttt{\detokenize{perm_to_seq (seq_to_perm w) = w}} for uniq, bounded seqs of size     \texttt{\detokenize{n}}.

\end{description}

% --------------------------------------------
\section*{\texttt{perm\_seq\_bridge.v}}
\addcontentsline{toc}{section}{\texttt{perm\_seq\_bridge.v}}

\begin{description}
\item[\href{https://github.com/LLM4Rocq/mathcomp-eulerian/blob/main/\detokenize{perm_seq_bridge.v}\#L44}{\texttt{\detokenize{char_mono_perm_to_seq}}}] \emph{Lemma, line 44}. \textsf{[aux]} \\
  \texttt{\detokenize{char_mono_perm_to_seq}} -- the seq-side \texttt{\detokenize{char_mono}} of     \texttt{\detokenize{perm_to_seq s}} equals the perm-side boolean vector of     \texttt{\detokenize{descent_set s}}. Bridges \texttt{\detokenize{psi_cdindex}} descent patterns to perm     descent sets.

\item[\href{https://github.com/LLM4Rocq/mathcomp-eulerian/blob/main/\detokenize{perm_seq_bridge.v}\#L70}{\texttt{\detokenize{psi_apply_psis_comm}}}] \emph{Lemma, line 70}. \textsf{[aux]} \\
  \texttt{\detokenize{psi_apply_psis_comm}} -- the single \texttt{\detokenize{psi i}} commutes with the     iterated \texttt{\detokenize{apply_psis ss}} (on uniq seqs), via \texttt{\detokenize{psi_comm}}.

\item[\href{https://github.com/LLM4Rocq/mathcomp-eulerian/blob/main/\detokenize{perm_seq_bridge.v}\#L83}{\texttt{\detokenize{apply_psis_rev}}}] \emph{Lemma, line 83}. \textsf{[aux]} \\
  \texttt{\detokenize{apply_psis_rev}} -- since each \texttt{\detokenize{psi i}} is involutive and they     commute, \texttt{\detokenize{apply_psis}} is independent of the order of the     operations: \texttt{\detokenize{rev ss}} gives the same result as \texttt{\detokenize{ss}}.

\item[\href{https://github.com/LLM4Rocq/mathcomp-eulerian/blob/main/\detokenize{perm_seq_bridge.v}\#L94}{\texttt{\detokenize{apply_psis_revK}}}] \emph{Lemma, line 94}. \textsf{[aux]} \\
  \texttt{\detokenize{apply_psis_revK}} -- applying \texttt{\detokenize{ss}} then \texttt{\detokenize{rev ss}} cancels back to     \texttt{\detokenize{w}} (each \texttt{\detokenize{psi i}} is involutive).

\item[\href{https://github.com/LLM4Rocq/mathcomp-eulerian/blob/main/\detokenize{perm_seq_bridge.v}\#L107}{\texttt{\detokenize{apply_psis_cancel}}}] \emph{Lemma, line 107}. \textsf{[aux]} \\
  \texttt{\detokenize{apply_psis_cancel}} -- \texttt{\detokenize{apply_psis ss}} is its own inverse:     applying \texttt{\detokenize{ss}} twice cancels back to \texttt{\detokenize{w}}.

\item[\href{https://github.com/LLM4Rocq/mathcomp-eulerian/blob/main/\detokenize{perm_seq_bridge.v}\#L119}{\texttt{\detokenize{powerset_internal_apply_psis}}}] \emph{Lemma, line 119}. \textsf{[aux]} \\
  \texttt{\detokenize{powerset_internal_apply_psis}} -- the internal-vertex powerset     used to enumerate the M-class is invariant under \texttt{\detokenize{apply_psis}};     consequence of \texttt{\detokenize{internal_vertices_apply_psis}}.

\item[\href{https://github.com/LLM4Rocq/mathcomp-eulerian/blob/main/\detokenize{perm_seq_bridge.v}\#L132}{\texttt{\detokenize{class_char_monos_uniq}}}] \emph{Lemma, line 132}. \textsf{[aux]} \\
  \texttt{\detokenize{class_char_monos_uniq}} -- the descent patterns of the M-class     members of \texttt{\detokenize{w}} are pairwise distinct (combines \texttt{\detokenize{fact3}} with     \texttt{\detokenize{uniq_expand_cde}}).

\item[\href{https://github.com/LLM4Rocq/mathcomp-eulerian/blob/main/\detokenize{perm_seq_bridge.v}\#L146}{\texttt{\detokenize{char_mono_class_inj}}}] \emph{Lemma, line 146}. \textsf{[aux]} \\
  \texttt{\detokenize{char_mono_class_inj}} -- M-class injectivity: within the M-class     of \texttt{\detokenize{w}}, two class members with the same descent pattern are     actually the same sequence. Follows from \texttt{\detokenize{fact3}} and     \texttt{\detokenize{uniq_expand_cde}}.

\item[\href{https://github.com/LLM4Rocq/mathcomp-eulerian/blob/main/\detokenize{perm_seq_bridge.v}\#L203}{\texttt{\detokenize{all_bnd_apply_psis}}}] \emph{Lemma, line 203}. \textsf{[aux]} \\
  \texttt{\detokenize{all_bnd_apply_psis}} -- the boundedness predicate     \texttt{\detokenize{all (< n)}} is preserved by \texttt{\detokenize{apply_psis}} (uses     \texttt{\detokenize{perm_eq_apply_psis}}).

\item[\href{https://github.com/LLM4Rocq/mathcomp-eulerian/blob/main/\detokenize{perm_seq_bridge.v}\#L213}{\texttt{\detokenize{apply_psis_size_eq}}}] \emph{Lemma, line 213}. \textsf{[aux]} \\
  \texttt{\detokenize{apply_psis_size_eq}} -- \texttt{\detokenize{apply_psis}} preserves size: if \texttt{\detokenize{size w =     n}} then \texttt{\detokenize{size (apply_psis ss w) = n}}.

\item[\href{https://github.com/LLM4Rocq/mathcomp-eulerian/blob/main/\detokenize{perm_seq_bridge.v}\#L223}{\texttt{\detokenize{uniq_expand_cde}}}] \emph{Lemma, line 223}. \textsf{[aux]} \\
  \texttt{\detokenize{uniq_expand_cde}} -- \texttt{\detokenize{expand_cde}} produces a list of pairwise     distinct boolean vectors (key for M-class injectivity).

\item[\href{https://github.com/LLM4Rocq/mathcomp-eulerian/blob/main/\detokenize{perm_seq_bridge.v}\#L243}{\texttt{\detokenize{nil_in_powerset_internal}}}] \emph{Lemma, line 243}. \textsf{[aux]} \\
  \texttt{\detokenize{nil_in_powerset_internal}} -- the empty list is always a member     of \texttt{\detokenize{powerset_internal w}} (the M-class always contains \texttt{\detokenize{w}}     itself).

\item[\href{https://github.com/LLM4Rocq/mathcomp-eulerian/blob/main/\detokenize{perm_seq_bridge.v}\#L256}{\texttt{\detokenize{char_mono_in_expand_cde}}}] \emph{Lemma, line 256}. \textsf{[aux]} \\
  \texttt{\detokenize{char_mono_in_expand_cde}} -- the descent pattern of \texttt{\detokenize{w}} itself     appears in \texttt{\detokenize{expand_cde (phi_w w)}} (corresponding to \texttt{\detokenize{ss = nil}}     in the \texttt{\detokenize{fact3}} enumeration).

\item[\href{https://github.com/LLM4Rocq/mathcomp-eulerian/blob/main/\detokenize{perm_seq_bridge.v}\#L276}{\texttt{\detokenize{find_ss}}}] \emph{Definition, line 276}. \textsf{[aux]} \\
  \texttt{\detokenize{find_ss w bv}} -- search for an \texttt{\detokenize{ss}} in \texttt{\detokenize{powerset_internal w}}     such that \texttt{\detokenize{apply_psis ss w}} has descent pattern \texttt{\detokenize{bv}}; returns the     first match (or \texttt{\detokenize{::}} if none).

\item[\href{https://github.com/LLM4Rocq/mathcomp-eulerian/blob/main/\detokenize{perm_seq_bridge.v}\#L282}{\texttt{\detokenize{find_ss_spec}}}] \emph{Lemma, line 282}. \textsf{[aux]} \\
  \texttt{\detokenize{find_ss_spec}} -- whenever \texttt{\detokenize{bv}} is a descent pattern realized in     the M-class of \texttt{\detokenize{w}}, \texttt{\detokenize{find_ss w bv}} returns a valid witness:     in \texttt{\detokenize{powerset_internal w}} and yielding the requested pattern.

\item[\href{https://github.com/LLM4Rocq/mathcomp-eulerian/blob/main/\detokenize{perm_seq_bridge.v}\#L312}{\texttt{\detokenize{class_map}}}] \emph{Definition, line 312}. \textsf{[aux]} \\
  \texttt{\detokenize{class_map bv sigma}} -- map a permutation \texttt{\detokenize{sigma}} to the M-class     member with descent pattern \texttt{\detokenize{bv}}, packaged back as a perm. The     core injection used in the proof of \texttt{\detokenize{omega_proper_beta_lt}}.

\item[\href{https://github.com/LLM4Rocq/mathcomp-eulerian/blob/main/\detokenize{perm_seq_bridge.v}\#L322}{\texttt{\detokenize{perm_to_seq_class_map}}}] \emph{Lemma, line 322}. \textsf{[aux]} \\
  \texttt{\detokenize{perm_to_seq_class_map}} -- \texttt{\detokenize{perm_to_seq}} of \texttt{\detokenize{class_map bv sigma}}     is exactly \texttt{\detokenize{apply_psis (find_ss ... bv) (perm_to_seq sigma)}}.

\item[\href{https://github.com/LLM4Rocq/mathcomp-eulerian/blob/main/\detokenize{perm_seq_bridge.v}\#L334}{\texttt{\detokenize{omega_seq_mem_eq}}}] \emph{Lemma, line 334}. \textsf{[aux]} \\
  \texttt{\detokenize{omega_seq_mem_eq}} -- the \texttt{\detokenize{psi_cdindex}} \texttt{\detokenize{omega_seq}} and the local     \texttt{\detokenize{omega_seq_local}} (in \texttt{\detokenize{beta_bridge}}) agree on memberships     (definitionally identical).

\item[\href{https://github.com/LLM4Rocq/mathcomp-eulerian/blob/main/\detokenize{perm_seq_bridge.v}\#L341}{\texttt{\detokenize{omega_set_seq_bridge_bounded}}}] \emph{Lemma, line 341}. \textsf{[aux]} \\
  \texttt{\detokenize{omega_set_seq_bridge_bounded}} -- bridge in the bounded form: for     \texttt{\detokenize{k < m}}, \texttt{k \textbackslash\{\}in omega\_seq (set\_to\_seq D)} iff \texttt{Ordinal Hkm \textbackslash\{\}in     omega\_set D}.

\item[\href{https://github.com/LLM4Rocq/mathcomp-eulerian/blob/main/\detokenize{perm_seq_bridge.v}\#L352}{\texttt{\detokenize{omega_seq_subset_bounded}}}] \emph{Lemma, line 352}. \textsf{[aux]} \\
  \texttt{\detokenize{omega_seq_subset_bounded}} -- a subset relation on omega-sets     transports to membership at the seq level for any list of     bounded indices.

\item[\href{https://github.com/LLM4Rocq/mathcomp-eulerian/blob/main/\detokenize{perm_seq_bridge.v}\#L373}{\texttt{\detokenize{index_lt_sorted}}}] \emph{Lemma, line 373}. \textsf{[aux]} \\
  \texttt{\detokenize{index_lt_sorted}} -- in a uniq sorted list, order on values     coincides with order on positions: \texttt{\detokenize{x < y}} iff \texttt{\detokenize{index x s <     index y s}}.

\item[\href{https://github.com/LLM4Rocq/mathcomp-eulerian/blob/main/\detokenize{perm_seq_bridge.v}\#L403}{\texttt{\detokenize{S_w_seq_all_lt}}}] \emph{Lemma, line 403}. \textsf{[aux]} \\
  \texttt{\detokenize{S_w_seq_all_lt}} -- elements of \texttt{\detokenize{S_w_seq w}} (the support indices     of the omega map) are bounded by \texttt{\detokenize{(size w).-2}}.

\item[\href{https://github.com/LLM4Rocq/mathcomp-eulerian/blob/main/\detokenize{perm_seq_bridge.v}\#L445}{\texttt{\detokenize{omega_proper_beta_lt}}}] \emph{Lemma, line 445}. \textsf{[Ch 8 §8.5]} \\
  \texttt{\detokenize{omega_proper_beta_lt}} -- Stanley EC1 (2nd ed.) Proposition 1.6.4     at the finset level: a strict inclusion \texttt{omega\_set D \textbackslash\{\}proper     omega\_set E} of omega-sets implies the strict inequality     \texttt{\detokenize{beta D < beta E}} of descent counts. Headline result of this     file; proof injects \{sigma | descent D\} into \{tau | descent E\}     via the M-class \texttt{\detokenize{class_map}} and exhibits a strict witness via     \texttt{\detokenize{strict_witness_exists}}.

\end{description}

% --------------------------------------------
\section*{\texttt{psi\_cdindex\_core.v}}
\addcontentsline{toc}{section}{\texttt{psi\_cdindex\_core.v}}

\begin{description}
\item[\href{https://github.com/LLM4Rocq/mathcomp-eulerian/blob/main/\detokenize{psi_cdindex_core.v}\#L20}{\texttt{\detokenize{window_size_apply_psis}}}] \emph{Lemma, line 20}. \textsf{[aux]} \\
  \texttt{\detokenize{apply_psis}} preserves \texttt{\detokenize{window_size}} at every index.

\item[\href{https://github.com/LLM4Rocq/mathcomp-eulerian/blob/main/\detokenize{psi_cdindex_core.v}\#L29}{\texttt{\detokenize{has_left_child_apply_psis}}}] \emph{Lemma, line 29}. \textsf{[aux]} \\
  \texttt{\detokenize{apply_psis}} preserves \texttt{\detokenize{has_left_child}} at every index.

\item[\href{https://github.com/LLM4Rocq/mathcomp-eulerian/blob/main/\detokenize{psi_cdindex_core.v}\#L38}{\texttt{\detokenize{is_internal_apply_psis}}}] \emph{Lemma, line 38}. \textsf{[aux]} \\
  \texttt{\detokenize{apply_psis}} preserves \texttt{\detokenize{is_internal}}.

\item[\href{https://github.com/LLM4Rocq/mathcomp-eulerian/blob/main/\detokenize{psi_cdindex_core.v}\#L47}{\texttt{\detokenize{internal_vertices_apply_psis}}}] \emph{Lemma, line 47}. \textsf{[aux]} \\
  \texttt{\detokenize{apply_psis}} leaves the list of internal vertices invariant.

\item[\href{https://github.com/LLM4Rocq/mathcomp-eulerian/blob/main/\detokenize{psi_cdindex_core.v}\#L58}{\texttt{\detokenize{phi_w_apply_psis}}}] \emph{Lemma, line 58}. \textsf{[aux]} \\
  Key invariance: the cd-index \texttt{\detokenize{phi_w}} is constant on the M-class     generated by repeated \texttt{\detokenize{psi}} applications.

\item[\href{https://github.com/LLM4Rocq/mathcomp-eulerian/blob/main/\detokenize{psi_cdindex_core.v}\#L72}{\texttt{\detokenize{phi_w_as_map}}}] \emph{Lemma, line 72}. \textsf{[aux]} \\
  \texttt{\detokenize{phi_w w}} expressed as a map over \texttt{\detokenize{internal_vertices w}}: each internal     vertex contributes \texttt{\detokenize{D_letter}} (if it has a left child) or \texttt{\detokenize{C_letter}}.

\item[\href{https://github.com/LLM4Rocq/mathcomp-eulerian/blob/main/\detokenize{psi_cdindex_core.v}\#L89}{\texttt{\detokenize{leq_seqb_total}}}] \emph{Lemma, line 89}. \textsf{[aux]} \\
  Lex order on \texttt{\detokenize{seq bool}} is total.

\item[\href{https://github.com/LLM4Rocq/mathcomp-eulerian/blob/main/\detokenize{psi_cdindex_core.v}\#L97}{\texttt{\detokenize{leq_seqb_trans}}}] \emph{Lemma, line 97}. \textsf{[aux]} \\
  Lex order on \texttt{\detokenize{seq bool}} is transitive.

\item[\href{https://github.com/LLM4Rocq/mathcomp-eulerian/blob/main/\detokenize{psi_cdindex_core.v}\#L105}{\texttt{\detokenize{leq_seqb_anti}}}] \emph{Lemma, line 105}. \textsf{[aux]} \\
  Lex order on \texttt{\detokenize{seq bool}} is antisymmetric.

\item[\href{https://github.com/LLM4Rocq/mathcomp-eulerian/blob/main/\detokenize{psi_cdindex_core.v}\#L116}{\texttt{\detokenize{sort_perm_eq_leq_seqb}}}] \emph{Lemma, line 116}. \textsf{[aux]} \\
  Sorting by \texttt{\detokenize{leq_seqb}} is a \texttt{\detokenize{perm_eq}}-invariant canonical form on     multisets of bit-strings, used to compare M-class char\_monos to     \texttt{\detokenize{expand_cde}} outputs.

\item[\href{https://github.com/LLM4Rocq/mathcomp-eulerian/blob/main/\detokenize{psi_cdindex_core.v}\#L131}{\texttt{\detokenize{descent_psi_effect}}}] \emph{Lemma, line 131}. \textsf{[aux]} \\
  Bit-level effect of \texttt{\detokenize{psi v}} on the descent sequence at any position \texttt{\detokenize{k}}:     R-vertex toggles bit \texttt{\detokenize{v}}; LR-vertex swaps bits \texttt{\detokenize{v.-1}} and \texttt{\detokenize{v}}.     Underlies the char\_mono-recovery proofs.

\item[\href{https://github.com/LLM4Rocq/mathcomp-eulerian/blob/main/\detokenize{psi_cdindex_core.v}\#L178}{\texttt{\detokenize{nth_char_mono}}}] \emph{Lemma, line 178}. \textsf{[aux]} \\
  \texttt{\detokenize{char_mono w}} indexed by \texttt{\detokenize{k}} returns the descent bit at position \texttt{\detokenize{k}}.

\item[\href{https://github.com/LLM4Rocq/mathcomp-eulerian/blob/main/\detokenize{psi_cdindex_core.v}\#L188}{\texttt{\detokenize{size_char_mono}}}] \emph{Lemma, line 188}. \textsf{[aux]} \\
  \texttt{\detokenize{char_mono}} has length \texttt{\detokenize{(size w).-1}} (one bit per descent position).

\item[\href{https://github.com/LLM4Rocq/mathcomp-eulerian/blob/main/\detokenize{psi_cdindex_core.v}\#L194}{\texttt{\detokenize{is_internal_lt}}}] \emph{Lemma, line 194}. \textsf{[aux]} \\
  Internal vertices live strictly inside the descent index range,     i.e. \texttt{\detokenize{v < (size w).-1}}.

\item[\href{https://github.com/LLM4Rocq/mathcomp-eulerian/blob/main/\detokenize{psi_cdindex_core.v}\#L207}{\texttt{\detokenize{char_mono_psi_effect}}}] \emph{Lemma, line 207}. \textsf{[aux]} \\
  Bit-level effect of \texttt{\detokenize{psi v}} on \texttt{\detokenize{char_mono}}:     a C-vertex (no left child) toggles bit at position \texttt{\detokenize{v}};     a D-vertex (has left child) swaps bits at positions \texttt{\detokenize{v.-1}} and \texttt{\detokenize{v}}.

\item[\href{https://github.com/LLM4Rocq/mathcomp-eulerian/blob/main/\detokenize{psi_cdindex_core.v}\#L235}{\texttt{\detokenize{index_rcons_eq_size}}}] \emph{Lemma, line 235}. \textsf{[aux]} \\
  Helper: if \texttt{\detokenize{index x (rcons s y) = size s}} then \texttt{x \textbackslash\{\}notin s}.

\item[\href{https://github.com/LLM4Rocq/mathcomp-eulerian/blob/main/\detokenize{psi_cdindex_core.v}\#L248}{\texttt{\detokenize{mm_pos_rcons_lt}}}] \emph{Lemma, line 248}. \textsf{[aux]} \\
  Helper: \texttt{\detokenize{mm_pos}} is never at the trailing position of a non-empty     prefix, because that would force the last element to be both min and     max, contradicting the presence of any other element.

\item[\href{https://github.com/LLM4Rocq/mathcomp-eulerian/blob/main/\detokenize{psi_cdindex_core.v}\#L344}{\texttt{\detokenize{has_left_child_t_last_valid}}}] \emph{Lemma, line 344}. \textsf{[aux]} \\
  Tree-level: in any valid mmtree, the last in-order position has no     left child.  Proved by structural induction on the tree.

\item[\href{https://github.com/LLM4Rocq/mathcomp-eulerian/blob/main/\detokenize{psi_cdindex_core.v}\#L379}{\texttt{\detokenize{has_left_child_last}}}] \emph{Lemma, line 379}. \textsf{[aux]} \\
  Seq-level corollary: the last vertex never has a left child.

\item[\href{https://github.com/LLM4Rocq/mathcomp-eulerian/blob/main/\detokenize{psi_cdindex_core.v}\#L393}{\texttt{\detokenize{has_left_child_last_fuel}}}] \emph{Lemma, line 393}. \textsf{[aux]} \\
  Fuel-bounded variant of \texttt{\detokenize{has_left_child_last}}.

\item[\href{https://github.com/LLM4Rocq/mathcomp-eulerian/blob/main/\detokenize{psi_cdindex_core.v}\#L425}{\texttt{\detokenize{has_left_child_is_internal}}}] \emph{Lemma, line 425}. \textsf{[aux]} \\
  \texttt{\detokenize{has_left_child}} implies \texttt{\detokenize{is_internal}} (window size > 1).  Tree     induction on the mmtree of \texttt{\detokenize{w}}.

\item[\href{https://github.com/LLM4Rocq/mathcomp-eulerian/blob/main/\detokenize{psi_cdindex_core.v}\#L470}{\texttt{\detokenize{size_powerset_of}}}] \emph{Lemma, line 470}. \textsf{[aux]} \\
  Generic powerset-by-fold has size \texttt{\detokenize{2 ^ size ivs}}.

\item[\href{https://github.com/LLM4Rocq/mathcomp-eulerian/blob/main/\detokenize{psi_cdindex_core.v}\#L488}{\texttt{\detokenize{size_powerset_internal}}}] \emph{Lemma, line 488}. \textsf{[aux]} \\
  \texttt{\detokenize{powerset_internal w}} has cardinality \texttt{\detokenize{2 ^ |internal_vertices w|}}.

\item[\href{https://github.com/LLM4Rocq/mathcomp-eulerian/blob/main/\detokenize{psi_cdindex_core.v}\#L493}{\texttt{\detokenize{uniq_powerset_of}}}] \emph{Lemma, line 493}. \textsf{[aux]} \\
  Generic powerset-by-fold is uniq when its input is uniq.

\item[\href{https://github.com/LLM4Rocq/mathcomp-eulerian/blob/main/\detokenize{psi_cdindex_core.v}\#L529}{\texttt{\detokenize{uniq_powerset_internal}}}] \emph{Lemma, line 529}. \textsf{[aux]} \\
  \texttt{\detokenize{powerset_internal w}} is duplicate-free.

\item[\href{https://github.com/LLM4Rocq/mathcomp-eulerian/blob/main/\detokenize{psi_cdindex_core.v}\#L533}{\texttt{\detokenize{subset_powerset_of}}}] \emph{Lemma, line 533}. \textsf{[aux]} \\
  Each element of the generic powerset is a subset of the input.

\item[\href{https://github.com/LLM4Rocq/mathcomp-eulerian/blob/main/\detokenize{psi_cdindex_core.v}\#L563}{\texttt{\detokenize{subset_powerset_internal}}}] \emph{Lemma, line 563}. \textsf{[aux]} \\
  Each element of \texttt{\detokenize{powerset_internal w}} is a subset of     \texttt{\detokenize{internal_vertices w}}.

\item[\href{https://github.com/LLM4Rocq/mathcomp-eulerian/blob/main/\detokenize{psi_cdindex_core.v}\#L568}{\texttt{\detokenize{subseq_powerset_of}}}] \emph{Lemma, line 568}. \textsf{[aux]} \\
  Each element of the generic powerset is a subseq of the input.

\item[\href{https://github.com/LLM4Rocq/mathcomp-eulerian/blob/main/\detokenize{psi_cdindex_core.v}\#L597}{\texttt{\detokenize{subseq_powerset_internal}}}] \emph{Lemma, line 597}. \textsf{[aux]} \\
  Each element of \texttt{\detokenize{powerset_internal w}} is a subseq of     \texttt{\detokenize{internal_vertices w}}, hence is uniq when \texttt{\detokenize{internal_vertices w}} is.

\item[\href{https://github.com/LLM4Rocq/mathcomp-eulerian/blob/main/\detokenize{psi_cdindex_core.v}\#L604}{\texttt{\detokenize{char_mono_apply_psis_C_bit}}}] \emph{Lemma, line 604}. \textsf{[aux]} \\
  Bit-level recovery for a C-vertex: bit \texttt{\detokenize{v}} of     \texttt{\detokenize{char_mono (apply_psis ss w)}} equals the original descent bit XOR the     indicator of \texttt{v \textbackslash\{\}in ss}.  Used to recover \texttt{\detokenize{ss}} from char\_monos.

\item[\href{https://github.com/LLM4Rocq/mathcomp-eulerian/blob/main/\detokenize{psi_cdindex_core.v}\#L650}{\texttt{\detokenize{char_mono_apply_psis_D_bit_pred}}}] \emph{Lemma, line 650}. \textsf{[aux]} \\
  Bit-level recovery for a D-vertex at the predecessor position \texttt{\detokenize{v.-1}}:     the bit equals the original descent bit at \texttt{\detokenize{v}} (if \texttt{v \textbackslash\{\}in ss}) or at     \texttt{\detokenize{v.-1}} (otherwise) -- encoding the swap behaviour.

\item[\href{https://github.com/LLM4Rocq/mathcomp-eulerian/blob/main/\detokenize{psi_cdindex_core.v}\#L723}{\texttt{\detokenize{char_mono_apply_psis_D_bit_self}}}] \emph{Lemma, line 723}. \textsf{[aux]} \\
  Bit-level recovery for a D-vertex at its self position \texttt{\detokenize{v}}: the bit     equals the original descent bit at \texttt{\detokenize{v.-1}} (if \texttt{v \textbackslash\{\}in ss}) or at \texttt{\detokenize{v}}     (otherwise) -- companion to \texttt{\detokenize{_D_bit_pred}}.

\item[\href{https://github.com/LLM4Rocq/mathcomp-eulerian/blob/main/\detokenize{psi_cdindex_core.v}\#L791}{\texttt{\detokenize{size_expand_cde}}}] \emph{Lemma, line 791}. \textsf{[aux]} \\
  \texttt{\detokenize{expand_cde m}} has \texttt{\detokenize{2 ^ k}} elements, where \texttt{\detokenize{k}} is the number of     non-\texttt{\detokenize{E}} letters in \texttt{\detokenize{m}}.

\item[\href{https://github.com/LLM4Rocq/mathcomp-eulerian/blob/main/\detokenize{psi_cdindex_core.v}\#L807}{\texttt{\detokenize{size_expand_cde_phi_w}}}] \emph{Lemma, line 807}. \textsf{[aux]} \\
  \texttt{\detokenize{expand_cde (phi_w w)}} has cardinality \texttt{\detokenize{2 ^ |internal_vertices w|}},     matching the size of \texttt{\detokenize{powerset_internal w}} for the M-class.

\item[\href{https://github.com/LLM4Rocq/mathcomp-eulerian/blob/main/\detokenize{psi_cdindex_core.v}\#L824}{\texttt{\detokenize{check_fact3}}}] \emph{Definition, line 824}. \textsf{[aux]} \\
  \texttt{\detokenize{check_fact3 w}} is the boolean Fact \#3 statement: the multiset of     M-class char\_monos (sorted) equals the multiset \texttt{\detokenize{expand_cde (phi_w w)}}     (sorted).

\item[\href{https://github.com/LLM4Rocq/mathcomp-eulerian/blob/main/\detokenize{psi_cdindex_core.v}\#L833}{\texttt{\detokenize{check_fact3P}}}] \emph{Lemma, line 833}. \textsf{[aux]} \\
  Reflection between \texttt{\detokenize{check_fact3 w}} and the equality of the sorted     multisets it encodes.

\item[\href{https://github.com/LLM4Rocq/mathcomp-eulerian/blob/main/\detokenize{psi_cdindex_core.v}\#L845}{\texttt{\detokenize{check_fact3_size1}}}] \emph{Lemma, line 845}. \textsf{[aux]} \\
  Base case for \texttt{\detokenize{check_fact3}}: sequences of size at most 1 satisfy it     trivially (no internal vertices, empty cd-index).

\end{description}

% --------------------------------------------
\section*{\texttt{psi\_cdindex\_defs.v}}
\addcontentsline{toc}{section}{\texttt{psi\_cdindex\_defs.v}}

\begin{description}
\item[\href{https://github.com/LLM4Rocq/mathcomp-eulerian/blob/main/\detokenize{psi_cdindex_defs.v}\#L26}{\texttt{\detokenize{is_internal}}}] \emph{Definition, line 26}. \textsf{[aux]} \\
  \texttt{\detokenize{is_internal i w}} holds iff vertex \texttt{\detokenize{i}} is non-endpoint in the     min-max tree of \texttt{\detokenize{w}} (its window has size > 1).

\item[\href{https://github.com/LLM4Rocq/mathcomp-eulerian/blob/main/\detokenize{psi_cdindex_defs.v}\#L31}{\texttt{\detokenize{apply_psis}}}] \emph{Definition, line 31}. \textsf{[aux]} \\
  \texttt{\detokenize{apply_psis ops w}} applies the sequence of \texttt{\detokenize{psi}} operators left-to-right,     threading \texttt{\detokenize{psi i}} for each \texttt{\detokenize{i}} in \texttt{\detokenize{ops}} through \texttt{\detokenize{w}}.

\item[\href{https://github.com/LLM4Rocq/mathcomp-eulerian/blob/main/\detokenize{psi_cdindex_defs.v}\#L36}{\texttt{\detokenize{char_mono}}}] \emph{Definition, line 36}. \textsf{[aux]} \\
  \texttt{\detokenize{char_mono w}} is the descent bit-string of \texttt{\detokenize{w}} (length \texttt{\detokenize{(size w).-1}}):     bit \texttt{\detokenize{k}} is \texttt{\detokenize{true}} iff there is a descent at position \texttt{\detokenize{k}}.

\item[\href{https://github.com/LLM4Rocq/mathcomp-eulerian/blob/main/\detokenize{psi_cdindex_defs.v}\#L41}{\texttt{\detokenize{cde}}}] \emph{Inductive, line 41}. \textsf{[aux]} \\
  The cd-alphabet for classifying vertices: \texttt{\detokenize{C_letter}} (right child only),     \texttt{\detokenize{D_letter}} (both children), \texttt{\detokenize{E_letter}} (endpoint).

\item[\href{https://github.com/LLM4Rocq/mathcomp-eulerian/blob/main/\detokenize{psi_cdindex_defs.v}\#L46}{\texttt{\detokenize{classify_vertex_cde}}}] \emph{Definition, line 46}. \textsf{[aux]} \\
  \texttt{\detokenize{classify_vertex_cde i w}} returns the cd-letter for vertex \texttt{\detokenize{i}} of     \texttt{\detokenize{w}}: \texttt{\detokenize{E_letter}} if endpoint, \texttt{\detokenize{D_letter}} if it has a left child, else     \texttt{\detokenize{C_letter}}.

\item[\href{https://github.com/LLM4Rocq/mathcomp-eulerian/blob/main/\detokenize{psi_cdindex_defs.v}\#L53}{\texttt{\detokenize{phi'_w}}}] \emph{Definition, line 53}. \textsf{[aux]} \\
  \texttt{\detokenize{phi'_w w}} is the full cde-classification string, one letter per vertex     of \texttt{\detokenize{w}} (length \texttt{\detokenize{size w}}).

\item[\href{https://github.com/LLM4Rocq/mathcomp-eulerian/blob/main/\detokenize{psi_cdindex_defs.v}\#L58}{\texttt{\detokenize{phi_w}}}] \emph{Definition, line 58}. \textsf{[aux]} \\
  \texttt{\detokenize{phi_w w}} is the cd-index of \texttt{\detokenize{w}}: \texttt{\detokenize{phi'_w w}} with all \texttt{\detokenize{E_letter}}     (endpoints) stripped.

\item[\href{https://github.com/LLM4Rocq/mathcomp-eulerian/blob/main/\detokenize{psi_cdindex_defs.v}\#L63}{\texttt{\detokenize{internal_vertices}}}] \emph{Definition, line 63}. \textsf{[aux]} \\
  \texttt{\detokenize{internal_vertices w}} enumerates internal-vertex positions of \texttt{\detokenize{w}}     in ascending order.

\item[\href{https://github.com/LLM4Rocq/mathcomp-eulerian/blob/main/\detokenize{psi_cdindex_defs.v}\#L69}{\texttt{\detokenize{expand_cde}}}] \emph{Fixpoint, line 69}. \textsf{[aux]} \\
  \texttt{\detokenize{expand_cde letters}} expands a cd-word to the multiset of bit-strings:     \texttt{\detokenize{c = a+b}} yields one bit, \texttt{\detokenize{d = ab+ba}} yields two bits in either order;     \texttt{\detokenize{E_letter}} contributes nothing.

\item[\href{https://github.com/LLM4Rocq/mathcomp-eulerian/blob/main/\detokenize{psi_cdindex_defs.v}\#L84}{\texttt{\detokenize{powerset_internal}}}] \emph{Definition, line 84}. \textsf{[aux]} \\
  \texttt{\detokenize{powerset_internal w}} enumerates all subsequences of \texttt{\detokenize{internal_vertices w}},     used to index the M-equivalence class of \texttt{\detokenize{w}}.

\item[\href{https://github.com/LLM4Rocq/mathcomp-eulerian/blob/main/\detokenize{psi_cdindex_defs.v}\#L90}{\texttt{\detokenize{leq_seqb}}}] \emph{Fixpoint, line 90}. \textsf{[aux]} \\
  \texttt{\detokenize{leq_seqb}} is the lexicographic order on \texttt{\detokenize{seq bool}}, used to canonicalize     the multiset of expanded cd-monomials by sorting.

\item[\href{https://github.com/LLM4Rocq/mathcomp-eulerian/blob/main/\detokenize{psi_cdindex_defs.v}\#L102}{\texttt{\detokenize{apply_psis_nil}}}] \emph{Lemma, line 102}. \textsf{[aux]} \\
  Empty operator list acts as identity.

\item[\href{https://github.com/LLM4Rocq/mathcomp-eulerian/blob/main/\detokenize{psi_cdindex_defs.v}\#L107}{\texttt{\detokenize{size_apply_psis}}}] \emph{Lemma, line 107}. \textsf{[aux]} \\
  \texttt{\detokenize{apply_psis}} preserves the length of \texttt{\detokenize{w}}.

\item[\href{https://github.com/LLM4Rocq/mathcomp-eulerian/blob/main/\detokenize{psi_cdindex_defs.v}\#L114}{\texttt{\detokenize{uniq_apply_psis}}}] \emph{Lemma, line 114}. \textsf{[aux]} \\
  \texttt{\detokenize{apply_psis}} preserves uniqueness of \texttt{\detokenize{w}}.

\item[\href{https://github.com/LLM4Rocq/mathcomp-eulerian/blob/main/\detokenize{psi_cdindex_defs.v}\#L122}{\texttt{\detokenize{perm_eq_apply_psis}}}] \emph{Lemma, line 122}. \textsf{[aux]} \\
  \texttt{\detokenize{apply_psis ops w}} is a permutation of \texttt{\detokenize{w}} (each \texttt{\detokenize{psi}} is a     transposition).

\item[\href{https://github.com/LLM4Rocq/mathcomp-eulerian/blob/main/\detokenize{psi_cdindex_defs.v}\#L129}{\texttt{\detokenize{apply_psis_cat}}}] \emph{Lemma, line 129}. \textsf{[aux]} \\
  Concatenation of operator lists corresponds to function composition.

\end{description}

% --------------------------------------------
\section*{\texttt{psi\_cdindex\_support.v}}
\addcontentsline{toc}{section}{\texttt{psi\_cdindex\_support.v}}

\begin{description}
\item[\href{https://github.com/LLM4Rocq/mathcomp-eulerian/blob/main/\detokenize{psi_cdindex_support.v}\#L41}{\texttt{\detokenize{S_w_seq_bound}}}] \emph{Lemma, line 41}. \textsf{[aux]} \\
  Every element of \texttt{\detokenize{S_w_seq w}} is bounded: \texttt{\detokenize{k < (size w).-2}}.  Holds     because the last vertex of \texttt{\detokenize{w}} is always an endpoint, hence cannot     contribute a \texttt{\detokenize{D}} letter.

\item[\href{https://github.com/LLM4Rocq/mathcomp-eulerian/blob/main/\detokenize{psi_cdindex_support.v}\#L71}{\texttt{\detokenize{classify_vertex_cde_psi}}}] \emph{Lemma, line 71}. \textsf{[aux]} \\
  Vertex classification is invariant under \texttt{\detokenize{psi j}}.

\item[\href{https://github.com/LLM4Rocq/mathcomp-eulerian/blob/main/\detokenize{psi_cdindex_support.v}\#L89}{\texttt{\detokenize{mm_pos_lt_pred}}}] \emph{Lemma, line 89}. \textsf{[aux]} \\
  \texttt{\detokenize{mm_pos s < (size s).-1}} for uniq \texttt{\detokenize{s}} of size at least 2.  Ensures     the root of the min-max tree is an internal vertex.

\item[\href{https://github.com/LLM4Rocq/mathcomp-eulerian/blob/main/\detokenize{psi_cdindex_support.v}\#L156}{\texttt{\detokenize{classify_vertex_left}}}] \emph{Lemma, line 156}. \textsf{[aux]} \\
  Classification splits across \texttt{\detokenize{mm_pos}}: for \texttt{\detokenize{i < j = mm_pos s}}, use     classification on the left subtree \texttt{\detokenize{take j s}}.

\item[\href{https://github.com/LLM4Rocq/mathcomp-eulerian/blob/main/\detokenize{psi_cdindex_support.v}\#L178}{\texttt{\detokenize{classify_vertex_right}}}] \emph{Lemma, line 178}. \textsf{[aux]} \\
  Classification on the right side of \texttt{\detokenize{mm_pos}}: for \texttt{\detokenize{i > j}}, shift     index by \texttt{\detokenize{j.+1}} and classify in the right subtree \texttt{\detokenize{drop j.+1 s}}.

\item[\href{https://github.com/LLM4Rocq/mathcomp-eulerian/blob/main/\detokenize{psi_cdindex_support.v}\#L207}{\texttt{\detokenize{classify_vertex_mm}}}] \emph{Lemma, line 207}. \textsf{[aux]} \\
  Classification at the root \texttt{\detokenize{j = mm_pos s}} of a uniq seq of size at     least 2: \texttt{\detokenize{D_letter}} if the root has a left subtree (\texttt{\detokenize{j > 0}}), else     \texttt{\detokenize{C_letter}}.

\item[\href{https://github.com/LLM4Rocq/mathcomp-eulerian/blob/main/\detokenize{psi_cdindex_support.v}\#L232}{\texttt{\detokenize{phi_w_cons_mm0}}}] \emph{Lemma, line 232}. \textsf{[aux]} \\
  Cons decomposition of \texttt{\detokenize{phi_w}} when \texttt{\detokenize{mm_pos = 0}} (root has no left     subtree): the root contributes a \texttt{\detokenize{C_letter}} and recursion peels into     the tail.

\item[\href{https://github.com/LLM4Rocq/mathcomp-eulerian/blob/main/\detokenize{psi_cdindex_support.v}\#L260}{\texttt{\detokenize{phi_w_decomp_mm}}}] \emph{Lemma, line 260}. \textsf{[aux]} \\
  Decomposition of \texttt{\detokenize{phi_w}} at \texttt{\detokenize{mm_pos > 0}}: \texttt{\detokenize{phi_w s}} equals the cd-index     of the left subtree concatenated with \texttt{\detokenize{D_letter}} then the cd-index of     the right subtree.

\item[\href{https://github.com/LLM4Rocq/mathcomp-eulerian/blob/main/\detokenize{psi_cdindex_support.v}\#L304}{\texttt{\detokenize{D_offsets_cons_C}}}] \emph{Lemma, line 304}. \textsf{[aux]} \\
  Cons decomposition of \texttt{\detokenize{D_offsets}} over a \texttt{\detokenize{C}} head: every offset     shifts by 1 bit.

\item[\href{https://github.com/LLM4Rocq/mathcomp-eulerian/blob/main/\detokenize{psi_cdindex_support.v}\#L323}{\texttt{\detokenize{D_offsets_cat_D}}}] \emph{Lemma, line 323}. \textsf{[aux]} \\
  Concatenation decomposition of \texttt{\detokenize{D_offsets}} across a \texttt{\detokenize{D}} inserted     between \texttt{\detokenize{m1}} and \texttt{\detokenize{m2}}: left offsets, then the \texttt{\detokenize{D}} at offset     \texttt{\detokenize{cde_total_width m1}}, then right offsets shifted by     \texttt{\detokenize{cde_total_width m1 + 2}}.

\item[\href{https://github.com/LLM4Rocq/mathcomp-eulerian/blob/main/\detokenize{psi_cdindex_support.v}\#L396}{\texttt{\detokenize{S_w_seq_decomp_mm}}}] \emph{Lemma, line 396}. \textsf{[aux]} \\
  Decomposition of \texttt{\detokenize{S_w_seq}} at \texttt{\detokenize{mm_pos > 0}}, parallel to     \texttt{\detokenize{phi_w_decomp_mm}}: left positions, then \texttt{\detokenize{j.-1}} (the D-position     contributed by the root), then right positions shifted by \texttt{\detokenize{j.+1}}.

\item[\href{https://github.com/LLM4Rocq/mathcomp-eulerian/blob/main/\detokenize{psi_cdindex_support.v}\#L481}{\texttt{\detokenize{cde_total_width_phi_w_all}}}] \emph{Lemma, line 481}. \textsf{[aux]} \\
  Total bit-width of \texttt{\detokenize{phi_w w}} equals \texttt{\detokenize{(size w).-1}} for any uniq \texttt{\detokenize{w}}     (proved by structural induction on size with \texttt{\detokenize{mm_pos}} decomposition).     Captures that \texttt{\detokenize{expand_cde (phi_w w)}} elements are descent-bit strings.

\item[\href{https://github.com/LLM4Rocq/mathcomp-eulerian/blob/main/\detokenize{psi_cdindex_support.v}\#L566}{\texttt{\detokenize{cde_total_width_phi_w}}}] \emph{Lemma, line 566}. \textsf{[aux]} \\
  Sized variant of \texttt{\detokenize{cde_total_width_phi_w_all}}; states the bit-width     identity under the usable hypothesis \texttt{\detokenize{2 <= size w}}.

\item[\href{https://github.com/LLM4Rocq/mathcomp-eulerian/blob/main/\detokenize{psi_cdindex_support.v}\#L581}{\texttt{\detokenize{D_offsets_phi_w_eq_S_w_seq}}}] \emph{Lemma, line 581}. \textsf{[aux]} \\
  Headline structural identity: the \texttt{\detokenize{D}}-offsets of \texttt{\detokenize{phi_w w}} coincide     with \texttt{\detokenize{S_w_seq w}}.  Proved by structural induction on size with     \texttt{\detokenize{mm_pos}} decomposition.

\item[\href{https://github.com/LLM4Rocq/mathcomp-eulerian/blob/main/\detokenize{psi_cdindex_support.v}\#L697}{\texttt{\detokenize{uniq_expand_cde}}}] \emph{Lemma, line 697}. \textsf{[aux]} \\
  \texttt{\detokenize{expand_cde m}} enumerates each cd-monomial expansion exactly once.

\item[\href{https://github.com/LLM4Rocq/mathcomp-eulerian/blob/main/\detokenize{psi_cdindex_support.v}\#L712}{\texttt{\detokenize{perm_eq_from_subset}}}] \emph{Lemma, line 712}. \textsf{[aux]} \\
  Two uniq sequences of the same size with one a subset of the other     are permutations of each other.

\item[\href{https://github.com/LLM4Rocq/mathcomp-eulerian/blob/main/\detokenize{psi_cdindex_support.v}\#L730}{\texttt{\detokenize{check_fact3_of_perm_eq}}}] \emph{Lemma, line 730}. \textsf{[aux]} \\
  \texttt{\detokenize{perm_eq}} of the M-class char\_monos with \texttt{\detokenize{expand_cde (phi_w w)}}     implies the boolean Fact \#3 statement.

\item[\href{https://github.com/LLM4Rocq/mathcomp-eulerian/blob/main/\detokenize{psi_cdindex_support.v}\#L743}{\texttt{\detokenize{D_vertex_descent_transition}}}] \emph{Lemma, line 743}. \textsf{[aux]} \\
  At a D-vertex (LR-internal, \texttt{\detokenize{has_left_child i w}}), the descent bits     at \texttt{\detokenize{i.-1}} and \texttt{\detokenize{i}} always differ -- the LR-swap behaviour of \texttt{\detokenize{psi}}.

\item[\href{https://github.com/LLM4Rocq/mathcomp-eulerian/blob/main/\detokenize{psi_cdindex_support.v}\#L759}{\texttt{\detokenize{char_mono_self_mem}}}] \emph{Lemma, line 759}. \textsf{[aux]} \\
  Self-support: \texttt{\detokenize{char_mono w}} is itself a member of \texttt{\detokenize{expand_cde (phi_w w)}},     using \texttt{\detokenize{phi_w_support_general}} applied to \texttt{\detokenize{X = char_mono w}}.

\item[\href{https://github.com/LLM4Rocq/mathcomp-eulerian/blob/main/\detokenize{psi_cdindex_support.v}\#L801}{\texttt{\detokenize{char_mono_apply_psis_mem}}}] \emph{Lemma, line 801}. \textsf{[aux]} \\
  Membership: every M-class element has its \texttt{\detokenize{char_mono}} inside     \texttt{\detokenize{expand_cde (phi_w w)}}, by combining \texttt{\detokenize{phi_w_apply_psis}} and     \texttt{\detokenize{char_mono_self_mem}}.

\item[\href{https://github.com/LLM4Rocq/mathcomp-eulerian/blob/main/\detokenize{psi_cdindex_support.v}\#L815}{\texttt{\detokenize{in_internal_vertices}}}] \emph{Lemma, line 815}. \textsf{[aux]} \\
  Membership in \texttt{\detokenize{internal_vertices}} is exactly \texttt{\detokenize{is_internal}}.

\item[\href{https://github.com/LLM4Rocq/mathcomp-eulerian/blob/main/\detokenize{psi_cdindex_support.v}\#L897}{\texttt{\detokenize{uniq_map_char_mono_powerset}}}] \emph{Lemma, line 897}. \textsf{[aux]} \\
  The list of M-class char\_monos (indexed by \texttt{\detokenize{powerset_internal w}}) is     duplicate-free, by \texttt{\detokenize{char_mono_apply_psis_inj}}.

\item[\href{https://github.com/LLM4Rocq/mathcomp-eulerian/blob/main/\detokenize{psi_cdindex_support.v}\#L911}{\texttt{\detokenize{check_fact3_true}}}] \emph{Lemma, line 911}. \textsf{[aux]} \\
  Boolean Fact \#3 holds for every uniq \texttt{\detokenize{w}}: combines uniq + same-size +     membership to derive \texttt{\detokenize{perm_eq}}, then the sorted equality.

\item[\href{https://github.com/LLM4Rocq/mathcomp-eulerian/blob/main/\detokenize{psi_cdindex_support.v}\#L932}{\texttt{\detokenize{fact3}}}] \emph{Lemma, line 932}. \textsf{[aux]} \\
  Headline Fact \#3 (Stanley EC1 \textbackslash{}S\{\}1.6.3, Theorem 1.6.3): the multiset of     char\_monos over the M-class of \texttt{\detokenize{w}}, canonicalised by \texttt{\detokenize{sort leq_seqb}},     equals the canonicalised expansion \texttt{\detokenize{expand_cde (phi_w w)}}.  Direct     consequence of \texttt{\detokenize{check_fact3_true}} via \texttt{\detokenize{check_fact3P}}.

\end{description}

% --------------------------------------------
\section*{\texttt{psi\_cdindex\_support\_defs.v}}
\addcontentsline{toc}{section}{\texttt{psi\_cdindex\_support\_defs.v}}

\begin{description}
\item[\href{https://github.com/LLM4Rocq/mathcomp-eulerian/blob/main/\detokenize{psi_cdindex_support_defs.v}\#L21}{\texttt{\detokenize{cde_width}}}] \emph{Definition, line 21}. \textsf{[aux]} \\
  \texttt{\detokenize{cde_width l}} is the number of bits a single cd-letter contributes to     \texttt{\detokenize{expand_cde}}: \texttt{\detokenize{C}} contributes 1, \texttt{\detokenize{D}} contributes 2, \texttt{\detokenize{E}} contributes 0.

\item[\href{https://github.com/LLM4Rocq/mathcomp-eulerian/blob/main/\detokenize{psi_cdindex_support_defs.v}\#L26}{\texttt{\detokenize{cde_total_width}}}] \emph{Definition, line 26}. \textsf{[aux]} \\
  \texttt{\detokenize{cde_total_width m}} is the bit-length of every element of     \texttt{\detokenize{expand_cde m}}; equals the sum of \texttt{\detokenize{cde_width}} over \texttt{\detokenize{m}}.

\item[\href{https://github.com/LLM4Rocq/mathcomp-eulerian/blob/main/\detokenize{psi_cdindex_support_defs.v}\#L31}{\texttt{\detokenize{cde_offset}}}] \emph{Definition, line 31}. \textsf{[aux]} \\
  \texttt{\detokenize{cde_offset m i}} is the cumulative bit-offset where the \texttt{\detokenize{i}}-th letter     of \texttt{\detokenize{m}} starts in any element of \texttt{\detokenize{expand_cde m}}.

\item[\href{https://github.com/LLM4Rocq/mathcomp-eulerian/blob/main/\detokenize{psi_cdindex_support_defs.v}\#L37}{\texttt{\detokenize{D_offsets}}}] \emph{Definition, line 37}. \textsf{[aux]} \\
  \texttt{\detokenize{D_offsets m}} enumerates the bit positions where each \texttt{\detokenize{D}} in \texttt{\detokenize{m}}     starts; these are exactly the indices at which an expansion must     have a transition (the two bits of \texttt{\detokenize{d}} differ).

\item[\href{https://github.com/LLM4Rocq/mathcomp-eulerian/blob/main/\detokenize{psi_cdindex_support_defs.v}\#L42}{\texttt{\detokenize{has_transition}}}] \emph{Definition, line 42}. \textsf{[aux]} \\
  \texttt{\detokenize{has_transition X k}} holds iff bits \texttt{\detokenize{k}} and \texttt{\detokenize{k.+1}} of \texttt{\detokenize{X}} differ.

\item[\href{https://github.com/LLM4Rocq/mathcomp-eulerian/blob/main/\detokenize{psi_cdindex_support_defs.v}\#L47}{\texttt{\detokenize{all_D_transitions}}}] \emph{Definition, line 47}. \textsf{[aux]} \\
  \texttt{\detokenize{all_D_transitions m X}} holds iff \texttt{\detokenize{X}} has a bit transition at every     \texttt{\detokenize{D}}-offset of \texttt{\detokenize{m}}; characterizes membership in \texttt{\detokenize{expand_cde m}}.

\item[\href{https://github.com/LLM4Rocq/mathcomp-eulerian/blob/main/\detokenize{psi_cdindex_support_defs.v}\#L51}{\texttt{\detokenize{size_in_expand_cde}}}] \emph{Lemma, line 51}. \textsf{[aux]} \\
  Every element of \texttt{\detokenize{expand_cde m}} has length \texttt{\detokenize{cde_total_width m}}.

\item[\href{https://github.com/LLM4Rocq/mathcomp-eulerian/blob/main/\detokenize{psi_cdindex_support_defs.v}\#L68}{\texttt{\detokenize{has_transition_cons}}}] \emph{Lemma, line 68}. \textsf{[aux]} \\
  Cons-shift for \texttt{\detokenize{has_transition}}: bumping the index by 1 corresponds     to dropping a leading bit.

\item[\href{https://github.com/LLM4Rocq/mathcomp-eulerian/blob/main/\detokenize{psi_cdindex_support_defs.v}\#L73}{\texttt{\detokenize{has_transition_cons2}}}] \emph{Lemma, line 73}. \textsf{[aux]} \\
  Cons-shift by 2 for \texttt{\detokenize{has_transition}}: dropping two leading bits.

\item[\href{https://github.com/LLM4Rocq/mathcomp-eulerian/blob/main/\detokenize{psi_cdindex_support_defs.v}\#L78}{\texttt{\detokenize{cde_offset_C_succ}}}] \emph{Lemma, line 78}. \textsf{[aux]} \\
  Cons-unfolding for \texttt{\detokenize{cde_offset}} over a \texttt{\detokenize{C}} head: shift by 1 bit.

\item[\href{https://github.com/LLM4Rocq/mathcomp-eulerian/blob/main/\detokenize{psi_cdindex_support_defs.v}\#L83}{\texttt{\detokenize{cde_offset_D_succ}}}] \emph{Lemma, line 83}. \textsf{[aux]} \\
  Cons-unfolding for \texttt{\detokenize{cde_offset}} over a \texttt{\detokenize{D}} head: shift by 2 bits.

\item[\href{https://github.com/LLM4Rocq/mathcomp-eulerian/blob/main/\detokenize{psi_cdindex_support_defs.v}\#L89}{\texttt{\detokenize{expand_cde_mem_transitions}}}] \emph{Lemma, line 89}. \textsf{[aux]} \\
  Forward direction of the transition characterisation: every \texttt{\detokenize{X}} in     \texttt{\detokenize{expand_cde m}} has the required bit transitions at all \texttt{\detokenize{D}}-offsets.

\item[\href{https://github.com/LLM4Rocq/mathcomp-eulerian/blob/main/\detokenize{psi_cdindex_support_defs.v}\#L141}{\texttt{\detokenize{transitions_expand_cde_mem}}}] \emph{Lemma, line 141}. \textsf{[aux]} \\
  Backward direction: any bit-string of the right length with \texttt{\detokenize{D}}-offset     transitions is a member of \texttt{\detokenize{expand_cde m}}.

\item[\href{https://github.com/LLM4Rocq/mathcomp-eulerian/blob/main/\detokenize{psi_cdindex_support_defs.v}\#L220}{\texttt{\detokenize{expand_cde_mem_iff}}}] \emph{Lemma, line 220}. \textsf{[aux]} \\
  Combined characterisation: when \texttt{\detokenize{X}} has the right length, membership     in \texttt{\detokenize{expand_cde m}} is equivalent to having a transition at every     \texttt{\detokenize{D}}-offset.  Pivot lemma used by \texttt{\detokenize{phi_w_support_general}}.

\item[\href{https://github.com/LLM4Rocq/mathcomp-eulerian/blob/main/\detokenize{psi_cdindex_support_defs.v}\#L235}{\texttt{\detokenize{mem_filter_iota_nth}}}] \emph{Lemma, line 235}. \textsf{[aux]} \\
  Membership in the descent set \texttt{\detokenize{filter (nth X) (iota 0 m)}} reduces to     the bit at that index.

\item[\href{https://github.com/LLM4Rocq/mathcomp-eulerian/blob/main/\detokenize{psi_cdindex_support_defs.v}\#L245}{\texttt{\detokenize{foldr_maxn_ge}}}] \emph{Lemma, line 245}. \textsf{[aux]} \\
  Every element of \texttt{\detokenize{s}} is bounded above by \texttt{\detokenize{foldr maxn 0 s}}.

\item[\href{https://github.com/LLM4Rocq/mathcomp-eulerian/blob/main/\detokenize{psi_cdindex_support_defs.v}\#L255}{\texttt{\detokenize{mem_omega_seq}}}] \emph{Lemma, line 255}. \textsf{[aux]} \\
  Bool form of \texttt{\detokenize{omega_seq}} membership: the XOR criterion plus the     bound \texttt{\detokenize{k <= max s}}.

\item[\href{https://github.com/LLM4Rocq/mathcomp-eulerian/blob/main/\detokenize{psi_cdindex_support_defs.v}\#L263}{\texttt{\detokenize{has_transition_omega_seq}}}] \emph{Lemma, line 263}. \textsf{[aux]} \\
  Bridge: a bit transition in \texttt{\detokenize{X}} at index \texttt{\detokenize{k < m.-1}} is equivalent to     \texttt{k \textbackslash\{\}in omega\_seq} of \texttt{\detokenize{X}}'s descent set.  Pivot used in the support     proof to translate D-offset transitions into \texttt{\detokenize{omega_seq}} membership.

\item[\href{https://github.com/LLM4Rocq/mathcomp-eulerian/blob/main/\detokenize{psi_cdindex_support_defs.v}\#L295}{\texttt{\detokenize{cde_total_width_cat}}}] \emph{Lemma, line 295}. \textsf{[aux]} \\
  \texttt{\detokenize{cde_total_width}} is additive over concatenation.

\end{description}

% --------------------------------------------
\section*{\texttt{psi\_cdindex\_tree.v}}
\addcontentsline{toc}{section}{\texttt{psi\_cdindex\_tree.v}}

\begin{description}
\item[\href{https://github.com/LLM4Rocq/mathcomp-eulerian/blob/main/\detokenize{psi_cdindex_tree.v}\#L44}{\texttt{\detokenize{window_size_last_fuel}}}] \emph{Lemma, line 44}. \textsf{[aux]} \\
  Fuel-bounded variant: \texttt{\detokenize{window_size}} at the last index of a non-empty     \texttt{\detokenize{w}} equals 1.  Used to prove the leaf-status of the last vertex.

\item[\href{https://github.com/LLM4Rocq/mathcomp-eulerian/blob/main/\detokenize{psi_cdindex_tree.v}\#L73}{\texttt{\detokenize{window_size_last}}}] \emph{Lemma, line 73}. \textsf{[aux]} \\
  Last vertex has window size 1, i.e. is always an endpoint.

\item[\href{https://github.com/LLM4Rocq/mathcomp-eulerian/blob/main/\detokenize{psi_cdindex_tree.v}\#L85}{\texttt{\detokenize{window_size_t}}}] \emph{Fixpoint, line 85}. \textsf{[aux]} \\
  \texttt{\detokenize{window_size_t i t}} is the tree-level analogue of \texttt{\detokenize{window_size}}:     recursion on the \texttt{\detokenize{mmtree}} \texttt{\detokenize{t}} using the in-order index \texttt{\detokenize{i}}.

\item[\href{https://github.com/LLM4Rocq/mathcomp-eulerian/blob/main/\detokenize{psi_cdindex_tree.v}\#L96}{\texttt{\detokenize{has_left_child_t}}}] \emph{Fixpoint, line 96}. \textsf{[aux]} \\
  \texttt{\detokenize{has_left_child_t i t}} is the tree-level analogue of \texttt{\detokenize{has_left_child}}:     direct structural recursion on the \texttt{\detokenize{mmtree}}.

\item[\href{https://github.com/LLM4Rocq/mathcomp-eulerian/blob/main/\detokenize{psi_cdindex_tree.v}\#L106}{\texttt{\detokenize{is_internal_t}}}] \emph{Definition, line 106}. \textsf{[aux]} \\
  Tree-level analogue of \texttt{\detokenize{is_internal}}.

\item[\href{https://github.com/LLM4Rocq/mathcomp-eulerian/blob/main/\detokenize{psi_cdindex_tree.v}\#L115}{\texttt{\detokenize{window_size_t_Node_lt}}}] \emph{Lemma, line 115}. \textsf{[aux]} \\
  Cons-style equation: index in left subtree forwards to \texttt{\detokenize{l}}.

\item[\href{https://github.com/LLM4Rocq/mathcomp-eulerian/blob/main/\detokenize{psi_cdindex_tree.v}\#L121}{\texttt{\detokenize{window_size_t_Node_eq}}}] \emph{Lemma, line 121}. \textsf{[aux]} \\
  At the root index, the window covers the root and the right subtree.

\item[\href{https://github.com/LLM4Rocq/mathcomp-eulerian/blob/main/\detokenize{psi_cdindex_tree.v}\#L127}{\texttt{\detokenize{window_size_t_Node_gt}}}] \emph{Lemma, line 127}. \textsf{[aux]} \\
  Index above the root forwards to the right subtree with shifted index.

\item[\href{https://github.com/LLM4Rocq/mathcomp-eulerian/blob/main/\detokenize{psi_cdindex_tree.v}\#L141}{\texttt{\detokenize{has_left_child_t_Node_lt}}}] \emph{Lemma, line 141}. \textsf{[aux]} \\
  Cons-style equation for \texttt{\detokenize{has_left_child_t}}: index in left subtree.

\item[\href{https://github.com/LLM4Rocq/mathcomp-eulerian/blob/main/\detokenize{psi_cdindex_tree.v}\#L147}{\texttt{\detokenize{has_left_child_t_Node_eq}}}] \emph{Lemma, line 147}. \textsf{[aux]} \\
  Root has a left child iff the left subtree is non-empty.

\item[\href{https://github.com/LLM4Rocq/mathcomp-eulerian/blob/main/\detokenize{psi_cdindex_tree.v}\#L153}{\texttt{\detokenize{has_left_child_t_Node_gt}}}] \emph{Lemma, line 153}. \textsf{[aux]} \\
  Index above the root forwards \texttt{\detokenize{has_left_child_t}} to right subtree.

\item[\href{https://github.com/LLM4Rocq/mathcomp-eulerian/blob/main/\detokenize{psi_cdindex_tree.v}\#L167}{\texttt{\detokenize{size_mmtree_to_seq_node}}}] \emph{Lemma, line 167}. \textsf{[aux]} \\
  In-order seq of \texttt{\detokenize{Node l x r}} has size \texttt{\detokenize{|l| + |r| + 1}}.

\item[\href{https://github.com/LLM4Rocq/mathcomp-eulerian/blob/main/\detokenize{psi_cdindex_tree.v}\#L173}{\texttt{\detokenize{has_left_child_t_0}}}] \emph{Lemma, line 173}. \textsf{[aux]} \\
  Position 0 (leftmost vertex) never has a left child.

\item[\href{https://github.com/LLM4Rocq/mathcomp-eulerian/blob/main/\detokenize{psi_cdindex_tree.v}\#L188}{\texttt{\detokenize{window_size_t_eq}}}] \emph{Lemma, line 188}. \textsf{[aux]} \\
  Bridge: tree-level \texttt{\detokenize{window_size_t}} agrees with seq-level \texttt{\detokenize{window_size}}     on the in-order traversal of any valid mmtree.

\item[\href{https://github.com/LLM4Rocq/mathcomp-eulerian/blob/main/\detokenize{psi_cdindex_tree.v}\#L207}{\texttt{\detokenize{has_left_child_t_eq}}}] \emph{Lemma, line 207}. \textsf{[aux]} \\
  Bridge: tree-level and seq-level \texttt{\detokenize{has_left_child}} agree.

\item[\href{https://github.com/LLM4Rocq/mathcomp-eulerian/blob/main/\detokenize{psi_cdindex_tree.v}\#L225}{\texttt{\detokenize{is_internal_t_eq}}}] \emph{Lemma, line 225}. \textsf{[aux]} \\
  Bridge: tree-level and seq-level \texttt{\detokenize{is_internal}} agree.

\item[\href{https://github.com/LLM4Rocq/mathcomp-eulerian/blob/main/\detokenize{psi_cdindex_tree.v}\#L237}{\texttt{\detokenize{endpoint_implies_next_t}}}] \emph{Lemma, line 237}. \textsf{[aux]} \\
  Tree-level: an endpoint at position \texttt{\detokenize{k}} is immediately followed (at     \texttt{\detokenize{k.+1}}) by a vertex with a left child.  Used to bound \texttt{\detokenize{S_w_seq}}.

\item[\href{https://github.com/LLM4Rocq/mathcomp-eulerian/blob/main/\detokenize{psi_cdindex_tree.v}\#L304}{\texttt{\detokenize{LR_pred_is_endpoint_t}}}] \emph{Lemma, line 304}. \textsf{[aux]} \\
  Tree-level: the predecessor of an LR-vertex (D-vertex) is an endpoint.     Key structural property used in the bit-recovery lemmas.

\item[\href{https://github.com/LLM4Rocq/mathcomp-eulerian/blob/main/\detokenize{psi_cdindex_tree.v}\#L394}{\texttt{\detokenize{endpoint_implies_next_has_left_child}}}] \emph{Lemma, line 394}. \textsf{[aux]} \\
  Seq-level corollary of \texttt{\detokenize{endpoint_implies_next_t}} via the mmtree bridge.

\item[\href{https://github.com/LLM4Rocq/mathcomp-eulerian/blob/main/\detokenize{psi_cdindex_tree.v}\#L412}{\texttt{\detokenize{LR_pred_is_endpoint}}}] \emph{Lemma, line 412}. \textsf{[aux]} \\
  Seq-level corollary of \texttt{\detokenize{LR_pred_is_endpoint_t}}: D-vertex predecessor is     always an endpoint.  Used by the bit-recovery proofs in     \texttt{\detokenize{psi_cdindex_core.v}}.

\end{description}

% --------------------------------------------
\section*{\texttt{psi\_cdindex\_tree\_shape.v}}
\addcontentsline{toc}{section}{\texttt{psi\_cdindex\_tree\_shape.v}}

\begin{description}
\item[\href{https://github.com/LLM4Rocq/mathcomp-eulerian/blob/main/\detokenize{psi_cdindex_tree_shape.v}\#L29}{\texttt{\detokenize{mmtree_shape_fuel}}}] \emph{Fixpoint, line 29}. \textsf{[aux]} \\
  \texttt{\detokenize{mmtree_shape_fuel fuel s}} is the fuel-bounded shape encoding of the     min-max tree of \texttt{\detokenize{s}}: head = \texttt{\detokenize{mm_pos s}}, tail = shape of left subtree     concatenated with shape of right subtree.

\item[\href{https://github.com/LLM4Rocq/mathcomp-eulerian/blob/main/\detokenize{psi_cdindex_tree_shape.v}\#L48}{\texttt{\detokenize{mmtree_shape}}}] \emph{Definition, line 48}. \textsf{[aux]} \\
  \texttt{\detokenize{mmtree_shape s}} is the canonical shape encoding of the min-max tree of     \texttt{\detokenize{s}}, obtained by saturating the fuel at \texttt{\detokenize{size s}}.  Two sequences with     equal shape have isomorphic min-max trees.

\item[\href{https://github.com/LLM4Rocq/mathcomp-eulerian/blob/main/\detokenize{psi_cdindex_tree_shape.v}\#L54}{\texttt{\detokenize{mmtree_shape_fuel_monotone}}}] \emph{Lemma, line 54}. \textsf{[aux]} \\
  \texttt{\detokenize{mmtree_shape_fuel}} is independent of the fuel once it is large enough.

\item[\href{https://github.com/LLM4Rocq/mathcomp-eulerian/blob/main/\detokenize{psi_cdindex_tree_shape.v}\#L78}{\texttt{\detokenize{mmtree_shape_nil}}}] \emph{Lemma, line 78}. \textsf{[aux]} \\
  Shape of the empty sequence is empty.

\item[\href{https://github.com/LLM4Rocq/mathcomp-eulerian/blob/main/\detokenize{psi_cdindex_tree_shape.v}\#L83}{\texttt{\detokenize{mmtree_shape_cons}}}] \emph{Lemma, line 83}. \textsf{[aux]} \\
  Cons unfolding: shape of \texttt{\detokenize{a :: s0}} = \texttt{\detokenize{mm_pos}} then shapes of left and     right subtrees, mirroring the tree-construction recursion.

\item[\href{https://github.com/LLM4Rocq/mathcomp-eulerian/blob/main/\detokenize{psi_cdindex_tree_shape.v}\#L106}{\texttt{\detokenize{size_mmtree_shape}}}] \emph{Lemma, line 106}. \textsf{[aux]} \\
  Shape encoding has the same length as the underlying sequence.

\item[\href{https://github.com/LLM4Rocq/mathcomp-eulerian/blob/main/\detokenize{psi_cdindex_tree_shape.v}\#L134}{\texttt{\detokenize{mmtree_shape_order_iso}}}] \emph{Lemma, line 134}. \textsf{[aux]} \\
  Headline structural lemma: order-isomorphic sequences (same size, both     uniq, same comparison pattern) produce the same tree shape.  This is the     single heavy proof underlying the per-property invariance lemmas below.

\item[\href{https://github.com/LLM4Rocq/mathcomp-eulerian/blob/main/\detokenize{psi_cdindex_tree_shape.v}\#L189}{\texttt{\detokenize{seq_cat_left_eq}}}] \emph{Lemma, line 189}. \textsf{[aux]} \\
  Left-cancellation for \texttt{\detokenize{++}} when prefixes have equal size.

\item[\href{https://github.com/LLM4Rocq/mathcomp-eulerian/blob/main/\detokenize{psi_cdindex_tree_shape.v}\#L200}{\texttt{\detokenize{seq_cat_right_eq}}}] \emph{Lemma, line 200}. \textsf{[aux]} \\
  Right-cancellation for \texttt{\detokenize{++}} when prefixes have equal size.

\item[\href{https://github.com/LLM4Rocq/mathcomp-eulerian/blob/main/\detokenize{psi_cdindex_tree_shape.v}\#L212}{\texttt{\detokenize{mmtree_shape_decompose}}}] \emph{Lemma, line 212}. \textsf{[aux]} \\
  Inverse of \texttt{\detokenize{mmtree_shape_cons}}: equal shapes split into equal roots and     equal subtree shapes, providing a structural recursion principle.

\item[\href{https://github.com/LLM4Rocq/mathcomp-eulerian/blob/main/\detokenize{psi_cdindex_tree_shape.v}\#L244}{\texttt{\detokenize{has_left_child_of_shape}}}] \emph{Lemma, line 244}. \textsf{[aux]} \\
  \texttt{\detokenize{has_left_child}} depends only on the tree shape: equal shapes give equal     \texttt{\detokenize{has_left_child}} at every index.

\item[\href{https://github.com/LLM4Rocq/mathcomp-eulerian/blob/main/\detokenize{psi_cdindex_tree_shape.v}\#L291}{\texttt{\detokenize{window_size_of_shape}}}] \emph{Lemma, line 291}. \textsf{[aux]} \\
  \texttt{\detokenize{window_size}} depends only on the tree shape.

\item[\href{https://github.com/LLM4Rocq/mathcomp-eulerian/blob/main/\detokenize{psi_cdindex_tree_shape.v}\#L345}{\texttt{\detokenize{has_left_child_order_iso}}}] \emph{Lemma, line 345}. \textsf{[aux]} \\
  Compatibility wrapper used downstream: order-isomorphic uniq sequences     give equal \texttt{\detokenize{has_left_child}} at every index.  Composes     \texttt{\detokenize{has_left_child_of_shape}} with \texttt{\detokenize{mmtree_shape_order_iso}}.

\item[\href{https://github.com/LLM4Rocq/mathcomp-eulerian/blob/main/\detokenize{psi_cdindex_tree_shape.v}\#L367}{\texttt{\detokenize{mmtree_shape_psi}}}] \emph{Lemma, line 367}. \textsf{[aux]} \\
  Heavy lemma: applying \texttt{\detokenize{psi j}} does not change the tree shape (for uniq     \texttt{\detokenize{w}}).  Proved by case-splitting on \texttt{\detokenize{j}} vs \texttt{\detokenize{m = mm_pos w}} and reducing     the root case to \texttt{\detokenize{mmtree_shape_order_iso}} on the rotated tail.

\item[\href{https://github.com/LLM4Rocq/mathcomp-eulerian/blob/main/\detokenize{psi_cdindex_tree_shape.v}\#L486}{\texttt{\detokenize{has_left_child_psi}}}] \emph{Lemma, line 486}. \textsf{[aux]} \\
  Headline corollary: \texttt{\detokenize{psi j}} preserves \texttt{\detokenize{has_left_child}}.  Five-line     proof composing \texttt{\detokenize{has_left_child_of_shape}} with \texttt{\detokenize{mmtree_shape_psi}}.

\end{description}

% --------------------------------------------
\section*{\texttt{psi\_cdindex\_witness.v}}
\addcontentsline{toc}{section}{\texttt{psi\_cdindex\_witness.v}}

\begin{description}
\item[\href{https://github.com/LLM4Rocq/mathcomp-eulerian/blob/main/\detokenize{psi_cdindex_witness.v}\#L30}{\texttt{\detokenize{omega_seq}}}] \emph{Definition, line 30}. \textsf{[aux]} \\
  \texttt{\detokenize{omega_seq s}} is the seq-level \texttt{\detokenize{omega}} map: positions \texttt{\detokenize{k}} such that     exactly one of \texttt{\detokenize{k}}, \texttt{\detokenize{k+1}} belongs to \texttt{\detokenize{s}}.  Mirrors \texttt{\detokenize{omega_set}} from     \texttt{\detokenize{beta_swap.v}} on raw \texttt{\detokenize{seq nat}} (no finset).

\item[\href{https://github.com/LLM4Rocq/mathcomp-eulerian/blob/main/\detokenize{psi_cdindex_witness.v}\#L39}{\texttt{\detokenize{is_D_letter}}}] \emph{Definition, line 39}. \textsf{[aux]} \\
  \texttt{\detokenize{is_D_letter l}} tests whether \texttt{\detokenize{l}} is the D cd-letter.

\item[\href{https://github.com/LLM4Rocq/mathcomp-eulerian/blob/main/\detokenize{psi_cdindex_witness.v}\#L45}{\texttt{\detokenize{S_w_seq}}}] \emph{Definition, line 45}. \textsf{[aux]} \\
  \texttt{\detokenize{S_w_seq w}} is the d-position set of \texttt{\detokenize{w}}: predecessors of D-vertices     in the cd-classification.  This is the support set \texttt{\detokenize{S_w}} from     Stanley's Prop 1.6.4.

\item[\href{https://github.com/LLM4Rocq/mathcomp-eulerian/blob/main/\detokenize{psi_cdindex_witness.v}\#L104}{\texttt{\detokenize{omega_monotone_class_count}}}] \emph{Lemma, line 104}. \textsf{[aux]} \\
  Weak monotonicity: if \texttt{\detokenize{omega_seq S}} is contained in \texttt{\detokenize{omega_seq T}}, then     every M-class whose d-positions are covered by \texttt{\detokenize{omega_seq S}} is also     covered by \texttt{\detokenize{omega_seq T}}.  Half of Stanley's Prop 1.6.4.

\item[\href{https://github.com/LLM4Rocq/mathcomp-eulerian/blob/main/\detokenize{psi_cdindex_witness.v}\#L125}{\texttt{\detokenize{witness_perm}}}] \emph{Definition, line 125}. \textsf{[aux]} \\
  \texttt{\detokenize{witness_perm n k}} is the permutation \texttt{\detokenize{1;..;k; k+2; k+1; k+3;..;n}}:     a single inversion at position \texttt{\detokenize{k}} embedded in a sorted backbone.     Realises a cd-word with \texttt{S\_w = \{k\}}, used in \texttt{\detokenize{strict_witness_exists}}.

\item[\href{https://github.com/LLM4Rocq/mathcomp-eulerian/blob/main/\detokenize{psi_cdindex_witness.v}\#L129}{\texttt{\detokenize{leq_maxn}}}] \emph{Lemma, line 129}. \textsf{[aux]} \\
  Local arithmetic helper: \texttt{\detokenize{maxn a b = a}} when \texttt{\detokenize{b <= a}}.

\item[\href{https://github.com/LLM4Rocq/mathcomp-eulerian/blob/main/\detokenize{psi_cdindex_witness.v}\#L137}{\texttt{\detokenize{geq_maxn}}}] \emph{Lemma, line 137}. \textsf{[aux]} \\
  Local arithmetic helper: \texttt{\detokenize{maxn a b = b}} when \texttt{\detokenize{a <= b}}.

\item[\href{https://github.com/LLM4Rocq/mathcomp-eulerian/blob/main/\detokenize{psi_cdindex_witness.v}\#L141}{\texttt{\detokenize{geq_minn}}}] \emph{Lemma, line 141}. \textsf{[aux]} \\
  Local arithmetic helper: \texttt{\detokenize{minn a b = a}} when \texttt{\detokenize{a <= b}}.

\item[\href{https://github.com/LLM4Rocq/mathcomp-eulerian/blob/main/\detokenize{psi_cdindex_witness.v}\#L149}{\texttt{\detokenize{path_leq_last}}}] \emph{Lemma, line 149}. \textsf{[aux]} \\
  A leq-path from \texttt{\detokenize{a}} over \texttt{\detokenize{s}} gives \texttt{\detokenize{a <= last a s}}.

\item[\href{https://github.com/LLM4Rocq/mathcomp-eulerian/blob/main/\detokenize{psi_cdindex_witness.v}\#L158}{\texttt{\detokenize{foldr_maxn_path}}}] \emph{Lemma, line 158}. \textsf{[aux]} \\
  On a sorted (leq-path) sequence, \texttt{\detokenize{foldr maxn}} reduces to its last     element.

\item[\href{https://github.com/LLM4Rocq/mathcomp-eulerian/blob/main/\detokenize{psi_cdindex_witness.v}\#L178}{\texttt{\detokenize{foldr_minn_ge}}}] \emph{Lemma, line 178}. \textsf{[aux]} \\
  \texttt{\detokenize{foldr minn a s = a}} when every element of \texttt{\detokenize{s}} is at least \texttt{\detokenize{a}}.

\item[\href{https://github.com/LLM4Rocq/mathcomp-eulerian/blob/main/\detokenize{psi_cdindex_witness.v}\#L187}{\texttt{\detokenize{all_le_iota}}}] \emph{Lemma, line 187}. \textsf{[aux]} \\
  Every element of \texttt{\detokenize{iota m.+1 n}} is at least \texttt{\detokenize{m}}.

\item[\href{https://github.com/LLM4Rocq/mathcomp-eulerian/blob/main/\detokenize{psi_cdindex_witness.v}\#L194}{\texttt{\detokenize{min_pos_iota}}}] \emph{Lemma, line 194}. \textsf{[aux]} \\
  \texttt{\detokenize{min_pos}} of an ascending \texttt{\detokenize{iota}} is 0 (minimum is the head).

\item[\href{https://github.com/LLM4Rocq/mathcomp-eulerian/blob/main/\detokenize{psi_cdindex_witness.v}\#L201}{\texttt{\detokenize{path_iota}}}] \emph{Lemma, line 201}. \textsf{[aux]} \\
  \texttt{\detokenize{iota m.+1 n}} is a leq-path with starting point \texttt{\detokenize{m}}.

\item[\href{https://github.com/LLM4Rocq/mathcomp-eulerian/blob/main/\detokenize{psi_cdindex_witness.v}\#L205}{\texttt{\detokenize{last_iota}}}] \emph{Lemma, line 205}. \textsf{[aux]} \\
  Last element of \texttt{\detokenize{m :: iota m.+1 n}} is \texttt{\detokenize{m + n}}.

\item[\href{https://github.com/LLM4Rocq/mathcomp-eulerian/blob/main/\detokenize{psi_cdindex_witness.v}\#L212}{\texttt{\detokenize{foldr_maxn_iota}}}] \emph{Lemma, line 212}. \textsf{[aux]} \\
  \texttt{\detokenize{foldr maxn}} of an ascending \texttt{\detokenize{iota}} equals its last element.

\item[\href{https://github.com/LLM4Rocq/mathcomp-eulerian/blob/main/\detokenize{psi_cdindex_witness.v}\#L222}{\texttt{\detokenize{max_pos_iota}}}] \emph{Lemma, line 222}. \textsf{[aux]} \\
  \texttt{\detokenize{max_pos}} of an ascending \texttt{\detokenize{iota}} is the last index.

\item[\href{https://github.com/LLM4Rocq/mathcomp-eulerian/blob/main/\detokenize{psi_cdindex_witness.v}\#L239}{\texttt{\detokenize{mm_pos_iota}}}] \emph{Lemma, line 239}. \textsf{[aux]} \\
  \texttt{\detokenize{mm_pos}} of an ascending \texttt{\detokenize{iota}} is 0 (root is the minimum).

\item[\href{https://github.com/LLM4Rocq/mathcomp-eulerian/blob/main/\detokenize{psi_cdindex_witness.v}\#L245}{\texttt{\detokenize{size_witness_perm}}}] \emph{Lemma, line 245}. \textsf{[aux]} \\
  \texttt{\detokenize{witness_perm n k}} has length \texttt{\detokenize{n}} when \texttt{\detokenize{k + 2 < n}}.

\item[\href{https://github.com/LLM4Rocq/mathcomp-eulerian/blob/main/\detokenize{psi_cdindex_witness.v}\#L256}{\texttt{\detokenize{witness_perm_uniq}}}] \emph{Lemma, line 256}. \textsf{[aux]} \\
  \texttt{\detokenize{witness_perm n k}} is duplicate-free when \texttt{\detokenize{k + 2 < n}}.

\item[\href{https://github.com/LLM4Rocq/mathcomp-eulerian/blob/main/\detokenize{psi_cdindex_witness.v}\#L297}{\texttt{\detokenize{check_strict_witness}}}] \emph{Definition, line 297}. \textsf{[aux]} \\
  Boolean witness: \texttt{\detokenize{witness_perm n k}} is uniq, of length \texttt{\detokenize{n}}, and has     \texttt{\detokenize{S_w_seq}} equal to \texttt{\detokenize{[:: k]}}.  Used for vm\_compute sanity checks.

\item[\href{https://github.com/LLM4Rocq/mathcomp-eulerian/blob/main/\detokenize{psi_cdindex_witness.v}\#L323}{\texttt{\detokenize{check_strict_witness_correct}}}] \emph{Lemma, line 323}. \textsf{[aux]} \\
  \texttt{\detokenize{check_strict_witness}} reflects \texttt{\detokenize{strict_witness_exists}} for that \texttt{\detokenize{n,k}}.

\item[\href{https://github.com/LLM4Rocq/mathcomp-eulerian/blob/main/\detokenize{psi_cdindex_witness.v}\#L335}{\texttt{\detokenize{has_left_child_iota}}}] \emph{Lemma, line 335}. \textsf{[aux]} \\
  Helper: ascending sequences (\texttt{\detokenize{iota}}) have no left children in their     min-max tree, since \texttt{\detokenize{mm_pos = 0}} at every recursion depth.

\item[\href{https://github.com/LLM4Rocq/mathcomp-eulerian/blob/main/\detokenize{psi_cdindex_witness.v}\#L348}{\texttt{\detokenize{foldr_minn_all_gt}}}] \emph{Lemma, line 348}. \textsf{[aux]} \\
  Helper: \texttt{\detokenize{foldr minn a s = a}} when \texttt{\detokenize{a}} is strictly less than every     element of \texttt{\detokenize{s}}.

\item[\href{https://github.com/LLM4Rocq/mathcomp-eulerian/blob/main/\detokenize{psi_cdindex_witness.v}\#L359}{\texttt{\detokenize{mm_pos_min_first}}}] \emph{Lemma, line 359}. \textsf{[aux]} \\
  \texttt{\detokenize{mm_pos (a :: s) = 0}} when \texttt{\detokenize{a}} is the strict minimum of \texttt{\detokenize{a :: s}}.

\item[\href{https://github.com/LLM4Rocq/mathcomp-eulerian/blob/main/\detokenize{psi_cdindex_witness.v}\#L372}{\texttt{\detokenize{min_pos_core}}}] \emph{Lemma, line 372}. \textsf{[aux]} \\
  \texttt{\detokenize{min_pos}} of the witness core \texttt{\detokenize{[k+2; k+1] ++ iota k.+3 m}} is 1     (the \texttt{\detokenize{k+1}} at position 1 is the unique minimum).

\item[\href{https://github.com/LLM4Rocq/mathcomp-eulerian/blob/main/\detokenize{psi_cdindex_witness.v}\#L386}{\texttt{\detokenize{max_pos_core_gt0}}}] \emph{Lemma, line 386}. \textsf{[aux]} \\
  \texttt{\detokenize{max_pos}} of the witness core is positive when the suffix \texttt{\detokenize{iota k.+3 m}}     is non-empty.

\item[\href{https://github.com/LLM4Rocq/mathcomp-eulerian/blob/main/\detokenize{psi_cdindex_witness.v}\#L401}{\texttt{\detokenize{mm_pos_core}}}] \emph{Lemma, line 401}. \textsf{[aux]} \\
  \texttt{\detokenize{mm_pos}} of the witness core is 1 (position of \texttt{\detokenize{k+1}}) when the suffix     is non-empty.

\item[\href{https://github.com/LLM4Rocq/mathcomp-eulerian/blob/main/\detokenize{psi_cdindex_witness.v}\#L418}{\texttt{\detokenize{hlc_core_not1}}}] \emph{Lemma, line 418}. \textsf{[aux]} \\
  In the witness core, only position 1 has a left child; all other     positions are leaves of the min-max tree.

\item[\href{https://github.com/LLM4Rocq/mathcomp-eulerian/blob/main/\detokenize{psi_cdindex_witness.v}\#L437}{\texttt{\detokenize{hlc_core_1}}}] \emph{Lemma, line 437}. \textsf{[aux]} \\
  Position 1 in the witness core has a left child (the \texttt{\detokenize{k+2}} vertex).

\item[\href{https://github.com/LLM4Rocq/mathcomp-eulerian/blob/main/\detokenize{psi_cdindex_witness.v}\#L450}{\texttt{\detokenize{ws_core_1}}}] \emph{Lemma, line 450}. \textsf{[aux]} \\
  \texttt{\detokenize{window_size}} at the \texttt{\detokenize{mm_pos}} position 1 in the witness core is \texttt{\detokenize{m.+1}}     (covers \texttt{\detokenize{k+1}} and the entire ascending suffix).

\item[\href{https://github.com/LLM4Rocq/mathcomp-eulerian/blob/main/\detokenize{psi_cdindex_witness.v}\#L466}{\texttt{\detokenize{S_w_seq_core}}}] \emph{Lemma, line 466}. \textsf{[aux]} \\
  \texttt{\detokenize{S_w_seq}} of the witness core is \texttt{\detokenize{[:: 0]}}: the only D-position is at     index 1, contributing \texttt{\detokenize{0}} to the d-position set.

\item[\href{https://github.com/LLM4Rocq/mathcomp-eulerian/blob/main/\detokenize{psi_cdindex_witness.v}\#L501}{\texttt{\detokenize{S_w_seq_witness_k0}}}] \emph{Lemma, line 501}. \textsf{[aux]} \\
  Base case \texttt{\detokenize{k = 0}}: \texttt{\detokenize{witness_perm n 0}} is exactly the witness core, so     \texttt{\detokenize{S_w_seq = [:: 0]}}.

\item[\href{https://github.com/LLM4Rocq/mathcomp-eulerian/blob/main/\detokenize{psi_cdindex_witness.v}\#L514}{\texttt{\detokenize{classify_skip_mm0}}}] \emph{Lemma, line 514}. \textsf{[aux]} \\
  When the head \texttt{\detokenize{a}} is at the root (\texttt{\detokenize{mm_pos = 0}}), classifying any     positive position \texttt{\detokenize{i}} in \texttt{\detokenize{a :: rest}} reduces to classifying \texttt{\detokenize{i-1}} in     the right subtree \texttt{\detokenize{rest}}.

\item[\href{https://github.com/LLM4Rocq/mathcomp-eulerian/blob/main/\detokenize{psi_cdindex_witness.v}\#L537}{\texttt{\detokenize{S_w_seq_shift}}}] \emph{Lemma, line 537}. \textsf{[aux]} \\
  Index-shift property: when the head is at the root, \texttt{\detokenize{S_w_seq}} of the     cons equals the \texttt{\detokenize{S_w_seq}} of the tail with every entry incremented.

\item[\href{https://github.com/LLM4Rocq/mathcomp-eulerian/blob/main/\detokenize{psi_cdindex_witness.v}\#L581}{\texttt{\detokenize{drop1_witness_map_succ}}}] \emph{Lemma, line 581}. \textsf{[aux]} \\
  Tail of \texttt{\detokenize{witness_perm n k.+1}} equals the \texttt{\detokenize{+1}}-shift of     \texttt{\detokenize{witness_perm n.-1 k}}; provides the inductive step for k.

\item[\href{https://github.com/LLM4Rocq/mathcomp-eulerian/blob/main/\detokenize{psi_cdindex_witness.v}\#L606}{\texttt{\detokenize{classify_map_succ}}}] \emph{Lemma, line 606}. \textsf{[aux]} \\
  \texttt{\detokenize{classify_vertex_cde}} is invariant under the \texttt{\detokenize{+1}}-shift on a uniq     sequence (since cd-classification depends only on order).

\item[\href{https://github.com/LLM4Rocq/mathcomp-eulerian/blob/main/\detokenize{psi_cdindex_witness.v}\#L634}{\texttt{\detokenize{S_w_seq_map_succ}}}] \emph{Lemma, line 634}. \textsf{[aux]} \\
  \texttt{\detokenize{S_w_seq}} is invariant under the \texttt{\detokenize{+1}}-shift on a uniq sequence.

\item[\href{https://github.com/LLM4Rocq/mathcomp-eulerian/blob/main/\detokenize{psi_cdindex_witness.v}\#L647}{\texttt{\detokenize{mm_pos_witness}}}] \emph{Lemma, line 647}. \textsf{[aux]} \\
  \texttt{\detokenize{mm_pos}} of \texttt{\detokenize{witness_perm n k}} is 0 for \texttt{\detokenize{k >= 1}}: the \texttt{\detokenize{1}} at the head     is the strict minimum of the permutation.

\item[\href{https://github.com/LLM4Rocq/mathcomp-eulerian/blob/main/\detokenize{psi_cdindex_witness.v}\#L698}{\texttt{\detokenize{strict_witness_exists}}}] \emph{Lemma, line 698}. \textsf{[aux]} \\
  Strict-witness existence (Stanley Prop 1.6.4, strict half): for every     \texttt{\detokenize{k < n.-2}} there is a uniq permutation \texttt{\detokenize{w}} of length \texttt{\detokenize{n}} with     \texttt{\detokenize{S_w_seq w = [:: k]}}.  The witness is \texttt{\detokenize{witness_perm n k}}.

\end{description}

% --------------------------------------------
\section*{\texttt{psi\_comm.v}}
\addcontentsline{toc}{section}{\texttt{psi\_comm.v}}

\begin{description}
\item[\href{https://github.com/LLM4Rocq/mathcomp-eulerian/blob/main/\detokenize{psi_comm.v}\#L15}{\texttt{\detokenize{sorted_uniq_nth_ltn}}}] \emph{Lemma, line 15}. \textsf{[aux]} \\
  \texttt{\detokenize{sorted_uniq_nth_ltn}} : on a uniq sorted sequence, value order at two     indices is the same as index order; ranking helper for \texttt{\detokenize{psi_comm}}.

\item[\href{https://github.com/LLM4Rocq/mathcomp-eulerian/blob/main/\detokenize{psi_comm.v}\#L36}{\texttt{\detokenize{shift_preserves_ltn}}}] \emph{Lemma, line 36}. \textsf{[aux]} \\
  \texttt{\detokenize{shift_preserves_ltn}} : the modular shift by \texttt{\detokenize{delta}} preserves the     \texttt{\detokenize{<}} relation on ranks within the appropriate non-extremum sub-range.     Modular-arithmetic core of interior order preservation.

\item[\href{https://github.com/LLM4Rocq/mathcomp-eulerian/blob/main/\detokenize{psi_comm.v}\#L79}{\texttt{\detokenize{rank_shift_preserves_interior_order}}}] \emph{Lemma, line 79}. \textsf{[aux]} \\
  \texttt{\detokenize{rank_shift_preserves_interior_order}} : at non-head positions \texttt{\detokenize{p, q}}     the rank-shift on \texttt{\detokenize{L}} is order-preserving (no wrap-around).  Used to     propagate descent comparisons through \texttt{\detokenize{psi}}.

\item[\href{https://github.com/LLM4Rocq/mathcomp-eulerian/blob/main/\detokenize{psi_comm.v}\#L163}{\texttt{\detokenize{window_size_psi}}}] \emph{Lemma, line 163}. \textsf{[aux]} \\
  \texttt{\detokenize{window_size_psi}} : applying \texttt{\detokenize{psi j}} does not change \texttt{\detokenize{window_size i}}     at any other position \texttt{\detokenize{i}} -- generalization of \texttt{\detokenize{window_size_psi_self}}     to cross positions, foundational for \texttt{\detokenize{psi_comm}}.

\item[\href{https://github.com/LLM4Rocq/mathcomp-eulerian/blob/main/\detokenize{psi_comm.v}\#L324}{\texttt{\detokenize{window_at_psi_disjoint}}}] \emph{Lemma, line 324}. \textsf{[aux]} \\
  \texttt{\detokenize{window_at_psi_disjoint}} : \texttt{\detokenize{psi j}} preserves \texttt{\detokenize{window_at i}} verbatim     when the windows at \texttt{\detokenize{i}} and \texttt{\detokenize{j}} are disjoint (i.e. \texttt{\detokenize{i + ws_i <= j}}).

\item[\href{https://github.com/LLM4Rocq/mathcomp-eulerian/blob/main/\detokenize{psi_comm.v}\#L354}{\texttt{\detokenize{psi_comm_disjoint_lr}}}] \emph{Lemma, line 354}. \textsf{[aux]} \\
  \texttt{\detokenize{psi_comm_disjoint_lr}} : disjoint commutativity in the directed form     \texttt{\detokenize{i + ws_i <= j}}; combinatorially trivial since the windows act on     non-overlapping slices of \texttt{\detokenize{w}}.

\item[\href{https://github.com/LLM4Rocq/mathcomp-eulerian/blob/main/\detokenize{psi_comm.v}\#L446}{\texttt{\detokenize{psi_comm_disjoint}}}] \emph{Lemma, line 446}. \textsf{[aux]} \\
  \texttt{\detokenize{psi_comm_disjoint}} : symmetric form of \texttt{\detokenize{psi_comm_disjoint_lr}} --     \texttt{\detokenize{psi i}} and \texttt{\detokenize{psi j}} commute whenever the two M-windows are disjoint.

\item[\href{https://github.com/LLM4Rocq/mathcomp-eulerian/blob/main/\detokenize{psi_comm.v}\#L479}{\texttt{\detokenize{window_size_psi_ancestor}}}] \emph{Lemma, line 479}. \textsf{[aux]} \\
  \texttt{\detokenize{window_size_psi_ancestor}} : nested specialization of     \texttt{\detokenize{window_size_psi}} -- when \texttt{\detokenize{W_j}} sits inside \texttt{\detokenize{W_i}}, applying \texttt{\detokenize{psi i}}     preserves \texttt{\detokenize{window_size j}}.

\item[\href{https://github.com/LLM4Rocq/mathcomp-eulerian/blob/main/\detokenize{psi_comm.v}\#L513}{\texttt{\detokenize{sort_map_mono}}}] \emph{Lemma, line 513}. \textsf{[aux]} \\
  \texttt{\detokenize{sort_map_mono}} : \texttt{\detokenize{sort leq}} commutes with a monotone injection on     the elements of \texttt{\detokenize{L}}; structural helper for \texttt{\detokenize{psi_map_comm}}.

\item[\href{https://github.com/LLM4Rocq/mathcomp-eulerian/blob/main/\detokenize{psi_comm.v}\#L551}{\texttt{\detokenize{index_map_inj_in}}}] \emph{Lemma, line 551}. \textsf{[aux]} \\
  \texttt{\detokenize{index_map_inj_in}} : \texttt{\detokenize{index}} commutes with a locally injective map     \texttt{\detokenize{f}}, a structural helper for the rank-shift transfer.

\item[\href{https://github.com/LLM4Rocq/mathcomp-eulerian/blob/main/\detokenize{psi_comm.v}\#L574}{\texttt{\detokenize{mono_inj_in}}}] \emph{Lemma, line 574}. \textsf{[aux]} \\
  \texttt{\detokenize{mono_inj_in}} : a \texttt{\detokenize{<}}-monotone function on a finite set is injective     there; trivial corollary used in \texttt{\detokenize{psi_map_comm}}.

\item[\href{https://github.com/LLM4Rocq/mathcomp-eulerian/blob/main/\detokenize{psi_comm.v}\#L590}{\texttt{\detokenize{rank_shift_map_comm}}}] \emph{Lemma, line 590}. \textsf{[aux]} \\
  \texttt{\detokenize{rank_shift_map_comm}} : \texttt{\detokenize{rank_shift_seq}} commutes with any monotone     injection on the elements of \texttt{\detokenize{L}}; algebraic core of \texttt{\detokenize{psi_map_comm}}.

\item[\href{https://github.com/LLM4Rocq/mathcomp-eulerian/blob/main/\detokenize{psi_comm.v}\#L657}{\texttt{\detokenize{psi_map_comm}}}] \emph{Lemma, line 657}. \textsf{[aux]} \\
  \texttt{\detokenize{psi_map_comm}} : \texttt{\detokenize{psi k}} commutes with any comparison-preserving     relabelling \texttt{\detokenize{f}}.  Reduces nested commutativity to interior-only     rank-shift commutativity.

\item[\href{https://github.com/LLM4Rocq/mathcomp-eulerian/blob/main/\detokenize{psi_comm.v}\#L720}{\texttt{\detokenize{rank_shift_psi_comm}}}] \emph{Lemma, line 720}. \textsf{[aux]} \\
  \texttt{\detokenize{rank_shift_psi_comm}} : when \texttt{\detokenize{d}} is rooted at the head (\texttt{\detokenize{mm_pos d = 0}})     and \texttt{\detokenize{k > 0}}, \texttt{\detokenize{rank_shift_seq}} commutes with \texttt{\detokenize{psi k}}.  Algebraic core     of \texttt{\detokenize{psi_comm_nested}}.

\item[\href{https://github.com/LLM4Rocq/mathcomp-eulerian/blob/main/\detokenize{psi_comm.v}\#L833}{\texttt{\detokenize{psi_comm_nested}}}] \emph{Lemma, line 833}. \textsf{[aux]} \\
  \texttt{\detokenize{psi_comm_nested}} : nested commutativity -- when \texttt{\detokenize{W_j}} sits inside     \texttt{\detokenize{W_i}}, \texttt{\detokenize{psi i}} and \texttt{\detokenize{psi j}} commute.  Combines window-stability and     \texttt{\detokenize{rank_shift_psi_comm}}.

\item[\href{https://github.com/LLM4Rocq/mathcomp-eulerian/blob/main/\detokenize{psi_comm.v}\#L1351}{\texttt{\detokenize{psi_comm}}}] \emph{Theorem, line 1351}. \textsf{[aux]} \\
  \texttt{\detokenize{psi_comm}} is the M3 commutativity theorem (Stanley Fact \#1, half 1):     the operators \texttt{\detokenize{psi i}} all pairwise commute on uniq inputs.  Proved by     \texttt{\detokenize{window_trichotomy}} dispatching to disjoint or nested commutativity.

\end{description}

% --------------------------------------------
\section*{\texttt{psi\_core.v}}
\addcontentsline{toc}{section}{\texttt{psi\_core.v}}

\begin{description}
\item[\href{https://github.com/LLM4Rocq/mathcomp-eulerian/blob/main/\detokenize{psi_core.v}\#L18}{\texttt{\detokenize{max_pos}}}] \emph{Definition, line 18}. \textsf{[aux]} \\
  \texttt{\detokenize{max_pos s}} is the least index in \texttt{\detokenize{s}} at which the maximum of \texttt{\detokenize{s}}     occurs, the dual of \texttt{\detokenize{min_pos}}.

\item[\href{https://github.com/LLM4Rocq/mathcomp-eulerian/blob/main/\detokenize{psi_core.v}\#L23}{\texttt{\detokenize{mm_pos}}}] \emph{Definition, line 23}. \textsf{[aux]} \\
  \texttt{\detokenize{mm_pos s}} is Stanley's min-or-max split index: the least index at     which either the minimum or the maximum of \texttt{\detokenize{s}} occurs.

\item[\href{https://github.com/LLM4Rocq/mathcomp-eulerian/blob/main/\detokenize{psi_core.v}\#L28}{\texttt{\detokenize{max_in}}}] \emph{Lemma, line 28}. \textsf{[aux]} \\
  \texttt{\detokenize{max_in}} is the dual of \texttt{\detokenize{min_in}}: \texttt{\detokenize{foldr maxn a s}} always lies in     \texttt{\detokenize{a :: s}}.

\item[\href{https://github.com/LLM4Rocq/mathcomp-eulerian/blob/main/\detokenize{psi_core.v}\#L39}{\texttt{\detokenize{max_pos_lt}}}] \emph{Lemma, line 39}. \textsf{[aux]} \\
  \texttt{\detokenize{max_pos_lt}} : on a nonempty sequence the index \texttt{\detokenize{max_pos s}} is in range.

\item[\href{https://github.com/LLM4Rocq/mathcomp-eulerian/blob/main/\detokenize{psi_core.v}\#L47}{\texttt{\detokenize{mm_pos_lt}}}] \emph{Lemma, line 47}. \textsf{[aux]} \\
  \texttt{\detokenize{mm_pos_lt}} : on a nonempty sequence the min-or-max split index is in     range, ensuring \texttt{\detokenize{mmtree_of_seq_mm_fuel}} recursion is well-defined.

\item[\href{https://github.com/LLM4Rocq/mathcomp-eulerian/blob/main/\detokenize{psi_core.v}\#L56}{\texttt{\detokenize{mmtree_of_seq_mm_fuel}}}] \emph{Fixpoint, line 56}. \textsf{[aux]} \\
  \texttt{\detokenize{mmtree_of_seq_mm_fuel fuel s}} is the Stanley-correct min-max tree     construction (M2 variant): splits \texttt{\detokenize{s}} at \texttt{\detokenize{mm_pos s}} rather than     \texttt{\detokenize{min_pos s}}, producing the alternating-extremum tree shape.

\item[\href{https://github.com/LLM4Rocq/mathcomp-eulerian/blob/main/\detokenize{psi_core.v}\#L72}{\texttt{\detokenize{mmtree_of_seq_mm}}}] \emph{Definition, line 72}. \textsf{[aux]} \\
  \texttt{\detokenize{mmtree_of_seq_mm s}} is Stanley's min-max tree \texttt{\detokenize{M(w)}} for \texttt{\detokenize{s = w}}: the     canonical M-class representative used to define \texttt{\detokenize{psi}}.

\item[\href{https://github.com/LLM4Rocq/mathcomp-eulerian/blob/main/\detokenize{psi_core.v}\#L77}{\texttt{\detokenize{mmtree_of_seq_mm_fuel_correct}}}] \emph{Lemma, line 77}. \textsf{[aux]} \\
  \texttt{\detokenize{mmtree_of_seq_mm_fuel_correct}} : the M2 fuel construction round-trips     in-order, the analogue of \texttt{\detokenize{mmtree_of_seq_fuel_correct}}.

\item[\href{https://github.com/LLM4Rocq/mathcomp-eulerian/blob/main/\detokenize{psi_core.v}\#L101}{\texttt{\detokenize{mmtree_of_seq_mmK}}}] \emph{Theorem, line 101}. \textsf{[aux]} \\
  \texttt{\detokenize{mmtree_of_seq_mmK}} is the M2 round-trip theorem: in-order traversal     inverts the Stanley min-max construction \texttt{\detokenize{mmtree_of_seq_mm}}.

\item[\href{https://github.com/LLM4Rocq/mathcomp-eulerian/blob/main/\detokenize{psi_core.v}\#L115}{\texttt{\detokenize{window_size_fuel}}}] \emph{Fixpoint, line 115}. \textsf{[aux]} \\
  \texttt{\detokenize{window_size_fuel fuel i s}} computes the size of the M-window at     in-order position \texttt{\detokenize{i}} (i.e. \texttt{\detokenize{1 + |right subtree of vertex i|}}) by     fuel-bounded recursion mirroring \texttt{\detokenize{mmtree_of_seq_mm_fuel}}.

\item[\href{https://github.com/LLM4Rocq/mathcomp-eulerian/blob/main/\detokenize{psi_core.v}\#L132}{\texttt{\detokenize{window_size}}}] \emph{Definition, line 132}. \textsf{[aux]} \\
  \texttt{\detokenize{window_size i s}} is \texttt{\detokenize{1 + |right subtree of vertex at in-order position i|}}     in \texttt{\detokenize{M(s)}}, i.e. the size of the M-window on which \texttt{\detokenize{psi i}} acts; \texttt{\detokenize{0}} when     \texttt{\detokenize{i >= size s}}.

\item[\href{https://github.com/LLM4Rocq/mathcomp-eulerian/blob/main/\detokenize{psi_core.v}\#L138}{\texttt{\detokenize{window_at}}}] \emph{Definition, line 138}. \textsf{[aux]} \\
  \texttt{\detokenize{window_at i w}} is the in-order labels of the M-window at position \texttt{\detokenize{i}}:     the contiguous slice \texttt{w\_i, w\_\{i+1\}, ..., w\_\{i+ws-1\}} where \texttt{\detokenize{ws =     window_size i w}}.

\item[\href{https://github.com/LLM4Rocq/mathcomp-eulerian/blob/main/\detokenize{psi_core.v}\#L149}{\texttt{\detokenize{window_size_fuel_bound}}}] \emph{Lemma, line 149}. \textsf{[aux]} \\
  \texttt{\detokenize{window_size_fuel_bound}} : combined invariant for the fuel version --     the window is positive iff \texttt{\detokenize{i}} is in range, and is bounded by the     available suffix length.

\item[\href{https://github.com/LLM4Rocq/mathcomp-eulerian/blob/main/\detokenize{psi_core.v}\#L212}{\texttt{\detokenize{window_size_gt0}}}] \emph{Lemma, line 212}. \textsf{[aux]} \\
  \texttt{\detokenize{window_size_gt0}} : the M-window at an in-range position is nonempty.

\item[\href{https://github.com/LLM4Rocq/mathcomp-eulerian/blob/main/\detokenize{psi_core.v}\#L221}{\texttt{\detokenize{window_size_bound}}}] \emph{Lemma, line 221}. \textsf{[aux]} \\
  \texttt{\detokenize{window_size_bound}} : the M-window at \texttt{\detokenize{i}} never exceeds the suffix     \texttt{\detokenize{size w - i}}, so \texttt{\detokenize{window_at i w}} is a slice of \texttt{\detokenize{w}}.

\item[\href{https://github.com/LLM4Rocq/mathcomp-eulerian/blob/main/\detokenize{psi_core.v}\#L228}{\texttt{\detokenize{window_size_oor}}}] \emph{Lemma, line 228}. \textsf{[aux]} \\
  \texttt{\detokenize{window_size_oor}} : out-of-range positions have empty window, by which     \texttt{\detokenize{psi i}} degenerates to the identity.

\item[\href{https://github.com/LLM4Rocq/mathcomp-eulerian/blob/main/\detokenize{psi_core.v}\#L245}{\texttt{\detokenize{rank_shift_seq}}}] \emph{Definition, line 245}. \textsf{[aux]} \\
  \texttt{\detokenize{rank_shift_seq L}} is the rank-shift on the M-window: replace the head     with the opposite extremum (max if head = min, else min) and rotate the     other ranks accordingly.  This is the algebraic core of \texttt{\detokenize{psi_i}}.

\item[\href{https://github.com/LLM4Rocq/mathcomp-eulerian/blob/main/\detokenize{psi_core.v}\#L253}{\texttt{\detokenize{size_rank_shift_seq}}}] \emph{Lemma, line 253}. \textsf{[aux]} \\
  \texttt{\detokenize{size_rank_shift_seq}} : \texttt{\detokenize{rank_shift_seq}} preserves length (definitional).

\item[\href{https://github.com/LLM4Rocq/mathcomp-eulerian/blob/main/\detokenize{psi_core.v}\#L265}{\texttt{\detokenize{psi}}}] \emph{Definition, line 265}. \textsf{[aux]} \\
  \texttt{\detokenize{psi i w}} is Stanley's \texttt{\detokenize{psi_i}} operator on the M-class representative     \texttt{\detokenize{w}}: it fixes the prefix of length \texttt{\detokenize{i}}, applies the rank-shift to the     M-window \texttt{\detokenize{window_at i w}}, and concatenates the unchanged suffix.  Total     function; identity outside the window or when \texttt{\detokenize{i >= size w}}.

\item[\href{https://github.com/LLM4Rocq/mathcomp-eulerian/blob/main/\detokenize{psi_core.v}\#L272}{\texttt{\detokenize{map_nth_iota_sorted}}}] \emph{Lemma, line 272}. \textsf{[aux]} \\
  \texttt{\detokenize{map_nth_iota_sorted}} : mapping \texttt{\detokenize{nth 0 sorted}} across \texttt{\detokenize{iota 0 (size sorted)}}     reproduces \texttt{\detokenize{sorted}}.  Helper for the rank-shift permutation analysis.

\item[\href{https://github.com/LLM4Rocq/mathcomp-eulerian/blob/main/\detokenize{psi_core.v}\#L286}{\texttt{\detokenize{map_mod_iota_rot}}}] \emph{Lemma, line 286}. \textsf{[aux]} \\
  \texttt{\detokenize{map_mod_iota_rot}} : the modular shift \texttt{r |-> (r + shift) \%\% k} on     \texttt{\detokenize{iota 0 k}} equals a rotation of \texttt{\detokenize{iota 0 k}}; key step in proving     \texttt{\detokenize{rank_shift_perm_eq}}.

\item[\href{https://github.com/LLM4Rocq/mathcomp-eulerian/blob/main/\detokenize{psi_core.v}\#L320}{\texttt{\detokenize{rank_shift_perm_eq}}}] \emph{Lemma, line 320}. \textsf{[aux]} \\
  \texttt{\detokenize{rank_shift_perm_eq}} : \texttt{\detokenize{rank_shift_seq L}} is a permutation of \texttt{\detokenize{L}}, so     \texttt{\detokenize{psi i}} preserves the multiset of labels in \texttt{\detokenize{w}}.

\item[\href{https://github.com/LLM4Rocq/mathcomp-eulerian/blob/main/\detokenize{psi_core.v}\#L362}{\texttt{\detokenize{sort_rank_shift_seq}}}] \emph{Lemma, line 362}. \textsf{[aux]} \\
  \texttt{\detokenize{sort_rank_shift_seq}} : the rank-shift commutes with \texttt{\detokenize{sort leq}}, a     consequence of \texttt{\detokenize{rank_shift_perm_eq}}.

\item[\href{https://github.com/LLM4Rocq/mathcomp-eulerian/blob/main/\detokenize{psi_core.v}\#L375}{\texttt{\detokenize{uniq_rank_shift_seq}}}] \emph{Lemma, line 375}. \textsf{[aux]} \\
  \texttt{\detokenize{uniq_rank_shift_seq}} : the rank-shift preserves uniqueness, by     permutation.

\item[\href{https://github.com/LLM4Rocq/mathcomp-eulerian/blob/main/\detokenize{psi_core.v}\#L382}{\texttt{\detokenize{size_rank_shift_seq2}}}] \emph{Lemma, line 382}. \textsf{[aux]} \\
  \texttt{\detokenize{size_rank_shift_seq2}} : size preservation derived from \texttt{\detokenize{perm_eq}},     convenient form for downstream rewrites.

\item[\href{https://github.com/LLM4Rocq/mathcomp-eulerian/blob/main/\detokenize{psi_core.v}\#L390}{\texttt{\detokenize{psi_perm_eq}}}] \emph{Lemma, line 390}. \textsf{[aux]} \\
  \texttt{\detokenize{psi_perm_eq}} is the M2 multiset preservation theorem: \texttt{\detokenize{psi i w}} is a     permutation of \texttt{\detokenize{w}} for every \texttt{\detokenize{i}} and \texttt{\detokenize{w}}.

\item[\href{https://github.com/LLM4Rocq/mathcomp-eulerian/blob/main/\detokenize{psi_core.v}\#L423}{\texttt{\detokenize{rank_shift_seqE}}}] \emph{Lemma, line 423}. \textsf{[aux]} \\
  \texttt{\detokenize{rank_shift_seqE}} : explicit map characterization of \texttt{\detokenize{rank_shift_seq L}}     in terms of the sort and the head's extremum role.

\item[\href{https://github.com/LLM4Rocq/mathcomp-eulerian/blob/main/\detokenize{psi_core.v}\#L439}{\texttt{\detokenize{nth_rank_shift_seq}}}] \emph{Lemma, line 439}. \textsf{[aux]} \\
  \texttt{\detokenize{nth_rank_shift_seq}} : positionwise formula for \texttt{\detokenize{rank_shift_seq L}},     used in the involutivity and interior-order proofs.

\item[\href{https://github.com/LLM4Rocq/mathcomp-eulerian/blob/main/\detokenize{psi_core.v}\#L452}{\texttt{\detokenize{head_rank_shift_seq}}}] \emph{Lemma, line 452}. \textsf{[aux]} \\
  \texttt{\detokenize{head_rank_shift_seq}} : the head of \texttt{\detokenize{rank_shift_seq L}} is the formula     above instantiated at \texttt{\detokenize{n = 0}}.

\item[\href{https://github.com/LLM4Rocq/mathcomp-eulerian/blob/main/\detokenize{psi_core.v}\#L467}{\texttt{\detokenize{rank_shift_seq_involutive}}}] \emph{Lemma, line 467}. \textsf{[aux]} \\
  \texttt{\detokenize{rank_shift_seq_involutive}} : applying the rank-shift twice is the     identity, provided the head of \texttt{\detokenize{L}} is the global min or max.  Algebraic     heart of the \texttt{\detokenize{psi_involutive}} theorem.

\item[\href{https://github.com/LLM4Rocq/mathcomp-eulerian/blob/main/\detokenize{psi_core.v}\#L540}{\texttt{\detokenize{window_size_fuel_monotone}}}] \emph{Lemma, line 540}. \textsf{[aux]} \\
  \texttt{\detokenize{window_size_fuel_monotone}} : the fuel-bounded \texttt{\detokenize{window_size}} is     independent of fuel as long as fuel covers the input size.

\item[\href{https://github.com/LLM4Rocq/mathcomp-eulerian/blob/main/\detokenize{psi_core.v}\#L568}{\texttt{\detokenize{window_size_cons}}}] \emph{Lemma, line 568}. \textsf{[aux]} \\
  \texttt{\detokenize{window_size_cons}} : structural unfolding of \texttt{\detokenize{window_size}} on a \texttt{\detokenize{cons}}     by the \texttt{\detokenize{mm_pos}} split, mirroring the recursion of \texttt{\detokenize{mmtree_of_seq_mm}}.

\item[\href{https://github.com/LLM4Rocq/mathcomp-eulerian/blob/main/\detokenize{psi_core.v}\#L592}{\texttt{\detokenize{psi_id_oor}}}] \emph{Lemma, line 592}. \textsf{[aux]} \\
  \texttt{\detokenize{psi_id_oor}} : \texttt{\detokenize{psi i w = w}} when \texttt{\detokenize{i}} is out of range, the     edge convention from M2\_PSI\_INFORMAL.md.

\item[\href{https://github.com/LLM4Rocq/mathcomp-eulerian/blob/main/\detokenize{psi_core.v}\#L603}{\texttt{\detokenize{psi_id_trivial}}}] \emph{Lemma, line 603}. \textsf{[aux]} \\
  \texttt{\detokenize{psi_id_trivial}} : \texttt{\detokenize{psi i w = w}} on a trivial window (size \texttt{\detokenize{<= 1}},     i.e. a leaf or single-vertex right subtree).

\item[\href{https://github.com/LLM4Rocq/mathcomp-eulerian/blob/main/\detokenize{psi_core.v}\#L624}{\texttt{\detokenize{window_at_cons}}}] \emph{Lemma, line 624}. \textsf{[aux]} \\
  \texttt{\detokenize{window_at_cons}} : structural unfolding of \texttt{\detokenize{window_at}} on a \texttt{\detokenize{cons}}, by     the \texttt{\detokenize{mm_pos}} split.  Companion of \texttt{\detokenize{window_size_cons}}.

\item[\href{https://github.com/LLM4Rocq/mathcomp-eulerian/blob/main/\detokenize{psi_core.v}\#L665}{\texttt{\detokenize{size_psi}}}] \emph{Lemma, line 665}. \textsf{[aux]} \\
  \texttt{\detokenize{size_psi}} : \texttt{\detokenize{psi i}} preserves length, derived from \texttt{\detokenize{psi_perm_eq}}.

\item[\href{https://github.com/LLM4Rocq/mathcomp-eulerian/blob/main/\detokenize{psi_core.v}\#L669}{\texttt{\detokenize{uniq_psi}}}] \emph{Lemma, line 669}. \textsf{[aux]} \\
  \texttt{\detokenize{uniq_psi}} : \texttt{\detokenize{psi i}} preserves uniqueness, derived from \texttt{\detokenize{psi_perm_eq}}.

\item[\href{https://github.com/LLM4Rocq/mathcomp-eulerian/blob/main/\detokenize{psi_core.v}\#L674}{\texttt{\detokenize{take_psi}}}] \emph{Lemma, line 674}. \textsf{[aux]} \\
  \texttt{\detokenize{take_psi}} : \texttt{\detokenize{psi j}} does not touch positions strictly before \texttt{\detokenize{j}}, so     any prefix of length \texttt{\detokenize{k <= j}} is unchanged.

\item[\href{https://github.com/LLM4Rocq/mathcomp-eulerian/blob/main/\detokenize{psi_core.v}\#L698}{\texttt{\detokenize{foldr_minn_aux}}}] \emph{Lemma, line 698}. \textsf{[aux]} \\
  \texttt{\detokenize{foldr_minn_aux}} : \texttt{\detokenize{foldr minn a s}} bounds \texttt{\detokenize{a}} and every element of     \texttt{\detokenize{s}} from above; technical helper for the extremum reasoning.

\item[\href{https://github.com/LLM4Rocq/mathcomp-eulerian/blob/main/\detokenize{psi_core.v}\#L711}{\texttt{\detokenize{foldr_minn_le}}}] \emph{Lemma, line 711}. \textsf{[aux]} \\
  \texttt{\detokenize{foldr_minn_le}} : convenient form of \texttt{\detokenize{foldr_minn_aux}} -- the fold-min     bounds every element of \texttt{\detokenize{a :: s}} from below.

\item[\href{https://github.com/LLM4Rocq/mathcomp-eulerian/blob/main/\detokenize{psi_core.v}\#L718}{\texttt{\detokenize{foldr_maxn_aux}}}] \emph{Lemma, line 718}. \textsf{[aux]} \\
  \texttt{\detokenize{foldr_maxn_aux}} : dual of \texttt{\detokenize{foldr_minn_aux}} for \texttt{\detokenize{maxn}}.

\item[\href{https://github.com/LLM4Rocq/mathcomp-eulerian/blob/main/\detokenize{psi_core.v}\#L730}{\texttt{\detokenize{foldr_maxn_ge}}}] \emph{Lemma, line 730}. \textsf{[aux]} \\
  \texttt{\detokenize{foldr_maxn_ge}} : dual of \texttt{\detokenize{foldr_minn_le}} for \texttt{\detokenize{maxn}}.

\item[\href{https://github.com/LLM4Rocq/mathcomp-eulerian/blob/main/\detokenize{psi_core.v}\#L738}{\texttt{\detokenize{min_val_perm_eq}}}] \emph{Lemma, line 738}. \textsf{[aux]} \\
  \texttt{\detokenize{min_val_perm_eq}} : the global min as computed by \texttt{\detokenize{foldr minn}} depends     only on the multiset, so \texttt{\detokenize{psi i}} preserves the minimum value.

\item[\href{https://github.com/LLM4Rocq/mathcomp-eulerian/blob/main/\detokenize{psi_core.v}\#L751}{\texttt{\detokenize{max_val_perm_eq}}}] \emph{Lemma, line 751}. \textsf{[aux]} \\
  \texttt{\detokenize{max_val_perm_eq}} : dual of \texttt{\detokenize{min_val_perm_eq}} for the global max.

\item[\href{https://github.com/LLM4Rocq/mathcomp-eulerian/blob/main/\detokenize{psi_core.v}\#L767}{\texttt{\detokenize{window_fits_left}}}] \emph{Lemma, line 767}. \textsf{[aux]} \\
  \texttt{\detokenize{window_fits_left}} : when \texttt{\detokenize{i}} sits in the left subtree of the root     (\texttt{\detokenize{i < mm_pos w}}), the M-window at \texttt{\detokenize{i}} does not cross the root.

\item[\href{https://github.com/LLM4Rocq/mathcomp-eulerian/blob/main/\detokenize{psi_core.v}\#L786}{\texttt{\detokenize{drop_psi}}}] \emph{Lemma, line 786}. \textsf{[aux]} \\
  \texttt{\detokenize{drop_psi}} : \texttt{\detokenize{psi i}} does not touch positions at or after the window     end, so the suffix from \texttt{\detokenize{i + window_size i w}} is unchanged.

\item[\href{https://github.com/LLM4Rocq/mathcomp-eulerian/blob/main/\detokenize{psi_core.v}\#L805}{\texttt{\detokenize{nth_psi_left}}}] \emph{Lemma, line 805}. \textsf{[aux]} \\
  \texttt{\detokenize{nth_psi_left}} : positionwise version of \texttt{\detokenize{take_psi}} -- entries before     the window are fixed.

\item[\href{https://github.com/LLM4Rocq/mathcomp-eulerian/blob/main/\detokenize{psi_core.v}\#L814}{\texttt{\detokenize{nth_psi_right}}}] \emph{Lemma, line 814}. \textsf{[aux]} \\
  \texttt{\detokenize{nth_psi_right}} : positionwise version of \texttt{\detokenize{drop_psi}} -- entries after     the window are fixed.

\item[\href{https://github.com/LLM4Rocq/mathcomp-eulerian/blob/main/\detokenize{psi_core.v}\#L824}{\texttt{\detokenize{nth_psi_inside}}}] \emph{Lemma, line 824}. \textsf{[aux]} \\
  \texttt{\detokenize{nth_psi_inside}} : within the window, \texttt{\detokenize{psi i w}} reads off     \texttt{\detokenize{rank_shift_seq (window_at i w)}} at the offset \texttt{\detokenize{k - i}}.

\item[\href{https://github.com/LLM4Rocq/mathcomp-eulerian/blob/main/\detokenize{psi_core.v}\#L849}{\texttt{\detokenize{take_psi_perm}}}] \emph{Lemma, line 849}. \textsf{[aux]} \\
  \texttt{\detokenize{take_psi_perm}} : when \texttt{\detokenize{j}} sits past the window end, the prefix     \texttt{\detokenize{take j (psi i w)}} is a permutation of \texttt{\detokenize{take j w}}; used in     \texttt{\detokenize{mm_pos_psi_eq}} to argue invariance of the global extremum positions.

\item[\href{https://github.com/LLM4Rocq/mathcomp-eulerian/blob/main/\detokenize{psi_core.v}\#L886}{\texttt{\detokenize{notin_take_mm}}}] \emph{Lemma, line 886}. \textsf{[aux]} \\
  \texttt{\detokenize{notin_take_mm}} : neither the global min nor the global max appears in     the prefix \texttt{\detokenize{take (mm_pos w) w}}, by minimality of \texttt{\detokenize{mm_pos}}.

\item[\href{https://github.com/LLM4Rocq/mathcomp-eulerian/blob/main/\detokenize{psi_core.v}\#L904}{\texttt{\detokenize{min_val_drop}}}] \emph{Lemma, line 904}. \textsf{[aux]} \\
  \texttt{\detokenize{min_val_drop}} : if the global min is not in \texttt{\detokenize{take j w}}, the min of     \texttt{\detokenize{drop j w}} equals the global min of \texttt{\detokenize{w}}.

\item[\href{https://github.com/LLM4Rocq/mathcomp-eulerian/blob/main/\detokenize{psi_core.v}\#L929}{\texttt{\detokenize{max_val_drop}}}] \emph{Lemma, line 929}. \textsf{[aux]} \\
  \texttt{\detokenize{max_val_drop}} : dual of \texttt{\detokenize{min_val_drop}} for the global max.

\item[\href{https://github.com/LLM4Rocq/mathcomp-eulerian/blob/main/\detokenize{psi_core.v}\#L958}{\texttt{\detokenize{mm_pos_char}}}] \emph{Lemma, line 958}. \textsf{[aux]} \\
  \texttt{\detokenize{mm_pos_char}} : the \texttt{\detokenize{mm_pos s}} index is uniquely characterized by:     no extremum sits in \texttt{\detokenize{take j s}} and the extremum sits at \texttt{\detokenize{nth 0 s j}}.     Used to prove \texttt{\detokenize{mm_pos}} is preserved by \texttt{\detokenize{psi}}.

\item[\href{https://github.com/LLM4Rocq/mathcomp-eulerian/blob/main/\detokenize{psi_core.v}\#L996}{\texttt{\detokenize{nth_w_mm_pos}}}] \emph{Lemma, line 996}. \textsf{[aux]} \\
  \texttt{\detokenize{nth_w_mm_pos}} : the entry at \texttt{\detokenize{mm_pos w}} is exactly the global min or     the global max of \texttt{\detokenize{w}}, by definition of \texttt{\detokenize{mm_pos}}.

\item[\href{https://github.com/LLM4Rocq/mathcomp-eulerian/blob/main/\detokenize{psi_core.v}\#L1011}{\texttt{\detokenize{sorted_head_le}}}] \emph{Lemma, line 1011}. \textsf{[aux]} \\
  \texttt{\detokenize{sorted_head_le}} : on a sorted sequence, the head bounds every element     from below.

\item[\href{https://github.com/LLM4Rocq/mathcomp-eulerian/blob/main/\detokenize{psi_core.v}\#L1024}{\texttt{\detokenize{sorted_last_ge}}}] \emph{Lemma, line 1024}. \textsf{[aux]} \\
  \texttt{\detokenize{sorted_last_ge}} : on a sorted sequence, the last element bounds every     element from above.

\item[\href{https://github.com/LLM4Rocq/mathcomp-eulerian/blob/main/\detokenize{psi_core.v}\#L1040}{\texttt{\detokenize{min_eq_nth_sort_0}}}] \emph{Lemma, line 1040}. \textsf{[aux]} \\
  \texttt{\detokenize{min_eq_nth_sort_0}} : the \texttt{\detokenize{foldr minn}} minimum equals the first entry     of the sorted list -- bridge from value-based to rank-based statements.

\item[\href{https://github.com/LLM4Rocq/mathcomp-eulerian/blob/main/\detokenize{psi_core.v}\#L1060}{\texttt{\detokenize{max_eq_nth_sort_last}}}] \emph{Lemma, line 1060}. \textsf{[aux]} \\
  \texttt{\detokenize{max_eq_nth_sort_last}} : dual of \texttt{\detokenize{min_eq_nth_sort_0}} -- the     \texttt{\detokenize{foldr maxn}} maximum equals the last entry of the sorted list.

\item[\href{https://github.com/LLM4Rocq/mathcomp-eulerian/blob/main/\detokenize{psi_core.v}\#L1087}{\texttt{\detokenize{rank_shift_head_min_to_max}}}] \emph{Lemma, line 1087}. \textsf{[aux]} \\
  \texttt{\detokenize{rank_shift_head_min_to_max}} : if the head of \texttt{\detokenize{L}} is the minimum, the     rank-shift sends it to the maximum; this is one half of the head-flip.

\item[\href{https://github.com/LLM4Rocq/mathcomp-eulerian/blob/main/\detokenize{psi_core.v}\#L1106}{\texttt{\detokenize{rank_shift_head_max_to_min}}}] \emph{Lemma, line 1106}. \textsf{[aux]} \\
  \texttt{\detokenize{rank_shift_head_max_to_min}} : the dual head-flip -- if the head of     \texttt{\detokenize{L}} is the maximum, the rank-shift sends it to the minimum.

\item[\href{https://github.com/LLM4Rocq/mathcomp-eulerian/blob/main/\detokenize{psi_core.v}\#L1137}{\texttt{\detokenize{mm_pos_psi_eq}}}] \emph{Lemma, line 1137}. \textsf{[aux]} \\
  \texttt{\detokenize{mm_pos_psi_eq}} : \texttt{\detokenize{psi i}} preserves the global root index \texttt{\detokenize{mm_pos}},     by the head-flip combined with extremum-preservation away from the     window.  Key step toward \texttt{\detokenize{psi_involutive}}.

\item[\href{https://github.com/LLM4Rocq/mathcomp-eulerian/blob/main/\detokenize{psi_core.v}\#L1266}{\texttt{\detokenize{ws_lt_size}}}] \emph{Lemma, line 1266}. \textsf{[aux]} \\
  \texttt{\detokenize{ws_lt_size}} : a nontrivial window forces \texttt{\detokenize{i < size w}}; used to     discharge in-range hypotheses cheaply.

\item[\href{https://github.com/LLM4Rocq/mathcomp-eulerian/blob/main/\detokenize{psi_core.v}\#L1276}{\texttt{\detokenize{take_mm_psi}}}] \emph{Lemma, line 1276}. \textsf{[aux]} \\
  \texttt{\detokenize{take_mm_psi}} : when \texttt{\detokenize{i}} is in the left subtree, \texttt{\detokenize{psi i}} commutes with     \texttt{\detokenize{take (mm_pos w)}}, i.e. acts on the left subtree's labels alone.     Recursive descent step for \texttt{\detokenize{window_size_psi_self}} and     \texttt{\detokenize{window_at_psi_self}}.

\item[\href{https://github.com/LLM4Rocq/mathcomp-eulerian/blob/main/\detokenize{psi_core.v}\#L1309}{\texttt{\detokenize{drop_mm_psi}}}] \emph{Lemma, line 1309}. \textsf{[aux]} \\
  \texttt{\detokenize{drop_mm_psi}} : dual of \texttt{\detokenize{take_mm_psi}} -- when \texttt{\detokenize{i}} is in the right     subtree, \texttt{\detokenize{psi i}} commutes with \texttt{\detokenize{drop (mm_pos w).+1}} and re-indexes     \texttt{\detokenize{i}} by subtracting the left-plus-root size.

\item[\href{https://github.com/LLM4Rocq/mathcomp-eulerian/blob/main/\detokenize{psi_core.v}\#L1358}{\texttt{\detokenize{window_size_psi_self}}}] \emph{Lemma, line 1358}. \textsf{[aux]} \\
  \texttt{\detokenize{window_size_psi_self}} : applying \texttt{\detokenize{psi i}} to \texttt{\detokenize{w}} does not change the     window size at the same position \texttt{\detokenize{i}}.  Proved by induction on size     using \texttt{\detokenize{take_mm_psi}}/\texttt{\detokenize{drop_mm_psi}}.

\item[\href{https://github.com/LLM4Rocq/mathcomp-eulerian/blob/main/\detokenize{psi_core.v}\#L1418}{\texttt{\detokenize{window_at_psi_self}}}] \emph{Lemma, line 1418}. \textsf{[aux]} \\
  \texttt{\detokenize{window_at_psi_self}} : the window at \texttt{\detokenize{i}} after one application of     \texttt{\detokenize{psi i}} is exactly the rank-shifted original window.  Required by the     second \texttt{\detokenize{psi i}} application in \texttt{\detokenize{psi_involutive}}.

\item[\href{https://github.com/LLM4Rocq/mathcomp-eulerian/blob/main/\detokenize{psi_core.v}\#L1495}{\texttt{\detokenize{window_head_extremum}}}] \emph{Lemma, line 1495}. \textsf{[aux]} \\
  \texttt{\detokenize{window_head_extremum}} : the head of \texttt{\detokenize{window_at i w}} is always an     extremum (min or max) of the window -- Stanley's structural fact (F2)     transferred to the window.  Discharge condition for     \texttt{\detokenize{rank_shift_seq_involutive}}.

\item[\href{https://github.com/LLM4Rocq/mathcomp-eulerian/blob/main/\detokenize{psi_core.v}\#L1559}{\texttt{\detokenize{psi_involutive}}}] \emph{Theorem, line 1559}. \textsf{[aux]} \\
  \texttt{\detokenize{psi_involutive}} is the M2 involutivity theorem: \texttt{\detokenize{psi i (psi i w) = w}}     for every \texttt{\detokenize{i}} and uniq \texttt{\detokenize{w}}.  Stanley fact about \texttt{\detokenize{psi_i}}.

\item[\href{https://github.com/LLM4Rocq/mathcomp-eulerian/blob/main/\detokenize{psi_core.v}\#L1617}{\texttt{\detokenize{window_trichotomy}}}] \emph{Lemma, line 1617}. \textsf{[aux]} \\
  \texttt{\detokenize{window_trichotomy}} : two M-windows in the same word are either     disjoint (left or right) or one is nested in the other -- the binary     tree geometry that drives the \texttt{\detokenize{psi_comm}} proof.

\item[\href{https://github.com/LLM4Rocq/mathcomp-eulerian/blob/main/\detokenize{psi_core.v}\#L1821}{\texttt{\detokenize{foldr_minn_le_nth}}}] \emph{Lemma, line 1821}. \textsf{[aux]} \\
  \texttt{\detokenize{foldr_minn_le_nth}} : positionwise version of \texttt{\detokenize{foldr_minn_le}}.

\item[\href{https://github.com/LLM4Rocq/mathcomp-eulerian/blob/main/\detokenize{psi_core.v}\#L1827}{\texttt{\detokenize{foldr_maxn_ge_nth}}}] \emph{Lemma, line 1827}. \textsf{[aux]} \\
  \texttt{\detokenize{foldr_maxn_ge_nth}} : positionwise version of \texttt{\detokenize{foldr_maxn_ge}}.

\item[\href{https://github.com/LLM4Rocq/mathcomp-eulerian/blob/main/\detokenize{psi_core.v}\#L1836}{\texttt{\detokenize{mm_pos_order_iso}}}] \emph{Lemma, line 1836}. \textsf{[aux]} \\
  \texttt{\detokenize{mm_pos_order_iso}} : \texttt{\detokenize{mm_pos}} depends only on the comparison structure     of the sequence -- two order-isomorphic sequences have the same     \texttt{\detokenize{mm_pos}}.  Foundational for transferring tree shape across the     rank-shift.

\item[\href{https://github.com/LLM4Rocq/mathcomp-eulerian/blob/main/\detokenize{psi_core.v}\#L1943}{\texttt{\detokenize{order_iso_take}}}] \emph{Lemma, line 1943}. \textsf{[aux]} \\
  \texttt{\detokenize{order_iso_take}} : the order-isomorphism property descends to \texttt{\detokenize{take m}}     of both sequences; helper for the recursive case in     \texttt{\detokenize{window_size_order_iso}}.

\item[\href{https://github.com/LLM4Rocq/mathcomp-eulerian/blob/main/\detokenize{psi_core.v}\#L1961}{\texttt{\detokenize{order_iso_drop}}}] \emph{Lemma, line 1961}. \textsf{[aux]} \\
  \texttt{\detokenize{order_iso_drop}} : the order-isomorphism property descends to     \texttt{\detokenize{drop m.+1}} of both sequences; companion of \texttt{\detokenize{order_iso_take}}.

\item[\href{https://github.com/LLM4Rocq/mathcomp-eulerian/blob/main/\detokenize{psi_core.v}\#L1978}{\texttt{\detokenize{window_size_order_iso}}}] \emph{Lemma, line 1978}. \textsf{[aux]} \\
  \texttt{\detokenize{window_size_order_iso}} : \texttt{\detokenize{window_size}} is a function only of     comparison structure -- two order-isomorphic sequences have the same     window sizes at every position.  Lifts \texttt{\detokenize{mm_pos_order_iso}} recursively.

\item[\href{https://github.com/LLM4Rocq/mathcomp-eulerian/blob/main/\detokenize{psi_core.v}\#L2045}{\texttt{\detokenize{mm_pos_drop_mm}}}] \emph{Lemma, line 2045}. \textsf{[aux]} \\
  \texttt{\detokenize{mm_pos_drop_mm}} : after dropping the prefix up to the root, the new     head is itself the root of its subtree, so its \texttt{\detokenize{mm_pos}} is \texttt{\detokenize{0}}.

\item[\href{https://github.com/LLM4Rocq/mathcomp-eulerian/blob/main/\detokenize{psi_core.v}\#L2077}{\texttt{\detokenize{psi_0_eq}}}] \emph{Lemma, line 2077}. \textsf{[aux]} \\
  \texttt{\detokenize{psi_0_eq}} : when \texttt{\detokenize{mm_pos s = 0}} (the head is the root), \texttt{\detokenize{psi 0 s}}     reduces to \texttt{\detokenize{rank_shift_seq s}} -- the action on the whole sequence.

\item[\href{https://github.com/LLM4Rocq/mathcomp-eulerian/blob/main/\detokenize{psi_core.v}\#L2102}{\texttt{\detokenize{drop_psi_above}}}] \emph{Lemma, line 2102}. \textsf{[aux]} \\
  \texttt{\detokenize{drop_psi_above}} : \texttt{\detokenize{psi j}} does not affect any suffix \texttt{\detokenize{drop k w}} when     \texttt{\detokenize{k}} is past the window end -- generalizes \texttt{\detokenize{drop_psi}}.

\end{description}

% --------------------------------------------
\section*{\texttt{psi\_descent\_thms.v}}
\addcontentsline{toc}{section}{\texttt{psi\_descent\_thms.v}}

\begin{description}
\item[\href{https://github.com/LLM4Rocq/mathcomp-eulerian/blob/main/\detokenize{psi_descent_thms.v}\#L14}{\texttt{\detokenize{window_at_uniq}}}] \emph{Lemma, line 14}. \textsf{[aux]} \\
  \texttt{\detokenize{window_at_uniq}} : uniqueness propagates from \texttt{\detokenize{w}} to \texttt{\detokenize{window_at i w}}.

\item[\href{https://github.com/LLM4Rocq/mathcomp-eulerian/blob/main/\detokenize{psi_descent_thms.v}\#L20}{\texttt{\detokenize{size_window_at}}}] \emph{Lemma, line 20}. \textsf{[aux]} \\
  \texttt{\detokenize{size_window_at}} : in-range, \texttt{\detokenize{size (window_at i w) = window_size i w}};     the slice fits exactly.

\item[\href{https://github.com/LLM4Rocq/mathcomp-eulerian/blob/main/\detokenize{psi_descent_thms.v}\#L31}{\texttt{\detokenize{nth_window_at}}}] \emph{Lemma, line 31}. \textsf{[aux]} \\
  \texttt{\detokenize{nth_window_at}} : indexing into \texttt{\detokenize{window_at i w}} is offset by \texttt{\detokenize{i}}     relative to indexing into \texttt{\detokenize{w}}.

\item[\href{https://github.com/LLM4Rocq/mathcomp-eulerian/blob/main/\detokenize{psi_descent_thms.v}\#L42}{\texttt{\detokenize{elem_in_range}}}] \emph{Lemma, line 42}. \textsf{[aux]} \\
  \texttt{\detokenize{elem_in_range}} : every entry of \texttt{\detokenize{L}} is bounded between the sorted     extremes of \texttt{\detokenize{L}}.  Tag-along helper for \texttt{\detokenize{cmp_out_of_range}}.

\item[\href{https://github.com/LLM4Rocq/mathcomp-eulerian/blob/main/\detokenize{psi_descent_thms.v}\#L81}{\texttt{\detokenize{head_min_not_descent}}}] \emph{Lemma, line 81}. \textsf{[aux]} \\
  \texttt{\detokenize{head_min_not_descent}} : if the window head is the window-min, then     \texttt{\detokenize{i}} is not a descent of \texttt{\detokenize{w}}.  Connects head-extremum to descent     polarity.

\item[\href{https://github.com/LLM4Rocq/mathcomp-eulerian/blob/main/\detokenize{psi_descent_thms.v}\#L104}{\texttt{\detokenize{head_max_is_descent}}}] \emph{Lemma, line 104}. \textsf{[aux]} \\
  \texttt{\detokenize{head_max_is_descent}} : dual of \texttt{\detokenize{head_min_not_descent}} -- if the     window head is the window-max, then \texttt{\detokenize{i}} is a descent of \texttt{\detokenize{w}}.

\item[\href{https://github.com/LLM4Rocq/mathcomp-eulerian/blob/main/\detokenize{psi_descent_thms.v}\#L147}{\texttt{\detokenize{elem_rs_in_range}}}] \emph{Lemma, line 147}. \textsf{[aux]} \\
  \texttt{\detokenize{elem_rs_in_range}} : every entry of \texttt{\detokenize{rank_shift_seq L}} is bounded     between the sorted extremes of \texttt{\detokenize{L}}; rank-shift counterpart of     \texttt{\detokenize{elem_in_range}}.

\item[\href{https://github.com/LLM4Rocq/mathcomp-eulerian/blob/main/\detokenize{psi_descent_thms.v}\#L198}{\texttt{\detokenize{rs_head_max_descent}}}] \emph{Lemma, line 198}. \textsf{[aux]} \\
  \texttt{\detokenize{rs_head_max_descent}} : after the rank-shift sends the head to the     max, position \texttt{\detokenize{0}} becomes a descent.  Used in the \texttt{\detokenize{psi_R_add}} and     \texttt{\detokenize{psi_LR_swap1}} cases of the descent-effect theorems.

\item[\href{https://github.com/LLM4Rocq/mathcomp-eulerian/blob/main/\detokenize{psi_descent_thms.v}\#L229}{\texttt{\detokenize{rs_head_min_no_descent}}}] \emph{Lemma, line 229}. \textsf{[aux]} \\
  \texttt{\detokenize{rs_head_min_no_descent}} : dual of \texttt{\detokenize{rs_head_max_descent}} -- after     rank-shift sends head to min, position \texttt{\detokenize{0}} is no longer a descent.

\item[\href{https://github.com/LLM4Rocq/mathcomp-eulerian/blob/main/\detokenize{psi_descent_thms.v}\#L261}{\texttt{\detokenize{cmp_out_of_range}}}] \emph{Lemma, line 261}. \textsf{[aux]} \\
  \texttt{\detokenize{cmp_out_of_range}} : when \texttt{\detokenize{v}} lies outside the sorted range of \texttt{\detokenize{L}},     the comparison \texttt{\detokenize{nth L j > v}} is invariant under replacing \texttt{\detokenize{L}} by     \texttt{\detokenize{rank_shift_seq L}}.  Used at right-boundary descent positions.

\item[\href{https://github.com/LLM4Rocq/mathcomp-eulerian/blob/main/\detokenize{psi_descent_thms.v}\#L285}{\texttt{\detokenize{cmp_out_of_range_left}}}] \emph{Lemma, line 285}. \textsf{[aux]} \\
  \texttt{\detokenize{cmp_out_of_range_left}} : the \texttt{\detokenize{cmp_out_of_range}} variant with \texttt{\detokenize{v}}     on the left of \texttt{\detokenize{>}}; used at the left boundary \texttt{\detokenize{i-1}}.

\item[\href{https://github.com/LLM4Rocq/mathcomp-eulerian/blob/main/\detokenize{psi_descent_thms.v}\#L325}{\texttt{\detokenize{descent_psi_interior}}}] \emph{Lemma, line 325}. \textsf{[aux]} \\
  \texttt{\detokenize{descent_psi_interior}} : at strictly interior window positions \texttt{\detokenize{k}},     \texttt{\detokenize{psi i}} preserves the descent value (interior order is preserved by     rank-shift).

\item[\href{https://github.com/LLM4Rocq/mathcomp-eulerian/blob/main/\detokenize{psi_descent_thms.v}\#L372}{\texttt{\detokenize{descent_psi_rboundary}}}] \emph{Lemma, line 372}. \textsf{[aux]} \\
  \texttt{\detokenize{descent_psi_rboundary}} : at the right boundary \texttt{\detokenize{k = i + ws - 1}},     \texttt{\detokenize{psi i}} preserves the descent value, by \texttt{\detokenize{post_window_extremum}} +     \texttt{\detokenize{cmp_out_of_range}}.

\item[\href{https://github.com/LLM4Rocq/mathcomp-eulerian/blob/main/\detokenize{psi_descent_thms.v}\#L418}{\texttt{\detokenize{descent_psi_lboundary_R}}}] \emph{Lemma, line 418}. \textsf{[aux]} \\
  \texttt{\detokenize{descent_psi_lboundary_R}} : at the left boundary \texttt{\detokenize{i-1}}, for an R     vertex, \texttt{\detokenize{psi i}} preserves the descent value.  Uses     \texttt{\detokenize{pre_window_extremum_R}} + \texttt{\detokenize{cmp_out_of_range_left}}.

\item[\href{https://github.com/LLM4Rocq/mathcomp-eulerian/blob/main/\detokenize{psi_descent_thms.v}\#L450}{\texttt{\detokenize{descent_psi_R_add}}}] \emph{Lemma, line 450}. \textsf{[aux]} \\
  \texttt{\detokenize{descent_psi_R_add}} : R-vertex case where \texttt{\detokenize{i}} is not a descent --     applying \texttt{\detokenize{psi i}} adds \texttt{\detokenize{i}} to the descent set, leaving all other     positions unchanged.  Half of Stanley's Fact \#2 (R case).

\item[\href{https://github.com/LLM4Rocq/mathcomp-eulerian/blob/main/\detokenize{psi_descent_thms.v}\#L502}{\texttt{\detokenize{descent_psi_R_remove}}}] \emph{Lemma, line 502}. \textsf{[aux]} \\
  \texttt{\detokenize{descent_psi_R_remove}} : R-vertex case where \texttt{\detokenize{i}} is a descent --     applying \texttt{\detokenize{psi i}} removes \texttt{\detokenize{i}} from the descent set.  Other half of     Stanley's Fact \#2 (R case).

\item[\href{https://github.com/LLM4Rocq/mathcomp-eulerian/blob/main/\detokenize{psi_descent_thms.v}\#L554}{\texttt{\detokenize{descent_psi_LR_swap1}}}] \emph{Lemma, line 554}. \textsf{[aux]} \\
  \texttt{\detokenize{descent_psi_LR_swap1}} : LR-vertex case where \texttt{\detokenize{i}} is not a descent     (hence \texttt{\detokenize{i-1}} is) -- applying \texttt{\detokenize{psi i}} swaps the descent from \texttt{\detokenize{i-1}}     to \texttt{\detokenize{i}}.  Stanley's Fact \#2 (LR case, swap1).

\item[\href{https://github.com/LLM4Rocq/mathcomp-eulerian/blob/main/\detokenize{psi_descent_thms.v}\#L623}{\texttt{\detokenize{descent_psi_LR_swap2}}}] \emph{Lemma, line 623}. \textsf{[aux]} \\
  \texttt{\detokenize{descent_psi_LR_swap2}} : LR-vertex case where \texttt{\detokenize{i}} is a descent     (hence \texttt{\detokenize{i-1}} is not) -- applying \texttt{\detokenize{psi i}} swaps the descent from \texttt{\detokenize{i}}     to \texttt{\detokenize{i-1}}.  Stanley's Fact \#2 (LR case, swap2).

\end{description}

% --------------------------------------------
\section*{\texttt{psi\_descent\_v2.v}}
\addcontentsline{toc}{section}{\texttt{psi\_descent\_v2.v}}

\begin{description}
\item[\href{https://github.com/LLM4Rocq/mathcomp-eulerian/blob/main/\detokenize{psi_descent_v2.v}\#L40}{\texttt{\detokenize{has_left_child_fuel}}}] \emph{Fixpoint, line 40}. \textsf{[aux]} \\
  \texttt{\detokenize{has_left_child_fuel fuel i s}} tests, by fuel-bounded recursion on     \texttt{\detokenize{mm_pos}}, whether the vertex at in-order position \texttt{\detokenize{i}} of \texttt{\detokenize{M(s)}} has a     left child (equivalent to it being an LR vertex).

\item[\href{https://github.com/LLM4Rocq/mathcomp-eulerian/blob/main/\detokenize{psi_descent_v2.v}\#L57}{\texttt{\detokenize{has_left_child}}}] \emph{Definition, line 57}. \textsf{[aux]} \\
  \texttt{\detokenize{has_left_child i w}} : the vertex at in-order position \texttt{\detokenize{i}} of \texttt{\detokenize{M(w)}}     has a left child.  Used to dispatch the descent-effect of \texttt{\detokenize{psi_i}}     into LR (with left child) and R (right-child only) cases.

\item[\href{https://github.com/LLM4Rocq/mathcomp-eulerian/blob/main/\detokenize{psi_descent_v2.v}\#L62}{\texttt{\detokenize{has_left_child_fuel_0}}}] \emph{Lemma, line 62}. \textsf{[aux]} \\
  \texttt{\detokenize{has_left_child_fuel_0}} : the vertex at in-order position \texttt{\detokenize{0}} never     has a left child (it is the leftmost vertex).

\item[\href{https://github.com/LLM4Rocq/mathcomp-eulerian/blob/main/\detokenize{psi_descent_v2.v}\#L71}{\texttt{\detokenize{has_left_child_0}}}] \emph{Lemma, line 71}. \textsf{[aux]} \\
  \texttt{\detokenize{has_left_child_0}} : top-level corollary of \texttt{\detokenize{has_left_child_fuel_0}}     -- position \texttt{\detokenize{0}} never has a left child.

\item[\href{https://github.com/LLM4Rocq/mathcomp-eulerian/blob/main/\detokenize{psi_descent_v2.v}\#L76}{\texttt{\detokenize{has_left_child_fuel_monotone}}}] \emph{Lemma, line 76}. \textsf{[aux]} \\
  \texttt{\detokenize{has_left_child_fuel_monotone}} : independence of fuel above the input     size, the \texttt{\detokenize{has_left_child}} analogue of \texttt{\detokenize{window_size_fuel_monotone}}.

\item[\href{https://github.com/LLM4Rocq/mathcomp-eulerian/blob/main/\detokenize{psi_descent_v2.v}\#L104}{\texttt{\detokenize{has_left_child_cons}}}] \emph{Lemma, line 104}. \textsf{[aux]} \\
  \texttt{\detokenize{has_left_child_cons}} : structural unfolding of \texttt{\detokenize{has_left_child}} on     a \texttt{\detokenize{cons}} by the \texttt{\detokenize{mm_pos}} split, mirroring \texttt{\detokenize{window_size_cons}}.

\item[\href{https://github.com/LLM4Rocq/mathcomp-eulerian/blob/main/\detokenize{psi_descent_v2.v}\#L152}{\texttt{\detokenize{valid_mm}}}] \emph{Fixpoint, line 152}. \textsf{[aux]} \\
  \texttt{\detokenize{valid_mm t}} : the tree \texttt{\detokenize{t}} is a valid min-max tree, i.e. at every     \texttt{\detokenize{Node l x r}} the root sits at \texttt{\detokenize{mm_pos}} of the in-order traversal.     Specifies the trees in the image of \texttt{\detokenize{mmtree_of_seq_mm}}.

\item[\href{https://github.com/LLM4Rocq/mathcomp-eulerian/blob/main/\detokenize{psi_descent_v2.v}\#L163}{\texttt{\detokenize{valid_mm_build}}}] \emph{Lemma, line 163}. \textsf{[aux]} \\
  \texttt{\detokenize{valid_mm_build}} : every output of \texttt{\detokenize{mmtree_of_seq_mm}} satisfies     \texttt{\detokenize{valid_mm}}, the structural invariant used by \texttt{\detokenize{tree_structure}}.

\item[\href{https://github.com/LLM4Rocq/mathcomp-eulerian/blob/main/\detokenize{psi_descent_v2.v}\#L217}{\texttt{\detokenize{take_mm_eq}}}] \emph{Lemma, line 217}. \textsf{[aux]} \\
  \texttt{\detokenize{take_mm_eq}} : at the \texttt{\detokenize{mm_pos}} split, \texttt{\detokenize{take (size sl) s}} recovers the     left-subtree in-order \texttt{\detokenize{sl}}.

\item[\href{https://github.com/LLM4Rocq/mathcomp-eulerian/blob/main/\detokenize{psi_descent_v2.v}\#L222}{\texttt{\detokenize{drop_mm_eq}}}] \emph{Lemma, line 222}. \textsf{[aux]} \\
  \texttt{\detokenize{drop_mm_eq}} : dual of \texttt{\detokenize{take_mm_eq}} -- \texttt{\detokenize{drop (size sl).+1 s}} is the     right-subtree in-order \texttt{\detokenize{sr}}.

\item[\href{https://github.com/LLM4Rocq/mathcomp-eulerian/blob/main/\detokenize{psi_descent_v2.v}\#L233}{\texttt{\detokenize{ws_bridge_left}}}] \emph{Lemma, line 233}. \textsf{[aux]} \\
  \texttt{\detokenize{ws_bridge_left}} : on a valid node \texttt{\detokenize{s = sl ++ x :: sr}}, when \texttt{\detokenize{i}} is     in the left subtree, \texttt{\detokenize{window_size i s}} reduces to \texttt{\detokenize{window_size i sl}}.     Bridge between flat and structural recursion.

\item[\href{https://github.com/LLM4Rocq/mathcomp-eulerian/blob/main/\detokenize{psi_descent_v2.v}\#L247}{\texttt{\detokenize{ws_bridge_root}}}] \emph{Lemma, line 247}. \textsf{[aux]} \\
  \texttt{\detokenize{ws_bridge_root}} : at the root of the subtree, the M-window has size     \texttt{\detokenize{size s - size sl}} (root + right subtree).

\item[\href{https://github.com/LLM4Rocq/mathcomp-eulerian/blob/main/\detokenize{psi_descent_v2.v}\#L260}{\texttt{\detokenize{ws_bridge_right}}}] \emph{Lemma, line 260}. \textsf{[aux]} \\
  \texttt{\detokenize{ws_bridge_right}} : when \texttt{\detokenize{i}} is in the right subtree, \texttt{\detokenize{window_size i s}}     reduces to \texttt{\detokenize{window_size (i - size sl - 1) sr}} -- the right-subtree     descent step.

\item[\href{https://github.com/LLM4Rocq/mathcomp-eulerian/blob/main/\detokenize{psi_descent_v2.v}\#L278}{\texttt{\detokenize{wa_bridge_left}}}] \emph{Lemma, line 278}. \textsf{[aux]} \\
  \texttt{\detokenize{wa_bridge_left}} : the \texttt{\detokenize{window_at}} analogue of \texttt{\detokenize{ws_bridge_left}}.

\item[\href{https://github.com/LLM4Rocq/mathcomp-eulerian/blob/main/\detokenize{psi_descent_v2.v}\#L292}{\texttt{\detokenize{wa_bridge_root}}}] \emph{Lemma, line 292}. \textsf{[aux]} \\
  \texttt{\detokenize{wa_bridge_root}} : at the root, \texttt{\detokenize{window_at}} is the suffix     \texttt{\detokenize{drop (size sl) s}} (root + right-subtree labels).

\item[\href{https://github.com/LLM4Rocq/mathcomp-eulerian/blob/main/\detokenize{psi_descent_v2.v}\#L303}{\texttt{\detokenize{wa_bridge_right}}}] \emph{Lemma, line 303}. \textsf{[aux]} \\
  \texttt{\detokenize{wa_bridge_right}} : the \texttt{\detokenize{window_at}} analogue of \texttt{\detokenize{ws_bridge_right}}.

\item[\href{https://github.com/LLM4Rocq/mathcomp-eulerian/blob/main/\detokenize{psi_descent_v2.v}\#L322}{\texttt{\detokenize{hlc_bridge_left}}}] \emph{Lemma, line 322}. \textsf{[aux]} \\
  \texttt{\detokenize{hlc_bridge_left}} : \texttt{\detokenize{has_left_child}} structural bridge analogous to     \texttt{\detokenize{ws_bridge_left}}.

\item[\href{https://github.com/LLM4Rocq/mathcomp-eulerian/blob/main/\detokenize{psi_descent_v2.v}\#L336}{\texttt{\detokenize{hlc_bridge_root}}}] \emph{Lemma, line 336}. \textsf{[aux]} \\
  \texttt{\detokenize{hlc_bridge_root}} : at the root, \texttt{\detokenize{has_left_child}} is true iff there     is a nonempty left subtree.

\item[\href{https://github.com/LLM4Rocq/mathcomp-eulerian/blob/main/\detokenize{psi_descent_v2.v}\#L348}{\texttt{\detokenize{hlc_bridge_right}}}] \emph{Lemma, line 348}. \textsf{[aux]} \\
  \texttt{\detokenize{hlc_bridge_right}} : \texttt{\detokenize{has_left_child}} structural bridge analogous to     \texttt{\detokenize{ws_bridge_right}}.

\item[\href{https://github.com/LLM4Rocq/mathcomp-eulerian/blob/main/\detokenize{psi_descent_v2.v}\#L525}{\texttt{\detokenize{tree_props}}}] \emph{Definition, line 525}. \textsf{[aux]} \\
  \texttt{\detokenize{tree_props i w}} is the conjunction of five structural properties at     in-order position \texttt{\detokenize{i}}: post-window extremum, two pre-window-vs-window     inequalities, the descent flip, and the no-left-child extremum.     Bundled to share a single structural induction.

\item[\href{https://github.com/LLM4Rocq/mathcomp-eulerian/blob/main/\detokenize{psi_descent_v2.v}\#L563}{\texttt{\detokenize{pred_sub_add}}}] \emph{Lemma, line 563}. \textsf{[aux]} \\
  \texttt{\detokenize{pred_sub_add}} : an arithmetic identity used to align in-order     indices when descending into the right subtree.

\item[\href{https://github.com/LLM4Rocq/mathcomp-eulerian/blob/main/\detokenize{psi_descent_v2.v}\#L584}{\texttt{\detokenize{tree_structure_via_tree}}}] \emph{Lemma, line 584}. \textsf{[aux]} \\
  \texttt{\detokenize{tree_structure_via_tree}} : \texttt{\detokenize{tree_props}} holds for every position \texttt{\detokenize{i}}     of the in-order traversal of any valid min-max tree.  Proved by     structural induction on \texttt{\detokenize{t}}, dispatching on \texttt{\detokenize{i}} vs the root index.

\item[\href{https://github.com/LLM4Rocq/mathcomp-eulerian/blob/main/\detokenize{psi_descent_v2.v}\#L1076}{\texttt{\detokenize{tree_structure}}}] \emph{Lemma, line 1076}. \textsf{[aux]} \\
  \texttt{\detokenize{tree_structure}} : \texttt{\detokenize{tree_props i w}} for every uniq \texttt{\detokenize{w}}, obtained by     re-routing through \texttt{\detokenize{mmtree_of_seq_mm}} and \texttt{\detokenize{tree_structure_via_tree}}.

\item[\href{https://github.com/LLM4Rocq/mathcomp-eulerian/blob/main/\detokenize{psi_descent_v2.v}\#L1090}{\texttt{\detokenize{post_window_extremum}}}] \emph{Lemma, line 1090}. \textsf{[aux]} \\
  \texttt{\detokenize{post_window_extremum}} : the entry just past the window is either     smaller than every window value or larger than every window value     -- the post-window position is a global extremum.  First projection     of \texttt{\detokenize{tree_props}}.

\item[\href{https://github.com/LLM4Rocq/mathcomp-eulerian/blob/main/\detokenize{psi_descent_v2.v}\#L1119}{\texttt{\detokenize{pre_window_lt_max_when_min_head}}}] \emph{Lemma, line 1119}. \textsf{[aux]} \\
  \texttt{\detokenize{pre_window_lt_max_when_min_head}} : at an LR vertex whose head is the     window-min, the pre-window entry \texttt{w\_\{i-1\}} is below the window-max.     Discharges the \texttt{\detokenize{head = min}} branch of \texttt{\detokenize{exactly_one_descent_LR}}.

\item[\href{https://github.com/LLM4Rocq/mathcomp-eulerian/blob/main/\detokenize{psi_descent_v2.v}\#L1137}{\texttt{\detokenize{pre_window_gt_min_when_max_head}}}] \emph{Lemma, line 1137}. \textsf{[aux]} \\
  \texttt{\detokenize{pre_window_gt_min_when_max_head}} : dual of     \texttt{\detokenize{pre_window_lt_max_when_min_head}} -- at an LR vertex whose head is the     window-max, \texttt{w\_\{i-1\}} is above the window-min.

\item[\href{https://github.com/LLM4Rocq/mathcomp-eulerian/blob/main/\detokenize{psi_descent_v2.v}\#L1166}{\texttt{\detokenize{exactly_one_descent_LR}}}] \emph{Lemma, line 1166}. \textsf{[aux]} \\
  \texttt{\detokenize{exactly_one_descent_LR}} : at an LR vertex with nontrivial window,     exactly one of positions \texttt{\detokenize{i-1, i}} is a descent of \texttt{\detokenize{w}}.  Used by the     descent-effect theorems in \texttt{\detokenize{psi_descent_thms.v}}.

\item[\href{https://github.com/LLM4Rocq/mathcomp-eulerian/blob/main/\detokenize{psi_descent_v2.v}\#L1192}{\texttt{\detokenize{pre_window_extremum_R}}}] \emph{Lemma, line 1192}. \textsf{[aux]} \\
  \texttt{\detokenize{pre_window_extremum_R}} : at an R vertex (no left child), the entry     \texttt{w\_\{i-1\}} is itself a global extremum relative to the window.  R-case     counterpart of the LR pre-window lemmas.

\end{description}

% --------------------------------------------
\section*{\texttt{qeul.v}}
\addcontentsline{toc}{section}{\texttt{qeul.v}}

\begin{description}
\item[\href{https://github.com/LLM4Rocq/mathcomp-eulerian/blob/main/\detokenize{qeul.v}\#L41}{\texttt{\detokenize{q_eul_pol}}}] \emph{Definition, line 41}. \textsf{[Ch 6 §6.2]} \\
  \texttt{\detokenize{q_eul_pol n}} is the q-Eulerian polynomial     \texttt{A\_n(q,t) = \textbackslash\{\}sum\_(sigma in S\_\{n+1\}) q\textasciicircum\{\}(maj sigma) * t\textasciicircum\{\}(des sigma)},     the bivariate joint generating function of \texttt{\detokenize{maj}} and \texttt{\detokenize{des}} over     permutations of \texttt{\detokenize{n+1}} letters.  Stanley EC1 \textbackslash{}S\{\}1.4.

\item[\href{https://github.com/LLM4Rocq/mathcomp-eulerian/blob/main/\detokenize{qeul.v}\#L50}{\texttt{\detokenize{q_eul_pol_t1}}}] \emph{Theorem, line 50}. \textsf{[Ch 6 §6.2]} \\
  Specialization \texttt{\detokenize{t = 1}} of \texttt{\detokenize{A_n(q,t)}} recovers the q-factorial:     \texttt{\detokenize{A_n(q, 1) = [n+1]_q!}}.  Direct consequence of \texttt{\detokenize{maj_q_fact}}.

\item[\href{https://github.com/LLM4Rocq/mathcomp-eulerian/blob/main/\detokenize{qeul.v}\#L68}{\texttt{\detokenize{eul_pol}}}] \emph{Definition, line 68}. \textsf{[Ch 6 §6.2]} \\
  \texttt{\detokenize{eul_pol n}} is the classical Eulerian polynomial     \texttt{\textbackslash\{\}sum\_(k <= n) eulerian(n,k) * t\textasciicircum\{\}k}, with coefficients the Eulerian     numbers counting permutations of \texttt{\detokenize{n+1}} letters by descent count.     Stanley EC1 \textbackslash{}S\{\}1.4.

\item[\href{https://github.com/LLM4Rocq/mathcomp-eulerian/blob/main/\detokenize{qeul.v}\#L74}{\texttt{\detokenize{sum_des_eq_eul_pol}}}] \emph{Lemma, line 74}. \textsf{[aux]} \\
  The descent generating function expressed as the Eulerian polynomial:     \texttt{\textbackslash\{\}sum\_(sigma in S\_\{n+1\}) X\textasciicircum\{\}(des sigma) = eul\_pol n}.  Obtained by     partitioning permutations according to their descent count.

\item[\href{https://github.com/LLM4Rocq/mathcomp-eulerian/blob/main/\detokenize{qeul.v}\#L96}{\texttt{\detokenize{q1_subst}}}] \emph{Definition, line 96}. \textsf{[Ch 6 §6.2]} \\
  \texttt{\detokenize{q1_subst}} is the substitution \texttt{\detokenize{q := 1}} on a bivariate polynomial in     \texttt{\{poly \{poly int\}\}}, implemented by evaluating each inner polynomial     coefficient at \texttt{\detokenize{1}}.

\item[\href{https://github.com/LLM4Rocq/mathcomp-eulerian/blob/main/\detokenize{qeul.v}\#L102}{\texttt{\detokenize{q_eul_pol_q1}}}] \emph{Theorem, line 102}. \textsf{[Ch 6 §6.2]} \\
  Specialization \texttt{\detokenize{q = 1}} of \texttt{\detokenize{A_n(q,t)}} recovers the classical Eulerian     polynomial: \texttt{\detokenize{A_n(1, t) = eul_pol n}}.  Each \texttt{\detokenize{q^(maj sigma)}} becomes \texttt{\detokenize{1}},     leaving the descent generating function.

\end{description}

% --------------------------------------------
\section*{\texttt{qfact.v}}
\addcontentsline{toc}{section}{\texttt{qfact.v}}

\begin{description}
\item[\href{https://github.com/LLM4Rocq/mathcomp-eulerian/blob/main/\detokenize{qfact.v}\#L28}{\texttt{\detokenize{q_int}}}] \emph{Definition, line 28}. \textsf{[Ch 6 §6.1]} \\
  \texttt{\detokenize{q_int n}} is the q-integer \texttt{\detokenize{n}}\_q = 1 + q + q\textasciicircum{}2 + ... + q\textasciicircum{}(n-1),     represented as the polynomial \texttt{\textbackslash\{\}sum\_(i < n) 'X\textasciicircum\{\}i} in \texttt{\{poly int\}}.     Stanley EC1 \textbackslash{}S\{\}1.4.

\item[\href{https://github.com/LLM4Rocq/mathcomp-eulerian/blob/main/\detokenize{qfact.v}\#L32}{\texttt{\detokenize{q_fact}}}] \emph{Definition, line 32}. \textsf{[Ch 6 §6.1]} \\
  \texttt{\detokenize{q_fact n}} is the q-factorial \texttt{\detokenize{n+1}}\_q! = \texttt{\detokenize{1}}\_q \texttt{\detokenize{2}}\_q ... \texttt{\detokenize{n+1}}\_q,     the product of \texttt{\detokenize{q_int k.+1}} for \texttt{\detokenize{k < n.+1}}.  Stanley EC1 \textbackslash{}S\{\}1.4.

\item[\href{https://github.com/LLM4Rocq/mathcomp-eulerian/blob/main/\detokenize{qfact.v}\#L41}{\texttt{\detokenize{sum_rev_X}}}] \emph{Lemma, line 41}. \textsf{[aux]} \\
  Reindexing identity: summing \texttt{\detokenize{X^(n - p)}} over \texttt{\detokenize{p < n.+1}} equals     \texttt{\detokenize{q_int n.+1}}.  Used to recognize the q-integer factor produced by     expanding \texttt{\detokenize{exprD}} over \texttt{\detokenize{val p}} in the \texttt{\detokenize{inv_q_fact}} inductive step.

\item[\href{https://github.com/LLM4Rocq/mathcomp-eulerian/blob/main/\detokenize{qfact.v}\#L70}{\texttt{\detokenize{is_inv_lift_pair}}}] \emph{Lemma, line 70}. \textsf{[aux]} \\
  Inversion correspondence on the "non-\texttt{\detokenize{p}} rows": a pair     \texttt{\detokenize{(lift p j1, lift p j2)}} is an inversion of \texttt{\detokenize{insert_max_perm t p}}     iff \texttt{\detokenize{(j1, j2)}} is an inversion of \texttt{\detokenize{t}}.

\item[\href{https://github.com/LLM4Rocq/mathcomp-eulerian/blob/main/\detokenize{qfact.v}\#L85}{\texttt{\detokenize{is_inv_lift_p}}}] \emph{Lemma, line 85}. \textsf{[aux]} \\
  Pairs \texttt{\detokenize{(lift p k, p)}} are never inversions of \texttt{\detokenize{insert_max_perm t p}}:     \texttt{\$\textbackslash\{\}sigma\$ p = ord\_max} is the maximum value, so it cannot be exceeded by any     earlier image \texttt{\$\textbackslash\{\}sigma\$ (lift p k)}.

\item[\href{https://github.com/LLM4Rocq/mathcomp-eulerian/blob/main/\detokenize{qfact.v}\#L98}{\texttt{\detokenize{sum_geq_p}}}] \emph{Lemma, line 98}. \textsf{[aux]} \\
  Counts the elements \texttt{\detokenize{k : 'I_n.+1}} with \texttt{\detokenize{p <= val k}}: there are     exactly \texttt{\detokenize{n.+1 - p}} of them.  Used to evaluate the \texttt{\detokenize{p}}-row contribution     in \texttt{\detokenize{inv_insert_max}}.

\item[\href{https://github.com/LLM4Rocq/mathcomp-eulerian/blob/main/\detokenize{qfact.v}\#L122}{\texttt{\detokenize{inv_via_sum}}}] \emph{Lemma, line 122}. \textsf{[aux]} \\
  Reformulates \texttt{\detokenize{inv s}} as the unconditional double sum     \texttt{\textbackslash\{\}sum\_i \textbackslash\{\}sum\_j is\_inv s i j}, unfolding the \texttt{\detokenize{#|...|}} counting form.

\item[\href{https://github.com/LLM4Rocq/mathcomp-eulerian/blob/main/\detokenize{qfact.v}\#L139}{\texttt{\detokenize{inv_insert_max}}}] \emph{Lemma, line 139}. \textsf{[aux]} \\
  Headline arithmetic identity for the \texttt{\detokenize{insert_max_perm}} bijection:     \texttt{\detokenize{inv (insert_max_perm t p) = inv t + (n.+1 - val p)}}.  The new image     \texttt{\detokenize{ord_max}} at position \texttt{\detokenize{p}} contributes exactly \texttt{\detokenize{n.+1 - val p}} new     inversions; the rest match \texttt{\detokenize{inv t}} under the lift.

\item[\href{https://github.com/LLM4Rocq/mathcomp-eulerian/blob/main/\detokenize{qfact.v}\#L204}{\texttt{\detokenize{inv_q_fact}}}] \emph{Theorem, line 204}. \textsf{[Ch 6 §6.1]} \\
  Headline q-factorial identity for the inversion statistic     (Stanley EC1 \textbackslash{}S\{\}1.4, Cor 1.3.10):     \texttt{\textbackslash\{\}sum\_(sigma in S\_\{n+1\}) X\textasciicircum\{\}(inv sigma) = [n+1]\_q!}.  Proved by induction     on \texttt{\detokenize{n}} using the \texttt{\detokenize{insert_max_perm}} bijection and \texttt{\detokenize{inv_insert_max}}.

\item[\href{https://github.com/LLM4Rocq/mathcomp-eulerian/blob/main/\detokenize{qfact.v}\#L240}{\texttt{\detokenize{maj_q_fact}}}] \emph{Theorem, line 240}. \textsf{[Ch 6 §6.1]} \\
  Headline q-factorial identity for the major-index statistic     (Stanley EC1 \textbackslash{}S\{\}1.4):     \texttt{\textbackslash\{\}sum\_(sigma in S\_\{n+1\}) X\textasciicircum\{\}(maj sigma) = [n+1]\_q!}.  One-line transfer     from \texttt{\detokenize{inv_q_fact}} via the Foata bijection \texttt{\detokenize{foata_perm}}, which sends     \texttt{\detokenize{inv}} to \texttt{\detokenize{maj}}.

\end{description}

% --------------------------------------------
\section*{\texttt{reflection.v}}
\addcontentsline{toc}{section}{\texttt{reflection.v}}

\begin{description}
\item[\href{https://github.com/LLM4Rocq/mathcomp-eulerian/blob/main/\detokenize{reflection.v}\#L28}{\texttt{\detokenize{euler}}}] \emph{Definition, line 28}. \textsf{[aux]} \\
  (no docstring)

\item[\href{https://github.com/LLM4Rocq/mathcomp-eulerian/blob/main/\detokenize{reflection.v}\#L32}{\texttt{\detokenize{eulerA}}}] \emph{Definition, line 32}. \textsf{[aux]} \\
  (no docstring)

\item[\href{https://github.com/LLM4Rocq/mathcomp-eulerian/blob/main/\detokenize{reflection.v}\#L39}{\texttt{\detokenize{eulerA_0}}}] \emph{Lemma, line 39}. \textsf{[aux]} \\
  (no docstring)

\item[\href{https://github.com/LLM4Rocq/mathcomp-eulerian/blob/main/\detokenize{reflection.v}\#L43}{\texttt{\detokenize{alt_desc_set_0}}}] \emph{Lemma, line 43}. \textsf{[aux]} \\
  (no docstring)

\item[\href{https://github.com/LLM4Rocq/mathcomp-eulerian/blob/main/\detokenize{reflection.v}\#L47}{\texttt{\detokenize{euler_0}}}] \emph{Lemma, line 47}. \textsf{[aux]} \\
  (no docstring)

\item[\href{https://github.com/LLM4Rocq/mathcomp-eulerian/blob/main/\detokenize{reflection.v}\#L50}{\texttt{\detokenize{eulerA_1}}}] \emph{Lemma, line 50}. \textsf{[aux]} \\
  (no docstring)

\item[\href{https://github.com/LLM4Rocq/mathcomp-eulerian/blob/main/\detokenize{reflection.v}\#L55}{\texttt{\detokenize{alt_desc_set_1}}}] \emph{Lemma, line 55}. \textsf{[aux]} \\
  (no docstring)

\item[\href{https://github.com/LLM4Rocq/mathcomp-eulerian/blob/main/\detokenize{reflection.v}\#L62}{\texttt{\detokenize{euler_1}}}] \emph{Lemma, line 62}. \textsf{[aux]} \\
  (no docstring)

\item[\href{https://github.com/LLM4Rocq/mathcomp-eulerian/blob/main/\detokenize{reflection.v}\#L65}{\texttt{\detokenize{eulerA_2}}}] \emph{Lemma, line 65}. \textsf{[aux]} \\
  (no docstring)

\item[\href{https://github.com/LLM4Rocq/mathcomp-eulerian/blob/main/\detokenize{reflection.v}\#L78}{\texttt{\detokenize{descent_set_insert_max_ord0}}}] \emph{Lemma, line 78}. \textsf{[aux]} \\
  (no docstring)

\item[\href{https://github.com/LLM4Rocq/mathcomp-eulerian/blob/main/\detokenize{reflection.v}\#L101}{\texttt{\detokenize{descent_set_insert_max_ord_max}}}] \emph{Lemma, line 101}. \textsf{[aux]} \\
  (no docstring)

\item[\href{https://github.com/LLM4Rocq/mathcomp-eulerian/blob/main/\detokenize{reflection.v}\#L151}{\texttt{\detokenize{descent_set_insert_max_interior}}}] \emph{Lemma, line 151}. \textsf{[aux]} \\
  (no docstring)

\item[\href{https://github.com/LLM4Rocq/mathcomp-eulerian/blob/main/\detokenize{reflection.v}\#L215}{\texttt{\detokenize{widen_inj}}}] \emph{Lemma, line 215}. \textsf{[aux]} \\
  (no docstring)

\item[\href{https://github.com/LLM4Rocq/mathcomp-eulerian/blob/main/\detokenize{reflection.v}\#L223}{\texttt{\detokenize{leqj_n1}}}] \emph{Definition, line 223}. \textsf{[aux]} \\
  (no docstring)

\item[\href{https://github.com/LLM4Rocq/mathcomp-eulerian/blob/main/\detokenize{reflection.v}\#L226}{\texttt{\detokenize{embed_left}}}] \emph{Definition, line 226}. \textsf{[aux]} \\
  (no docstring)

\item[\href{https://github.com/LLM4Rocq/mathcomp-eulerian/blob/main/\detokenize{reflection.v}\#L228}{\texttt{\detokenize{embed_left_inj}}}] \emph{Lemma, line 228}. \textsf{[aux]} \\
  (no docstring)

\item[\href{https://github.com/LLM4Rocq/mathcomp-eulerian/blob/main/\detokenize{reflection.v}\#L232}{\texttt{\detokenize{image_left}}}] \emph{Definition, line 232}. \textsf{[aux]} \\
  (no docstring)

\item[\href{https://github.com/LLM4Rocq/mathcomp-eulerian/blob/main/\detokenize{reflection.v}\#L235}{\texttt{\detokenize{t_embed_inj}}}] \emph{Lemma, line 235}. \textsf{[aux]} \\
  (no docstring)

\item[\href{https://github.com/LLM4Rocq/mathcomp-eulerian/blob/main/\detokenize{reflection.v}\#L238}{\texttt{\detokenize{card_image_left}}}] \emph{Lemma, line 238}. \textsf{[aux]} \\
  (no docstring)

\item[\href{https://github.com/LLM4Rocq/mathcomp-eulerian/blob/main/\detokenize{reflection.v}\#L244}{\texttt{\detokenize{mem_image_left}}}] \emph{Lemma, line 244}. \textsf{[aux]} \\
  (no docstring)

\item[\href{https://github.com/LLM4Rocq/mathcomp-eulerian/blob/main/\detokenize{reflection.v}\#L269}{\texttt{\detokenize{rank_left}}}] \emph{Definition, line 269}. \textsf{[aux]} \\
  (no docstring)

\item[\href{https://github.com/LLM4Rocq/mathcomp-eulerian/blob/main/\detokenize{reflection.v}\#L274}{\texttt{\detokenize{cast_to_j}}}] \emph{Definition, line 274}. \textsf{[aux]} \\
  (no docstring)

\item[\href{https://github.com/LLM4Rocq/mathcomp-eulerian/blob/main/\detokenize{reflection.v}\#L277}{\texttt{\detokenize{cast_to_j_inj}}}] \emph{Lemma, line 277}. \textsf{[aux]} \\
  (no docstring)

\item[\href{https://github.com/LLM4Rocq/mathcomp-eulerian/blob/main/\detokenize{reflection.v}\#L281}{\texttt{\detokenize{perm_left_fun}}}] \emph{Definition, line 281}. \textsf{[aux]} \\
  (no docstring)

\item[\href{https://github.com/LLM4Rocq/mathcomp-eulerian/blob/main/\detokenize{reflection.v}\#L285}{\texttt{\detokenize{perm_left_fun_inj}}}] \emph{Lemma, line 285}. \textsf{[aux]} \\
  (no docstring)

\item[\href{https://github.com/LLM4Rocq/mathcomp-eulerian/blob/main/\detokenize{reflection.v}\#L294}{\texttt{\detokenize{perm_left}}}] \emph{Definition, line 294}. \textsf{[aux]} \\
  (no docstring)

\item[\href{https://github.com/LLM4Rocq/mathcomp-eulerian/blob/main/\detokenize{reflection.v}\#L297}{\texttt{\detokenize{perm_leftE}}}] \emph{Lemma, line 297}. \textsf{[aux]} \\
  (no docstring)

\item[\href{https://github.com/LLM4Rocq/mathcomp-eulerian/blob/main/\detokenize{reflection.v}\#L302}{\texttt{\detokenize{enum_val_perm_left}}}] \emph{Lemma, line 302}. \textsf{[aux]} \\
  (no docstring)

\item[\href{https://github.com/LLM4Rocq/mathcomp-eulerian/blob/main/\detokenize{reflection.v}\#L317}{\texttt{\detokenize{perm_left_witness_indep}}}] \emph{Lemma, line 317}. \textsf{[aux]} \\
  (no docstring)

\item[\href{https://github.com/LLM4Rocq/mathcomp-eulerian/blob/main/\detokenize{reflection.v}\#L331}{\texttt{\detokenize{image_left_witness_pos}}}] \emph{Lemma, line 331}. \textsf{[aux]} \\
  (no docstring)

\item[\href{https://github.com/LLM4Rocq/mathcomp-eulerian/blob/main/\detokenize{reflection.v}\#L347}{\texttt{\detokenize{embed_right}}}] \emph{Definition, line 347}. \textsf{[aux]} \\
  (no docstring)

\item[\href{https://github.com/LLM4Rocq/mathcomp-eulerian/blob/main/\detokenize{reflection.v}\#L350}{\texttt{\detokenize{embed_right_inj}}}] \emph{Lemma, line 350}. \textsf{[aux]} \\
  (no docstring)

\item[\href{https://github.com/LLM4Rocq/mathcomp-eulerian/blob/main/\detokenize{reflection.v}\#L356}{\texttt{\detokenize{image_right}}}] \emph{Definition, line 356}. \textsf{[aux]} \\
  (no docstring)

\item[\href{https://github.com/LLM4Rocq/mathcomp-eulerian/blob/main/\detokenize{reflection.v}\#L359}{\texttt{\detokenize{t_embed_right_inj}}}] \emph{Lemma, line 359}. \textsf{[aux]} \\
  (no docstring)

\item[\href{https://github.com/LLM4Rocq/mathcomp-eulerian/blob/main/\detokenize{reflection.v}\#L362}{\texttt{\detokenize{card_image_right}}}] \emph{Lemma, line 362}. \textsf{[aux]} \\
  (no docstring)

\item[\href{https://github.com/LLM4Rocq/mathcomp-eulerian/blob/main/\detokenize{reflection.v}\#L367}{\texttt{\detokenize{mem_image_right}}}] \emph{Lemma, line 367}. \textsf{[aux]} \\
  (no docstring)

\item[\href{https://github.com/LLM4Rocq/mathcomp-eulerian/blob/main/\detokenize{reflection.v}\#L380}{\texttt{\detokenize{cast_to_subj}}}] \emph{Definition, line 380}. \textsf{[aux]} \\
  (no docstring)

\item[\href{https://github.com/LLM4Rocq/mathcomp-eulerian/blob/main/\detokenize{reflection.v}\#L383}{\texttt{\detokenize{cast_to_subj_inj}}}] \emph{Lemma, line 383}. \textsf{[aux]} \\
  (no docstring)

\item[\href{https://github.com/LLM4Rocq/mathcomp-eulerian/blob/main/\detokenize{reflection.v}\#L386}{\texttt{\detokenize{perm_right_fun}}}] \emph{Definition, line 386}. \textsf{[aux]} \\
  (no docstring)

\item[\href{https://github.com/LLM4Rocq/mathcomp-eulerian/blob/main/\detokenize{reflection.v}\#L390}{\texttt{\detokenize{perm_right_fun_inj}}}] \emph{Lemma, line 390}. \textsf{[aux]} \\
  (no docstring)

\item[\href{https://github.com/LLM4Rocq/mathcomp-eulerian/blob/main/\detokenize{reflection.v}\#L399}{\texttt{\detokenize{perm_right}}}] \emph{Definition, line 399}. \textsf{[aux]} \\
  (no docstring)

\item[\href{https://github.com/LLM4Rocq/mathcomp-eulerian/blob/main/\detokenize{reflection.v}\#L402}{\texttt{\detokenize{perm_rightE}}}] \emph{Lemma, line 402}. \textsf{[aux]} \\
  (no docstring)

\item[\href{https://github.com/LLM4Rocq/mathcomp-eulerian/blob/main/\detokenize{reflection.v}\#L407}{\texttt{\detokenize{enum_val_perm_right}}}] \emph{Lemma, line 407}. \textsf{[aux]} \\
  (no docstring)

\item[\href{https://github.com/LLM4Rocq/mathcomp-eulerian/blob/main/\detokenize{reflection.v}\#L418}{\texttt{\detokenize{perm_right_witness_indep}}}] \emph{Lemma, line 418}. \textsf{[aux]} \\
  (no docstring)

\item[\href{https://github.com/LLM4Rocq/mathcomp-eulerian/blob/main/\detokenize{reflection.v}\#L435}{\texttt{\detokenize{sorted_ltn_val_enum}}}] \emph{Lemma, line 435}. \textsf{[aux]} \\
  (no docstring)

\item[\href{https://github.com/LLM4Rocq/mathcomp-eulerian/blob/main/\detokenize{reflection.v}\#L448}{\texttt{\detokenize{enum_val_ltn}}}] \emph{Lemma, line 448}. \textsf{[aux]} \\
  (no docstring)

\item[\href{https://github.com/LLM4Rocq/mathcomp-eulerian/blob/main/\detokenize{reflection.v}\#L477}{\texttt{\detokenize{enum_rank_in_val_ltn}}}] \emph{Lemma, line 477}. \textsf{[aux]} \\
  (no docstring)

\item[\href{https://github.com/LLM4Rocq/mathcomp-eulerian/blob/main/\detokenize{reflection.v}\#L493}{\texttt{\detokenize{perm_left_lt_iff}}}] \emph{Lemma, line 493}. \textsf{[aux]} \\
  (no docstring)

\item[\href{https://github.com/LLM4Rocq/mathcomp-eulerian/blob/main/\detokenize{reflection.v}\#L506}{\texttt{\detokenize{perm_right_lt_iff}}}] \emph{Lemma, line 506}. \textsf{[aux]} \\
  (no docstring)

\item[\href{https://github.com/LLM4Rocq/mathcomp-eulerian/blob/main/\detokenize{reflection.v}\#L520}{\texttt{\detokenize{leqj_pred_n}}}] \emph{Lemma, line 520}. \textsf{[aux]} \\
  (no docstring)

\item[\href{https://github.com/LLM4Rocq/mathcomp-eulerian/blob/main/\detokenize{reflection.v}\#L546}{\texttt{\detokenize{embed_desc_left}}}] \emph{Definition, line 546}. \textsf{[aux]} \\
  (no docstring)

\item[\href{https://github.com/LLM4Rocq/mathcomp-eulerian/blob/main/\detokenize{reflection.v}\#L548}{\texttt{\detokenize{is_descent_perm_left}}}] \emph{Lemma, line 548}. \textsf{[aux]} \\
  (no docstring)

\item[\href{https://github.com/LLM4Rocq/mathcomp-eulerian/blob/main/\detokenize{reflection.v}\#L573}{\texttt{\detokenize{sub_succ}}}] \emph{Lemma, line 573}. \textsf{[aux]} \\
  (no docstring)

\item[\href{https://github.com/LLM4Rocq/mathcomp-eulerian/blob/main/\detokenize{reflection.v}\#L581}{\texttt{\detokenize{j_plus_lt_n}}}] \emph{Lemma, line 581}. \textsf{[aux]} \\
  (no docstring)

\item[\href{https://github.com/LLM4Rocq/mathcomp-eulerian/blob/main/\detokenize{reflection.v}\#L587}{\texttt{\detokenize{embed_desc_right}}}] \emph{Definition, line 587}. \textsf{[aux]} \\
  (no docstring)

\item[\href{https://github.com/LLM4Rocq/mathcomp-eulerian/blob/main/\detokenize{reflection.v}\#L589}{\texttt{\detokenize{is_descent_perm_right}}}] \emph{Lemma, line 589}. \textsf{[aux]} \\
  (no docstring)

\item[\href{https://github.com/LLM4Rocq/mathcomp-eulerian/blob/main/\detokenize{reflection.v}\#L614}{\texttt{\detokenize{image_left_right_disjoint}}}] \emph{Lemma, line 614}. \textsf{[aux]} \\
  (no docstring)

\item[\href{https://github.com/LLM4Rocq/mathcomp-eulerian/blob/main/\detokenize{reflection.v}\#L628}{\texttt{\detokenize{image_left_right_cover}}}] \emph{Lemma, line 628}. \textsf{[aux]} \\
  (no docstring)

\item[\href{https://github.com/LLM4Rocq/mathcomp-eulerian/blob/main/\detokenize{reflection.v}\#L640}{\texttt{\detokenize{image_left_card_plus_right}}}] \emph{Lemma, line 640}. \textsf{[aux]} \\
  (no docstring)

\item[\href{https://github.com/LLM4Rocq/mathcomp-eulerian/blob/main/\detokenize{reflection.v}\#L722}{\texttt{\detokenize{descent_left_of_t}}}] \emph{Lemma, line 722}. \textsf{[aux]} \\
  (no docstring)

\item[\href{https://github.com/LLM4Rocq/mathcomp-eulerian/blob/main/\detokenize{reflection.v}\#L739}{\texttt{\detokenize{descent_right_of_t}}}] \emph{Lemma, line 739}. \textsf{[aux]} \\
  (no docstring)

\item[\href{https://github.com/LLM4Rocq/mathcomp-eulerian/blob/main/\detokenize{reflection.v}\#L770}{\texttt{\detokenize{descent_left_part}}}] \emph{Definition, line 770}. \textsf{[aux]} \\
  (no docstring)

\item[\href{https://github.com/LLM4Rocq/mathcomp-eulerian/blob/main/\detokenize{reflection.v}\#L773}{\texttt{\detokenize{descent_boundary_part}}}] \emph{Definition, line 773}. \textsf{[aux]} \\
  (no docstring)

\item[\href{https://github.com/LLM4Rocq/mathcomp-eulerian/blob/main/\detokenize{reflection.v}\#L776}{\texttt{\detokenize{descent_right_part}}}] \emph{Definition, line 776}. \textsf{[aux]} \\
  (no docstring)

\item[\href{https://github.com/LLM4Rocq/mathcomp-eulerian/blob/main/\detokenize{reflection.v}\#L779}{\texttt{\detokenize{descent_set_decomp_partition}}}] \emph{Lemma, line 779}. \textsf{[aux]} \\
  (no docstring)

\item[\href{https://github.com/LLM4Rocq/mathcomp-eulerian/blob/main/\detokenize{reflection.v}\#L788}{\texttt{\detokenize{descent_left_boundary_disjoint}}}] \emph{Lemma, line 788}. \textsf{[aux]} \\
  (no docstring)

\item[\href{https://github.com/LLM4Rocq/mathcomp-eulerian/blob/main/\detokenize{reflection.v}\#L796}{\texttt{\detokenize{descent_left_right_disjoint}}}] \emph{Lemma, line 796}. \textsf{[aux]} \\
  (no docstring)

\item[\href{https://github.com/LLM4Rocq/mathcomp-eulerian/blob/main/\detokenize{reflection.v}\#L804}{\texttt{\detokenize{descent_boundary_right_disjoint}}}] \emph{Lemma, line 804}. \textsf{[aux]} \\
  (no docstring)

\item[\href{https://github.com/LLM4Rocq/mathcomp-eulerian/blob/main/\detokenize{reflection.v}\#L813}{\texttt{\detokenize{descent_left_part_image}}}] \emph{Lemma, line 813}. \textsf{[aux]} \\
  (no docstring)

\item[\href{https://github.com/LLM4Rocq/mathcomp-eulerian/blob/main/\detokenize{reflection.v}\#L841}{\texttt{\detokenize{descent_right_part_R}}}] \emph{Definition, line 841}. \textsf{[aux]} \\
  (no docstring)

\item[\href{https://github.com/LLM4Rocq/mathcomp-eulerian/blob/main/\detokenize{reflection.v}\#L844}{\texttt{\detokenize{descent_right_part_R_image}}}] \emph{Lemma, line 844}. \textsf{[aux]} \\
  (no docstring)

\item[\href{https://github.com/LLM4Rocq/mathcomp-eulerian/blob/main/\detokenize{reflection.v}\#L891}{\texttt{\detokenize{cardCL_eq}}}] \emph{Lemma, line 891}. \textsf{[aux]} \\
  (no docstring)

\item[\href{https://github.com/LLM4Rocq/mathcomp-eulerian/blob/main/\detokenize{reflection.v}\#L896}{\texttt{\detokenize{castL}}}] \emph{Definition, line 896}. \textsf{[aux]} \\
  (no docstring)

\item[\href{https://github.com/LLM4Rocq/mathcomp-eulerian/blob/main/\detokenize{reflection.v}\#L897}{\texttt{\detokenize{castR}}}] \emph{Definition, line 897}. \textsf{[aux]} \\
  (no docstring)

\item[\href{https://github.com/LLM4Rocq/mathcomp-eulerian/blob/main/\detokenize{reflection.v}\#L899}{\texttt{\detokenize{leqj}}}] \emph{Definition, line 899}. \textsf{[aux]} \\
  (no docstring)

\item[\href{https://github.com/LLM4Rocq/mathcomp-eulerian/blob/main/\detokenize{reflection.v}\#L902}{\texttt{\detokenize{split_pos}}}] \emph{Definition, line 902}. \textsf{[aux]} \\
  (no docstring)

\item[\href{https://github.com/LLM4Rocq/mathcomp-eulerian/blob/main/\detokenize{reflection.v}\#L908}{\texttt{\detokenize{assemble_fun}}}] \emph{Definition, line 908}. \textsf{[aux]} \\
  (no docstring)

\item[\href{https://github.com/LLM4Rocq/mathcomp-eulerian/blob/main/\detokenize{reflection.v}\#L914}{\texttt{\detokenize{split_pos_inj}}}] \emph{Lemma, line 914}. \textsf{[aux]} \\
  (no docstring)

\item[\href{https://github.com/LLM4Rocq/mathcomp-eulerian/blob/main/\detokenize{reflection.v}\#L928}{\texttt{\detokenize{assemble_fun_inj}}}] \emph{Lemma, line 928}. \textsf{[aux]} \\
  (no docstring)

\item[\href{https://github.com/LLM4Rocq/mathcomp-eulerian/blob/main/\detokenize{reflection.v}\#L950}{\texttt{\detokenize{assemble_perm}}}] \emph{Definition, line 950}. \textsf{[aux]} \\
  (no docstring)

\item[\href{https://github.com/LLM4Rocq/mathcomp-eulerian/blob/main/\detokenize{reflection.v}\#L952}{\texttt{\detokenize{assemble_permE}}}] \emph{Lemma, line 952}. \textsf{[aux]} \\
  (no docstring)

\item[\href{https://github.com/LLM4Rocq/mathcomp-eulerian/blob/main/\detokenize{reflection.v}\#L969}{\texttt{\detokenize{assemble_left}}}] \emph{Lemma, line 969}. \textsf{[aux]} \\
  (no docstring)

\item[\href{https://github.com/LLM4Rocq/mathcomp-eulerian/blob/main/\detokenize{reflection.v}\#L987}{\texttt{\detokenize{assemble_right}}}] \emph{Lemma, line 987}. \textsf{[aux]} \\
  (no docstring)

\item[\href{https://github.com/LLM4Rocq/mathcomp-eulerian/blob/main/\detokenize{reflection.v}\#L1007}{\texttt{\detokenize{assemble_image_left}}}] \emph{Lemma, line 1007}. \textsf{[aux]} \\
  (no docstring)

\item[\href{https://github.com/LLM4Rocq/mathcomp-eulerian/blob/main/\detokenize{reflection.v}\#L1024}{\texttt{\detokenize{assemble_image_right}}}] \emph{Lemma, line 1024}. \textsf{[aux]} \\
  (no docstring)

\item[\href{https://github.com/LLM4Rocq/mathcomp-eulerian/blob/main/\detokenize{reflection.v}\#L1053}{\texttt{\detokenize{assemble_perm_left}}}] \emph{Lemma, line 1053}. \textsf{[aux]} \\
  (no docstring)

\item[\href{https://github.com/LLM4Rocq/mathcomp-eulerian/blob/main/\detokenize{reflection.v}\#L1082}{\texttt{\detokenize{assemble_perm_right}}}] \emph{Lemma, line 1082}. \textsf{[aux]} \\
  (no docstring)

\item[\href{https://github.com/LLM4Rocq/mathcomp-eulerian/blob/main/\detokenize{reflection.v}\#L1129}{\texttt{\detokenize{card_image_left_eq}}}] \emph{Definition, line 1129}. \textsf{[aux]} \\
  (no docstring)

\item[\href{https://github.com/LLM4Rocq/mathcomp-eulerian/blob/main/\detokenize{reflection.v}\#L1135}{\texttt{\detokenize{assemble_decomp_inverse_left}}}] \emph{Lemma, line 1135}. \textsf{[aux]} \\
  (no docstring)

\item[\href{https://github.com/LLM4Rocq/mathcomp-eulerian/blob/main/\detokenize{reflection.v}\#L1154}{\texttt{\detokenize{image_right_eq_compl}}}] \emph{Lemma, line 1154}. \textsf{[aux]} \\
  (no docstring)

\item[\href{https://github.com/LLM4Rocq/mathcomp-eulerian/blob/main/\detokenize{reflection.v}\#L1178}{\texttt{\detokenize{assemble_decomp_inverse_right}}}] \emph{Lemma, line 1178}. \textsf{[aux]} \\
  (no docstring)

\item[\href{https://github.com/LLM4Rocq/mathcomp-eulerian/blob/main/\detokenize{reflection.v}\#L1206}{\texttt{\detokenize{split_embed}}}] \emph{Lemma, line 1206}. \textsf{[aux]} \\
  (no docstring)

\item[\href{https://github.com/LLM4Rocq/mathcomp-eulerian/blob/main/\detokenize{reflection.v}\#L1220}{\texttt{\detokenize{assemble_decomp_inverse}}}] \emph{Lemma, line 1220}. \textsf{[aux]} \\
  (no docstring)

\item[\href{https://github.com/LLM4Rocq/mathcomp-eulerian/blob/main/\detokenize{reflection.v}\#L1246}{\texttt{\detokenize{assemble_decomp_unique}}}] \emph{Lemma, line 1246}. \textsf{[aux]} \\
  (no docstring)

\item[\href{https://github.com/LLM4Rocq/mathcomp-eulerian/blob/main/\detokenize{reflection.v}\#L1262}{\texttt{\detokenize{eulerA_S2}}}] \emph{Lemma, line 1262}. \textsf{[aux]} \\
  (no docstring)

\item[\href{https://github.com/LLM4Rocq/mathcomp-eulerian/blob/main/\detokenize{reflection.v}\#L1267}{\texttt{\detokenize{two_eulerA_split}}}] \emph{Lemma, line 1267}. \textsf{[aux]} \\
  (no docstring)

\item[\href{https://github.com/LLM4Rocq/mathcomp-eulerian/blob/main/\detokenize{reflection.v}\#L1283}{\texttt{\detokenize{beta_eq_pair_sum}}}] \emph{Lemma, line 1283}. \textsf{[aux]} \\
  (no docstring)

\item[\href{https://github.com/LLM4Rocq/mathcomp-eulerian/blob/main/\detokenize{reflection.v}\#L1294}{\texttt{\detokenize{beta_eq_double_sum}}}] \emph{Lemma, line 1294}. \textsf{[aux]} \\
  (no docstring)

\item[\href{https://github.com/LLM4Rocq/mathcomp-eulerian/blob/main/\detokenize{reflection.v}\#L1306}{\texttt{\detokenize{beta_eq_triple_split}}}] \emph{Lemma, line 1306}. \textsf{[aux]} \\
  (no docstring)

\item[\href{https://github.com/LLM4Rocq/mathcomp-eulerian/blob/main/\detokenize{reflection.v}\#L1328}{\texttt{\detokenize{sum_partition_image_left}}}] \emph{Lemma, line 1328}. \textsf{[aux]} \\
  (no docstring)

\item[\href{https://github.com/LLM4Rocq/mathcomp-eulerian/blob/main/\detokenize{reflection.v}\#L1340}{\texttt{\detokenize{assemble_perm_pirr}}}] \emph{Lemma, line 1340}. \textsf{[aux]} \\
  (no docstring)

\item[\href{https://github.com/LLM4Rocq/mathcomp-eulerian/blob/main/\detokenize{reflection.v}\#L1347}{\texttt{\detokenize{assemble_decomp_inverse_gen}}}] \emph{Lemma, line 1347}. \textsf{[aux]} \\
  (no docstring)

\item[\href{https://github.com/LLM4Rocq/mathcomp-eulerian/blob/main/\detokenize{reflection.v}\#L1359}{\texttt{\detokenize{sub_gt0_n_j}}}] \emph{Lemma, line 1359}. \textsf{[aux]} \\
  (no docstring)

\item[\href{https://github.com/LLM4Rocq/mathcomp-eulerian/blob/main/\detokenize{reflection.v}\#L1375}{\texttt{\detokenize{fwdAss}}}] \emph{Definition, line 1375}. \textsf{[aux]} \\
  (no docstring)

\item[\href{https://github.com/LLM4Rocq/mathcomp-eulerian/blob/main/\detokenize{reflection.v}\#L1378}{\texttt{\detokenize{bwdAss}}}] \emph{Definition, line 1378}. \textsf{[aux]} \\
  (no docstring)

\item[\href{https://github.com/LLM4Rocq/mathcomp-eulerian/blob/main/\detokenize{reflection.v}\#L1381}{\texttt{\detokenize{bwd_fwd}}}] \emph{Lemma, line 1381}. \textsf{[aux]} \\
  (no docstring)

\item[\href{https://github.com/LLM4Rocq/mathcomp-eulerian/blob/main/\detokenize{reflection.v}\#L1387}{\texttt{\detokenize{fwd_bwd}}}] \emph{Lemma, line 1387}. \textsf{[aux]} \\
  (no docstring)

\item[\href{https://github.com/LLM4Rocq/mathcomp-eulerian/blob/main/\detokenize{reflection.v}\#L1396}{\texttt{\detokenize{sum_reindex_inner}}}] \emph{Lemma, line 1396}. \textsf{[aux]} \\
  (no docstring)

\item[\href{https://github.com/LLM4Rocq/mathcomp-eulerian/blob/main/\detokenize{reflection.v}\#L1414}{\texttt{\detokenize{imset_lift_eq}}}] \emph{Lemma, line 1414}. \textsf{[aux]} \\
  (no docstring)

\item[\href{https://github.com/LLM4Rocq/mathcomp-eulerian/blob/main/\detokenize{reflection.v}\#L1438}{\texttt{\detokenize{interior_descent_set_eq}}}] \emph{Lemma, line 1438}. \textsf{[aux]} \\
  (no docstring)

\item[\href{https://github.com/LLM4Rocq/mathcomp-eulerian/blob/main/\detokenize{reflection.v}\#L1623}{\texttt{\detokenize{alt_desc_set_is_alt}}}] \emph{Lemma, line 1623}. \textsf{[aux]} \\
  (no docstring)

\item[\href{https://github.com/LLM4Rocq/mathcomp-eulerian/blob/main/\detokenize{reflection.v}\#L1633}{\texttt{\detokenize{set_is_alt_indicator}}}] \emph{Lemma, line 1633}. \textsf{[aux]} \\
  (no docstring)

\item[\href{https://github.com/LLM4Rocq/mathcomp-eulerian/blob/main/\detokenize{reflection.v}\#L1664}{\texttt{\detokenize{alt_plus_nalt_as_set_is_alt_sum}}}] \emph{Lemma, line 1664}. \textsf{[aux]} \\
  (no docstring)

\item[\href{https://github.com/LLM4Rocq/mathcomp-eulerian/blob/main/\detokenize{reflection.v}\#L1680}{\texttt{\detokenize{sum_descent_set_indicator}}}] \emph{Lemma, line 1680}. \textsf{[aux]} \\
  (no docstring)

\item[\href{https://github.com/LLM4Rocq/mathcomp-eulerian/blob/main/\detokenize{reflection.v}\#L1688}{\texttt{\detokenize{setU1_imset_lift_eq}}}] \emph{Lemma, line 1688}. \textsf{[aux]} \\
  (no docstring)

\item[\href{https://github.com/LLM4Rocq/mathcomp-eulerian/blob/main/\detokenize{reflection.v}\#L1707}{\texttt{\detokenize{nat_of_andb}}}] \emph{Lemma, line 1707}. \textsf{[aux]} \\
  (no docstring)

\item[\href{https://github.com/LLM4Rocq/mathcomp-eulerian/blob/main/\detokenize{reflection.v}\#L1716}{\texttt{\detokenize{sum_set_is_alt_ord0}}}] \emph{Lemma, line 1716}. \textsf{[aux]} \\
  (no docstring)

\item[\href{https://github.com/LLM4Rocq/mathcomp-eulerian/blob/main/\detokenize{reflection.v}\#L1738}{\texttt{\detokenize{sum_set_is_alt_ord_max}}}] \emph{Lemma, line 1738}. \textsf{[aux]} \\
  (no docstring)

\item[\href{https://github.com/LLM4Rocq/mathcomp-eulerian/blob/main/\detokenize{reflection.v}\#L1770}{\texttt{\detokenize{andre_target}}}] \emph{Definition, line 1770}. \textsf{[aux]} \\
  (no docstring)

\item[\href{https://github.com/LLM4Rocq/mathcomp-eulerian/blob/main/\detokenize{reflection.v}\#L1773}{\texttt{\detokenize{interior_set_is_alt}}}] \emph{Lemma, line 1773}. \textsf{[aux]} \\
  (no docstring)

\item[\href{https://github.com/LLM4Rocq/mathcomp-eulerian/blob/main/\detokenize{reflection.v}\#L1801}{\texttt{\detokenize{andre_split_lt}}}] \emph{Lemma, line 1801}. \textsf{[aux]} \\
  (no docstring)

\item[\href{https://github.com/LLM4Rocq/mathcomp-eulerian/blob/main/\detokenize{reflection.v}\#L1806}{\texttt{\detokenize{andre_s_ltn}}}] \emph{Lemma, line 1806}. \textsf{[aux]} \\
  (no docstring)

\item[\href{https://github.com/LLM4Rocq/mathcomp-eulerian/blob/main/\detokenize{reflection.v}\#L1810}{\texttt{\detokenize{andre_target_lt}}}] \emph{Lemma, line 1810}. \textsf{[aux]} \\
  (no docstring)

\item[\href{https://github.com/LLM4Rocq/mathcomp-eulerian/blob/main/\detokenize{reflection.v}\#L1814}{\texttt{\detokenize{andre_target_at}}}] \emph{Lemma, line 1814}. \textsf{[aux]} \\
  (no docstring)

\item[\href{https://github.com/LLM4Rocq/mathcomp-eulerian/blob/main/\detokenize{reflection.v}\#L1817}{\texttt{\detokenize{andre_target_gt}}}] \emph{Lemma, line 1817}. \textsf{[aux]} \\
  (no docstring)

\item[\href{https://github.com/LLM4Rocq/mathcomp-eulerian/blob/main/\detokenize{reflection.v}\#L1825}{\texttt{\detokenize{andre_union_eq_split}}}] \emph{Lemma, line 1825}. \textsf{[aux]} \\
  (no docstring)

\item[\href{https://github.com/LLM4Rocq/mathcomp-eulerian/blob/main/\detokenize{reflection.v}\#L1872}{\texttt{\detokenize{beta_andre_left}}}] \emph{Lemma, line 1872}. \textsf{[aux]} \\
  (no docstring)

\item[\href{https://github.com/LLM4Rocq/mathcomp-eulerian/blob/main/\detokenize{reflection.v}\#L1883}{\texttt{\detokenize{beta_andre_right}}}] \emph{Lemma, line 1883}. \textsf{[aux]} \\
  (no docstring)

\item[\href{https://github.com/LLM4Rocq/mathcomp-eulerian/blob/main/\detokenize{reflection.v}\#L1893}{\texttt{\detokenize{andre_interior_count}}}] \emph{Lemma, line 1893}. \textsf{[aux]} \\
  (no docstring)

\item[\href{https://github.com/LLM4Rocq/mathcomp-eulerian/blob/main/\detokenize{reflection.v}\#L1947}{\texttt{\detokenize{sum_set_is_alt_eq_andre_sum}}}] \emph{Lemma, line 1947}. \textsf{[aux]} \\
  (no docstring)

\item[\href{https://github.com/LLM4Rocq/mathcomp-eulerian/blob/main/\detokenize{reflection.v}\#L1980}{\texttt{\detokenize{boundary_cancellation_alt}}}] \emph{Lemma, line 1980}. \textsf{[aux]} \\
  (no docstring)

\item[\href{https://github.com/LLM4Rocq/mathcomp-eulerian/blob/main/\detokenize{reflection.v}\#L1988}{\texttt{\detokenize{euler_rec}}}] \emph{Theorem, line 1988}. \textsf{[aux]} \\
  (no docstring)

\end{description}

% --------------------------------------------
\section*{\texttt{stanley\_egf.v}}
\addcontentsline{toc}{section}{\texttt{stanley\_egf.v}}

\begin{description}
\item[\href{https://github.com/LLM4Rocq/mathcomp-eulerian/blob/main/\detokenize{stanley_egf.v}\#L35}{\texttt{\detokenize{rat_char0}}}] \emph{Lemma, line 35}. \textsf{[aux]} \\
  rat has characteristic 0.

\item[\href{https://github.com/LLM4Rocq/mathcomp-eulerian/blob/main/\detokenize{stanley_egf.v}\#L39}{\texttt{\detokenize{euler_egf}}}] \emph{Definition, line 39}. \textsf{[aux]} \\
  The exponential generating function of the Euler numbers, over rat.

\item[\href{https://github.com/LLM4Rocq/mathcomp-eulerian/blob/main/\detokenize{stanley_egf.v}\#L41}{\texttt{\detokenize{euler_egf0}}}] \emph{Lemma, line 41}. \textsf{[aux]} \\
  (no docstring)

\item[\href{https://github.com/LLM4Rocq/mathcomp-eulerian/blob/main/\detokenize{stanley_egf.v}\#L46}{\texttt{\detokenize{euler_egf_ode}}}] \emph{Lemma, line 46}. \textsf{[aux]} \\
  (no docstring)

\item[\href{https://github.com/LLM4Rocq/mathcomp-eulerian/blob/main/\detokenize{stanley_egf.v}\#L66}{\texttt{\detokenize{stanley_1_6_1}}}] \emph{Theorem, line 66}. \textsf{[aux]} \\
  (no docstring)

\item[\href{https://github.com/LLM4Rocq/mathcomp-eulerian/blob/main/\detokenize{stanley_egf.v}\#L73}{\texttt{\detokenize{stanley_1_6_1_coef}}}] \emph{Corollary, line 73}. \textsf{[aux]} \\
  Coefficient form: E\_n = n! * \texttt{\detokenize{x^n}} (sec x + tan x).

\end{description}

% --------------------------------------------
\section*{\texttt{stanley\_ogf.v}}
\addcontentsline{toc}{section}{\texttt{stanley\_ogf.v}}

\begin{description}
\item[\href{https://github.com/LLM4Rocq/mathcomp-eulerian/blob/main/\detokenize{stanley_ogf.v}\#L29}{\texttt{\detokenize{pow_ogf}}}] \emph{Definition, line 29}. \textsf{[aux]} \\
  The OGF of the (n+1)-st powers: coefficient m is (m+1)\textasciicircum{}(n+1).

\item[\href{https://github.com/LLM4Rocq/mathcomp-eulerian/blob/main/\detokenize{stanley_ogf.v}\#L39}{\texttt{\detokenize{stanley_1_4_div}}}] \emph{Theorem, line 39}. \textsf{[aux]} \\
  (no docstring)

\item[\href{https://github.com/LLM4Rocq/mathcomp-eulerian/blob/main/\detokenize{stanley_ogf.v}\#L67}{\texttt{\detokenize{stanley_1_4}}}] \emph{Theorem, line 67}. \textsf{[aux]} \\
  (no docstring)

\item[\href{https://github.com/LLM4Rocq/mathcomp-eulerian/blob/main/\detokenize{stanley_ogf.v}\#L75}{\texttt{\detokenize{stanley_1_4_inv}}}] \emph{Corollary, line 75}. \textsf{[aux]} \\
  Display form with the explicit inverse of 1 - x.

\end{description}

% --------------------------------------------
\section*{\texttt{stirling\_fiber.v}}
\addcontentsline{toc}{section}{\texttt{stirling\_fiber.v}}

\begin{description}
\item[\href{https://github.com/LLM4Rocq/mathcomp-eulerian/blob/main/\detokenize{stirling_fiber.v}\#L34}{\texttt{\detokenize{insert_after_fun}}}] \emph{Definition, line 34}. \textsf{[aux]} \\
  \texttt{\detokenize{insert_after_fun j0 s}} is the underlying function of \texttt{\detokenize{insert_after     j0 s}}: it splices \texttt{\detokenize{ord_max}} into the \texttt{\detokenize{j0}}-cycle of \texttt{\detokenize{s}} right after     \texttt{\detokenize{j0}}.  See \texttt{\detokenize{insert_after}} for the perm wrapping.

\item[\href{https://github.com/LLM4Rocq/mathcomp-eulerian/blob/main/\detokenize{stirling_fiber.v}\#L44}{\texttt{\detokenize{insert_after_fun_inj}}}] \emph{Lemma, line 44}. \textsf{[aux]} \\
  \texttt{\detokenize{insert_after_fun j0 s}} is injective on \texttt{\detokenize{'I_n.+1}}.  Used to package it     as a \texttt{\detokenize{perm}} in \texttt{\detokenize{insert_after}}.

\item[\href{https://github.com/LLM4Rocq/mathcomp-eulerian/blob/main/\detokenize{stirling_fiber.v}\#L82}{\texttt{\detokenize{insert_after}}}] \emph{Definition, line 82}. \textsf{[aux]} \\
  \texttt{insert\_after j0 s : \{perm 'I\_n.+1\}} is the permutation that splices     \texttt{\detokenize{ord_max}} into the \texttt{\detokenize{j0}}-cycle of \texttt{\detokenize{s}} right after \texttt{\detokenize{j0}}.  An \texttt{\detokenize{s}}-cycle     \texttt{\detokenize{p1 -> ... -> j0 -> s j0 -> ...}} becomes     \texttt{\detokenize{p1 -> ... -> j0 -> ord_max -> s j0 -> ...}}.

\item[\href{https://github.com/LLM4Rocq/mathcomp-eulerian/blob/main/\detokenize{stirling_fiber.v}\#L86}{\texttt{\detokenize{insert_afterE}}}] \emph{Lemma, line 86}. \textsf{[aux]} \\
  \texttt{\detokenize{insert_after}} unfolds to \texttt{\detokenize{insert_after_fun}} pointwise.

\item[\href{https://github.com/LLM4Rocq/mathcomp-eulerian/blob/main/\detokenize{stirling_fiber.v}\#L92}{\texttt{\detokenize{insert_after_ord_max}}}] \emph{Lemma, line 92}. \textsf{[aux]} \\
  Defining equation: \texttt{\detokenize{insert_after j0 s}} sends \texttt{\detokenize{ord_max}} to     \texttt{\detokenize{lift ord_max (s j0)}}.

\item[\href{https://github.com/LLM4Rocq/mathcomp-eulerian/blob/main/\detokenize{stirling_fiber.v}\#L98}{\texttt{\detokenize{insert_after_j0}}}] \emph{Lemma, line 98}. \textsf{[aux]} \\
  Defining equation: \texttt{\detokenize{insert_after j0 s}} sends \texttt{\detokenize{lift ord_max j0}} to     \texttt{\detokenize{ord_max}} (the spliced-in element follows \texttt{\detokenize{j0}} in its cycle).

\item[\href{https://github.com/LLM4Rocq/mathcomp-eulerian/blob/main/\detokenize{stirling_fiber.v}\#L104}{\texttt{\detokenize{insert_after_other}}}] \emph{Lemma, line 104}. \textsf{[aux]} \\
  Defining equation: on points \texttt{\detokenize{lift ord_max k}} with \texttt{\detokenize{k != j0}},     \texttt{\detokenize{insert_after j0 s}} acts like the \texttt{\detokenize{lift ord_max}}-extension of \texttt{\detokenize{s}}.

\item[\href{https://github.com/LLM4Rocq/mathcomp-eulerian/blob/main/\detokenize{stirling_fiber.v}\#L113}{\texttt{\detokenize{insert_after_ord_max_neq}}}] \emph{Lemma, line 113}. \textsf{[aux]} \\
  \texttt{\detokenize{insert_after j0 s}} does not fix \texttt{\detokenize{ord_max}}.  Used to characterise     the image of \texttt{\detokenize{insert_after}} in the bijection of \textbackslash{}S\{\}D.

\item[\href{https://github.com/LLM4Rocq/mathcomp-eulerian/blob/main/\detokenize{stirling_fiber.v}\#L149}{\texttt{\detokenize{insert_after_decomp}}}] \emph{Lemma, line 149}. \textsf{[aux]} \\
  Decomposition: \texttt{\detokenize{insert_after j0 s}} factors as the transposition     \texttt{\detokenize{tperm ord_max (lift ord_max j0)}} composed with \texttt{\detokenize{perm.lift_perm     ord_max ord_max s}}.  Used in \texttt{\detokenize{cycle_count_insert_after}} to apply     \texttt{\detokenize{porbits_mul_tperm}}.

\item[\href{https://github.com/LLM4Rocq/mathcomp-eulerian/blob/main/\detokenize{stirling_fiber.v}\#L169}{\texttt{\detokenize{cycle_count_insert_after}}}] \emph{Lemma, line 169}. \textsf{[aux]} \\
  Key invariant: \texttt{\detokenize{insert_after j0 s}} preserves the cycle count of \texttt{\detokenize{s}}.     Splicing \texttt{\detokenize{ord_max}} into the \texttt{\detokenize{j0}}-cycle widens that cycle by one     element without creating or destroying cycles.

\item[\href{https://github.com/LLM4Rocq/mathcomp-eulerian/blob/main/\detokenize{stirling_fiber.v}\#L201}{\texttt{\detokenize{insert_after_inj_pair}}}] \emph{Lemma, line 201}. \textsf{[aux]} \\
  \texttt{\detokenize{insert_after}} is injective in the pair \texttt{\detokenize{(j0, s)}}: distinct     \texttt{\detokenize{(j0, s)}} yield distinct permutations of \texttt{\detokenize{'I_n.+1}}.

\item[\href{https://github.com/LLM4Rocq/mathcomp-eulerian/blob/main/\detokenize{stirling_fiber.v}\#L228}{\texttt{\detokenize{insert_after_surj}}}] \emph{Lemma, line 228}. \textsf{[aux]} \\
  Surjectivity onto the non-fixing class: every \texttt{sigma : \{perm 'I\_n.+1\}}     with \texttt{\detokenize{sigma ord_max != ord_max}} arises as \texttt{\detokenize{insert_after j0 s}} for     some \texttt{\detokenize{j0}} and \texttt{\detokenize{s}}.  Constructed by conjugating \texttt{\detokenize{sigma}} with a     transposition to land in the fixed-\texttt{\detokenize{ord_max}} class, then applying     \texttt{\detokenize{drop_perm}}.

\item[\href{https://github.com/LLM4Rocq/mathcomp-eulerian/blob/main/\detokenize{stirling_fiber.v}\#L266}{\texttt{\detokenize{card_neq_ord_max_fiber}}}] \emph{Lemma, line 266}. \textsf{[aux]} \\
  Count of permutations \texttt{\detokenize{sigma}} of \texttt{\detokenize{'I_n.+1}} with \texttt{\detokenize{sigma ord_max !=     ord_max}} and exactly \texttt{\detokenize{k.+1}} cycles equals \texttt{\detokenize{n * stirling_c n k.+1}}.     This is the \texttt{\detokenize{j != ord_max}} half of the Stirling fiber bijection.

\item[\href{https://github.com/LLM4Rocq/mathcomp-eulerian/blob/main/\detokenize{stirling_fiber.v}\#L309}{\texttt{\detokenize{card_eq_ord_max_fiber}}}] \emph{Lemma, line 309}. \textsf{[aux]} \\
  Count of permutations \texttt{\detokenize{sigma}} of \texttt{\detokenize{'I_n.+1}} fixing \texttt{\detokenize{ord_max}} with     exactly \texttt{\detokenize{k.+1}} cycles equals \texttt{\detokenize{stirling_c n k}}.  Bijection via     \texttt{\detokenize{perm.lift_perm ord_max ord_max}} and \texttt{\detokenize{drop_perm}}.  This is the     \texttt{\detokenize{j = ord_max}} half of the Stirling fiber bijection.

\item[\href{https://github.com/LLM4Rocq/mathcomp-eulerian/blob/main/\detokenize{stirling_fiber.v}\#L364}{\texttt{\detokenize{stirling_c_rec}}}] \emph{Lemma, line 364}. \textsf{[Ch 2 §2.2]} \\
  Stirling cycle-number recurrence:     \texttt{\detokenize{c(n+1, k+1) = n * c(n, k+1) + c(n, k)}} (Stanley EC1 \textbackslash{}S\{\}1.3.2).     Final assembly from \texttt{\detokenize{card_neq_ord_max_fiber}} and     \texttt{\detokenize{card_eq_ord_max_fiber}}.

\end{description}

% --------------------------------------------
\section*{\texttt{worpitzky.v}}
\addcontentsline{toc}{section}{\texttt{worpitzky.v}}

\begin{description}
\item[\href{https://github.com/LLM4Rocq/mathcomp-eulerian/blob/main/\detokenize{worpitzky.v}\#L33}{\texttt{\detokenize{worpitzky_bin_step}}}] \emph{Lemma, line 33}. \textsf{[aux]} \\
  (no docstring)

\item[\href{https://github.com/LLM4Rocq/mathcomp-eulerian/blob/main/\detokenize{worpitzky.v}\#L56}{\texttt{\detokenize{worpitzky}}}] \emph{Theorem, line 56}. \textsf{[aux]} \\
  (no docstring)

\item[\href{https://github.com/LLM4Rocq/mathcomp-eulerian/blob/main/\detokenize{worpitzky.v}\#L92}{\texttt{\detokenize{coef_eul_pol}}}] \emph{Lemma, line 92}. \textsf{[aux]} \\
  (no docstring)

\end{description}
